%% Generated by Sphinx.
\def\sphinxdocclass{jupyterBook}
\documentclass[letterpaper,10pt,english]{jupyterBook}
\ifdefined\pdfpxdimen
   \let\sphinxpxdimen\pdfpxdimen\else\newdimen\sphinxpxdimen
\fi \sphinxpxdimen=.75bp\relax
\ifdefined\pdfimageresolution
    \pdfimageresolution= \numexpr \dimexpr1in\relax/\sphinxpxdimen\relax
\fi
%% let collapsible pdf bookmarks panel have high depth per default
\PassOptionsToPackage{bookmarksdepth=5}{hyperref}
%% turn off hyperref patch of \index as sphinx.xdy xindy module takes care of
%% suitable \hyperpage mark-up, working around hyperref-xindy incompatibility
\PassOptionsToPackage{hyperindex=false}{hyperref}
%% memoir class requires extra handling
\makeatletter\@ifclassloaded{memoir}
{\ifdefined\memhyperindexfalse\memhyperindexfalse\fi}{}\makeatother

\PassOptionsToPackage{booktabs}{sphinx}
\PassOptionsToPackage{colorrows}{sphinx}

\PassOptionsToPackage{warn}{textcomp}

\catcode`^^^^00a0\active\protected\def^^^^00a0{\leavevmode\nobreak\ }
\usepackage{cmap}
\usepackage{fontspec}
\defaultfontfeatures[\rmfamily,\sffamily,\ttfamily]{}
\usepackage{amsmath,amssymb,amstext}
\usepackage{polyglossia}
\setmainlanguage{english}



\setmainfont{FreeSerif}[
  Extension      = .otf,
  UprightFont    = *,
  ItalicFont     = *Italic,
  BoldFont       = *Bold,
  BoldItalicFont = *BoldItalic
]
\setsansfont{FreeSans}[
  Extension      = .otf,
  UprightFont    = *,
  ItalicFont     = *Oblique,
  BoldFont       = *Bold,
  BoldItalicFont = *BoldOblique,
]
\setmonofont{FreeMono}[
  Extension      = .otf,
  UprightFont    = *,
  ItalicFont     = *Oblique,
  BoldFont       = *Bold,
  BoldItalicFont = *BoldOblique,
]



\usepackage[Bjarne]{fncychap}
\usepackage[,numfigreset=2,mathnumfig]{sphinx}

\fvset{fontsize=\small}
\usepackage{geometry}


% Include hyperref last.
\usepackage{hyperref}
% Fix anchor placement for figures with captions.
\usepackage{hypcap}% it must be loaded after hyperref.
% Set up styles of URL: it should be placed after hyperref.
\urlstyle{same}


\usepackage{sphinxmessages}



        % Start of preamble defined in sphinx-jupyterbook-latex %
         \usepackage[Latin,Greek]{ucharclasses}
        \usepackage{unicode-math}
        % fixing title of the toc
        \addto\captionsenglish{\renewcommand{\contentsname}{Contents}}
        \hypersetup{
            pdfencoding=auto,
            psdextra
        }
        % End of preamble defined in sphinx-jupyterbook-latex %
        

\title{PSSP Virtual Notebook}
\date{Oct 01, 2025}
\release{}
\author{Matteo Salvalaglio}
\newcommand{\sphinxlogo}{\vbox{}}
\renewcommand{\releasename}{}
\makeindex
\begin{document}

\pagestyle{empty}
\sphinxmaketitle
\pagestyle{plain}
\sphinxtableofcontents
\pagestyle{normal}
\phantomsection\label{\detokenize{index::doc}}


\sphinxAtStartPar
This is a collection of notes and exercises associated with the \sphinxhref{https://moodle.ucl.ac.uk/course/view.php?id=1191}{Particulate Systems and Separations II module}, part of the second year ChemEng curriculum at UCL{]}(\sphinxurl{https://www.ucl.ac.uk/chemical-engineering/}).

\sphinxAtStartPar
These note are intended a complement to the lecture material available to UCL students.

\begin{sphinxadmonition}{note}{Contribute to evolve these notes!}

\sphinxAtStartPar
You can:
\begin{itemize}
\item {} 
\sphinxAtStartPar
Submit an issue (using the GitHub button, on the top)

\item {} 
\sphinxAtStartPar
Suggest an edit (also using the GitHub button, on the top)

\item {} 
\sphinxAtStartPar
Add a new notebook solving an exercise or developing a new topic (using GitHub and/or getting in touch via email at \sphinxhref{mailto:m.salvalaglio@ucl.ac.uk}{m.salvalaglio@ucl.ac.uk})

\end{itemize}
\end{sphinxadmonition}

\sphinxstepscope


\section{1. Material Balances in a Gas Permeation Unit}
\label{\detokenize{Gas_permeation:material-balances-in-a-gas-permeation-unit}}\label{\detokenize{Gas_permeation::doc}}

\subsection{1.1 Operating Equation}
\label{\detokenize{Gas_permeation:operating-equation}}
\sphinxAtStartPar
Consider a binary gas phase mixture fed into a perfectly mixed membrane separation unit with molar flow rate \(F_{in}\), and composition \(\vec{z}\). The retentate stream is characterised by molar flow rate \(F_{r}\), and composition \(\vec{x}\), while  permeate is characterised by molar flow rate \(F_{p}\)  and composition \(\vec{y}\).

\sphinxAtStartPar
Since we are dealing with a binary mixture, in all streams the compositions of A and B complement to 1 (e.g. for the permeate stream \(y_B=1-y_A\)). Hence we drop unnecessary subscripts and indicate with \(z\), \(x\), and \(y\) the molar fraction of component A in the mixture.

\sphinxAtStartPar
The global material balance on the unit reads: \(F_{in}=F_p+F_r\)

\sphinxAtStartPar
The material balance for component A reads: \(F_{in}z=F_py+F_rx\)

\sphinxAtStartPar
Combining these two equations:
\begin{equation*}
\begin{split}
F_{in}z=F_py+(F_{in}-F_p)x
\end{split}
\end{equation*}\begin{equation}\label{equation:Gas_permeation:GPeq1}
\begin{split}
z=\frac{F_p}{F_{in}}y+(1-\frac{F_p}{F_{in}})x
\end{split}
\end{equation}
\sphinxAtStartPar
the quantity \(\frac{F_p}{F_{in}}\) is commonly defined as cut and indicated as \(\theta\).

\sphinxAtStartPar
Solving Eq.\eqref{equation:Gas_permeation:GPeq1} for \(y\) leads to the expression:
\begin{equation}\label{equation:Gas_permeation:GPeq:Operating}
\begin{split}
y=\frac{\left(\theta-1\right)}{\theta}x+\frac{1}{\theta}z
\end{split}
\end{equation}
\sphinxAtStartPar
Eq. \eqref{equation:Gas_permeation:GPeq:Operating} is typically referred to as the Operating Equation of a perfectly mixed gas phase membrane module.


\subsubsection{Graphical Representation}
\label{\detokenize{Gas_permeation:graphical-representation}}
\begin{sphinxuseclass}{cell}\begin{sphinxVerbatimInput}

\begin{sphinxuseclass}{cell_input}
\begin{sphinxVerbatim}[commandchars=\\\{\}]
\PYG{k+kn}{import}\PYG{+w}{ }\PYG{n+nn}{numpy}\PYG{+w}{ }\PYG{k}{as}\PYG{+w}{ }\PYG{n+nn}{np}
\PYG{k+kn}{import}\PYG{+w}{ }\PYG{n+nn}{matplotlib}\PYG{n+nn}{.}\PYG{n+nn}{pyplot}\PYG{+w}{ }\PYG{k}{as}\PYG{+w}{ }\PYG{n+nn}{plt} 

\PYG{c+c1}{\PYGZsh{} Parameters: }
\PYG{c+c1}{\PYGZsh{} cut}
\PYG{n}{theta} \PYG{o}{=} \PYG{l+m+mf}{0.7}
\PYG{c+c1}{\PYGZsh{} molar fraction in the feed}
\PYG{n}{z} \PYG{o}{=} \PYG{l+m+mf}{0.3}
\PYG{n}{N} \PYG{o}{=} \PYG{l+m+mi}{100} \PYG{c+c1}{\PYGZsh{}number of points}
\PYG{n}{x\PYGZus{}op} \PYG{o}{=} \PYG{n}{np}\PYG{o}{.}\PYG{n}{linspace}\PYG{p}{(}\PYG{l+m+mi}{0}\PYG{p}{,} \PYG{l+m+mi}{1}\PYG{p}{,} \PYG{n}{N}\PYG{p}{)}

\PYG{c+c1}{\PYGZsh{}Operating Equation}
\PYG{n}{y\PYGZus{}op} \PYG{o}{=} \PYG{p}{(}\PYG{p}{(}\PYG{n}{theta} \PYG{o}{\PYGZhy{}} \PYG{l+m+mi}{1}\PYG{p}{)} \PYG{o}{/} \PYG{n}{theta}\PYG{p}{)} \PYG{o}{*} \PYG{n}{x\PYGZus{}op} \PYG{o}{+} \PYG{p}{(}\PYG{n}{z} \PYG{o}{/} \PYG{n}{theta}\PYG{p}{)} 

\PYG{c+c1}{\PYGZsh{}Plotting}
\PYG{n}{figure}\PYG{o}{=}\PYG{n}{plt}\PYG{o}{.}\PYG{n}{figure}\PYG{p}{(}\PYG{p}{)}
\PYG{n}{axes} \PYG{o}{=} \PYG{n}{figure}\PYG{o}{.}\PYG{n}{add\PYGZus{}axes}\PYG{p}{(}\PYG{p}{[}\PYG{l+m+mf}{0.1}\PYG{p}{,}\PYG{l+m+mf}{0.1}\PYG{p}{,}\PYG{l+m+mf}{0.8}\PYG{p}{,}\PYG{l+m+mf}{0.8}\PYG{p}{]}\PYG{p}{)}
\PYG{n}{plt}\PYG{o}{.}\PYG{n}{xticks}\PYG{p}{(}\PYG{n}{fontsize}\PYG{o}{=}\PYG{l+m+mi}{14}\PYG{p}{)}
\PYG{n}{plt}\PYG{o}{.}\PYG{n}{yticks}\PYG{p}{(}\PYG{n}{fontsize}\PYG{o}{=}\PYG{l+m+mi}{14}\PYG{p}{)}
\PYG{n}{axes}\PYG{o}{.}\PYG{n}{plot}\PYG{p}{(}\PYG{n}{x\PYGZus{}op}\PYG{p}{,}\PYG{n}{y\PYGZus{}op}\PYG{p}{,} \PYG{n}{marker}\PYG{o}{=}\PYG{l+s+s1}{\PYGZsq{}}\PYG{l+s+s1}{ }\PYG{l+s+s1}{\PYGZsq{}} \PYG{p}{,} \PYG{n}{color}\PYG{o}{=}\PYG{l+s+s1}{\PYGZsq{}}\PYG{l+s+s1}{r}\PYG{l+s+s1}{\PYGZsq{}}\PYG{p}{)}
\PYG{n}{axes}\PYG{o}{.}\PYG{n}{set\PYGZus{}xlim}\PYG{p}{(}\PYG{p}{[}\PYG{l+m+mi}{0}\PYG{p}{,}\PYG{l+m+mi}{1}\PYG{p}{]}\PYG{p}{)}
\PYG{n}{axes}\PYG{o}{.}\PYG{n}{set\PYGZus{}ylim}\PYG{p}{(}\PYG{p}{[}\PYG{l+m+mi}{0}\PYG{p}{,}\PYG{l+m+mi}{1}\PYG{p}{]}\PYG{p}{)}\PYG{p}{;}
\PYG{n}{plt}\PYG{o}{.}\PYG{n}{title}\PYG{p}{(}\PYG{l+s+s1}{\PYGZsq{}}\PYG{l+s+s1}{Operating Equation}\PYG{l+s+s1}{\PYGZsq{}}\PYG{p}{,} \PYG{n}{fontsize}\PYG{o}{=}\PYG{l+m+mi}{18}\PYG{p}{)}\PYG{p}{;}
\PYG{n}{axes}\PYG{o}{.}\PYG{n}{set\PYGZus{}xlabel}\PYG{p}{(}\PYG{l+s+s1}{\PYGZsq{}}\PYG{l+s+s1}{Molar Fraction in the Retentate}\PYG{l+s+s1}{\PYGZsq{}}\PYG{p}{,} \PYG{n}{fontsize}\PYG{o}{=}\PYG{l+m+mi}{14}\PYG{p}{)}\PYG{p}{;}
\PYG{n}{axes}\PYG{o}{.}\PYG{n}{set\PYGZus{}ylabel}\PYG{p}{(}\PYG{l+s+s1}{\PYGZsq{}}\PYG{l+s+s1}{Molar Fraction in the Permeate}\PYG{l+s+s1}{\PYGZsq{}}\PYG{p}{,}\PYG{n}{fontsize}\PYG{o}{=}\PYG{l+m+mi}{14}\PYG{p}{)}\PYG{p}{;}
\end{sphinxVerbatim}

\end{sphinxuseclass}\end{sphinxVerbatimInput}
\begin{sphinxVerbatimOutput}

\begin{sphinxuseclass}{cell_output}
\noindent\sphinxincludegraphics{{621c69282ab6f654ded6b9f906d37c56be8017640c6a94743bdd10837856632c}.png}

\end{sphinxuseclass}\end{sphinxVerbatimOutput}

\end{sphinxuseclass}

\subsection{1.2 Rate Transfer Equation}
\label{\detokenize{Gas_permeation:rate-transfer-equation}}
\sphinxAtStartPar
Let us begin by considering that the permeate flow rate of component A is given by its molar flux through the membrane multiplied by the membrane total area:
\begin{equation}\label{equation:Gas_permeation:GPeq2}
\begin{split}
F_py=J\rho_AA_M
\end{split}
\end{equation}
\sphinxAtStartPar
where \(J\) is the volumetric flux across the membrane, \(\rho_A\) is the molar density of component A, and \(A_M\) is the total area of the membrane. Applying the solution/diffusion model to define the volumetric flux Eq. \eqref{equation:Gas_permeation:GPeq2} becomes:
\begin{equation}\label{equation:Gas_permeation:GPeq3}
\begin{split}
F_py=\frac{P_A}{l}\rho_AA_M\left(xp_r-yp_p\right)
\end{split}
\end{equation}
\sphinxAtStartPar
where \(P_A\) is the permeability of the membrane to component \(A\), \(l\) is the thickness of the membrane, \(p_r\) is the pressure on the retentate side, \(p_p\) is the pressure on the permeate side.

\sphinxAtStartPar
The same expression can be written for component \(B\):
\begin{equation}\label{equation:Gas_permeation:GPeq4}
\begin{split}
F_p(1-y)=\frac{P_B}{l}\rho_BA_M\left((1-x)p_r-(1-y)p_p\right)
\end{split}
\end{equation}
\sphinxAtStartPar
Both Eq. \eqref{equation:Gas_permeation:GPeq3} and \eqref{equation:Gas_permeation:GPeq4} can be solved for \(F_p\) and equated, leading to:
\begin{equation}\label{equation:Gas_permeation:GPeq5}
\begin{split}
\frac{P_A}{y}\rho_A\left(xp_r-yp_p\right)=\frac{P_B}{(1-y)}\rho_B\left((1-x)p_r-(1-y)p_p\right)
\end{split}
\end{equation}
\sphinxAtStartPar
Eq. \eqref{equation:Gas_permeation:GPeq5} can be solved either for \(y\) or for \(x\), the latter being simpler. After some algebraic manipulation the solution for \(x\) is:
\begin{equation}\label{equation:Gas_permeation:GPeq6}
\begin{split}
x=\frac{y\left[1+\frac{p_p}{p_r}(1-y)\left(\alpha_{AB}\frac{\rho_A}{\rho_B}-1\right)\right]}{\alpha_{AB}\frac{\rho_A}{\rho_B}-\left(\alpha_{AB}\frac{\rho_A}{\rho_B}-1\right)y}
\end{split}
\end{equation}
\sphinxAtStartPar
where \(\alpha_{A,B}=\frac{P_A}{P_B}\) is the ideal separation factor.

\sphinxAtStartPar
Eq. \eqref{equation:Gas_permeation:GPeq6} can be simplified when the gas phase can be considered ideal. In such case \(\rho_A=\rho_B=P/RT\), yielding:
\begin{equation}\label{equation:Gas_permeation:GPeq7}
\begin{split}
x=\frac{y\left[1+\frac{p_p}{p_r}(1-y)\left(\alpha_{AB}-1\right)\right]}{\alpha_{AB}-\left(\alpha_{AB}-1\right)y}
\end{split}
\end{equation}
\sphinxAtStartPar
Eq. \eqref{equation:Gas_permeation:GPeq7} is called \sphinxstyleemphasis{rate transfer equation} and together with th operating equation identify the membrane operating conditions in composition space.


\subsubsection{Graphical Representation}
\label{\detokenize{Gas_permeation:id1}}
\begin{sphinxuseclass}{cell}\begin{sphinxVerbatimInput}

\begin{sphinxuseclass}{cell_input}
\begin{sphinxVerbatim}[commandchars=\\\{\}]
\PYG{c+c1}{\PYGZsh{} Parameters: }
\PYG{c+c1}{\PYGZsh{} Ratio between the permeate and retentate side pressures}
\PYG{n}{Pp\PYGZus{}Pr} \PYG{o}{=} \PYG{l+m+mf}{0.3}
\PYG{c+c1}{\PYGZsh{} ideal separation factor}
\PYG{n}{alpha} \PYG{o}{=} \PYG{l+m+mi}{100}
\PYG{n}{N} \PYG{o}{=} \PYG{l+m+mi}{100} \PYG{c+c1}{\PYGZsh{}number of points}
\PYG{n}{y\PYGZus{}rt} \PYG{o}{=} \PYG{n}{np}\PYG{o}{.}\PYG{n}{linspace}\PYG{p}{(}\PYG{l+m+mi}{0}\PYG{p}{,} \PYG{l+m+mi}{1}\PYG{p}{,} \PYG{n}{N}\PYG{p}{)}

\PYG{c+c1}{\PYGZsh{}Rate Transfer}
\PYG{n}{x\PYGZus{}rt} \PYG{o}{=} \PYG{n}{y\PYGZus{}rt} \PYG{o}{*} \PYG{p}{(}\PYG{l+m+mi}{1} \PYG{o}{+} \PYG{n}{Pp\PYGZus{}Pr} \PYG{o}{*} \PYG{p}{(}\PYG{l+m+mi}{1}\PYG{o}{\PYGZhy{}}\PYG{n}{y\PYGZus{}rt}\PYG{p}{)} \PYG{o}{*} \PYG{p}{(}\PYG{n}{alpha} \PYG{o}{\PYGZhy{}} \PYG{l+m+mi}{1}\PYG{p}{)}\PYG{p}{)} \PYG{o}{/} \PYG{p}{(}\PYG{n}{alpha} \PYG{o}{\PYGZhy{}} \PYG{p}{(}\PYG{n}{alpha} \PYG{o}{\PYGZhy{}} \PYG{l+m+mi}{1}\PYG{p}{)} \PYG{o}{*} \PYG{n}{y\PYGZus{}rt}\PYG{p}{)}

\PYG{c+c1}{\PYGZsh{}Plotting}
\PYG{n}{figure}\PYG{o}{=}\PYG{n}{plt}\PYG{o}{.}\PYG{n}{figure}\PYG{p}{(}\PYG{p}{)}
\PYG{n}{axes} \PYG{o}{=} \PYG{n}{figure}\PYG{o}{.}\PYG{n}{add\PYGZus{}axes}\PYG{p}{(}\PYG{p}{[}\PYG{l+m+mf}{0.1}\PYG{p}{,}\PYG{l+m+mf}{0.1}\PYG{p}{,}\PYG{l+m+mf}{1.0}\PYG{p}{,}\PYG{l+m+mf}{1.0}\PYG{p}{]}\PYG{p}{)}
\PYG{n}{plt}\PYG{o}{.}\PYG{n}{xticks}\PYG{p}{(}\PYG{n}{fontsize}\PYG{o}{=}\PYG{l+m+mi}{14}\PYG{p}{)}
\PYG{n}{plt}\PYG{o}{.}\PYG{n}{yticks}\PYG{p}{(}\PYG{n}{fontsize}\PYG{o}{=}\PYG{l+m+mi}{14}\PYG{p}{)}
\PYG{n}{axes}\PYG{o}{.}\PYG{n}{plot}\PYG{p}{(}\PYG{n}{x\PYGZus{}op}\PYG{p}{,}\PYG{n}{y\PYGZus{}op}\PYG{p}{,} \PYG{n}{marker}\PYG{o}{=}\PYG{l+s+s1}{\PYGZsq{}}\PYG{l+s+s1}{ }\PYG{l+s+s1}{\PYGZsq{}} \PYG{p}{,} \PYG{n}{color}\PYG{o}{=}\PYG{l+s+s1}{\PYGZsq{}}\PYG{l+s+s1}{r}\PYG{l+s+s1}{\PYGZsq{}}\PYG{p}{)}
\PYG{n}{axes}\PYG{o}{.}\PYG{n}{plot}\PYG{p}{(}\PYG{n}{x\PYGZus{}rt}\PYG{p}{,}\PYG{n}{y\PYGZus{}rt}\PYG{p}{,} \PYG{n}{marker}\PYG{o}{=}\PYG{l+s+s1}{\PYGZsq{}}\PYG{l+s+s1}{ }\PYG{l+s+s1}{\PYGZsq{}} \PYG{p}{,} \PYG{n}{color}\PYG{o}{=}\PYG{l+s+s1}{\PYGZsq{}}\PYG{l+s+s1}{b}\PYG{l+s+s1}{\PYGZsq{}}\PYG{p}{)}
\PYG{n}{axes}\PYG{o}{.}\PYG{n}{set\PYGZus{}xlim}\PYG{p}{(}\PYG{p}{[}\PYG{l+m+mi}{0}\PYG{p}{,}\PYG{l+m+mi}{1}\PYG{p}{]}\PYG{p}{)}
\PYG{n}{axes}\PYG{o}{.}\PYG{n}{set\PYGZus{}ylim}\PYG{p}{(}\PYG{p}{[}\PYG{l+m+mi}{0}\PYG{p}{,}\PYG{l+m+mi}{1}\PYG{p}{]}\PYG{p}{)}
\PYG{n}{plt}\PYG{o}{.}\PYG{n}{title}\PYG{p}{(}\PYG{l+s+s1}{\PYGZsq{}}\PYG{l+s+s1}{Operating and Rate Transfer Equations}\PYG{l+s+s1}{\PYGZsq{}}\PYG{p}{,} \PYG{n}{fontsize}\PYG{o}{=}\PYG{l+m+mi}{18}\PYG{p}{)}\PYG{p}{;}
\PYG{n}{axes}\PYG{o}{.}\PYG{n}{set\PYGZus{}xlabel}\PYG{p}{(}\PYG{l+s+s1}{\PYGZsq{}}\PYG{l+s+s1}{Molar Fraction in the Retentate}\PYG{l+s+s1}{\PYGZsq{}}\PYG{p}{,} \PYG{n}{fontsize}\PYG{o}{=}\PYG{l+m+mi}{14}\PYG{p}{)}\PYG{p}{;}
\PYG{n}{axes}\PYG{o}{.}\PYG{n}{set\PYGZus{}ylabel}\PYG{p}{(}\PYG{l+s+s1}{\PYGZsq{}}\PYG{l+s+s1}{Molar Fraction in the Permeate}\PYG{l+s+s1}{\PYGZsq{}}\PYG{p}{,}\PYG{n}{fontsize}\PYG{o}{=}\PYG{l+m+mi}{14}\PYG{p}{)}\PYG{p}{;}
\end{sphinxVerbatim}

\end{sphinxuseclass}\end{sphinxVerbatimInput}
\begin{sphinxVerbatimOutput}

\begin{sphinxuseclass}{cell_output}
\noindent\sphinxincludegraphics{{bdb53b9430c1f929f803826956edb2f8069637c7d81f3ca477582090a7c082ce}.png}

\end{sphinxuseclass}\end{sphinxVerbatimOutput}

\end{sphinxuseclass}
\sphinxAtStartPar
The operating conditions correspond to the interasection of the RT and OP equations, i.e. the solution of a system of equations defined by the operating and rate transfer equations.


\subsubsection{Numerical Solution}
\label{\detokenize{Gas_permeation:numerical-solution}}
\begin{sphinxuseclass}{cell}\begin{sphinxVerbatimInput}

\begin{sphinxuseclass}{cell_input}
\begin{sphinxVerbatim}[commandchars=\\\{\}]
\PYG{k+kn}{from}\PYG{+w}{ }\PYG{n+nn}{scipy}\PYG{n+nn}{.}\PYG{n+nn}{optimize}\PYG{+w}{ }\PYG{k+kn}{import} \PYG{n}{fsolve}

\PYG{k}{def}\PYG{+w}{ }\PYG{n+nf}{equations}\PYG{p}{(}\PYG{n+nb}{vars}\PYG{p}{)}\PYG{p}{:}
    \PYG{n}{x}\PYG{p}{,} \PYG{n}{y} \PYG{o}{=} \PYG{n+nb}{vars}
    \PYG{n}{eq1} \PYG{o}{=} \PYG{p}{(}\PYG{p}{(}\PYG{n}{theta} \PYG{o}{\PYGZhy{}} \PYG{l+m+mi}{1}\PYG{p}{)} \PYG{o}{/} \PYG{n}{theta}\PYG{p}{)} \PYG{o}{*} \PYG{n}{x} \PYG{o}{+} \PYG{p}{(}\PYG{n}{z} \PYG{o}{/} \PYG{n}{theta}\PYG{p}{)} \PYG{o}{\PYGZhy{}} \PYG{n}{y}
    \PYG{n}{eq2} \PYG{o}{=} \PYG{n}{y} \PYG{o}{*} \PYG{p}{(}\PYG{l+m+mi}{1} \PYG{o}{+} \PYG{n}{Pp\PYGZus{}Pr} \PYG{o}{*} \PYG{p}{(}\PYG{l+m+mi}{1}\PYG{o}{\PYGZhy{}}\PYG{n}{y}\PYG{p}{)} \PYG{o}{*} \PYG{p}{(}\PYG{n}{alpha} \PYG{o}{\PYGZhy{}} \PYG{l+m+mi}{1}\PYG{p}{)}\PYG{p}{)} \PYG{o}{/} \PYG{p}{(}\PYG{n}{alpha} \PYG{o}{\PYGZhy{}} \PYG{p}{(}\PYG{n}{alpha} \PYG{o}{\PYGZhy{}} \PYG{l+m+mi}{1}\PYG{p}{)} \PYG{o}{*} \PYG{n}{y}\PYG{p}{)} \PYG{o}{\PYGZhy{}} \PYG{n}{x}
    \PYG{k}{return} \PYG{p}{[}\PYG{n}{eq1}\PYG{p}{,} \PYG{n}{eq2}\PYG{p}{]}

\PYG{n}{x\PYGZus{}set}\PYG{p}{,} \PYG{n}{y\PYGZus{}set} \PYG{o}{=}  \PYG{n}{fsolve}\PYG{p}{(}\PYG{n}{equations}\PYG{p}{,} \PYG{p}{(}\PYG{l+m+mi}{1}\PYG{p}{,} \PYG{l+m+mi}{1}\PYG{p}{)}\PYG{p}{)}


\PYG{c+c1}{\PYGZsh{}Plotting}
\PYG{n}{figure}\PYG{o}{=}\PYG{n}{plt}\PYG{o}{.}\PYG{n}{figure}\PYG{p}{(}\PYG{p}{)}
\PYG{n}{axes} \PYG{o}{=} \PYG{n}{figure}\PYG{o}{.}\PYG{n}{add\PYGZus{}axes}\PYG{p}{(}\PYG{p}{[}\PYG{l+m+mf}{0.1}\PYG{p}{,}\PYG{l+m+mf}{0.1}\PYG{p}{,}\PYG{l+m+mf}{1.0}\PYG{p}{,}\PYG{l+m+mf}{1.0}\PYG{p}{]}\PYG{p}{)}
\PYG{n}{plt}\PYG{o}{.}\PYG{n}{xticks}\PYG{p}{(}\PYG{n}{fontsize}\PYG{o}{=}\PYG{l+m+mi}{14}\PYG{p}{)}
\PYG{n}{plt}\PYG{o}{.}\PYG{n}{yticks}\PYG{p}{(}\PYG{n}{fontsize}\PYG{o}{=}\PYG{l+m+mi}{14}\PYG{p}{)}
\PYG{n}{axes}\PYG{o}{.}\PYG{n}{plot}\PYG{p}{(}\PYG{n}{x\PYGZus{}op}\PYG{p}{,}\PYG{n}{y\PYGZus{}op}\PYG{p}{,} \PYG{n}{marker}\PYG{o}{=}\PYG{l+s+s1}{\PYGZsq{}}\PYG{l+s+s1}{ }\PYG{l+s+s1}{\PYGZsq{}} \PYG{p}{,} \PYG{n}{color}\PYG{o}{=}\PYG{l+s+s1}{\PYGZsq{}}\PYG{l+s+s1}{r}\PYG{l+s+s1}{\PYGZsq{}}\PYG{p}{)}
\PYG{n}{axes}\PYG{o}{.}\PYG{n}{plot}\PYG{p}{(}\PYG{n}{x\PYGZus{}rt}\PYG{p}{,}\PYG{n}{y\PYGZus{}rt}\PYG{p}{,} \PYG{n}{marker}\PYG{o}{=}\PYG{l+s+s1}{\PYGZsq{}}\PYG{l+s+s1}{ }\PYG{l+s+s1}{\PYGZsq{}} \PYG{p}{,} \PYG{n}{color}\PYG{o}{=}\PYG{l+s+s1}{\PYGZsq{}}\PYG{l+s+s1}{b}\PYG{l+s+s1}{\PYGZsq{}}\PYG{p}{)}
\PYG{n}{axes}\PYG{o}{.}\PYG{n}{plot}\PYG{p}{(}\PYG{n}{x\PYGZus{}set}\PYG{p}{,}\PYG{n}{y\PYGZus{}set}\PYG{p}{,} \PYG{n}{marker}\PYG{o}{=}\PYG{l+s+s1}{\PYGZsq{}}\PYG{l+s+s1}{o}\PYG{l+s+s1}{\PYGZsq{}} \PYG{p}{,} \PYG{n}{color}\PYG{o}{=}\PYG{l+s+s1}{\PYGZsq{}}\PYG{l+s+s1}{lime}\PYG{l+s+s1}{\PYGZsq{}}\PYG{p}{,} \PYG{n}{markersize}\PYG{o}{=}\PYG{l+m+mi}{14}\PYG{p}{)}
\PYG{n}{axes}\PYG{o}{.}\PYG{n}{set\PYGZus{}xlim}\PYG{p}{(}\PYG{p}{[}\PYG{l+m+mi}{0}\PYG{p}{,}\PYG{l+m+mi}{1}\PYG{p}{]}\PYG{p}{)}
\PYG{n}{axes}\PYG{o}{.}\PYG{n}{set\PYGZus{}ylim}\PYG{p}{(}\PYG{p}{[}\PYG{l+m+mi}{0}\PYG{p}{,}\PYG{l+m+mi}{1}\PYG{p}{]}\PYG{p}{)}
\PYG{n}{plt}\PYG{o}{.}\PYG{n}{title}\PYG{p}{(}\PYG{l+s+s1}{\PYGZsq{}}\PYG{l+s+s1}{Operating and Rate Transfer Equations}\PYG{l+s+s1}{\PYGZsq{}}\PYG{p}{,} \PYG{n}{fontsize}\PYG{o}{=}\PYG{l+m+mi}{18}\PYG{p}{)}\PYG{p}{;}
\PYG{n}{axes}\PYG{o}{.}\PYG{n}{set\PYGZus{}xlabel}\PYG{p}{(}\PYG{l+s+s1}{\PYGZsq{}}\PYG{l+s+s1}{Molar Fraction in the Retentate}\PYG{l+s+s1}{\PYGZsq{}}\PYG{p}{,} \PYG{n}{fontsize}\PYG{o}{=}\PYG{l+m+mi}{14}\PYG{p}{)}\PYG{p}{;}
\PYG{n}{axes}\PYG{o}{.}\PYG{n}{set\PYGZus{}ylabel}\PYG{p}{(}\PYG{l+s+s1}{\PYGZsq{}}\PYG{l+s+s1}{Molar Fraction in the Permeate}\PYG{l+s+s1}{\PYGZsq{}}\PYG{p}{,}\PYG{n}{fontsize}\PYG{o}{=}\PYG{l+m+mi}{14}\PYG{p}{)}\PYG{p}{;}


\PYG{n+nb}{print}\PYG{p}{(}\PYG{l+s+s2}{\PYGZdq{}}\PYG{l+s+s2}{Operating point \PYGZhy{}\PYGZgt{} x = }\PYG{l+s+s2}{\PYGZdq{}}\PYG{p}{,} \PYG{n}{x\PYGZus{}set}\PYG{p}{,} \PYG{l+s+s2}{\PYGZdq{}}\PYG{l+s+s2}{ y = }\PYG{l+s+s2}{\PYGZdq{}}\PYG{p}{,} \PYG{n}{y\PYGZus{}set}\PYG{p}{)} 
\end{sphinxVerbatim}

\end{sphinxuseclass}\end{sphinxVerbatimInput}
\begin{sphinxVerbatimOutput}

\begin{sphinxuseclass}{cell_output}
\begin{sphinxVerbatim}[commandchars=\\\{\}]
Operating point \PYGZhy{}\PYGZgt{} x =  0.11767288696786887  y =  0.3781401912994848
\end{sphinxVerbatim}

\noindent\sphinxincludegraphics{{a12be8a24878f5bcadb65178a1cd9a92bc70e0329cd3a7c37947650c00f8021e}.png}

\end{sphinxuseclass}\end{sphinxVerbatimOutput}

\end{sphinxuseclass}

\subsection{Contributions}
\label{\detokenize{Gas_permeation:contributions}}\begin{itemize}
\item {} 
\sphinxAtStartPar
Radoslav Lukanov, 16 Feb 2021

\item {} 
\sphinxAtStartPar
Garv Mittal, 4 Feb 2024

\end{itemize}

\sphinxstepscope


\section{2. Pervaporation}
\label{\detokenize{Pervaporation:pervaporation}}\label{\detokenize{Pervaporation::doc}}

\subsection{2.1 Vapour\sphinxhyphen{}Liquid Equilibrium in a Pervaporation unit}
\label{\detokenize{Pervaporation:vapour-liquid-equilibrium-in-a-pervaporation-unit}}
\sphinxAtStartPar
In pervaporation processes the feed/retentate is liquid, while the permeate is in the vapour phase.
In this process, evaporation and permeation of different chemical through a semipermeable membrane take place simultaneously.

\sphinxAtStartPar
In order to tame the complexity of the processes taking place in this unit, a typical modelling assumption is that Vapour liquid equilibrium is established on the feed retentate side and permeation through the membrane takes place in the vapour phase.

\sphinxAtStartPar
Hence, pervaporation can be treated as a gas permeation problem by introducing a \sphinxstyleemphasis{virtual vapour phase} in equilibrium with the liquid on the retentate side of the membrane.

\sphinxAtStartPar
The driving force for a specie \(i\) can thus be written as the difference in the partial pressure for component \(i\) in this virtual vapour phase, and in the real vapour phase on the permeate side.


\subsection{2.2 Vapour\sphinxhyphen{}Liquid Equilibrium conditions}
\label{\detokenize{Pervaporation:vapour-liquid-equilibrium-conditions}}
\sphinxAtStartPar
To express the partial pressure of component \(i\) in a vapour phase in equilibrium with the liquid on the retentate side we shall begin by expressing the fundamental condition that has to be fulfilled by all species at thermodynamic equilibrium:
\begin{equation}\label{equation:Pervaporation:PVAPeq:fund}
\begin{split}
\mu_i^L(T,P,\vec{x})=\mu_i^V(T,P,\vec{y})
\end{split}
\end{equation}
\sphinxAtStartPar
where \(\mu_i^L(T,P,\vec{x})\) and \(\mu_i^V(T,P,\vec{y})\) are the chemical potentials for specie \(i\) in the liquid and vapour phase, respectively.

\sphinxAtStartPar
For mixtures of ideal gases the chemical potential is conveniently expressed in differential form as a function of the partial pressure of component \(i\), \(p_i\) as \(d\mu_i=RTd\ln{p_i}\).\\
In analogy with this expression in the case of real gases fugacity \(f_i\) is introduced, leading to the differential expression \(d\mu_i=RTd\ln{f_i}\).
Integrating this expression between the phases in equilibrium yields:
\begin{equation*}
\begin{split}
\int_{L,T,P,\vec{x}}^{V,T,P,\vec{y}}d\mu_i=\int_{L,T,P,\vec{x}}^{V,T,P,\vec{y}}RTd\ln{f_i}
\end{split}
\end{equation*}\begin{equation}\label{equation:Pervaporation:PVAPeq:chempot}
\begin{split}
\mu_i^V(T,P,\vec{y})-\mu_i^L(T,P,\vec{x})=RT\ln\frac{f_i^V(T,P,\vec{y})}{f_i^L(T,P,\vec{x})}
\end{split}
\end{equation}
\sphinxAtStartPar
Hence, the equilibrium condition expressed by Eq. \eqref{equation:Pervaporation:PVAPeq:fund} can be conveniently stated as:
\begin{equation}\label{equation:Pervaporation:PVAPeq:EQUI}
\begin{split}
f_i^V(T,P,\vec{y})=f_i^L(T,P,\vec{x})
\end{split}
\end{equation}
\sphinxAtStartPar
The left\sphinxhyphen{}hand term can be written as:
\begin{equation}\label{equation:Pervaporation:PVAPeq:vapour}
\begin{split}
f_i^V(T,P,\vec{y})=\phi(T,P,\vec{y})Py_i
\end{split}
\end{equation}
\sphinxAtStartPar
where \(\phi(T,P,\vec{y})\) is the fugacity coefficient, \(P\) the pressure and \(y_i\) the molar fraction in the vapour phase.

\sphinxAtStartPar
The right hand side term can instead be written as:
\begin{equation}\label{equation:Pervaporation:PVAPeq:liquid}
\begin{split}
f_i^L(T,P,\vec{x})=\gamma_i(T,P,\vec{x})x_i\phi(T,P^o(T))P^o(T)e^{\frac{v_i(P-P^o(T))}{RT}}
\end{split}
\end{equation}
\sphinxAtStartPar
where \(x_i\) is the molar fraction of component \(i\) in the liquid phase, \(\gamma_i(T,P,\vec{x})\) is the activity coefficient for component i in the liquid mixture, \(P^o(T)\) is the equilibrium vapour pressure of the pure component \(i\), \(\phi(T,P^o(T))\) is the fugacity coefficient for the pure component \(i\) at \(T\) and \(P^o(T)\), and the term \(e^{\frac{v_i(P-P^o(T))}{RT}}\) is the Poynting correction. For a discussion on the origin of the Poynting correction please refer to the notes on osmosis.

\sphinxAtStartPar
Introducing Eq.\eqref{equation:Pervaporation:PVAPeq:vapour} and Eq. \eqref{equation:Pervaporation:PVAPeq:liquid} in Eq.\eqref{equation:Pervaporation:PVAPeq:EQUI}, while considering negligible the Poynting correction, and ideal gas approximation applicable yields:
\begin{equation}\label{equation:Pervaporation:PVAPeq:final}
\begin{split}
Py_i=\gamma_i(T,P,\vec{x})P^o(T)x_i
\end{split}
\end{equation}
\sphinxAtStartPar
It should be noted that introducing the further hypothesis of ideal liquid mixture yields the Raoult law:
\begin{equation}\label{equation:Pervaporation:PVAPeq:Raoult}
\begin{split}
Py_i=P^o(T)x_i
\end{split}
\end{equation}

\subsection{2.3 Flux Equation in a pervaporation Unit}
\label{\detokenize{Pervaporation:flux-equation-in-a-pervaporation-unit}}
\sphinxAtStartPar
Considering now a binary mixture of components \(i\) and \(j\), the flux equation for species \(i\) in a pervaporation unit is:
\begin{equation}\label{equation:Pervaporation:PVAPeq:flux}
\begin{split}
J_i=\frac{P_i}{l}\left(\gamma_i(T,P,\vec{x})x^r_ip_i^o(T)-p_py^p_i\right)
\end{split}
\end{equation}
\sphinxAtStartPar
where \(P_i\) is the permeability of species \(i\) through the membrane, \(l\) the thickness of the membrane, \(\gamma_i(T,P,\vec{x})x^r_ip_i^o(T)\) is the partial pressure in the virtual vapour phase in equilibrium with the liquid on the retentate side, \(p_p\) the pressure on the permeate side, and \(y^p_i\) the molar fraction of component \(i\) in the real vapour phase on the permeate side.

\sphinxAtStartPar
The flux equation of component \(j\) can be expressed as a function of the molar fractions of component \(i\) taking advantage of the stoichiometric constraints \(\sum_{i=1}^{N}x^r_i=1\), \(\sum_{i=1}^{N}y^p_i=1\).
\begin{equation}\label{equation:Pervaporation:PVAPeq:flux2}
\begin{split}
J_j=\frac{P_j}{l}\left(\gamma_j(T,P,\vec{x})(1-x^r_i)p_j^o(T)-p_p(1-y^p_i)\right)
\end{split}
\end{equation}
\sphinxstepscope


\section{3. The Osmotic Pressure}
\label{\detokenize{Osmosis:the-osmotic-pressure}}\label{\detokenize{Osmosis::doc}}

\subsection{3.1 Osmosis}
\label{\detokenize{Osmosis:osmosis}}
\sphinxAtStartPar
Osmosis is a mass transfer process that takes place spontaneously when two compartments, containing fluid mixtures of different composition are separated by a semipermeable membrane, i.e. a membrane that restricts the transfer of one of the components across compartments.

\sphinxAtStartPar
For the sake of simplicity, let us consider for now an ideal semipermeable membrane that completely restricts the transfer of the solute, while instead allowing the permeation of a solvent.

\sphinxAtStartPar
Let us consider an initial state characterised by two compartments (1 and 2), at the same temperature T and pressure P.
Compartment 1 contains a binary mixture of component B (solute) in component A (solvent), e.g. NaCl in water. Compartment 2 instead contains a pure component A (e.g. pure water).

\sphinxAtStartPar
If the membrane is semipermeable and allows for the passage of the solvent alone the process of osmosis will lead to an equilibrium condition in which:
\begin{itemize}
\item {} 
\sphinxAtStartPar
Component A has transferred from compartment 2 to compartment 1.

\item {} 
\sphinxAtStartPar
An pressure difference between compartments is established, where : \(P_1>P_2\).

\item {} 
\sphinxAtStartPar
The pressure difference between compartments at equilibrium is the \sphinxstyleemphasis{osmotic pressure} \(\pi=P_1-P_2\).

\end{itemize}

\sphinxAtStartPar
In the next paragraphs we will see what is the thermodynamic origin of the osmotic pressure.


\subsection{3.2 Equilibrium Conditions}
\label{\detokenize{Osmosis:equilibrium-conditions}}
\sphinxAtStartPar
In order to understand the origin of the osmotic pressure, i.e. pressure difference induced by a composition difference across compartments separated by a semipermeable membrane, we begin by defining the the equilibrium state for the system described in the previous paragraph.
To this aim we can rely on the thermodynamic definition of equilibrium for component A (e.g. water), i.e. the only component that can transfer across the membrane.
Equilibrium is defined when \(A\) has the same chemical potential in both compartments
\begin{equation}\label{equation:Osmosis:OSMeq0}
\begin{split}
\mu_{A,1}=\mu_{A,2}
\end{split}
\end{equation}
\sphinxAtStartPar
The equilibrium condition is conveniently written introducing fugacity.
Fugacity and chemical potential are related by the expression \(\mu_i=\mu_0+RT\ln{f_i}\), where \(\mu_0\) is the chemical potential of an arbitrary reference state.

\sphinxAtStartPar
Equilibrium can conditions are thus equivalently formulated as:
\begin{equation}\label{equation:Osmosis:OSMeq:Equilibrium}
\begin{split}
f_{A,1}=f_{A,2}
\end{split}
\end{equation}
\sphinxAtStartPar
\sphinxstyleemphasis{In compartment 1} we have a \sphinxstyleemphasis{non\sphinxhyphen{}ideal} mixture of two components. The fugacity of component A in compartment 1 is thus written as:
\begin{equation}\label{equation:Osmosis:OSMeq:leftside}
\begin{split}
f_{A,1}=f_A(T,P_1,x_A)=\gamma_A(T,P_1,x_A)x_A f_A(T,P_1)
\end{split}
\end{equation}
\sphinxAtStartPar
where:
\begin{itemize}
\item {} 
\sphinxAtStartPar
\(\gamma_A(T,P_1,x_A)\) is the activity coefficient of specie A in the mixture

\item {} 
\sphinxAtStartPar
\(x_A\) is the molar fraction of component A

\item {} 
\sphinxAtStartPar
\(f_A(T,P_1)\) is the fugacity of the \sphinxstyleemphasis{pure component A} at temperature \(T\) and pressure \(P_1\)

\end{itemize}

\sphinxAtStartPar
\sphinxstyleemphasis{In compartment 2} we have instewad a pure component A liquid, at temperature \(T\) and pressure \(P_2\). The fugacity of component A in compartment 2 thus simply corresponds to the fugacity of the pure component A at \(T\), \(P_2\):
\begin{equation}\label{equation:Osmosis:OSMeq:rightside}
\begin{split}
f_{A,2}=f_A(T,P_2)
\end{split}
\end{equation}
\sphinxAtStartPar
\sphinxstylestrong{The Poynting correction} In order to define a common reference state between the two compartments we rewrite fugacity of the pure component A at pressure \(P_2\) (\(f_A(T,P_2)\)) as a function of the fugacity of the pure component A at pressure \(P_1\) and of the pressure difference \((P_2-P_1)\).

\sphinxAtStartPar
To do that we express the difference in chemical potential \(\Delta\mu_{P_1\rightarrow{P_2}}\) between pure liquid A at temperature \(T\) and pressure \(P_1\) and pure liquid A at temperature \(T\) and pressure \(P_2\):
\begin{equation}\label{equation:Osmosis:OSMeq1}
\begin{split}
\Delta\mu_{P_1\rightarrow{P_2}}=\int_{T,P_1}^{T,P_2}d\mu=\int_{T,P_1}^{T,P_2}vdP=v(P_2-P_1)
\end{split}
\end{equation}
\sphinxAtStartPar
where:
\begin{itemize}
\item {} 
\sphinxAtStartPar
the expression introduced for the definition of the \(\Delta\mu\) is the Gibbs Duhem equation (\(d{\mu}=-sdT+vdP\))

\item {} 
\sphinxAtStartPar
\(v\) is the molar volume of the system

\end{itemize}

\sphinxAtStartPar
Using the definition of fugacity introduced earlier one can thus write:
\begin{equation}\label{equation:Osmosis:OSMeq2}
\begin{split}
\Delta\mu_{P_1\rightarrow{P_2}}=RT\ln\frac{f_A(T,P_2)}{f_A(T,P_1)}=v(P_2-P_1)
\end{split}
\end{equation}
\sphinxAtStartPar
The fugacity of a pure A liquid at pressure \(P_2\) can thus be written as:
\begin{equation}\label{equation:Osmosis:OSMeq3}
\begin{split}
f_A(T,P_2)={f_A(T,P_1)}\left[\exp\left(\frac{v(P_2-P_1)}{RT}\right)\right]
\label{eq:Poynting}
\end{split}
\end{equation}
\sphinxAtStartPar
where the term \(\left[\exp\left(\frac{v(P_2-P_1)}{RT}\right)\right]\) is the Poynting correction.


\subsection{3.3 The Osmotic Pressure}
\label{\detokenize{Osmosis:id1}}
\sphinxAtStartPar
Given the definition of equilibrium conditions, the expressions of the fugacity of component A in the two compartments 1 and 2, and the Poynting correction given in the previous paragraph we can write:
\begin{equation}\label{equation:Osmosis:OSMeq}
\begin{split}
\gamma_A(T,P_1,x_A)x_A f_A(T,P_1)={f_A(T,P_1)}\left[\exp\left(\frac{v(P_2-P_1)}{RT}\right)\right]
\end{split}
\end{equation}
\sphinxAtStartPar
where \(f_A(T,P_1)\) is the fugacity of the pure component A at pressure \(P_1\), now appearing on both sides of the equation.

\sphinxAtStartPar
We can thus divide both sides by \(f_A(T,P_1)\) to obtain the expression:
\begin{equation}\label{equation:Osmosis:OSMeq4}
\begin{split}
\gamma_A(T,P_1,x_A)x_A =\left[\exp\left(\frac{v(P_2-P_1)}{RT}\right)\right]
\end{split}
\end{equation}
\sphinxAtStartPar
Rearranging this expression one obtains:
\begin{equation}\label{equation:Osmosis:OSMeq5}
\begin{split}
P_2-P_1=\frac{RT}{v}\ln\left({\gamma_A(T,P_1,x_A)x_A}\right)
\end{split}
\end{equation}
\sphinxAtStartPar
Recalling that the Osmotic pressure \(\pi=P_1-P_2\) we get:
\begin{equation}\label{equation:Osmosis:OSMeq6}
\begin{split}
\pi=-\frac{RT}{v}\ln\left({\gamma_A(T,P_1,x_A)x_A}\right)
\end{split}
\end{equation}
\sphinxAtStartPar
This expression \sphinxstylestrong{relates the composition of compartment 1}, containing a mixture of components A and B (e.g NaCl in water),  \sphinxstylestrong{to the osmotic pressure} generated at equilibrium conditions when such compartment is in contact with a \sphinxstylestrong{semipermeable} membrane that has on the opposite side a compartment filled by pure component A liquid.

\sphinxAtStartPar
Please note that in this expression \(x_A\) is the molar fraction of the solvent (e.g. water). Since we are considering a binary system and from an engineering standpoint what is typically measured is the solute concentration we can rewrite the expression as follows:
\begin{equation}\label{equation:Osmosis:OSMeq7}
\begin{split}
\pi=-\frac{RT}{v}\ln\left({\gamma_A(T,P_1,x_A)(1-x_B)}\right)
\end{split}
\end{equation}
\sphinxAtStartPar
Where now \(x_B\) is the molar fraction of component \(B\) (e.g. NaCl).
When the solution in compartment 1 is diluted (\(x_B\) small), \(\gamma_A(T,P_1,x_A)\rightarrow{1}\),  \(\ln{\left({(1-x_B)}\right)}\rightarrow{-x_B}\) thus:
\begin{equation}\label{equation:Osmosis:OSMeq8}
\begin{split}
\pi=-\frac{RT}{v}\ln\left({(1-x_B)}\right)={RT\frac{x_B}{v}}\simeq{RTc_B}
\end{split}
\end{equation}
\sphinxAtStartPar
Where \(c_B\simeq{x_B}v^{-1}\) is the molar concentration of component B, the solute.


\subsubsection{3.3.1 Dependence of the Osmotic pressure on the composition of the system}
\label{\detokenize{Osmosis:dependence-of-the-osmotic-pressure-on-the-composition-of-the-system}}
\sphinxAtStartPar
Note that the expression \(\pi={RTc_B}\) is only valid for small molar fractions of component B, the solute.

\begin{sphinxuseclass}{cell}\begin{sphinxVerbatimInput}

\begin{sphinxuseclass}{cell_input}
\begin{sphinxVerbatim}[commandchars=\\\{\}]
\PYG{k+kn}{import}\PYG{+w}{ }\PYG{n+nn}{numpy}\PYG{+w}{ }\PYG{k}{as}\PYG{+w}{ }\PYG{n+nn}{np}
\PYG{k+kn}{import}\PYG{+w}{ }\PYG{n+nn}{matplotlib}\PYG{n+nn}{.}\PYG{n+nn}{pyplot}\PYG{+w}{ }\PYG{k}{as}\PYG{+w}{ }\PYG{n+nn}{plt} 

\PYG{c+c1}{\PYGZsh{} Parameters: }
\PYG{c+c1}{\PYGZsh{} molar fraction in the feed}
\PYG{n}{N} \PYG{o}{=} \PYG{l+m+mi}{100} \PYG{c+c1}{\PYGZsh{}number of points}
\PYG{n}{x\PYGZus{}B} \PYG{o}{=} \PYG{n}{np}\PYG{o}{.}\PYG{n}{linspace}\PYG{p}{(}\PYG{l+m+mi}{0}\PYG{p}{,} \PYG{l+m+mf}{0.999}\PYG{p}{,} \PYG{n}{N}\PYG{p}{)}

\PYG{c+c1}{\PYGZsh{}Operating Equation}
\PYG{n}{pi} \PYG{o}{=} \PYG{o}{\PYGZhy{}} \PYG{n}{np}\PYG{o}{.}\PYG{n}{log}\PYG{p}{(}\PYG{l+m+mi}{1}\PYG{o}{\PYGZhy{}}\PYG{n}{x\PYGZus{}B}\PYG{p}{)}
\PYG{n}{pi\PYGZus{}simple} \PYG{o}{=} \PYG{n}{x\PYGZus{}B}

\PYG{c+c1}{\PYGZsh{}Plotting}
\PYG{n}{figure}\PYG{o}{=}\PYG{n}{plt}\PYG{o}{.}\PYG{n}{figure}\PYG{p}{(}\PYG{p}{)}
\PYG{n}{axes} \PYG{o}{=} \PYG{n}{figure}\PYG{o}{.}\PYG{n}{add\PYGZus{}axes}\PYG{p}{(}\PYG{p}{[}\PYG{l+m+mf}{0.1}\PYG{p}{,}\PYG{l+m+mf}{0.1}\PYG{p}{,}\PYG{l+m+mf}{0.8}\PYG{p}{,}\PYG{l+m+mf}{0.8}\PYG{p}{]}\PYG{p}{)}
\PYG{n}{plt}\PYG{o}{.}\PYG{n}{xticks}\PYG{p}{(}\PYG{n}{fontsize}\PYG{o}{=}\PYG{l+m+mi}{14}\PYG{p}{)}
\PYG{n}{plt}\PYG{o}{.}\PYG{n}{yticks}\PYG{p}{(}\PYG{n}{fontsize}\PYG{o}{=}\PYG{l+m+mi}{14}\PYG{p}{)}
\PYG{n}{axes}\PYG{o}{.}\PYG{n}{plot}\PYG{p}{(}\PYG{n}{x\PYGZus{}B}\PYG{p}{,}\PYG{n}{pi}\PYG{p}{,} \PYG{n}{marker}\PYG{o}{=}\PYG{l+s+s1}{\PYGZsq{}}\PYG{l+s+s1}{ }\PYG{l+s+s1}{\PYGZsq{}} \PYG{p}{,} \PYG{n}{color}\PYG{o}{=}\PYG{l+s+s1}{\PYGZsq{}}\PYG{l+s+s1}{r}\PYG{l+s+s1}{\PYGZsq{}}\PYG{p}{)}
\PYG{n}{axes}\PYG{o}{.}\PYG{n}{plot}\PYG{p}{(}\PYG{n}{x\PYGZus{}B}\PYG{p}{,}\PYG{n}{pi\PYGZus{}simple}\PYG{p}{,} \PYG{n}{marker}\PYG{o}{=}\PYG{l+s+s1}{\PYGZsq{}}\PYG{l+s+s1}{ }\PYG{l+s+s1}{\PYGZsq{}} \PYG{p}{,} \PYG{n}{color}\PYG{o}{=}\PYG{l+s+s1}{\PYGZsq{}}\PYG{l+s+s1}{b}\PYG{l+s+s1}{\PYGZsq{}}\PYG{p}{)}
\PYG{n}{axes}\PYG{o}{.}\PYG{n}{set\PYGZus{}xlim}\PYG{p}{(}\PYG{p}{[}\PYG{l+m+mi}{0}\PYG{p}{,}\PYG{l+m+mi}{1}\PYG{p}{]}\PYG{p}{)}\PYG{p}{;}
\PYG{n}{axes}\PYG{o}{.}\PYG{n}{set\PYGZus{}ylim}\PYG{p}{(}\PYG{p}{[}\PYG{l+m+mi}{0}\PYG{p}{,}\PYG{l+m+mi}{1}\PYG{p}{]}\PYG{p}{)}\PYG{p}{;}
\PYG{n}{plt}\PYG{o}{.}\PYG{n}{title}\PYG{p}{(}\PYG{l+s+s1}{\PYGZsq{}}\PYG{l+s+s1}{Osmotic Pressure vs Solute Concentration}\PYG{l+s+s1}{\PYGZsq{}}\PYG{p}{,} \PYG{n}{fontsize}\PYG{o}{=}\PYG{l+m+mi}{18}\PYG{p}{)}\PYG{p}{;}
\PYG{n}{axes}\PYG{o}{.}\PYG{n}{set\PYGZus{}xlabel}\PYG{p}{(}\PYG{l+s+s1}{\PYGZsq{}}\PYG{l+s+s1}{Molar Fraction of component B}\PYG{l+s+s1}{\PYGZsq{}}\PYG{p}{,} \PYG{n}{fontsize}\PYG{o}{=}\PYG{l+m+mi}{14}\PYG{p}{)}\PYG{p}{;}
\PYG{n}{axes}\PYG{o}{.}\PYG{n}{set\PYGZus{}ylabel}\PYG{p}{(}\PYG{l+s+s1}{\PYGZsq{}}\PYG{l+s+s1}{Dimensionless Osmotic Pressure}\PYG{l+s+s1}{\PYGZsq{}}\PYG{p}{,}\PYG{n}{fontsize}\PYG{o}{=}\PYG{l+m+mi}{14}\PYG{p}{)}\PYG{p}{;}
\end{sphinxVerbatim}

\end{sphinxuseclass}\end{sphinxVerbatimInput}
\begin{sphinxVerbatimOutput}

\begin{sphinxuseclass}{cell_output}
\noindent\sphinxincludegraphics{{ffdd361f32380a99b954fd669c524b492fa392d7610a5f7f32e42a04eed6de0a}.png}

\end{sphinxuseclass}\end{sphinxVerbatimOutput}

\end{sphinxuseclass}
\sphinxAtStartPar
Note that the dimensionless osmotic pressure plotted abotve is obtained as \(\pi{v}/RT\).


\subsection{3.4 Flux of the solvent during Osmosis}
\label{\detokenize{Osmosis:flux-of-the-solvent-during-osmosis}}
\sphinxAtStartPar
In osmosis the flux of the solvent (component A) through a membrane is driven by the difference between the \(\Delta{P}=P_1-P_2\) across compartments and the osmotic pressure \(\pi\):
\begin{equation}\label{equation:Osmosis:OSMeq9}
\begin{split}
J_A=p_A\left( \Delta{P} - \pi \right)
\end{split}
\end{equation}
\sphinxAtStartPar
where \(p_A\) is the permeability of the membrane with respect to component A.

\sphinxAtStartPar
At equilibrium, when \(\Delta{P}=\pi\) the net flux across compartments is null. Without adding any external pressure the spontaneous direction of the flux is from compartment 2 to compartment 1 (\(J<0\)).

\sphinxAtStartPar
If we add a pressure, such that \(\Delta{P}>\pi\) we induce the opposite flux, from compartment 1 to compartment 2. This process is named reverse osmosis.

\begin{sphinxuseclass}{cell}\begin{sphinxVerbatimInput}

\begin{sphinxuseclass}{cell_input}
\begin{sphinxVerbatim}[commandchars=\\\{\}]
\PYG{k+kn}{import}\PYG{+w}{ }\PYG{n+nn}{numpy}\PYG{+w}{ }\PYG{k}{as}\PYG{+w}{ }\PYG{n+nn}{np}
\PYG{k+kn}{import}\PYG{+w}{ }\PYG{n+nn}{pandas}\PYG{+w}{ }\PYG{k}{as}\PYG{+w}{ }\PYG{n+nn}{pd}
\PYG{k+kn}{from}\PYG{+w}{ }\PYG{n+nn}{matplotlib}\PYG{+w}{ }\PYG{k+kn}{import} \PYG{n}{cm}
\PYG{k+kn}{import}\PYG{+w}{ }\PYG{n+nn}{matplotlib}\PYG{n+nn}{.}\PYG{n+nn}{pyplot}\PYG{+w}{ }\PYG{k}{as}\PYG{+w}{ }\PYG{n+nn}{plt} 
\PYG{k+kn}{from}\PYG{+w}{ }\PYG{n+nn}{mpl\PYGZus{}toolkits}\PYG{n+nn}{.}\PYG{n+nn}{mplot3d}\PYG{+w}{ }\PYG{k+kn}{import} \PYG{n}{Axes3D} 

\PYG{c+c1}{\PYGZsh{} Parameters: }
\PYG{c+c1}{\PYGZsh{} molar fraction in the feed}
\PYG{n}{N} \PYG{o}{=} \PYG{l+m+mi}{100} \PYG{c+c1}{\PYGZsh{}number of points}
\PYG{n}{DeltaP} \PYG{o}{=} \PYG{n}{np}\PYG{o}{.}\PYG{n}{linspace}\PYG{p}{(}\PYG{l+m+mi}{0}\PYG{p}{,} \PYG{l+m+mi}{1}\PYG{p}{,} \PYG{n}{N}\PYG{p}{)}
\PYG{n}{x\PYGZus{}B} \PYG{o}{=} \PYG{n}{np}\PYG{o}{.}\PYG{n}{linspace}\PYG{p}{(}\PYG{l+m+mi}{0}\PYG{p}{,} \PYG{l+m+mf}{0.8}\PYG{p}{,} \PYG{n}{N}\PYG{p}{)}
\PYG{n}{P\PYGZus{}A}\PYG{o}{=} \PYG{l+m+mi}{1} 
\PYG{n}{DeltaP}\PYG{p}{,} \PYG{n}{x\PYGZus{}B} \PYG{o}{=} \PYG{n}{np}\PYG{o}{.}\PYG{n}{meshgrid}\PYG{p}{(}\PYG{n}{DeltaP}\PYG{p}{,} \PYG{n}{x\PYGZus{}B}\PYG{p}{)} 

\PYG{n}{pi} \PYG{o}{=} \PYG{o}{\PYGZhy{}} \PYG{n}{np}\PYG{o}{.}\PYG{n}{log}\PYG{p}{(}\PYG{l+m+mi}{1}\PYG{o}{\PYGZhy{}}\PYG{n}{x\PYGZus{}B}\PYG{p}{)}

\PYG{c+c1}{\PYGZsh{}Flux Equation}
\PYG{n}{J\PYGZus{}A} \PYG{o}{=} \PYG{n}{P\PYGZus{}A} \PYG{o}{*} \PYG{p}{(}\PYG{n}{DeltaP} \PYG{o}{\PYGZhy{}} \PYG{n}{pi}\PYG{p}{)}


\PYG{c+c1}{\PYGZsh{}Plotting}

\PYG{n}{figure}\PYG{o}{=}\PYG{n}{plt}\PYG{o}{.}\PYG{n}{figure}\PYG{p}{(}\PYG{n}{figsize}\PYG{o}{=}\PYG{p}{(}\PYG{l+m+mi}{10}\PYG{p}{,} \PYG{l+m+mi}{8}\PYG{p}{)}\PYG{p}{)}
\PYG{c+c1}{\PYGZsh{}axes = figure.add\PYGZus{}axes([0.1,0.1,0.1,1.2,1.2,1.2])}
\PYG{n}{axes} \PYG{o}{=} \PYG{n}{figure}\PYG{o}{.}\PYG{n}{add\PYGZus{}subplot}\PYG{p}{(}\PYG{n}{projection}\PYG{o}{=}\PYG{l+s+s1}{\PYGZsq{}}\PYG{l+s+s1}{3d}\PYG{l+s+s1}{\PYGZsq{}}\PYG{p}{)}


\PYG{n}{plt}\PYG{o}{.}\PYG{n}{xticks}\PYG{p}{(}\PYG{n}{fontsize}\PYG{o}{=}\PYG{l+m+mi}{14}\PYG{p}{)}
\PYG{n}{plt}\PYG{o}{.}\PYG{n}{yticks}\PYG{p}{(}\PYG{n}{fontsize}\PYG{o}{=}\PYG{l+m+mi}{14}\PYG{p}{)}

\PYG{n}{surf}\PYG{o}{=}\PYG{n}{axes}\PYG{o}{.}\PYG{n}{plot\PYGZus{}surface}\PYG{p}{(}\PYG{n}{x\PYGZus{}B}\PYG{p}{,}\PYG{n}{DeltaP}\PYG{p}{,}\PYG{n}{J\PYGZus{}A}\PYG{p}{,}\PYG{n}{cmap}\PYG{o}{=}\PYG{n}{cm}\PYG{o}{.}\PYG{n}{coolwarm}\PYG{p}{,}
                       \PYG{n}{linewidth}\PYG{o}{=}\PYG{l+m+mi}{0}\PYG{p}{,} \PYG{n}{antialiased}\PYG{o}{=}\PYG{k+kc}{False}\PYG{p}{,}\PYG{n}{alpha}\PYG{o}{=}\PYG{l+m+mf}{0.5}\PYG{p}{)}

\PYG{n}{contour}\PYG{o}{=}\PYG{n}{axes}\PYG{o}{.}\PYG{n}{contour}\PYG{p}{(}\PYG{n}{x\PYGZus{}B}\PYG{p}{,}\PYG{n}{DeltaP}\PYG{p}{,}\PYG{n}{J\PYGZus{}A}\PYG{p}{,}\PYG{p}{[}\PYG{l+m+mi}{0}\PYG{p}{]}\PYG{p}{)}

\PYG{n}{plt}\PYG{o}{.}\PYG{n}{title}\PYG{p}{(}\PYG{l+s+s1}{\PYGZsq{}}\PYG{l+s+s1}{Osmotic Flux vs \PYGZdl{}}\PYG{l+s+s1}{\PYGZbs{}}\PYG{l+s+s1}{Delta}\PYG{l+s+si}{\PYGZob{}P\PYGZcb{}}\PYG{l+s+s1}{\PYGZdl{} / \PYGZdl{}x\PYGZus{}B\PYGZdl{}}\PYG{l+s+s1}{\PYGZsq{}}\PYG{p}{,} \PYG{n}{fontsize}\PYG{o}{=}\PYG{l+m+mi}{18}\PYG{p}{)}\PYG{p}{;}
\PYG{n}{axes}\PYG{o}{.}\PYG{n}{set\PYGZus{}xlabel}\PYG{p}{(}\PYG{l+s+s1}{\PYGZsq{}}\PYG{l+s+s1}{\PYGZdl{}x\PYGZus{}B\PYGZdl{}}\PYG{l+s+s1}{\PYGZsq{}}\PYG{p}{,} \PYG{n}{fontsize}\PYG{o}{=}\PYG{l+m+mi}{14}\PYG{p}{)}\PYG{p}{;}
\PYG{n}{axes}\PYG{o}{.}\PYG{n}{set\PYGZus{}ylabel}\PYG{p}{(}\PYG{l+s+s1}{\PYGZsq{}}\PYG{l+s+s1}{\PYGZdl{}}\PYG{l+s+s1}{\PYGZbs{}}\PYG{l+s+s1}{Delta}\PYG{l+s+si}{\PYGZob{}P\PYGZcb{}}\PYG{l+s+s1}{\PYGZdl{}}\PYG{l+s+s1}{\PYGZsq{}}\PYG{p}{,}\PYG{n}{fontsize}\PYG{o}{=}\PYG{l+m+mi}{14}\PYG{p}{)}\PYG{p}{;}
\PYG{n}{figure}\PYG{o}{.}\PYG{n}{colorbar}\PYG{p}{(}\PYG{n}{surf}\PYG{p}{,} \PYG{n}{shrink}\PYG{o}{=}\PYG{l+m+mf}{0.5}\PYG{p}{,} \PYG{n}{aspect}\PYG{o}{=}\PYG{l+m+mi}{5}\PYG{p}{)}\PYG{p}{;}
\end{sphinxVerbatim}

\end{sphinxuseclass}\end{sphinxVerbatimInput}
\begin{sphinxVerbatimOutput}

\begin{sphinxuseclass}{cell_output}
\noindent\sphinxincludegraphics{{7781d59ff5b57abc2b6493ccad995a7670ff98e119e135d32ba1a425d46e5d57}.png}

\end{sphinxuseclass}\end{sphinxVerbatimOutput}

\end{sphinxuseclass}
\sphinxstepscope


\section{4. Membrane Processes Layouts}
\label{\detokenize{Process_layouts:membrane-processes-layouts}}\label{\detokenize{Process_layouts::doc}}
\sphinxAtStartPar
In this set of notes we will use an example ultrafiltration process to introduce material balances in different process configurations. We will consider a Batch, semi\sphinxhyphen{}Batch (feed and bleed) and continuous membrane processes to develop general ideas around material balances in membrane units. Every example will be introduced by a slight variant of a common problem statement.


\subsection{4.1 Batch}
\label{\detokenize{Process_layouts:batch}}
\sphinxAtStartPar
\sphinxstylestrong{Problem Statement}

\sphinxAtStartPar
500 l of fruit juice are concentrated from an initial solid\sphinxhyphen{}residue content of 0.05 kg/l to a solid\sphinxhyphen{}residue content of 0.2 kg/l through a batch microfiltration process. The total area of the membrane is of 20 \(\mathrm{m^2}\) and the flux of pure water through the membrane is captured by the empirical expression: \(J=BC^{-1} [m/h]\) with \(B=0.1\) where C is the solid residue concentration in kg/l. Compute the process time necessary to reach the target concentration.

\begin{figure}[htbp]
\centering
\capstart

\noindent\sphinxincludegraphics[height=300\sphinxpxdimen]{{batch}.png}
\caption{Schematics of a Batch membrane process.}\label{\detokenize{Process_layouts:scheme}}\end{figure}

\sphinxAtStartPar
\sphinxstylestrong{Solution trace}

\sphinxAtStartPar
The global differential material balance (volume basis) reads:
\begin{equation}\label{equation:Process_layouts:PLYeq1}
\begin{split}
\frac{dV}{dt}=-J\,A
\end{split}
\end{equation}
\sphinxAtStartPar
where \(V\) is the volume in the concentrate loop, \(J\) the permeate flux per unit area and \(A\) the total membrane area.
The permeate flux is defined by the expression:
\begin{equation}\label{equation:Process_layouts:PLYeq2}
\begin{split}
J=BC^{-1}
\end{split}
\end{equation}
\sphinxAtStartPar
with \(B=0.1\).

\sphinxAtStartPar
Since the solid residue never leaves the retentate loop,  the differential material balance on the solid residue reads:
\begin{equation}\label{equation:Process_layouts:PLYeq3}
\begin{split}
\frac{dm}{dt}=0 
\end{split}
\end{equation}
\sphinxAtStartPar
Which means that the total mass of solid residue is constant and equal to its initial mass:
\begin{equation}\label{equation:Process_layouts:PLYeq4}
\begin{split}
m=m_0=V_0C_0 
\end{split}
\end{equation}
\sphinxAtStartPar
The global differential material balance can thus be rewritten as:
\begin{equation}\label{equation:Process_layouts:PLYeq5}
\begin{split}
\frac{dV}{dt}=-J\,A=-\frac{B}{V_0C_0 }AV
\end{split}
\end{equation}
\sphinxAtStartPar
In this simple case it can be integrated analytically to compute the volume as a function of time:
\begin{equation}\label{equation:Process_layouts:PLYeq6}
\begin{split}
\int_{V_0}^{V^{\prime}}\frac{dV}{V}=\int_0^{t^{\prime}}-\frac{B}{V_0C_0 }Adt
\end{split}
\end{equation}\begin{equation}\label{equation:Process_layouts:PLYeq7}
\begin{split}
\int_{V_0}^{V^{\prime}}\frac{dV}{V}=\left[\ln{V}\right]^{V^{\prime}}_{V_0}=\ln\left(\frac{V^\prime}{V_0}\right)
\end{split}
\end{equation}\begin{equation}\label{equation:Process_layouts:PLYeq8}
\begin{split}
\int_0^{t^{\prime}}-\frac{B}{V_0C_0 }Adt=-\frac{B}{V_0C_0 }At^\prime
\end{split}
\end{equation}
\sphinxAtStartPar
Putting together both sides of the equation and substituting \(V^\prime\) with \(V\) for the sake of simplicity in the notation we get:
\begin{equation}\label{equation:Process_layouts:PLYeq9}
\begin{split}
V(t)=V_0e^{-\frac{B}{V_0C_0 }At}
\end{split}
\end{equation}
\sphinxAtStartPar
which describes the time dependence of the volume in the batch membrane separator.

\sphinxAtStartPar
The concentration can thus be calculated as:
\begin{equation}\label{equation:Process_layouts:PLYeq10}
\begin{split}
C(t)=\frac{m}{V(t)}=\frac{C_0V_0}{V_0}e^{\frac{B}{V_0C_0 }At}={C_0}e^{\frac{B}{V_0C_0 }At}
\end{split}
\end{equation}
\sphinxAtStartPar
The process time necessary to obtain a solid residue concentration of 0.2 \([kg/l]\) is of 17.3 \([h]\)

\sphinxAtStartPar
\sphinxstylestrong{Numerical solution}

\begin{sphinxuseclass}{cell}\begin{sphinxVerbatimInput}

\begin{sphinxuseclass}{cell_input}
\begin{sphinxVerbatim}[commandchars=\\\{\}]
\PYG{k+kn}{import}\PYG{+w}{ }\PYG{n+nn}{numpy}\PYG{+w}{ }\PYG{k}{as}\PYG{+w}{ }\PYG{n+nn}{np}
\PYG{k+kn}{import}\PYG{+w}{ }\PYG{n+nn}{matplotlib}\PYG{n+nn}{.}\PYG{n+nn}{pyplot}\PYG{+w}{ }\PYG{k}{as}\PYG{+w}{ }\PYG{n+nn}{plt} 
\PYG{k+kn}{from}\PYG{+w}{ }\PYG{n+nn}{scipy}\PYG{n+nn}{.}\PYG{n+nn}{optimize}\PYG{+w}{ }\PYG{k+kn}{import} \PYG{n}{fsolve}

\PYG{c+c1}{\PYGZsh{} Parameters: }
\PYG{n}{N} \PYG{o}{=} \PYG{l+m+mi}{500} \PYG{c+c1}{\PYGZsh{}number of points}
\PYG{n}{time} \PYG{o}{=} \PYG{n}{np}\PYG{o}{.}\PYG{n}{linspace}\PYG{p}{(}\PYG{l+m+mi}{0}\PYG{p}{,} \PYG{l+m+mi}{30}\PYG{p}{,} \PYG{n}{N}\PYG{p}{)}

\PYG{c+c1}{\PYGZsh{}Every length in m}

\PYG{n}{A}\PYG{o}{=}\PYG{l+m+mi}{20}\PYG{p}{;}  \PYG{c+c1}{\PYGZsh{}m\PYGZca{}2}
\PYG{n}{B}\PYG{o}{=}\PYG{l+m+mf}{0.1}\PYG{p}{;}
\PYG{n}{C0}\PYG{o}{=}\PYG{l+m+mf}{0.05}\PYG{o}{*}\PYG{l+m+mf}{1E3}\PYG{p}{;} \PYG{c+c1}{\PYGZsh{}kg /l * dm\PYGZca{}3/m\PYGZca{}3}
\PYG{n}{V0}\PYG{o}{=}\PYG{l+m+mi}{500}\PYG{o}{*}\PYG{l+m+mf}{1E\PYGZhy{}3}\PYG{p}{;} \PYG{c+c1}{\PYGZsh{}l * m\PYGZca{}3/dm\PYGZca{}3}

\PYG{n}{C\PYGZus{}specific} \PYG{o}{=} \PYG{l+m+mf}{0.2}\PYG{o}{*}\PYG{l+m+mf}{1E3}

\PYG{c+c1}{\PYGZsh{}Operating Equation}
\PYG{n}{C} \PYG{o}{=} \PYG{n}{C0}\PYG{o}{*}\PYG{n}{np}\PYG{o}{.}\PYG{n}{exp}\PYG{p}{(}\PYG{n}{B}\PYG{o}{*}\PYG{n}{A}\PYG{o}{/}\PYG{n}{V0}\PYG{o}{/}\PYG{n}{C0}\PYG{o}{*}\PYG{n}{time}\PYG{p}{)} 
\PYG{n}{V} \PYG{o}{=} \PYG{n}{V0}\PYG{o}{*}\PYG{n}{np}\PYG{o}{.}\PYG{n}{exp}\PYG{p}{(}\PYG{o}{\PYGZhy{}}\PYG{n}{B}\PYG{o}{*}\PYG{n}{A}\PYG{o}{/}\PYG{n}{V0}\PYG{o}{/}\PYG{n}{C0}\PYG{o}{*}\PYG{n}{time}\PYG{p}{)} 

\PYG{k}{def}\PYG{+w}{ }\PYG{n+nf}{equation}\PYG{p}{(}\PYG{n}{proc\PYGZus{}time}\PYG{p}{)}\PYG{p}{:}
    \PYG{n}{eq1} \PYG{o}{=} \PYG{n}{C0}\PYG{o}{*}\PYG{n}{np}\PYG{o}{.}\PYG{n}{exp}\PYG{p}{(}\PYG{n}{B}\PYG{o}{*}\PYG{n}{A}\PYG{o}{/}\PYG{n}{V0}\PYG{o}{/}\PYG{n}{C0}\PYG{o}{*}\PYG{n}{proc\PYGZus{}time}\PYG{p}{)} \PYG{o}{\PYGZhy{}} \PYG{n}{C\PYGZus{}specific}
    \PYG{k}{return} \PYG{n}{eq1}

\PYG{n}{process\PYGZus{}time} \PYG{o}{=} \PYG{n}{fsolve}\PYG{p}{(}\PYG{n}{equation}\PYG{p}{,}\PYG{l+m+mi}{1}\PYG{p}{)}

\PYG{c+c1}{\PYGZsh{}Plotting}
\PYG{n}{figure}\PYG{o}{=}\PYG{n}{plt}\PYG{o}{.}\PYG{n}{figure}\PYG{p}{(}\PYG{p}{)}
\PYG{n}{axes} \PYG{o}{=} \PYG{n}{figure}\PYG{o}{.}\PYG{n}{add\PYGZus{}axes}\PYG{p}{(}\PYG{p}{[}\PYG{l+m+mf}{0.1}\PYG{p}{,}\PYG{l+m+mf}{0.1}\PYG{p}{,}\PYG{l+m+mf}{0.8}\PYG{p}{,}\PYG{l+m+mf}{0.8}\PYG{p}{]}\PYG{p}{)}
\PYG{n}{plt}\PYG{o}{.}\PYG{n}{xticks}\PYG{p}{(}\PYG{n}{fontsize}\PYG{o}{=}\PYG{l+m+mi}{14}\PYG{p}{)}
\PYG{n}{plt}\PYG{o}{.}\PYG{n}{yticks}\PYG{p}{(}\PYG{n}{fontsize}\PYG{o}{=}\PYG{l+m+mi}{14}\PYG{p}{)}
\PYG{n}{axes}\PYG{o}{.}\PYG{n}{plot}\PYG{p}{(}\PYG{n}{time}\PYG{p}{,}\PYG{n}{V}\PYG{p}{,} \PYG{n}{marker}\PYG{o}{=}\PYG{l+s+s1}{\PYGZsq{}}\PYG{l+s+s1}{ }\PYG{l+s+s1}{\PYGZsq{}} \PYG{p}{,} \PYG{n}{color}\PYG{o}{=}\PYG{l+s+s1}{\PYGZsq{}}\PYG{l+s+s1}{b}\PYG{l+s+s1}{\PYGZsq{}}\PYG{p}{)}

\PYG{n}{plt}\PYG{o}{.}\PYG{n}{title}\PYG{p}{(}\PYG{l+s+s1}{\PYGZsq{}}\PYG{l+s+s1}{Suspension Volume}\PYG{l+s+s1}{\PYGZsq{}}\PYG{p}{,} \PYG{n}{fontsize}\PYG{o}{=}\PYG{l+m+mi}{18}\PYG{p}{)}\PYG{p}{;}
\PYG{n}{axes}\PYG{o}{.}\PYG{n}{set\PYGZus{}xlabel}\PYG{p}{(}\PYG{l+s+s1}{\PYGZsq{}}\PYG{l+s+s1}{time [h]}\PYG{l+s+s1}{\PYGZsq{}}\PYG{p}{,} \PYG{n}{fontsize}\PYG{o}{=}\PYG{l+m+mi}{14}\PYG{p}{)}\PYG{p}{;}
\PYG{n}{axes}\PYG{o}{.}\PYG{n}{set\PYGZus{}ylabel}\PYG{p}{(}\PYG{l+s+s1}{\PYGZsq{}}\PYG{l+s+s1}{V [m\PYGZdl{}\PYGZca{}3\PYGZdl{}]}\PYG{l+s+s1}{\PYGZsq{}}\PYG{p}{,}\PYG{n}{fontsize}\PYG{o}{=}\PYG{l+m+mi}{14}\PYG{p}{)}\PYG{p}{;}

\PYG{n}{figure}\PYG{o}{=}\PYG{n}{plt}\PYG{o}{.}\PYG{n}{figure}\PYG{p}{(}\PYG{p}{)}
\PYG{n}{axes} \PYG{o}{=} \PYG{n}{figure}\PYG{o}{.}\PYG{n}{add\PYGZus{}axes}\PYG{p}{(}\PYG{p}{[}\PYG{l+m+mf}{0.1}\PYG{p}{,}\PYG{l+m+mf}{0.1}\PYG{p}{,}\PYG{l+m+mf}{0.8}\PYG{p}{,}\PYG{l+m+mf}{0.8}\PYG{p}{]}\PYG{p}{)}
\PYG{n}{plt}\PYG{o}{.}\PYG{n}{xticks}\PYG{p}{(}\PYG{n}{fontsize}\PYG{o}{=}\PYG{l+m+mi}{14}\PYG{p}{)}
\PYG{n}{plt}\PYG{o}{.}\PYG{n}{yticks}\PYG{p}{(}\PYG{n}{fontsize}\PYG{o}{=}\PYG{l+m+mi}{14}\PYG{p}{)}
\PYG{n}{axes}\PYG{o}{.}\PYG{n}{plot}\PYG{p}{(}\PYG{n}{time}\PYG{p}{,}\PYG{n}{C}\PYG{p}{,} \PYG{n}{marker}\PYG{o}{=}\PYG{l+s+s1}{\PYGZsq{}}\PYG{l+s+s1}{ }\PYG{l+s+s1}{\PYGZsq{}} \PYG{p}{,} \PYG{n}{color}\PYG{o}{=}\PYG{l+s+s1}{\PYGZsq{}}\PYG{l+s+s1}{r}\PYG{l+s+s1}{\PYGZsq{}}\PYG{p}{)}
\PYG{n}{axes}\PYG{o}{.}\PYG{n}{plot}\PYG{p}{(}\PYG{n}{np}\PYG{o}{.}\PYG{n}{array}\PYG{p}{(}\PYG{p}{[}\PYG{n}{process\PYGZus{}time}\PYG{p}{[}\PYG{l+m+mi}{0}\PYG{p}{]}\PYG{p}{,}\PYG{n}{process\PYGZus{}time}\PYG{p}{[}\PYG{l+m+mi}{0}\PYG{p}{]}\PYG{p}{]}\PYG{p}{)}\PYG{p}{,}\PYG{n}{np}\PYG{o}{.}\PYG{n}{array}\PYG{p}{(}\PYG{p}{[}\PYG{l+m+mi}{0}\PYG{p}{,} \PYG{n}{C\PYGZus{}specific}\PYG{p}{]}\PYG{p}{)}\PYG{p}{,} \PYG{n}{marker}\PYG{o}{=}\PYG{l+s+s1}{\PYGZsq{}}\PYG{l+s+s1}{ }\PYG{l+s+s1}{\PYGZsq{}} \PYG{p}{,} \PYG{n}{color}\PYG{o}{=}\PYG{l+s+s1}{\PYGZsq{}}\PYG{l+s+s1}{lime}\PYG{l+s+s1}{\PYGZsq{}}\PYG{p}{,} \PYG{n}{markersize}\PYG{o}{=}\PYG{l+m+mi}{3}\PYG{p}{)}
\PYG{n}{axes}\PYG{o}{.}\PYG{n}{plot}\PYG{p}{(}\PYG{n}{np}\PYG{o}{.}\PYG{n}{array}\PYG{p}{(}\PYG{p}{[}\PYG{l+m+mi}{0}\PYG{p}{,}\PYG{n}{process\PYGZus{}time}\PYG{p}{[}\PYG{l+m+mi}{0}\PYG{p}{]}\PYG{p}{]}\PYG{p}{)}\PYG{p}{,}\PYG{n}{np}\PYG{o}{.}\PYG{n}{array}\PYG{p}{(}\PYG{p}{[}\PYG{n}{C\PYGZus{}specific}\PYG{p}{,} \PYG{n}{C\PYGZus{}specific}\PYG{p}{]}\PYG{p}{)}\PYG{p}{,} \PYG{n}{marker}\PYG{o}{=}\PYG{l+s+s1}{\PYGZsq{}}\PYG{l+s+s1}{ }\PYG{l+s+s1}{\PYGZsq{}} \PYG{p}{,} \PYG{n}{color}\PYG{o}{=}\PYG{l+s+s1}{\PYGZsq{}}\PYG{l+s+s1}{lime}\PYG{l+s+s1}{\PYGZsq{}}\PYG{p}{,} \PYG{n}{markersize}\PYG{o}{=}\PYG{l+m+mi}{3}\PYG{p}{)}


\PYG{n}{axes}\PYG{o}{.}\PYG{n}{plot}\PYG{p}{(}\PYG{p}{[}\PYG{n}{process\PYGZus{}time}\PYG{p}{,}\PYG{n}{process\PYGZus{}time}\PYG{p}{]}\PYG{p}{,}\PYG{p}{[}\PYG{l+m+mi}{0}\PYG{p}{,} \PYG{n}{C\PYGZus{}specific}\PYG{p}{]}\PYG{p}{,} \PYG{n}{marker}\PYG{o}{=}\PYG{l+s+s1}{\PYGZsq{}}\PYG{l+s+s1}{ }\PYG{l+s+s1}{\PYGZsq{}} \PYG{p}{,} \PYG{n}{color}\PYG{o}{=}\PYG{l+s+s1}{\PYGZsq{}}\PYG{l+s+s1}{lime}\PYG{l+s+s1}{\PYGZsq{}}\PYG{p}{,} \PYG{n}{markersize}\PYG{o}{=}\PYG{l+m+mi}{3}\PYG{p}{)}


\PYG{n}{axes}\PYG{o}{.}\PYG{n}{plot}\PYG{p}{(}\PYG{n}{process\PYGZus{}time}\PYG{p}{,}\PYG{n}{C\PYGZus{}specific}\PYG{p}{,} \PYG{n}{marker}\PYG{o}{=}\PYG{l+s+s1}{\PYGZsq{}}\PYG{l+s+s1}{o}\PYG{l+s+s1}{\PYGZsq{}} \PYG{p}{,} \PYG{n}{color}\PYG{o}{=}\PYG{l+s+s1}{\PYGZsq{}}\PYG{l+s+s1}{lime}\PYG{l+s+s1}{\PYGZsq{}}\PYG{p}{,} \PYG{n}{markersize}\PYG{o}{=}\PYG{l+m+mi}{10}\PYG{p}{)}

\PYG{n}{plt}\PYG{o}{.}\PYG{n}{title}\PYG{p}{(}\PYG{l+s+s1}{\PYGZsq{}}\PYG{l+s+s1}{Solute concentration}\PYG{l+s+s1}{\PYGZsq{}}\PYG{p}{,} \PYG{n}{fontsize}\PYG{o}{=}\PYG{l+m+mi}{18}\PYG{p}{)}\PYG{p}{;}
\PYG{n}{axes}\PYG{o}{.}\PYG{n}{set\PYGZus{}xlabel}\PYG{p}{(}\PYG{l+s+s1}{\PYGZsq{}}\PYG{l+s+s1}{time [h]}\PYG{l+s+s1}{\PYGZsq{}}\PYG{p}{,} \PYG{n}{fontsize}\PYG{o}{=}\PYG{l+m+mi}{14}\PYG{p}{)}\PYG{p}{;}
\PYG{n}{axes}\PYG{o}{.}\PYG{n}{set\PYGZus{}ylabel}\PYG{p}{(}\PYG{l+s+s1}{\PYGZsq{}}\PYG{l+s+s1}{concentration [kg/m\PYGZdl{}\PYGZca{}3\PYGZdl{}]}\PYG{l+s+s1}{\PYGZsq{}}\PYG{p}{,}\PYG{n}{fontsize}\PYG{o}{=}\PYG{l+m+mi}{14}\PYG{p}{)}\PYG{p}{;}

\PYG{n+nb}{print}\PYG{p}{(}\PYG{l+s+s2}{\PYGZdq{}}\PYG{l+s+s2}{The process time necessary to obtain a solid residue concentration of}\PYG{l+s+s2}{\PYGZdq{}}\PYG{p}{,} \PYG{n}{C\PYGZus{}specific}\PYG{p}{,} \PYG{l+s+s2}{\PYGZdq{}}\PYG{l+s+s2}{[kg/l] is:}\PYG{l+s+s2}{\PYGZdq{}}\PYG{p}{,} \PYG{n}{process\PYGZus{}time}\PYG{p}{,} \PYG{l+s+s2}{\PYGZdq{}}\PYG{l+s+s2}{[h]}\PYG{l+s+s2}{\PYGZdq{}}\PYG{p}{)} 
\end{sphinxVerbatim}

\end{sphinxuseclass}\end{sphinxVerbatimInput}
\begin{sphinxVerbatimOutput}

\begin{sphinxuseclass}{cell_output}
\begin{sphinxVerbatim}[commandchars=\\\{\}]
The process time necessary to obtain a solid residue concentration of 200.0 [kg/l] is: [17.32867951] [h]
\end{sphinxVerbatim}

\noindent\sphinxincludegraphics{{4f23a8e94ff9c50f28bfd2ec62217f4a2f486df7083cc5783a264151ff4c814a}.png}

\noindent\sphinxincludegraphics{{18eb7f3bca806d6744eb1cbd1fa0beed088e73fbdeedd19dc60d88342c767e76}.png}

\end{sphinxuseclass}\end{sphinxVerbatimOutput}

\end{sphinxuseclass}

\subsection{4.2 Feed and Bleed}
\label{\detokenize{Process_layouts:feed-and-bleed}}
\sphinxAtStartPar
\sphinxstylestrong{Feed and Bleed:  Problem Statement}

\sphinxAtStartPar
500 l of fruit juice are concentrated in 5 h of steady state operation from an initial solid\sphinxhyphen{}residue content of 0.05 \(kg/l\) through a feed and bleed process. The total area of the membrane is of 20 \(m^2\) and the flux of pure water through the membrane is captured by the empirical expression: \(J=BC^{-1} [m/h]\) with \(B=0.1\) where C is the solid residue concentration in kg/l. Is the steady state concentration of the retentate compatible with the specifics of 0.2 \(kg/l\)?

\begin{figure}[htbp]
\centering
\capstart

\noindent\sphinxincludegraphics[height=300\sphinxpxdimen]{{feedandbleed}.png}
\caption{Feed and Bleed process configuration scheme}\label{\detokenize{Process_layouts:id1}}\end{figure}

\sphinxAtStartPar
\sphinxstylestrong{Feed and Bleed: Solution trace}

\sphinxAtStartPar
Also in this case the global differential material balance and the solid\sphinxhyphen{}residue material balance can be written as follows:
\begin{equation}\label{equation:Process_layouts:PLYeq11}
\begin{split}
\frac{dV}{dt}=Q_{IN}-J\,A-Q_{OUT}=0
\end{split}
\end{equation}\begin{equation}\label{equation:Process_layouts:PLYeq12}
\begin{split}
\frac{dm}{dt}=Q_{IN}C_{IN}-Q_{OUT}C_{OUT}=0
\end{split}
\end{equation}
\sphinxAtStartPar
At steady state \({dV}/{dt}=0\) as well as \({dm}/{dt}=0\). The two ODEs become then two algebraic equations that should be solved together to compute the steady state concentration.
From the global material balance we get:
\begin{equation}\label{equation:Process_layouts:PLYeq13}
\begin{split}
Q_{OUT}=Q_{IN}-J\,A=Q_{IN}-\frac{B}{C_{OUT}}A
\end{split}
\end{equation}
\sphinxAtStartPar
and then, with some manipulations:
\begin{equation}\label{equation:Process_layouts:PLYeq14}
\begin{split}
C_{OUT}=C_{IN}+\frac{AB}{Q_{IN}}
\end{split}
\end{equation}
\begin{sphinxuseclass}{cell}\begin{sphinxVerbatimInput}

\begin{sphinxuseclass}{cell_input}
\begin{sphinxVerbatim}[commandchars=\\\{\}]
\PYG{k+kn}{import}\PYG{+w}{ }\PYG{n+nn}{numpy}\PYG{+w}{ }\PYG{k}{as}\PYG{+w}{ }\PYG{n+nn}{np}
\PYG{k+kn}{import}\PYG{+w}{ }\PYG{n+nn}{matplotlib}\PYG{n+nn}{.}\PYG{n+nn}{pyplot}\PYG{+w}{ }\PYG{k}{as}\PYG{+w}{ }\PYG{n+nn}{plt} 
\PYG{k+kn}{from}\PYG{+w}{ }\PYG{n+nn}{scipy}\PYG{n+nn}{.}\PYG{n+nn}{optimize}\PYG{+w}{ }\PYG{k+kn}{import} \PYG{n}{fsolve}

\PYG{c+c1}{\PYGZsh{} data: }
\PYG{n}{A}\PYG{o}{=}\PYG{l+m+mi}{20}\PYG{p}{;}  \PYG{c+c1}{\PYGZsh{}[m\PYGZca{}2]}
\PYG{n}{B}\PYG{o}{=}\PYG{l+m+mf}{0.1}\PYG{p}{;}
\PYG{n}{C0}\PYG{o}{=}\PYG{l+m+mf}{0.05}\PYG{o}{*}\PYG{l+m+mf}{1E3}\PYG{p}{;} \PYG{c+c1}{\PYGZsh{}kg /l * dm\PYGZca{}3/m\PYGZca{}3}
\PYG{n}{V0}\PYG{o}{=}\PYG{l+m+mi}{500}\PYG{o}{*}\PYG{l+m+mf}{1E\PYGZhy{}3}\PYG{p}{;} \PYG{c+c1}{\PYGZsh{}l * m\PYGZca{}3/dm\PYGZca{}3}
\PYG{n}{process\PYGZus{}time}\PYG{o}{=}\PYG{l+m+mi}{5}\PYG{p}{;} \PYG{c+c1}{\PYGZsh{} [h]}

\PYG{n}{Cout}\PYG{o}{=}\PYG{n}{C0}\PYG{o}{+}\PYG{n}{A}\PYG{o}{*}\PYG{n}{B}\PYG{o}{*}\PYG{n}{process\PYGZus{}time}\PYG{o}{/}\PYG{n}{V0}


\PYG{n+nb}{print}\PYG{p}{(}\PYG{l+s+s2}{\PYGZdq{}}\PYG{l+s+s2}{The steady state concentration is}\PYG{l+s+s2}{\PYGZdq{}}\PYG{p}{,} \PYG{n}{Cout}\PYG{p}{,} \PYG{l+s+s2}{\PYGZdq{}}\PYG{l+s+s2}{[kg/m\PYGZca{}3]}\PYG{l+s+s2}{\PYGZdq{}}\PYG{p}{)} 
\end{sphinxVerbatim}

\end{sphinxuseclass}\end{sphinxVerbatimInput}
\begin{sphinxVerbatimOutput}

\begin{sphinxuseclass}{cell_output}
\begin{sphinxVerbatim}[commandchars=\\\{\}]
The steady state concentration is 70.0 [kg/m\PYGZca{}3]
\end{sphinxVerbatim}

\end{sphinxuseclass}\end{sphinxVerbatimOutput}

\end{sphinxuseclass}

\subsection{4.3 Cascade configuration}
\label{\detokenize{Process_layouts:cascade-configuration}}
\sphinxAtStartPar
\sphinxstylestrong{Problem Statement}

\sphinxAtStartPar
500 l of fruit juice are concentrated in 5 h of steady state operation from an initial solid\sphinxhyphen{}residue content of 0.05 kg/l through four membrane separation units characterised by a total membrane area of 20 \(\mathrm{m^2}\) each.
The flux of pure water through the membrane is captured by the empirical expression: \(J=BC^{-1} [m/h]\) with \(B=0.1\) where C is the solid residue concentration in kg/l.

\sphinxAtStartPar
Is it more efficient to design a single stage configuration with four units in parallel or a cascade configuration with four units in series?

\begin{figure}[htbp]
\centering
\capstart

\noindent\sphinxincludegraphics[height=300\sphinxpxdimen]{{cascade}.png}
\caption{Cascade configuration schematic representation.}\label{\detokenize{Process_layouts:id2}}\end{figure}

\sphinxAtStartPar
\sphinxstylestrong{Cascade: Solution trace}

\sphinxAtStartPar
Each stage can be treated like a single unit in which the feed stream corresponds to the retentate stream from the previous unit.
We can thus define:
\begin{equation}\label{equation:Process_layouts:PLYeq15}
\begin{split}
C_{OUT,i}=C_{IN,i+1}
\end{split}
\end{equation}\begin{equation}\label{equation:Process_layouts:PLYeq16}
\begin{split}
Q_{OUT,i}=Q_{IN,i+1}
\end{split}
\end{equation}
\sphinxAtStartPar
where the index \(i\) identifies the stage.

\sphinxAtStartPar
The material balances for each stage at steady state can thus be written as:
\begin{equation}\label{equation:Process_layouts:PLYeq17}
\begin{split}
\frac{dV}{dt}=Q_{IN,i}-n_iJ(C_{IN,i+1})\,A-Q_{IN,i+1}=0
\end{split}
\end{equation}\begin{equation}\label{equation:Process_layouts:PLYeq18}
\begin{split}
\frac{dm}{dt}=Q_{IN,i}C_{IN}-Q_{IN,i+1}C_{IN,i+1}=0
\end{split}
\end{equation}
\sphinxAtStartPar
where \(n_i\) is the number of modules used in stage \(i\).

\sphinxAtStartPar
For each stage we can thus compute the steady state concentration and concentrate volumetric flow by solving sequentially the following equations:
\begin{equation}\label{equation:Process_layouts:PLYeq19}
\begin{split}
C_{IN,i+1}=C_{IN,i}+n_i\frac{AB}{Q_{IN,i}}
\end{split}
\end{equation}\begin{equation}\label{equation:Process_layouts:PLYeq20}
\begin{split}
Q_{IN,i+1}=\frac{C_{IN,i}{Q_{IN,i}}}{C_{IN,i+1}}
\end{split}
\end{equation}
\sphinxAtStartPar
with \(i=1...N\) with \(N\) is the total number of stages.

\sphinxAtStartPar
\sphinxstylestrong{Numerical Solution}

\sphinxAtStartPar
The solution of a cascade composed of any number of stages, each formed by an arbitrary number of modules in parallel can be tackled sequentially through a simple cycle similar to the following, let’s for example consider a system that reflects the sketch above i.e. with 4 stages implementing 4, 3, 2, and 1 modules each:

\begin{sphinxuseclass}{cell}\begin{sphinxVerbatimInput}

\begin{sphinxuseclass}{cell_input}
\begin{sphinxVerbatim}[commandchars=\\\{\}]
\PYG{k+kn}{import}\PYG{+w}{ }\PYG{n+nn}{numpy}\PYG{+w}{ }\PYG{k}{as}\PYG{+w}{ }\PYG{n+nn}{np}
\PYG{k+kn}{import}\PYG{+w}{ }\PYG{n+nn}{matplotlib}\PYG{n+nn}{.}\PYG{n+nn}{pyplot}\PYG{+w}{ }\PYG{k}{as}\PYG{+w}{ }\PYG{n+nn}{plt} 
\PYG{k+kn}{from}\PYG{+w}{ }\PYG{n+nn}{scipy}\PYG{n+nn}{.}\PYG{n+nn}{optimize}\PYG{+w}{ }\PYG{k+kn}{import} \PYG{n}{fsolve}

\PYG{c+c1}{\PYGZsh{} data: }
\PYG{n}{A}\PYG{o}{=}\PYG{l+m+mi}{20}\PYG{p}{;}  
\PYG{n}{B}\PYG{o}{=}\PYG{l+m+mf}{0.1}\PYG{p}{;}
\PYG{n}{C0}\PYG{o}{=}\PYG{l+m+mf}{0.05}\PYG{o}{*}\PYG{l+m+mf}{1E3}\PYG{p}{;} \PYG{c+c1}{\PYGZsh{}kg /l * dm\PYGZca{}3/m\PYGZca{}3}
\PYG{n}{V0}\PYG{o}{=}\PYG{l+m+mi}{500}\PYG{o}{*}\PYG{l+m+mf}{1E\PYGZhy{}3}\PYG{p}{;} \PYG{c+c1}{\PYGZsh{}l * m\PYGZca{}3/dm\PYGZca{}3}
\PYG{n}{process\PYGZus{}time}\PYG{o}{=}\PYG{l+m+mi}{5}\PYG{p}{;} \PYG{c+c1}{\PYGZsh{} [h]}


\PYG{c+c1}{\PYGZsh{} The number of elements of this array corresponds to the number of stages. }
\PYG{c+c1}{\PYGZsh{} The value in each element is the number of modules per stage. }
\PYG{n}{n}\PYG{o}{=}\PYG{n}{np}\PYG{o}{.}\PYG{n}{array}\PYG{p}{(}\PYG{p}{[}\PYG{l+m+mi}{4}\PYG{p}{,}\PYG{l+m+mi}{3}\PYG{p}{,}\PYG{l+m+mi}{2}\PYG{p}{,}\PYG{l+m+mi}{1}\PYG{p}{]}\PYG{p}{)}\PYG{p}{;}

\PYG{n}{CIN}\PYG{o}{=}\PYG{n}{np}\PYG{o}{.}\PYG{n}{append}\PYG{p}{(}\PYG{n}{C0}\PYG{p}{,}\PYG{n}{np}\PYG{o}{.}\PYG{n}{zeros}\PYG{p}{(}\PYG{n}{np}\PYG{o}{.}\PYG{n}{size}\PYG{p}{(}\PYG{n}{n}\PYG{p}{)}\PYG{o}{\PYGZhy{}}\PYG{l+m+mi}{1}\PYG{p}{)}\PYG{p}{)}
\PYG{n}{FIN}\PYG{o}{=}\PYG{n}{np}\PYG{o}{.}\PYG{n}{append}\PYG{p}{(}\PYG{n}{V0}\PYG{o}{/}\PYG{n}{process\PYGZus{}time}\PYG{p}{,} \PYG{n}{np}\PYG{o}{.}\PYG{n}{zeros}\PYG{p}{(}\PYG{n}{np}\PYG{o}{.}\PYG{n}{size}\PYG{p}{(}\PYG{n}{n}\PYG{p}{)}\PYG{o}{\PYGZhy{}}\PYG{l+m+mi}{1}\PYG{p}{)}\PYG{p}{)}\PYG{p}{;}

\PYG{c+c1}{\PYGZsh{} Input to the intermediate stages}
\PYG{k}{for} \PYG{n}{i} \PYG{o+ow}{in} \PYG{n+nb}{range}\PYG{p}{(}\PYG{l+m+mi}{1}\PYG{p}{,}\PYG{n}{np}\PYG{o}{.}\PYG{n}{size}\PYG{p}{(}\PYG{n}{n}\PYG{p}{)}\PYG{p}{)}\PYG{p}{:} 
    \PYG{n}{CIN}\PYG{p}{[}\PYG{n}{i}\PYG{p}{]}\PYG{o}{=}\PYG{n}{CIN}\PYG{p}{[}\PYG{n}{i}\PYG{o}{\PYGZhy{}}\PYG{l+m+mi}{1}\PYG{p}{]}\PYG{o}{+}\PYG{n}{n}\PYG{p}{[}\PYG{n}{i}\PYG{o}{\PYGZhy{}}\PYG{l+m+mi}{1}\PYG{p}{]}\PYG{o}{*}\PYG{n}{A}\PYG{o}{*}\PYG{n}{B}\PYG{o}{/}\PYG{n}{FIN}\PYG{p}{[}\PYG{n}{i}\PYG{o}{\PYGZhy{}}\PYG{l+m+mi}{1}\PYG{p}{]}\PYG{p}{;}
    \PYG{n}{FIN}\PYG{p}{[}\PYG{n}{i}\PYG{p}{]}\PYG{o}{=}\PYG{n}{CIN}\PYG{p}{[}\PYG{n}{i}\PYG{o}{\PYGZhy{}}\PYG{l+m+mi}{1}\PYG{p}{]}\PYG{o}{*}\PYG{n}{FIN}\PYG{p}{[}\PYG{n}{i}\PYG{o}{\PYGZhy{}}\PYG{l+m+mi}{1}\PYG{p}{]}\PYG{o}{/}\PYG{n}{CIN}\PYG{p}{[}\PYG{n}{i}\PYG{p}{]}\PYG{p}{;}

\PYG{c+c1}{\PYGZsh{} Output concentration}
\PYG{n}{Cout}\PYG{o}{=}\PYG{n}{CIN}\PYG{p}{[}\PYG{n}{np}\PYG{o}{.}\PYG{n}{size}\PYG{p}{(}\PYG{n}{n}\PYG{p}{)}\PYG{o}{\PYGZhy{}}\PYG{l+m+mi}{1}\PYG{p}{]}\PYG{o}{+}\PYG{n}{n}\PYG{p}{[}\PYG{n}{np}\PYG{o}{.}\PYG{n}{size}\PYG{p}{(}\PYG{n}{n}\PYG{p}{)}\PYG{o}{\PYGZhy{}}\PYG{l+m+mi}{1}\PYG{p}{]}\PYG{o}{*}\PYG{n}{A}\PYG{o}{*}\PYG{n}{B}\PYG{o}{/}\PYG{n}{FIN}\PYG{p}{[}\PYG{n}{np}\PYG{o}{.}\PYG{n}{size}\PYG{p}{(}\PYG{n}{n}\PYG{p}{)}\PYG{o}{\PYGZhy{}}\PYG{l+m+mi}{1}\PYG{p}{]}\PYG{p}{;}

\PYG{n+nb}{print}\PYG{p}{(}\PYG{l+s+s2}{\PYGZdq{}}\PYG{l+s+s2}{The steady state concentration is}\PYG{l+s+s2}{\PYGZdq{}}\PYG{p}{,} \PYG{n}{Cout}\PYG{p}{,} \PYG{l+s+s2}{\PYGZdq{}}\PYG{l+s+s2}{[kg/m\PYGZca{}3]}\PYG{l+s+s2}{\PYGZdq{}}\PYG{p}{)} 
\end{sphinxVerbatim}

\end{sphinxuseclass}\end{sphinxVerbatimInput}
\begin{sphinxVerbatimOutput}

\begin{sphinxuseclass}{cell_output}
\begin{sphinxVerbatim}[commandchars=\\\{\}]
The steady state concentration is 720.7199999999999 [kg/m\PYGZca{}3]
\end{sphinxVerbatim}

\end{sphinxuseclass}\end{sphinxVerbatimOutput}

\end{sphinxuseclass}
\sphinxAtStartPar
In order to answer the problem request one should solve the system for two different configurations.
In the first, representing a single\sphinxhyphen{}stage configuration with four membrane units in parallel, the number of stages should be set to \(N=1\) and the number of units in the first stage \(n_1\) to 4.

\begin{sphinxuseclass}{cell}\begin{sphinxVerbatimInput}

\begin{sphinxuseclass}{cell_input}
\begin{sphinxVerbatim}[commandchars=\\\{\}]
\PYG{c+c1}{\PYGZsh{} The number of elements of this array corresponds to the number of stages. }
\PYG{c+c1}{\PYGZsh{} The value in each element is the number of modules per stage. }
\PYG{n}{n}\PYG{o}{=}\PYG{n}{np}\PYG{o}{.}\PYG{n}{array}\PYG{p}{(}\PYG{p}{[}\PYG{l+m+mi}{4}\PYG{p}{]}\PYG{p}{)}\PYG{p}{;}

\PYG{n}{CIN}\PYG{o}{=}\PYG{n}{np}\PYG{o}{.}\PYG{n}{append}\PYG{p}{(}\PYG{n}{C0}\PYG{p}{,}\PYG{n}{np}\PYG{o}{.}\PYG{n}{zeros}\PYG{p}{(}\PYG{n}{np}\PYG{o}{.}\PYG{n}{size}\PYG{p}{(}\PYG{n}{n}\PYG{p}{)}\PYG{o}{\PYGZhy{}}\PYG{l+m+mi}{1}\PYG{p}{)}\PYG{p}{)}
\PYG{n}{FIN}\PYG{o}{=}\PYG{n}{np}\PYG{o}{.}\PYG{n}{append}\PYG{p}{(}\PYG{n}{V0}\PYG{o}{/}\PYG{n}{process\PYGZus{}time}\PYG{p}{,} \PYG{n}{np}\PYG{o}{.}\PYG{n}{zeros}\PYG{p}{(}\PYG{n}{np}\PYG{o}{.}\PYG{n}{size}\PYG{p}{(}\PYG{n}{n}\PYG{p}{)}\PYG{o}{\PYGZhy{}}\PYG{l+m+mi}{1}\PYG{p}{)}\PYG{p}{)}\PYG{p}{;}

\PYG{c+c1}{\PYGZsh{} Input to the intermediate stages}
\PYG{k}{for} \PYG{n}{i} \PYG{o+ow}{in} \PYG{n+nb}{range}\PYG{p}{(}\PYG{l+m+mi}{1}\PYG{p}{,}\PYG{n}{np}\PYG{o}{.}\PYG{n}{size}\PYG{p}{(}\PYG{n}{n}\PYG{p}{)}\PYG{p}{)}\PYG{p}{:} 
    \PYG{n}{CIN}\PYG{p}{[}\PYG{n}{i}\PYG{p}{]}\PYG{o}{=}\PYG{n}{CIN}\PYG{p}{[}\PYG{n}{i}\PYG{o}{\PYGZhy{}}\PYG{l+m+mi}{1}\PYG{p}{]}\PYG{o}{+}\PYG{n}{n}\PYG{p}{[}\PYG{n}{i}\PYG{o}{\PYGZhy{}}\PYG{l+m+mi}{1}\PYG{p}{]}\PYG{o}{*}\PYG{n}{A}\PYG{o}{*}\PYG{n}{B}\PYG{o}{/}\PYG{n}{FIN}\PYG{p}{[}\PYG{n}{i}\PYG{o}{\PYGZhy{}}\PYG{l+m+mi}{1}\PYG{p}{]}\PYG{p}{;}
    \PYG{n}{FIN}\PYG{p}{[}\PYG{n}{i}\PYG{p}{]}\PYG{o}{=}\PYG{n}{CIN}\PYG{p}{[}\PYG{n}{i}\PYG{o}{\PYGZhy{}}\PYG{l+m+mi}{1}\PYG{p}{]}\PYG{o}{*}\PYG{n}{FIN}\PYG{p}{[}\PYG{n}{i}\PYG{o}{\PYGZhy{}}\PYG{l+m+mi}{1}\PYG{p}{]}\PYG{o}{/}\PYG{n}{CIN}\PYG{p}{[}\PYG{n}{i}\PYG{p}{]}\PYG{p}{;}

\PYG{c+c1}{\PYGZsh{} Output concentration}
\PYG{n}{Cout}\PYG{o}{=}\PYG{n}{CIN}\PYG{p}{[}\PYG{n}{np}\PYG{o}{.}\PYG{n}{size}\PYG{p}{(}\PYG{n}{n}\PYG{p}{)}\PYG{o}{\PYGZhy{}}\PYG{l+m+mi}{1}\PYG{p}{]}\PYG{o}{+}\PYG{n}{n}\PYG{p}{[}\PYG{n}{np}\PYG{o}{.}\PYG{n}{size}\PYG{p}{(}\PYG{n}{n}\PYG{p}{)}\PYG{o}{\PYGZhy{}}\PYG{l+m+mi}{1}\PYG{p}{]}\PYG{o}{*}\PYG{n}{A}\PYG{o}{*}\PYG{n}{B}\PYG{o}{/}\PYG{n}{FIN}\PYG{p}{[}\PYG{n}{np}\PYG{o}{.}\PYG{n}{size}\PYG{p}{(}\PYG{n}{n}\PYG{p}{)}\PYG{o}{\PYGZhy{}}\PYG{l+m+mi}{1}\PYG{p}{]}\PYG{p}{;}

\PYG{n+nb}{print}\PYG{p}{(}\PYG{l+s+s2}{\PYGZdq{}}\PYG{l+s+s2}{The steady state concentration is}\PYG{l+s+s2}{\PYGZdq{}}\PYG{p}{,} \PYG{n}{Cout}\PYG{p}{,} \PYG{l+s+s2}{\PYGZdq{}}\PYG{l+s+s2}{[kg/m\PYGZca{}3]}\PYG{l+s+s2}{\PYGZdq{}}\PYG{p}{)} 
\end{sphinxVerbatim}

\end{sphinxuseclass}\end{sphinxVerbatimInput}
\begin{sphinxVerbatimOutput}

\begin{sphinxuseclass}{cell_output}
\begin{sphinxVerbatim}[commandchars=\\\{\}]
The steady state concentration is 130.0 [kg/m\PYGZca{}3]
\end{sphinxVerbatim}

\end{sphinxuseclass}\end{sphinxVerbatimOutput}

\end{sphinxuseclass}
\sphinxAtStartPar
In the second the number of stages should be set to \(N=4\), each of the stages being assembled as a single unit (\(n_i=1\) for \(i=[1, 4]\)).

\begin{sphinxuseclass}{cell}\begin{sphinxVerbatimInput}

\begin{sphinxuseclass}{cell_input}
\begin{sphinxVerbatim}[commandchars=\\\{\}]
\PYG{c+c1}{\PYGZsh{} The number of elements of this array corresponds to the number of stages. }
\PYG{c+c1}{\PYGZsh{} The value in each element is the number of modules per stage. }
\PYG{n}{n}\PYG{o}{=}\PYG{n}{np}\PYG{o}{.}\PYG{n}{array}\PYG{p}{(}\PYG{p}{[}\PYG{l+m+mi}{1}\PYG{p}{,} \PYG{l+m+mi}{1}\PYG{p}{,} \PYG{l+m+mi}{1}\PYG{p}{,} \PYG{l+m+mi}{1}\PYG{p}{]}\PYG{p}{)}\PYG{p}{;}

\PYG{n}{CIN}\PYG{o}{=}\PYG{n}{np}\PYG{o}{.}\PYG{n}{append}\PYG{p}{(}\PYG{n}{C0}\PYG{p}{,}\PYG{n}{np}\PYG{o}{.}\PYG{n}{zeros}\PYG{p}{(}\PYG{n}{np}\PYG{o}{.}\PYG{n}{size}\PYG{p}{(}\PYG{n}{n}\PYG{p}{)}\PYG{o}{\PYGZhy{}}\PYG{l+m+mi}{1}\PYG{p}{)}\PYG{p}{)}
\PYG{n}{FIN}\PYG{o}{=}\PYG{n}{np}\PYG{o}{.}\PYG{n}{append}\PYG{p}{(}\PYG{n}{V0}\PYG{o}{/}\PYG{n}{process\PYGZus{}time}\PYG{p}{,} \PYG{n}{np}\PYG{o}{.}\PYG{n}{zeros}\PYG{p}{(}\PYG{n}{np}\PYG{o}{.}\PYG{n}{size}\PYG{p}{(}\PYG{n}{n}\PYG{p}{)}\PYG{o}{\PYGZhy{}}\PYG{l+m+mi}{1}\PYG{p}{)}\PYG{p}{)}\PYG{p}{;}

\PYG{c+c1}{\PYGZsh{} Input to the intermediate stages}
\PYG{k}{for} \PYG{n}{i} \PYG{o+ow}{in} \PYG{n+nb}{range}\PYG{p}{(}\PYG{l+m+mi}{1}\PYG{p}{,}\PYG{n}{np}\PYG{o}{.}\PYG{n}{size}\PYG{p}{(}\PYG{n}{n}\PYG{p}{)}\PYG{p}{)}\PYG{p}{:} 
    \PYG{n}{CIN}\PYG{p}{[}\PYG{n}{i}\PYG{p}{]}\PYG{o}{=}\PYG{n}{CIN}\PYG{p}{[}\PYG{n}{i}\PYG{o}{\PYGZhy{}}\PYG{l+m+mi}{1}\PYG{p}{]}\PYG{o}{+}\PYG{n}{n}\PYG{p}{[}\PYG{n}{i}\PYG{o}{\PYGZhy{}}\PYG{l+m+mi}{1}\PYG{p}{]}\PYG{o}{*}\PYG{n}{A}\PYG{o}{*}\PYG{n}{B}\PYG{o}{/}\PYG{n}{FIN}\PYG{p}{[}\PYG{n}{i}\PYG{o}{\PYGZhy{}}\PYG{l+m+mi}{1}\PYG{p}{]}\PYG{p}{;}
    \PYG{n}{FIN}\PYG{p}{[}\PYG{n}{i}\PYG{p}{]}\PYG{o}{=}\PYG{n}{CIN}\PYG{p}{[}\PYG{n}{i}\PYG{o}{\PYGZhy{}}\PYG{l+m+mi}{1}\PYG{p}{]}\PYG{o}{*}\PYG{n}{FIN}\PYG{p}{[}\PYG{n}{i}\PYG{o}{\PYGZhy{}}\PYG{l+m+mi}{1}\PYG{p}{]}\PYG{o}{/}\PYG{n}{CIN}\PYG{p}{[}\PYG{n}{i}\PYG{p}{]}\PYG{p}{;}

\PYG{c+c1}{\PYGZsh{} Output concentration}
\PYG{n}{Cout}\PYG{o}{=}\PYG{n}{CIN}\PYG{p}{[}\PYG{n}{np}\PYG{o}{.}\PYG{n}{size}\PYG{p}{(}\PYG{n}{n}\PYG{p}{)}\PYG{o}{\PYGZhy{}}\PYG{l+m+mi}{1}\PYG{p}{]}\PYG{o}{+}\PYG{n}{n}\PYG{p}{[}\PYG{n}{np}\PYG{o}{.}\PYG{n}{size}\PYG{p}{(}\PYG{n}{n}\PYG{p}{)}\PYG{o}{\PYGZhy{}}\PYG{l+m+mi}{1}\PYG{p}{]}\PYG{o}{*}\PYG{n}{A}\PYG{o}{*}\PYG{n}{B}\PYG{o}{/}\PYG{n}{FIN}\PYG{p}{[}\PYG{n}{np}\PYG{o}{.}\PYG{n}{size}\PYG{p}{(}\PYG{n}{n}\PYG{p}{)}\PYG{o}{\PYGZhy{}}\PYG{l+m+mi}{1}\PYG{p}{]}\PYG{p}{;}

\PYG{n+nb}{print}\PYG{p}{(}\PYG{l+s+s2}{\PYGZdq{}}\PYG{l+s+s2}{The steady state concentration is}\PYG{l+s+s2}{\PYGZdq{}}\PYG{p}{,} \PYG{n}{Cout}\PYG{p}{,} \PYG{l+s+s2}{\PYGZdq{}}\PYG{l+s+s2}{[kg/m\PYGZca{}3]}\PYG{l+s+s2}{\PYGZdq{}}\PYG{p}{)} 
\end{sphinxVerbatim}

\end{sphinxuseclass}\end{sphinxVerbatimInput}
\begin{sphinxVerbatimOutput}

\begin{sphinxuseclass}{cell_output}
\begin{sphinxVerbatim}[commandchars=\\\{\}]
The steady state concentration is 192.07999999999998 [kg/m\PYGZca{}3]
\end{sphinxVerbatim}

\end{sphinxuseclass}\end{sphinxVerbatimOutput}

\end{sphinxuseclass}
\sphinxAtStartPar
The cascade configuration allows for a more efficient process since it allows to obtain a larger concentration with the same number of modules.

\sphinxstepscope


\section{5. The Langmuir Isotherm}
\label{\detokenize{Langmuir_isotherm:the-langmuir-isotherm}}\label{\detokenize{Langmuir_isotherm::doc}}
\sphinxAtStartPar
The Langmuir isotherm provides a relationship between the compositions of the fluid phase and of the adsorbed phase at equilibrium.

\sphinxAtStartPar
In the following we shall note with \(C\) the concentration of the component that undergoes adsorption in the fluid phase, and with \(q\) its  molar fraction in the adsorbed phase.

\sphinxAtStartPar
The Langmuir isother is based on a set of assumptions:
\begin{itemize}
\item {} 
\sphinxAtStartPar
The kinetic constants of the adsorption and the desorption processes, as well as the equilibrium constant for the adsorption process, are independent from the fraction of free adsorption sites.

\item {} 
\sphinxAtStartPar
Interactions between adsorbed molecules are neglibile, or in other words the adsorption process does not depend on the local environment around an adsorption site.

\item {} 
\sphinxAtStartPar
The surface of the adsorbent material can host a single monolayer of adsorbed molecules, hence its fraction of free sites \(\theta_0\) and the fraction of occupied sites \(\Gamma_1\) are related through the stoichiometric relation \(\Gamma_0+\Gamma_1=1\).

\end{itemize}

\sphinxAtStartPar
In order to derive the langmuir adsorption isotherm from this set of hypotheses let’s start by defining expressions for the adsorption and desorption rates.

\sphinxAtStartPar
The adsorption process takes place when a molecule of the component that is undergoing adsorption finds a free adsorption site. The rate of an adsorption event can thus be considered analogous to a bimolecular elementary reaction step with rate:
\begin{equation}\label{equation:Langmuir_isotherm:LANGeq1}
\begin{split}
R_A=k_A\Gamma_0{C}
\end{split}
\end{equation}
\sphinxAtStartPar
Where \(k_A\) is the kinetic constant, \(\Gamma_0\) is the fraction of free adsorption sites,  and \(x\) is the molar fraction of the adsorbing component in the fluid phase.

\sphinxAtStartPar
The rate of desorption instead can be considered analogous to a monomolecular reaction, in which a (non\sphinxhyphen{}covalent) bond is broken and an adsorbed molecule is released to the fluid phase, leaving one free site. Its rate can be written as:
\begin{equation}\label{equation:Langmuir_isotherm:LANGeq2}
\begin{split}
R_D=k_D\Gamma_1
\end{split}
\end{equation}
\sphinxAtStartPar
At equilibrium, the rates of adsorption and desorption are equal and the net flux of molecules to and from the adsorbed phase is null.
We can thus write:
\begin{equation}\label{equation:Langmuir_isotherm:LANGeq3}
\begin{split}
k_A\Gamma_0{C}=k_D\Gamma_1
\end{split}
\end{equation}
\sphinxAtStartPar
Recalling the stoichiometric relation between the fraction of accupied and free sites Eq. \textbackslash{}eqref\{eq:3\} becomes:
\begin{equation}\label{equation:Langmuir_isotherm:LANGeq4}
\begin{split}
k_A(1-\Gamma_1){C}=k_D\Gamma_1
\end{split}
\end{equation}
\sphinxAtStartPar
which can be rearranged to provide an explicit expression of \(\Gamma_1\), the fraction of accupied sites, as a function of the concentration \(C\) in the fluid phase:
\begin{equation}\label{equation:Langmuir_isotherm:LANGeq5}
\begin{split}
\Gamma_1=\frac{k_A{C}}{k_D+k_A{C}}
\end{split}
\end{equation}
\sphinxAtStartPar
To obtain a typical expression of the Langmuir isotherm we shall introduce as a single parameter for this equation the adsorption equilibrium constant \(B_0=k_A/k_D\):
\begin{equation}\label{equation:Langmuir_isotherm:LANGeq6}
\begin{split}
\Gamma_1=\frac{B_0{C}}{1+B_0{C}}
\end{split}
\end{equation}

\subsection{5.1 Isotherm dependence on \protect\(B_0\protect\)}
\label{\detokenize{Langmuir_isotherm:isotherm-dependence-on-b-0}}
\sphinxAtStartPar
In order to understand the role of the adsorption equilibrium constant let us plot the fraction of occupied adsorption sites as a function of \(B_0\)

\begin{sphinxuseclass}{cell}\begin{sphinxVerbatimInput}

\begin{sphinxuseclass}{cell_input}
\begin{sphinxVerbatim}[commandchars=\\\{\}]
\PYG{k+kn}{import}\PYG{+w}{ }\PYG{n+nn}{numpy}\PYG{+w}{ }\PYG{k}{as}\PYG{+w}{ }\PYG{n+nn}{np}
\PYG{k+kn}{import}\PYG{+w}{ }\PYG{n+nn}{matplotlib}\PYG{n+nn}{.}\PYG{n+nn}{pyplot}\PYG{+w}{ }\PYG{k}{as}\PYG{+w}{ }\PYG{n+nn}{plt} 
\PYG{k+kn}{from}\PYG{+w}{ }\PYG{n+nn}{matplotlib}\PYG{n+nn}{.}\PYG{n+nn}{pyplot}\PYG{+w}{ }\PYG{k+kn}{import} \PYG{n}{cm}

\PYG{c+c1}{\PYGZsh{}Plotting}
\PYG{n}{figure}\PYG{o}{=}\PYG{n}{plt}\PYG{o}{.}\PYG{n}{figure}\PYG{p}{(}\PYG{p}{)}
\PYG{n}{axes} \PYG{o}{=} \PYG{n}{figure}\PYG{o}{.}\PYG{n}{add\PYGZus{}axes}\PYG{p}{(}\PYG{p}{[}\PYG{l+m+mf}{0.1}\PYG{p}{,}\PYG{l+m+mf}{0.1}\PYG{p}{,}\PYG{l+m+mf}{1.2}\PYG{p}{,}\PYG{l+m+mf}{1.2}\PYG{p}{]}\PYG{p}{)}
\PYG{n}{plt}\PYG{o}{.}\PYG{n}{xticks}\PYG{p}{(}\PYG{n}{fontsize}\PYG{o}{=}\PYG{l+m+mi}{14}\PYG{p}{)}
\PYG{n}{plt}\PYG{o}{.}\PYG{n}{yticks}\PYG{p}{(}\PYG{n}{fontsize}\PYG{o}{=}\PYG{l+m+mi}{14}\PYG{p}{)}


\PYG{c+c1}{\PYGZsh{} molar fraction in the feed}
\PYG{n}{z} \PYG{o}{=} \PYG{l+m+mf}{0.3}
\PYG{n}{N} \PYG{o}{=} \PYG{l+m+mi}{500} \PYG{c+c1}{\PYGZsh{}number of points}
\PYG{n}{C} \PYG{o}{=} \PYG{n}{np}\PYG{o}{.}\PYG{n}{linspace}\PYG{p}{(}\PYG{l+m+mi}{0}\PYG{p}{,} \PYG{l+m+mi}{1}\PYG{p}{,} \PYG{n}{N}\PYG{p}{)}
\PYG{n}{B0} \PYG{o}{=} \PYG{n}{np}\PYG{o}{.}\PYG{n}{array}\PYG{p}{(}\PYG{p}{[}\PYG{l+m+mf}{0.1}\PYG{p}{,} \PYG{l+m+mi}{1}\PYG{p}{,} \PYG{l+m+mi}{2}\PYG{p}{,} \PYG{l+m+mi}{5}\PYG{p}{,} \PYG{l+m+mi}{10}\PYG{p}{,} \PYG{l+m+mi}{50}\PYG{p}{,} \PYG{l+m+mi}{100}\PYG{p}{,} \PYG{l+m+mi}{1000}\PYG{p}{]}\PYG{p}{)}\PYG{p}{;}

\PYG{n}{color}\PYG{o}{=}\PYG{n+nb}{iter}\PYG{p}{(}\PYG{n}{cm}\PYG{o}{.}\PYG{n}{gist\PYGZus{}heat}\PYG{p}{(}\PYG{n}{np}\PYG{o}{.}\PYG{n}{linspace}\PYG{p}{(}\PYG{l+m+mi}{0}\PYG{p}{,}\PYG{l+m+mi}{1}\PYG{p}{,}\PYG{n}{np}\PYG{o}{.}\PYG{n}{size}\PYG{p}{(}\PYG{n}{B0}\PYG{p}{)}\PYG{p}{)}\PYG{p}{)}\PYG{p}{)}

\PYG{k}{for} \PYG{n}{i} \PYG{o+ow}{in} \PYG{n+nb}{range}\PYG{p}{(}\PYG{l+m+mi}{0}\PYG{p}{,}\PYG{n}{np}\PYG{o}{.}\PYG{n}{size}\PYG{p}{(}\PYG{n}{B0}\PYG{p}{)}\PYG{p}{)}\PYG{p}{:}     
    \PYG{n}{c}\PYG{o}{=}\PYG{n+nb}{next}\PYG{p}{(}\PYG{n}{color}\PYG{p}{)}
    \PYG{c+c1}{\PYGZsh{}Langmuir isotherm}
    \PYG{n}{Gamma1} \PYG{o}{=} \PYG{n}{B0}\PYG{p}{[}\PYG{n}{i}\PYG{p}{]}\PYG{o}{*}\PYG{n}{C} \PYG{o}{/} \PYG{p}{(}\PYG{l+m+mi}{1} \PYG{o}{+} \PYG{n}{B0}\PYG{p}{[}\PYG{n}{i}\PYG{p}{]}\PYG{o}{*}\PYG{n}{C}\PYG{p}{)} 
    \PYG{n}{axes}\PYG{o}{.}\PYG{n}{plot}\PYG{p}{(}\PYG{n}{C}\PYG{p}{,}\PYG{n}{Gamma1}\PYG{p}{,} \PYG{n}{marker}\PYG{o}{=}\PYG{l+s+s1}{\PYGZsq{}}\PYG{l+s+s1}{ }\PYG{l+s+s1}{\PYGZsq{}} \PYG{p}{,} \PYG{n}{c}\PYG{o}{=}\PYG{n}{c}\PYG{p}{)}

    
\PYG{n}{plt}\PYG{o}{.}\PYG{n}{title}\PYG{p}{(}\PYG{l+s+s1}{\PYGZsq{}}\PYG{l+s+s1}{Langmuir Isotherm}\PYG{l+s+s1}{\PYGZsq{}}\PYG{p}{,} \PYG{n}{fontsize}\PYG{o}{=}\PYG{l+m+mi}{18}\PYG{p}{)}\PYG{p}{;}
\PYG{n}{axes}\PYG{o}{.}\PYG{n}{set\PYGZus{}xlabel}\PYG{p}{(}\PYG{l+s+s1}{\PYGZsq{}}\PYG{l+s+s1}{C [mol l\PYGZdl{}\PYGZca{}}\PYG{l+s+s1}{\PYGZob{}}\PYG{l+s+s1}{\PYGZhy{}1\PYGZcb{}\PYGZdl{}]}\PYG{l+s+s1}{\PYGZsq{}}\PYG{p}{,} \PYG{n}{fontsize}\PYG{o}{=}\PYG{l+m+mi}{14}\PYG{p}{)}\PYG{p}{;}
\PYG{n}{axes}\PYG{o}{.}\PYG{n}{set\PYGZus{}ylabel}\PYG{p}{(}\PYG{l+s+s1}{\PYGZsq{}}\PYG{l+s+s1}{\PYGZdl{}}\PYG{l+s+s1}{\PYGZbs{}}\PYG{l+s+s1}{Gamma\PYGZus{}1\PYGZdl{}}\PYG{l+s+s1}{\PYGZsq{}}\PYG{p}{,}\PYG{n}{fontsize}\PYG{o}{=}\PYG{l+m+mi}{18}\PYG{p}{)}\PYG{p}{;}
\end{sphinxVerbatim}

\end{sphinxuseclass}\end{sphinxVerbatimInput}
\begin{sphinxVerbatimOutput}

\begin{sphinxuseclass}{cell_output}
\begin{sphinxVerbatim}[commandchars=\\\{\}]
\PYGZlt{}\PYGZgt{}:29: SyntaxWarning: invalid escape sequence \PYGZsq{}\PYGZbs{}G\PYGZsq{}
\PYGZlt{}\PYGZgt{}:29: SyntaxWarning: invalid escape sequence \PYGZsq{}\PYGZbs{}G\PYGZsq{}
/var/folders/4v/jxb311rn5nq2q59y91zzgxd40000gp/T/ipykernel\PYGZus{}85059/3365319449.py:29: SyntaxWarning: invalid escape sequence \PYGZsq{}\PYGZbs{}G\PYGZsq{}
  axes.set\PYGZus{}ylabel(\PYGZsq{}\PYGZdl{}\PYGZbs{}Gamma\PYGZus{}1\PYGZdl{}\PYGZsq{},fontsize=18);
/var/folders/4v/jxb311rn5nq2q59y91zzgxd40000gp/T/ipykernel\PYGZus{}85059/3365319449.py:29: SyntaxWarning: invalid escape sequence \PYGZsq{}\PYGZbs{}G\PYGZsq{}
  axes.set\PYGZus{}ylabel(\PYGZsq{}\PYGZdl{}\PYGZbs{}Gamma\PYGZus{}1\PYGZdl{}\PYGZsq{},fontsize=18);
\end{sphinxVerbatim}

\begin{sphinxVerbatim}[commandchars=\\\{\}]
\PYG{g+gt}{\PYGZhy{}\PYGZhy{}\PYGZhy{}\PYGZhy{}\PYGZhy{}\PYGZhy{}\PYGZhy{}\PYGZhy{}\PYGZhy{}\PYGZhy{}\PYGZhy{}\PYGZhy{}\PYGZhy{}\PYGZhy{}\PYGZhy{}\PYGZhy{}\PYGZhy{}\PYGZhy{}\PYGZhy{}\PYGZhy{}\PYGZhy{}\PYGZhy{}\PYGZhy{}\PYGZhy{}\PYGZhy{}\PYGZhy{}\PYGZhy{}\PYGZhy{}\PYGZhy{}\PYGZhy{}\PYGZhy{}\PYGZhy{}\PYGZhy{}\PYGZhy{}\PYGZhy{}\PYGZhy{}\PYGZhy{}\PYGZhy{}\PYGZhy{}\PYGZhy{}\PYGZhy{}\PYGZhy{}\PYGZhy{}\PYGZhy{}\PYGZhy{}\PYGZhy{}\PYGZhy{}\PYGZhy{}\PYGZhy{}\PYGZhy{}\PYGZhy{}\PYGZhy{}\PYGZhy{}\PYGZhy{}\PYGZhy{}\PYGZhy{}\PYGZhy{}\PYGZhy{}\PYGZhy{}\PYGZhy{}\PYGZhy{}\PYGZhy{}\PYGZhy{}\PYGZhy{}\PYGZhy{}\PYGZhy{}\PYGZhy{}\PYGZhy{}\PYGZhy{}\PYGZhy{}\PYGZhy{}\PYGZhy{}\PYGZhy{}\PYGZhy{}\PYGZhy{}}
\PYG{n+ne}{ModuleNotFoundError}\PYG{g+gWhitespace}{                       }Traceback (most recent call last)
\PYG{n}{Cell} \PYG{n}{In}\PYG{p}{[}\PYG{l+m+mi}{1}\PYG{p}{]}\PYG{p}{,} \PYG{n}{line} \PYG{l+m+mi}{1}
\PYG{n+ne}{\PYGZhy{}\PYGZhy{}\PYGZhy{}\PYGZhy{}\PYGZgt{} }\PYG{l+m+mi}{1} \PYG{k+kn}{import}\PYG{+w}{ }\PYG{n+nn}{numpy}\PYG{+w}{ }\PYG{k}{as}\PYG{+w}{ }\PYG{n+nn}{np}
\PYG{g+gWhitespace}{      }\PYG{l+m+mi}{2} \PYG{k+kn}{import}\PYG{+w}{ }\PYG{n+nn}{matplotlib}\PYG{n+nn}{.}\PYG{n+nn}{pyplot}\PYG{+w}{ }\PYG{k}{as}\PYG{+w}{ }\PYG{n+nn}{plt} 
\PYG{g+gWhitespace}{      }\PYG{l+m+mi}{3} \PYG{k+kn}{from}\PYG{+w}{ }\PYG{n+nn}{matplotlib}\PYG{n+nn}{.}\PYG{n+nn}{pyplot}\PYG{+w}{ }\PYG{k+kn}{import} \PYG{n}{cm}

\PYG{n+ne}{ModuleNotFoundError}: No module named \PYGZsq{}numpy\PYGZsq{}
\end{sphinxVerbatim}

\end{sphinxuseclass}\end{sphinxVerbatimOutput}

\end{sphinxuseclass}

\subsection{5.2 Concentration in the adsorbed phase.}
\label{\detokenize{Langmuir_isotherm:concentration-in-the-adsorbed-phase}}
\sphinxAtStartPar
The Langmuir isotherm as discussed so far provides a relationship between the fractional occupation of adsorption sites \(\Gamma_1\) and the concentration in the fluid phase \(C\). This relationship is only dependent on the thermodynamics of adsorption, captured by the value of the adsorption equilibrium coinstant \(B_0\). In order to compute the concentration adsorbed by a specific material we need one additional parameter, the maximum concentration \(q_M\) in the adsorbed phase.
The concentration in the adsorbed phase is thus computed as: \(q=\Gamma_1{q_M}\), hence:
\begin{equation}\label{equation:Langmuir_isotherm:LANGeq7}
\begin{split}
q=q_M\frac{B_0{C}}{1+B_0{C}}
\end{split}
\end{equation}
\sphinxAtStartPar
The parameter \(q_M\) depends on the adsorbent material, and in particular on the amount of surface area per unit volume.

\sphinxAtStartPar
In the following we report the concentration in the adsorbed phase, in equilibrium with an arbitrary concetration \(C'\) in the fluid phase as a function of the affinity towards the stationary phase (captured by parameter \(B_0\)) and of its surface area (captured by parameter \(q_M\)).

\begin{sphinxuseclass}{cell}\begin{sphinxVerbatimInput}

\begin{sphinxuseclass}{cell_input}
\begin{sphinxVerbatim}[commandchars=\\\{\}]
\PYG{k+kn}{import}\PYG{+w}{ }\PYG{n+nn}{numpy}\PYG{+w}{ }\PYG{k}{as}\PYG{+w}{ }\PYG{n+nn}{np}
\PYG{k+kn}{import}\PYG{+w}{ }\PYG{n+nn}{pandas}\PYG{+w}{ }\PYG{k}{as}\PYG{+w}{ }\PYG{n+nn}{pd}
\PYG{k+kn}{from}\PYG{+w}{ }\PYG{n+nn}{matplotlib}\PYG{+w}{ }\PYG{k+kn}{import} \PYG{n}{cm}
\PYG{k+kn}{import}\PYG{+w}{ }\PYG{n+nn}{matplotlib}\PYG{n+nn}{.}\PYG{n+nn}{pyplot}\PYG{+w}{ }\PYG{k}{as}\PYG{+w}{ }\PYG{n+nn}{plt} 
\PYG{k+kn}{from}\PYG{+w}{ }\PYG{n+nn}{mpl\PYGZus{}toolkits}\PYG{n+nn}{.}\PYG{n+nn}{mplot3d}\PYG{+w}{ }\PYG{k+kn}{import} \PYG{n}{Axes3D} 

\PYG{c+c1}{\PYGZsh{} Parameters: }
\PYG{c+c1}{\PYGZsh{} molar fraction in the feed}
\PYG{n}{N} \PYG{o}{=} \PYG{l+m+mi}{100} \PYG{c+c1}{\PYGZsh{}number of points}
\PYG{n}{B0} \PYG{o}{=} \PYG{n}{np}\PYG{o}{.}\PYG{n}{logspace}\PYG{p}{(}\PYG{o}{\PYGZhy{}}\PYG{l+m+mi}{3}\PYG{p}{,} \PYG{l+m+mf}{1.5}\PYG{p}{,} \PYG{n}{N}\PYG{p}{)}
\PYG{n}{q\PYGZus{}M} \PYG{o}{=} \PYG{n}{np}\PYG{o}{.}\PYG{n}{logspace}\PYG{p}{(}\PYG{l+m+mi}{1}\PYG{p}{,} \PYG{l+m+mi}{3}\PYG{p}{,} \PYG{n}{N}\PYG{p}{)}
\PYG{n}{C} \PYG{o}{=} \PYG{l+m+mi}{1} \PYG{p}{;} \PYG{c+c1}{\PYGZsh{} arbitrary value of C\PYGZsq{}}
\PYG{n}{B0}\PYG{p}{,} \PYG{n}{q\PYGZus{}M} \PYG{o}{=} \PYG{n}{np}\PYG{o}{.}\PYG{n}{meshgrid}\PYG{p}{(}\PYG{n}{B0}\PYG{p}{,} \PYG{n}{q\PYGZus{}M}\PYG{p}{)} 

\PYG{n}{q}\PYG{o}{=}\PYG{n}{q\PYGZus{}M} \PYG{o}{*} \PYG{n}{B0} \PYG{o}{*} \PYG{n}{C} \PYG{o}{/} \PYG{p}{(}\PYG{l+m+mi}{1}\PYG{o}{+}\PYG{n}{B0} \PYG{o}{*} \PYG{n}{C}\PYG{p}{)}


\PYG{c+c1}{\PYGZsh{}Plotting}

\PYG{n}{figure}\PYG{o}{=}\PYG{n}{plt}\PYG{o}{.}\PYG{n}{figure}\PYG{p}{(}\PYG{n}{figsize}\PYG{o}{=}\PYG{p}{(}\PYG{l+m+mi}{10}\PYG{p}{,} \PYG{l+m+mi}{8}\PYG{p}{)}\PYG{p}{)}
\PYG{n}{axes} \PYG{o}{=} \PYG{n}{figure}\PYG{o}{.}\PYG{n}{gca}\PYG{p}{(}\PYG{n}{projection} \PYG{o}{=}\PYG{l+s+s1}{\PYGZsq{}}\PYG{l+s+s1}{3d}\PYG{l+s+s1}{\PYGZsq{}}\PYG{p}{)} 

\PYG{n}{plt}\PYG{o}{.}\PYG{n}{xticks}\PYG{p}{(}\PYG{n}{fontsize}\PYG{o}{=}\PYG{l+m+mi}{14}\PYG{p}{)}
\PYG{n}{plt}\PYG{o}{.}\PYG{n}{yticks}\PYG{p}{(}\PYG{n}{fontsize}\PYG{o}{=}\PYG{l+m+mi}{14}\PYG{p}{)}


\PYG{n}{surf}\PYG{o}{=}\PYG{n}{axes}\PYG{o}{.}\PYG{n}{plot\PYGZus{}surface}\PYG{p}{(}\PYG{n}{q\PYGZus{}M}\PYG{p}{,}\PYG{n}{B0}\PYG{p}{,}\PYG{n}{q}\PYG{p}{,}\PYG{n}{cmap}\PYG{o}{=}\PYG{n}{cm}\PYG{o}{.}\PYG{n}{coolwarm}\PYG{p}{,}
                       \PYG{n}{linewidth}\PYG{o}{=}\PYG{l+m+mi}{0}\PYG{p}{,} \PYG{n}{antialiased}\PYG{o}{=}\PYG{k+kc}{False}\PYG{p}{,}\PYG{n}{alpha}\PYG{o}{=}\PYG{l+m+mf}{0.2}\PYG{p}{)}

\PYG{n}{contour}\PYG{o}{=}\PYG{n}{axes}\PYG{o}{.}\PYG{n}{contour}\PYG{p}{(}\PYG{n}{q\PYGZus{}M}\PYG{p}{,}\PYG{n}{B0}\PYG{p}{,}\PYG{n}{q}\PYG{p}{,}\PYG{n}{np}\PYG{o}{.}\PYG{n}{linspace}\PYG{p}{(}\PYG{l+m+mi}{100}\PYG{p}{,} \PYG{l+m+mi}{800}\PYG{p}{,} \PYG{l+m+mi}{10}\PYG{p}{)}\PYG{p}{,}\PYG{n}{cmap}\PYG{o}{=}\PYG{n}{cm}\PYG{o}{.}\PYG{n}{coolwarm}\PYG{p}{)}

\PYG{n}{plt}\PYG{o}{.}\PYG{n}{title}\PYG{p}{(}\PYG{l+s+s1}{\PYGZsq{}}\PYG{l+s+s1}{Adsorbed concentration vs \PYGZdl{}B\PYGZus{}0\PYGZdl{} / \PYGZdl{}q\PYGZus{}M\PYGZdl{}}\PYG{l+s+s1}{\PYGZsq{}}\PYG{p}{,} \PYG{n}{fontsize}\PYG{o}{=}\PYG{l+m+mi}{18}\PYG{p}{)}\PYG{p}{;}
\PYG{n}{axes}\PYG{o}{.}\PYG{n}{set\PYGZus{}xlabel}\PYG{p}{(}\PYG{l+s+s1}{\PYGZsq{}}\PYG{l+s+s1}{\PYGZdl{}q\PYGZus{}M\PYGZdl{}}\PYG{l+s+s1}{\PYGZsq{}}\PYG{p}{,} \PYG{n}{fontsize}\PYG{o}{=}\PYG{l+m+mi}{14}\PYG{p}{)}\PYG{p}{;}
\PYG{n}{axes}\PYG{o}{.}\PYG{n}{set\PYGZus{}ylabel}\PYG{p}{(}\PYG{l+s+s1}{\PYGZsq{}}\PYG{l+s+s1}{\PYGZdl{}B\PYGZus{}0\PYGZdl{}}\PYG{l+s+s1}{\PYGZsq{}}\PYG{p}{,}\PYG{n}{fontsize}\PYG{o}{=}\PYG{l+m+mi}{14}\PYG{p}{)}\PYG{p}{;}
\PYG{n}{axes}\PYG{o}{.}\PYG{n}{set\PYGZus{}zlabel}\PYG{p}{(}\PYG{l+s+s1}{\PYGZsq{}}\PYG{l+s+s1}{\PYGZdl{}q\PYGZdl{}}\PYG{l+s+s1}{\PYGZsq{}}\PYG{p}{,}\PYG{n}{fontsize}\PYG{o}{=}\PYG{l+m+mi}{14}\PYG{p}{)}\PYG{p}{;}
\PYG{n}{figure}\PYG{o}{.}\PYG{n}{colorbar}\PYG{p}{(}\PYG{n}{surf}\PYG{p}{,} \PYG{n}{shrink}\PYG{o}{=}\PYG{l+m+mf}{0.5}\PYG{p}{,} \PYG{n}{aspect}\PYG{o}{=}\PYG{l+m+mi}{5}\PYG{p}{)}\PYG{p}{;}
\end{sphinxVerbatim}

\end{sphinxuseclass}\end{sphinxVerbatimInput}
\begin{sphinxVerbatimOutput}

\begin{sphinxuseclass}{cell_output}
\noindent\sphinxincludegraphics{{6359911cf974e034b0410bca1b89e8e1d2e4f3ec60daf819fd9839cc5cd39759}.png}

\end{sphinxuseclass}\end{sphinxVerbatimOutput}

\end{sphinxuseclass}

\subsection{5.3 Langmuir isotherm in the gas phase.}
\label{\detokenize{Langmuir_isotherm:langmuir-isotherm-in-the-gas-phase}}
\sphinxAtStartPar
When describing the thermodynamics of adsorption from the gas phase it is convenient to introduce explicitly the pressure \(P\) in the Langmuir isotherm.
In this case the expression of the monocomponent Langmuir isotherm becomes:
\begin{equation}\label{equation:Langmuir_isotherm:LANGeq8}
\begin{split}
q=q_M\frac{B_1{P}}{1+B_1{P}}
\end{split}
\end{equation}
\sphinxAtStartPar
where \(B_1=B_0/RT\) and \(P\) is pressure.


\subsection{5.4 Multicomponent, competitive adsorption}
\label{\detokenize{Langmuir_isotherm:multicomponent-competitive-adsorption}}
\sphinxAtStartPar
In multicomponent systems, in which \(N\) components compete for the same adsorption sites we can define the adsorption equilibrium constant of component \(i\) as:
\begin{equation}\label{equation:Langmuir_isotherm:LANGeq9}
\begin{split}
B_{1,i}=\frac{\Gamma_i}{P_i\Gamma_0}
\end{split}
\end{equation}
\sphinxAtStartPar
In a multicomponent system, the single layer hypothesis bocomes:
\begin{equation}\label{equation:Langmuir_isotherm:LANGeq10}
\begin{split}
\Gamma_0+\sum_{i=1}^N{\Gamma_i}=1 
\end{split}
\end{equation}
\sphinxAtStartPar
Substituting in Eq. \eqref{equation:Langmuir_isotherm:LANGeq10} the definition of adsorption equilibrium constant given in Eq. \eqref{equation:Langmuir_isotherm:LANGeq9} we get:
\begin{equation}\label{equation:Langmuir_isotherm:LANGeq11}
\begin{split}
\Gamma_0^{-1}=1+\sum_{i=1}^N{B_{1,i}P_i}
\end{split}
\end{equation}
\sphinxAtStartPar
finally, substituting again \(\Gamma_0^ {-1}\) in Eq. \eqref{equation:Langmuir_isotherm:LANGeq9}, and rearranging we get an expression of the fraction of sites occupied by specie \(i\) as a function of the partial pressure (concentration) of \(i\) in the fluid phase.
\begin{equation}\label{equation:Langmuir_isotherm:LANGeq12}
\begin{split}
\Gamma_i=\frac{B_{1,i}P_i}{1+\sum_{i=1}^NB_{1,i}{P_i}}
\end{split}
\end{equation}
\sphinxAtStartPar
Finally, inserting the density of specie \(i\) in the adsorbed phase we get:
\begin{equation}\label{equation:Langmuir_isotherm:LANGeq13}
\begin{split}
q_i=q_M\frac{B_{1,i}{P_i}}{1+\sum_{i=1}^NB_{1,i}{P_i}}
\end{split}
\end{equation}
\sphinxAtStartPar
where \(N\) is the total number of components, \(P_i\) is the partial pressure component \(i\) in the gas phase, and \(B_{1,i}=B_{0,i}/RT\) where \(B_{0,i}\) is the adsorption equilibrium constant for component \(i\).

\sphinxAtStartPar
In the following we report the competitive adsorption isotherms of two components adsorbing from a binary mixture.

\begin{sphinxuseclass}{cell}\begin{sphinxVerbatimInput}

\begin{sphinxuseclass}{cell_input}
\begin{sphinxVerbatim}[commandchars=\\\{\}]
\PYG{k+kn}{import}\PYG{+w}{ }\PYG{n+nn}{numpy}\PYG{+w}{ }\PYG{k}{as}\PYG{+w}{ }\PYG{n+nn}{np}
\PYG{k+kn}{import}\PYG{+w}{ }\PYG{n+nn}{matplotlib}\PYG{n+nn}{.}\PYG{n+nn}{pyplot}\PYG{+w}{ }\PYG{k}{as}\PYG{+w}{ }\PYG{n+nn}{plt} 
\PYG{k+kn}{from}\PYG{+w}{ }\PYG{n+nn}{matplotlib}\PYG{n+nn}{.}\PYG{n+nn}{pyplot}\PYG{+w}{ }\PYG{k+kn}{import} \PYG{n}{cm}

\PYG{c+c1}{\PYGZsh{}Plotting}
\PYG{n}{figure}\PYG{o}{=}\PYG{n}{plt}\PYG{o}{.}\PYG{n}{figure}\PYG{p}{(}\PYG{p}{)}
\PYG{n}{axes} \PYG{o}{=} \PYG{n}{figure}\PYG{o}{.}\PYG{n}{add\PYGZus{}axes}\PYG{p}{(}\PYG{p}{[}\PYG{l+m+mf}{0.1}\PYG{p}{,}\PYG{l+m+mf}{0.1}\PYG{p}{,}\PYG{l+m+mf}{1.2}\PYG{p}{,}\PYG{l+m+mf}{1.2}\PYG{p}{]}\PYG{p}{)}
\PYG{n}{plt}\PYG{o}{.}\PYG{n}{xticks}\PYG{p}{(}\PYG{n}{fontsize}\PYG{o}{=}\PYG{l+m+mi}{14}\PYG{p}{)}
\PYG{n}{plt}\PYG{o}{.}\PYG{n}{yticks}\PYG{p}{(}\PYG{n}{fontsize}\PYG{o}{=}\PYG{l+m+mi}{14}\PYG{p}{)}


\PYG{c+c1}{\PYGZsh{} molar fraction in the feed}
\PYG{n}{z} \PYG{o}{=} \PYG{l+m+mf}{0.3}
\PYG{n}{N} \PYG{o}{=} \PYG{l+m+mi}{500} \PYG{c+c1}{\PYGZsh{}number of points}
\PYG{n}{P} \PYG{o}{=} \PYG{n}{np}\PYG{o}{.}\PYG{n}{linspace}\PYG{p}{(}\PYG{l+m+mi}{0}\PYG{p}{,}\PYG{l+m+mi}{50}\PYG{p}{,} \PYG{n}{N}\PYG{p}{)}
\PYG{n}{y} \PYG{o}{=} \PYG{l+m+mf}{0.1} \PYG{c+c1}{\PYGZsh{}component 1 is in a smaller proportion}
\PYG{n}{B11} \PYG{o}{=} \PYG{l+m+mi}{10}\PYG{p}{;} \PYG{c+c1}{\PYGZsh{}component 1 affinity for the stationary phase is higher than that of component 2}
\PYG{n}{B12} \PYG{o}{=} \PYG{l+m+mi}{1}\PYG{p}{;}

\PYG{n}{color}\PYG{o}{=}\PYG{n+nb}{iter}\PYG{p}{(}\PYG{n}{cm}\PYG{o}{.}\PYG{n}{rainbow}\PYG{p}{(}\PYG{n}{np}\PYG{o}{.}\PYG{n}{linspace}\PYG{p}{(}\PYG{l+m+mi}{0}\PYG{p}{,}\PYG{l+m+mi}{1}\PYG{p}{,}\PYG{l+m+mi}{3}\PYG{p}{)}\PYG{p}{)}\PYG{p}{)}
\PYG{n}{c}\PYG{o}{=}\PYG{n+nb}{next}\PYG{p}{(}\PYG{n}{color}\PYG{p}{)}
\PYG{n}{Gamma11} \PYG{o}{=} \PYG{n}{B11}\PYG{o}{*}\PYG{n}{P}\PYG{o}{*}\PYG{n}{y} \PYG{o}{/} \PYG{p}{(}\PYG{l+m+mi}{1} \PYG{o}{+} \PYG{n}{B11}\PYG{o}{*}\PYG{n}{P}\PYG{o}{*}\PYG{n}{y} \PYG{o}{+} \PYG{n}{B12}\PYG{o}{*}\PYG{n}{P}\PYG{o}{*}\PYG{p}{(}\PYG{l+m+mi}{1}\PYG{o}{\PYGZhy{}}\PYG{n}{y}\PYG{p}{)}\PYG{p}{)} 
\PYG{n}{axes}\PYG{o}{.}\PYG{n}{plot}\PYG{p}{(}\PYG{n}{P}\PYG{p}{,}\PYG{n}{Gamma11}\PYG{p}{,} \PYG{n}{marker}\PYG{o}{=}\PYG{l+s+s1}{\PYGZsq{}}\PYG{l+s+s1}{ }\PYG{l+s+s1}{\PYGZsq{}} \PYG{p}{,} \PYG{n}{c}\PYG{o}{=}\PYG{n}{c}\PYG{p}{)}
\PYG{n}{c}\PYG{o}{=}\PYG{n+nb}{next}\PYG{p}{(}\PYG{n}{color}\PYG{p}{)}
\PYG{n}{Gamma12} \PYG{o}{=} \PYG{n}{B12}\PYG{o}{*}\PYG{n}{P}\PYG{o}{*}\PYG{n}{y} \PYG{o}{/} \PYG{p}{(}\PYG{l+m+mi}{1} \PYG{o}{+} \PYG{n}{B11}\PYG{o}{*}\PYG{n}{P}\PYG{o}{*}\PYG{n}{y} \PYG{o}{+} \PYG{n}{B12}\PYG{o}{*}\PYG{n}{P}\PYG{o}{*}\PYG{p}{(}\PYG{l+m+mi}{1}\PYG{o}{\PYGZhy{}}\PYG{n}{y}\PYG{p}{)}\PYG{p}{)} 
\PYG{n}{axes}\PYG{o}{.}\PYG{n}{plot}\PYG{p}{(}\PYG{n}{P}\PYG{p}{,}\PYG{n}{Gamma12}\PYG{p}{,} \PYG{n}{marker}\PYG{o}{=}\PYG{l+s+s1}{\PYGZsq{}}\PYG{l+s+s1}{ }\PYG{l+s+s1}{\PYGZsq{}} \PYG{p}{,} \PYG{n}{c}\PYG{o}{=}\PYG{n}{c}\PYG{p}{)}
\PYG{n}{c}\PYG{o}{=}\PYG{n+nb}{next}\PYG{p}{(}\PYG{n}{color}\PYG{p}{)}
\PYG{n}{Gamma0} \PYG{o}{=} \PYG{l+m+mi}{1} \PYG{o}{\PYGZhy{}}  \PYG{n}{Gamma11} \PYG{o}{\PYGZhy{}} \PYG{n}{Gamma12}
\PYG{n}{axes}\PYG{o}{.}\PYG{n}{plot}\PYG{p}{(}\PYG{n}{P}\PYG{p}{,}\PYG{n}{Gamma0}\PYG{p}{,} \PYG{n}{marker}\PYG{o}{=}\PYG{l+s+s1}{\PYGZsq{}}\PYG{l+s+s1}{ }\PYG{l+s+s1}{\PYGZsq{}} \PYG{p}{,} \PYG{n}{c}\PYG{o}{=}\PYG{n}{c}\PYG{p}{)}

\PYG{n}{axes}\PYG{o}{.}\PYG{n}{set\PYGZus{}xlabel}\PYG{p}{(}\PYG{l+s+s1}{\PYGZsq{}}\PYG{l+s+s1}{P [bar]}\PYG{l+s+s1}{\PYGZsq{}}\PYG{p}{,} \PYG{n}{fontsize}\PYG{o}{=}\PYG{l+m+mi}{14}\PYG{p}{)}\PYG{p}{;}
\PYG{n}{axes}\PYG{o}{.}\PYG{n}{set\PYGZus{}ylabel}\PYG{p}{(}\PYG{l+s+s1}{\PYGZsq{}}\PYG{l+s+s1}{\PYGZdl{}}\PYG{l+s+s1}{\PYGZbs{}}\PYG{l+s+s1}{Gamma\PYGZus{}0 / }\PYG{l+s+s1}{\PYGZbs{}}\PYG{l+s+s1}{Gamma\PYGZus{}}\PYG{l+s+s1}{\PYGZob{}}\PYG{l+s+s1}{1,i\PYGZcb{}\PYGZdl{}}\PYG{l+s+s1}{\PYGZsq{}}\PYG{p}{,} \PYG{n}{fontsize}\PYG{o}{=}\PYG{l+m+mi}{18}\PYG{p}{)}\PYG{p}{;}
\PYG{n}{axes}\PYG{o}{.}\PYG{n}{legend}\PYG{p}{(}\PYG{p}{[}\PYG{l+s+s1}{\PYGZsq{}}\PYG{l+s+s1}{\PYGZdl{}}\PYG{l+s+s1}{\PYGZbs{}}\PYG{l+s+s1}{Gamma\PYGZus{}}\PYG{l+s+s1}{\PYGZob{}}\PYG{l+s+s1}{1,1\PYGZcb{}\PYGZdl{}}\PYG{l+s+s1}{\PYGZsq{}}\PYG{p}{,}\PYG{l+s+s1}{\PYGZsq{}}\PYG{l+s+s1}{\PYGZdl{}}\PYG{l+s+s1}{\PYGZbs{}}\PYG{l+s+s1}{Gamma\PYGZus{}}\PYG{l+s+s1}{\PYGZob{}}\PYG{l+s+s1}{1,2\PYGZcb{}\PYGZdl{}}\PYG{l+s+s1}{\PYGZsq{}}\PYG{p}{,}\PYG{l+s+s1}{\PYGZsq{}}\PYG{l+s+s1}{\PYGZdl{}}\PYG{l+s+s1}{\PYGZbs{}}\PYG{l+s+s1}{Gamma\PYGZus{}}\PYG{l+s+si}{\PYGZob{}0\PYGZcb{}}\PYG{l+s+s1}{\PYGZdl{}}\PYG{l+s+s1}{\PYGZsq{}}\PYG{p}{]}\PYG{p}{,} \PYG{n}{fontsize}\PYG{o}{=}\PYG{l+m+mi}{18}\PYG{p}{)}\PYG{p}{;}
\end{sphinxVerbatim}

\end{sphinxuseclass}\end{sphinxVerbatimInput}
\begin{sphinxVerbatimOutput}

\begin{sphinxuseclass}{cell_output}
\noindent\sphinxincludegraphics{{43e198fa0121b9df31be02d5530b4c510023f9d82d64a27af55536e00bfc7fe9}.png}

\end{sphinxuseclass}\end{sphinxVerbatimOutput}

\end{sphinxuseclass}
\sphinxstepscope


\section{6. The BET Isotherm}
\label{\detokenize{BET_isotherm:the-bet-isotherm}}\label{\detokenize{BET_isotherm::doc}}
\sphinxAtStartPar
The BET (Brunauer, Emmet and Teller) isotherm allows to extend ideas and hypotheses that are adopted to obtain the Langmuir isotherm to describe the adsorption of a component from a gas phase in cases in which it can form multiple layers on the surface of the stationary phase, up to the limit in which it condenses into a macroscopic liquid phase.

\sphinxAtStartPar
The hypotheses at the foundation of the BET theory are:
\begin{itemize}
\item {} 
\sphinxAtStartPar
Lateral interactions between adsorbed molecules (i.e. interactions between molecules adsorbed in the same layer) are neglibile. In other words the adsorption process does not depend on the local environment around an adsorption site. This hypothesis is common to the derivation of the Langmuir isotherm.

\item {} 
\sphinxAtStartPar
Molecules can adsorb in an infinite number of layers.

\item {} 
\sphinxAtStartPar
The adsorption enthalpy on the first layer is constant and independent from the fractional occupation of that layer.

\item {} 
\sphinxAtStartPar
The adsorption enthalpy in layers from the second onwards is constant and corresponds to the condensation enthalpy.

\end{itemize}

\sphinxAtStartPar
Following an approach similar to the one adopted for the derivation of the Langmuir isotherm let us begin by defining the rates of adsorption and desorption. In the case of th BET isotherm the rates will depend on the layer considered and will therefore have subscript \(i\) that indicate the layer on which adsorption takes place.

\sphinxAtStartPar
The adsorption rate on layer \(i\) can be written as:
\begin{equation}\label{equation:BET_isotherm:BETeq1}
\begin{split}
R_{A,i}=k_{A,i}\Gamma_iP
\end{split}
\end{equation}
\sphinxAtStartPar
The desorption rate instead, explicitly introducing the dependence on the adsorption enthalpy through an Arrhenius\sphinxhyphen{}like functional form, is written as:
\begin{equation}\label{equation:BET_isotherm:BETeq2}
\begin{split}
R_{D,i}=k_{D,i}\Gamma_i=A_ie^{\frac{-E_i}{RT}}\Gamma_i
\end{split}
\end{equation}
\sphinxAtStartPar
where \(E_i\) is the adsorption energy. Denoting with \(E_0\) the adsorption energy on a free adsorption site, the adsorption energy on all successive layers is \(E_2=E_3=...E_n=\lambda\), where \(\lambda\) is the condensation enthalpy.

\sphinxAtStartPar
Following the same criterion also the adsorption constant \(k_{A,i}\), and the pre\sphinxhyphen{}exponential factor of the desorption constant \(A_i\) are independent from the layer index \sphinxhyphen{} except for those related to the adsorption on a free surface and the desorption from the first layer:
\begin{equation}\label{equation:BET_isotherm:BETeq3}
\begin{split}
A_{i\geq2}=const=A
\end{split}
\end{equation}\begin{equation}\label{equation:BET_isotherm:BETeq4}
\begin{split}
k_{A,i\geq1}=const=k_A
\end{split}
\end{equation}
\sphinxAtStartPar
Having defined the rates of adsorption and desorption one can write population balances for every layer of the adsorbed phase.
In the following \(\Gamma_0\) represents the fraction of free adsorption sites, \(\Gamma_i\) represents the fractional occupation of the \(i^{th}\) layer.
\begin{equation}\label{equation:BET_isotherm:BETeq5}
\begin{split}
\begin{cases}
 \frac{d\Gamma_0}{dt}&=-k_{A,0}\Gamma_0P +k_{D,1}\Gamma_1 \\
 \frac{d\Gamma_1}{dt}&=+k_{A,0}\Gamma_0P -k_{D,1}\Gamma_1-k_{A,1}\Gamma_1P+k_{D,2}\Gamma_2 \\
 \frac{d\Gamma_2}{dt}&=+k_{A,1}\Gamma_1P -k_{D,2}\Gamma_2-k_{A,2}\Gamma_2P+k_{D,3}\Gamma_3 \\
 \vdots \\
 \frac{d\Gamma_n}{dt}&=+k_{A,n-1}\Gamma_{n-1}P -k_{D,n}\Gamma_n-k_{A,n}\Gamma_nP+k_{D,n+1}\Gamma_{n+1} \\
 \vdots\\
 \end{cases}
\end{split}
\end{equation}
\sphinxAtStartPar
At equilibrium the balances reduce to:
\begin{equation}\label{equation:BET_isotherm:BETeq6}
\begin{split}
\begin{cases}
 0=-k_{A,0}\Gamma_0P +k_{D,1}\Gamma_1 \\
 0=+k_{A,0}\Gamma_0P -k_{D,1}\Gamma_1-k_{A,1}\Gamma_1P+k_{D,2}\Gamma_2 \\
 0=+k_{A,1}\Gamma_1P -k_{D,2}\Gamma_2-k_{A,2}\Gamma_2P+k_{D,3}\Gamma_3 \\
 \vdots \\
 0=+k_{A,n-1}\Gamma_{n-1}P -k_{D,n}\Gamma_n-k_{A,n}\Gamma_nP+k_{D,n+1}\Gamma_{n+1} \\
 \vdots\\
 \end{cases}
\end{split}
\end{equation}
\sphinxAtStartPar
From the steady state balance on free sites we get:
\begin{equation}\label{equation:BET_isotherm:BETeq7}
\begin{split}
k_{A,0}\Gamma_0P=k_{D,1}\Gamma_1
\end{split}
\end{equation}
\sphinxAtStartPar
hence, inserting this equality into the balance on the fraction of sites occupied by a single layer, \(\Gamma_1\) we get:
\begin{equation}\label{equation:BET_isotherm:BETeq8}
\begin{split}
k_{A,1}\Gamma_1P=k_{D,2}\Gamma_2
\end{split}
\end{equation}
\sphinxAtStartPar
Introducing the Arrhenius expression for the adsorption rates we get:
\begin{equation}\label{equation:BET_isotherm:BETeq9}
\begin{split}
\Gamma_1=\frac{k_{A,0}}{k_{D,1}}\Gamma_0P=\frac{k_{A,0}}{A_{1}}e^{\frac{E_1}{RT}}P\Gamma_0
\end{split}
\end{equation}
\sphinxAtStartPar
and analogously for \(\Gamma_2\):
\begin{equation}\label{equation:BET_isotherm:BETeq10}
\begin{split}
\Gamma_2=\frac{k_{A,1}}{k_{D,2}}\Gamma_1P=\frac{k_{A,1}}{A_{2}}e^{\frac{E_2}{RT}}P\Gamma_1
\end{split}
\end{equation}
\sphinxAtStartPar
a similar procedure can now be followed recursively to obtain a compact expression for the fraction of sites occupied by an arbitrary number of layers \(n\). To this aim we can introduce two constants, \(\alpha\) and \(\beta\) \sphinxhyphen{} defined as:
\begin{equation}\label{equation:BET_isotherm:BETeq11}
\begin{split}
\alpha=\frac{k_{A,0}}{A_{1}}e^{\frac{E_1}{RT}}P
\end{split}
\end{equation}\begin{equation}\label{equation:BET_isotherm:BETeq12}
\begin{split}
\beta=\frac{k_{A}}{A}e^{\frac{\lambda}{RT}}P
\end{split}
\end{equation}
\sphinxAtStartPar
At this point the fraction of sites occupied by \(n\) layers can be obtained as:
\begin{equation}\label{equation:BET_isotherm:BETeq13}
\begin{split}
\Gamma_n=\beta\Gamma_{n-1}=\beta^2\Gamma_{n-2}=...=\beta^{n-1}\Gamma_{1}=\beta^n\frac{\alpha}{\beta}\Gamma_0
\end{split}
\end{equation}
\sphinxAtStartPar
Deining for the sake of compactness the constant \(B_2=\alpha/\beta\) and introducing the volume of adsorbate per unit area  \(v_s^1\) one can express the total volume adsorbed as:
\begin{equation}\label{equation:BET_isotherm:BETeq14}
\begin{split}
v_s=v_s^1\sum_{i=1}^n{iB_2\beta^i\Gamma_0}=v_s^1B_2\beta\sum_{i=1}^ni\beta^{i-1}=v_s^1B_2\beta\frac{d\sum_{i=1}^n\beta^{i}}{d\beta}
\end{split}
\end{equation}
\sphinxAtStartPar
By introducing the notable result for the term \(\sum_{i=1}^n\beta^{i}\):
\begin{equation}\label{equation:BET_isotherm:BETeq15}
\begin{split}
\sum_{i=1}^n\beta^{i}=\left(\frac{1-\beta^n}{1-\beta}\right)\beta
\end{split}
\end{equation}
\sphinxAtStartPar
one gets:
\begin{equation}\label{equation:BET_isotherm:BETeq16}
\begin{split}
\frac{v_s}{v_s^1}=\frac{B_2\beta}{1-\beta}\frac{\left[1-(n+1)\beta^n+n\beta^{n+1}\right]}{\left[1+(B_2-1)\beta-B_2\beta^{n+1}\right]}
\label{eq:16}
\end{split}
\end{equation}
\sphinxAtStartPar
In order to obtain the final formulation of the BET isotherm it is useful to consider the limit corresponding to macroscopic condensation, which corresponds to a buildup of an infinite number of adsorbed layers, which corresponds to \(\beta=1\). Since this condition is realised for \(P=P^o\), it can be shown that \(\beta=\frac{k_{A}}{A}e^{\frac{\lambda}{RT}}P=P/P^o\).

\sphinxAtStartPar
Introducing this equality in \sphinxcode{\sphinxupquote{eq16}} one gets the so called limited form of the BET isotherm:
\begin{equation}\label{equation:BET_isotherm:BETeq17}
\begin{split}
\frac{v_s}{v_s^1}=\frac{B_2\left(\frac{P}{P^o}\right)}{1-\left(\frac{P}{P^o}\right)}\frac{\left[1-(n+1)\left(\frac{P}{P^o}\right)^n+n\left(\frac{P}{P^o}\right)^{n+1}\right]}{\left[1+(B_2-1)\left(\frac{P}{P^o}\right)-B_2\left(\frac{P}{P^o}\right)^{n+1}\right]}
\end{split}
\end{equation}
\sphinxAtStartPar
note that the left hand side of this equation can be expressed in terms of unit mass of adsorbent, instead of unit surface.


\subsection{6.1 Many isotherm shapes from the same expression}
\label{\detokenize{BET_isotherm:many-isotherm-shapes-from-the-same-expression}}
\sphinxAtStartPar
Depending on the parameters introduced in the BET isotherm its functional form can adapt to capture all known shapes of adsorption isotherm. In the following we shall demonstrate combination of parameters that produce isotherms belonging to all five of the classes introduced by Brunauer (Brunauer et al. JACS 1940).

\begin{sphinxuseclass}{cell}\begin{sphinxVerbatimInput}

\begin{sphinxuseclass}{cell_input}
\begin{sphinxVerbatim}[commandchars=\\\{\}]
\PYG{k+kn}{import}\PYG{+w}{ }\PYG{n+nn}{numpy}\PYG{+w}{ }\PYG{k}{as}\PYG{+w}{ }\PYG{n+nn}{np}
\PYG{k+kn}{import}\PYG{+w}{ }\PYG{n+nn}{matplotlib}\PYG{n+nn}{.}\PYG{n+nn}{pyplot}\PYG{+w}{ }\PYG{k}{as}\PYG{+w}{ }\PYG{n+nn}{plt} 
\PYG{k+kn}{from}\PYG{+w}{ }\PYG{n+nn}{matplotlib}\PYG{n+nn}{.}\PYG{n+nn}{pyplot}\PYG{+w}{ }\PYG{k+kn}{import} \PYG{n}{cm}

\PYG{c+c1}{\PYGZsh{}Plotting}
\PYG{n}{figure}\PYG{o}{=}\PYG{n}{plt}\PYG{o}{.}\PYG{n}{figure}\PYG{p}{(}\PYG{p}{)}
\PYG{n}{axes} \PYG{o}{=} \PYG{n}{figure}\PYG{o}{.}\PYG{n}{add\PYGZus{}axes}\PYG{p}{(}\PYG{p}{[}\PYG{l+m+mf}{0.1}\PYG{p}{,}\PYG{l+m+mf}{0.1}\PYG{p}{,}\PYG{l+m+mf}{1.2}\PYG{p}{,}\PYG{l+m+mf}{1.2}\PYG{p}{]}\PYG{p}{)}
\PYG{n}{plt}\PYG{o}{.}\PYG{n}{xticks}\PYG{p}{(}\PYG{n}{fontsize}\PYG{o}{=}\PYG{l+m+mi}{14}\PYG{p}{)}
\PYG{n}{plt}\PYG{o}{.}\PYG{n}{yticks}\PYG{p}{(}\PYG{n}{fontsize}\PYG{o}{=}\PYG{l+m+mi}{14}\PYG{p}{)}

\PYG{n}{N}\PYG{o}{=}\PYG{l+m+mi}{100}\PYG{p}{;}

\PYG{c+c1}{\PYGZsh{}Parameters}
\PYG{n}{P0} \PYG{o}{=} \PYG{n}{np}\PYG{o}{.}\PYG{n}{array}\PYG{p}{(}\PYG{p}{[}\PYG{l+m+mi}{5}\PYG{p}{,} \PYG{l+m+mi}{12}\PYG{p}{,} \PYG{l+m+mi}{12}\PYG{p}{,} \PYG{l+m+mi}{5}\PYG{p}{,} \PYG{l+m+mi}{11}\PYG{p}{,} \PYG{l+m+mi}{11}\PYG{p}{]}\PYG{p}{)}\PYG{p}{;}
\PYG{n}{B2} \PYG{o}{=} \PYG{n}{np}\PYG{o}{.}\PYG{n}{array}\PYG{p}{(}\PYG{p}{[}\PYG{l+m+mi}{10}\PYG{p}{,} \PYG{l+m+mi}{50}\PYG{p}{,} \PYG{l+m+mi}{2}\PYG{p}{,} \PYG{l+m+mi}{50}\PYG{p}{,} \PYG{l+m+mi}{2}\PYG{p}{]}\PYG{p}{)}\PYG{p}{;}
\PYG{n}{n}  \PYG{o}{=} \PYG{n}{np}\PYG{o}{.}\PYG{n}{array}\PYG{p}{(}\PYG{p}{[}\PYG{l+m+mi}{1}\PYG{p}{,}\PYG{l+m+mi}{500}\PYG{p}{,} \PYG{l+m+mi}{500}\PYG{p}{,} \PYG{l+m+mi}{5}\PYG{p}{,} \PYG{l+m+mi}{5}\PYG{p}{]}\PYG{p}{)}\PYG{p}{;}

\PYG{n}{P} \PYG{o}{=} \PYG{n}{np}\PYG{o}{.}\PYG{n}{linspace}\PYG{p}{(}\PYG{l+m+mi}{0}\PYG{p}{,} \PYG{l+m+mi}{10}\PYG{p}{,} \PYG{n}{N}\PYG{p}{)}
\PYG{n}{vv1} \PYG{o}{=} \PYG{n}{np}\PYG{o}{.}\PYG{n}{zeros}\PYG{p}{(}\PYG{n}{np}\PYG{o}{.}\PYG{n}{size}\PYG{p}{(}\PYG{n}{P}\PYG{p}{)}\PYG{p}{)}
   


\PYG{n}{color}\PYG{o}{=}\PYG{n+nb}{iter}\PYG{p}{(}\PYG{n}{cm}\PYG{o}{.}\PYG{n}{rainbow}\PYG{p}{(}\PYG{n}{np}\PYG{o}{.}\PYG{n}{linspace}\PYG{p}{(}\PYG{l+m+mi}{0}\PYG{p}{,}\PYG{l+m+mi}{1}\PYG{p}{,}\PYG{n}{np}\PYG{o}{.}\PYG{n}{size}\PYG{p}{(}\PYG{n}{B2}\PYG{p}{)}\PYG{p}{)}\PYG{p}{)}\PYG{p}{)}

\PYG{k}{for} \PYG{n}{i} \PYG{o+ow}{in} \PYG{n+nb}{range}\PYG{p}{(}\PYG{l+m+mi}{0}\PYG{p}{,}\PYG{n}{np}\PYG{o}{.}\PYG{n}{size}\PYG{p}{(}\PYG{n}{B2}\PYG{p}{)}\PYG{p}{)}\PYG{p}{:}     
    \PYG{n}{S}\PYG{o}{=}\PYG{n}{P}\PYG{o}{/}\PYG{n}{P0}\PYG{p}{[}\PYG{n}{i}\PYG{p}{]}
    \PYG{n}{c}\PYG{o}{=}\PYG{n+nb}{next}\PYG{p}{(}\PYG{n}{color}\PYG{p}{)}
    
    \PYG{c+c1}{\PYGZsh{}Langmuir isotherm   }
    \PYG{n}{vv1}\PYG{o}{=}\PYG{n}{B2}\PYG{p}{[}\PYG{n}{i}\PYG{p}{]}\PYG{o}{*}\PYG{n}{S}\PYG{o}{/}\PYG{p}{(}\PYG{l+m+mi}{1}\PYG{o}{\PYGZhy{}}\PYG{n}{S}\PYG{p}{)}\PYG{o}{*}\PYG{p}{(}\PYG{l+m+mi}{1}\PYG{o}{\PYGZhy{}}\PYG{p}{(}\PYG{n}{n}\PYG{p}{[}\PYG{n}{i}\PYG{p}{]}\PYG{o}{+}\PYG{l+m+mi}{1}\PYG{p}{)}\PYG{o}{*}\PYG{n}{np}\PYG{o}{.}\PYG{n}{power}\PYG{p}{(}\PYG{n}{S}\PYG{p}{,}\PYG{n}{n}\PYG{p}{[}\PYG{n}{i}\PYG{p}{]}\PYG{p}{)}\PYG{o}{+}\PYG{n}{n}\PYG{p}{[}\PYG{n}{i}\PYG{p}{]}\PYG{o}{*}\PYG{n}{np}\PYG{o}{.}\PYG{n}{power}\PYG{p}{(}\PYG{n}{S}\PYG{p}{,}\PYG{n}{n}\PYG{p}{[}\PYG{n}{i}\PYG{p}{]}\PYG{o}{+}\PYG{l+m+mi}{1}\PYG{p}{)}\PYG{p}{)}\PYG{o}{/}\PYG{p}{(}\PYG{l+m+mi}{1}\PYG{o}{+}\PYG{p}{(}\PYG{n}{B2}\PYG{p}{[}\PYG{n}{i}\PYG{p}{]}\PYG{o}{\PYGZhy{}}\PYG{l+m+mi}{1}\PYG{p}{)}\PYG{o}{*}\PYG{n}{S}\PYG{o}{\PYGZhy{}}\PYG{n}{B2}\PYG{p}{[}\PYG{n}{i}\PYG{p}{]}\PYG{o}{*}\PYG{n}{np}\PYG{o}{.}\PYG{n}{power}\PYG{p}{(}\PYG{n}{S}\PYG{p}{,}\PYG{n}{n}\PYG{p}{[}\PYG{n}{i}\PYG{p}{]}\PYG{o}{+}\PYG{l+m+mi}{1}\PYG{p}{)}\PYG{p}{)}\PYG{p}{;}
    
    \PYG{n}{axes}\PYG{o}{.}\PYG{n}{plot}\PYG{p}{(}\PYG{n}{P}\PYG{p}{,}\PYG{n}{vv1}\PYG{p}{,} \PYG{n}{marker}\PYG{o}{=}\PYG{l+s+s1}{\PYGZsq{}}\PYG{l+s+s1}{ }\PYG{l+s+s1}{\PYGZsq{}} \PYG{p}{,} \PYG{n}{c}\PYG{o}{=}\PYG{n}{c}\PYG{p}{)}

    
\PYG{n}{plt}\PYG{o}{.}\PYG{n}{title}\PYG{p}{(}\PYG{l+s+s1}{\PYGZsq{}}\PYG{l+s+s1}{BET isotherm types}\PYG{l+s+s1}{\PYGZsq{}}\PYG{p}{,} \PYG{n}{fontsize}\PYG{o}{=}\PYG{l+m+mi}{18}\PYG{p}{)}\PYG{p}{;}
\PYG{n}{axes}\PYG{o}{.}\PYG{n}{set\PYGZus{}xlabel}\PYG{p}{(}\PYG{l+s+s1}{\PYGZsq{}}\PYG{l+s+s1}{P [bar]}\PYG{l+s+s1}{\PYGZsq{}}\PYG{p}{,} \PYG{n}{fontsize}\PYG{o}{=}\PYG{l+m+mi}{18}\PYG{p}{)}\PYG{p}{;}
\PYG{n}{axes}\PYG{o}{.}\PYG{n}{set\PYGZus{}ylabel}\PYG{p}{(}\PYG{l+s+s1}{\PYGZsq{}}\PYG{l+s+s1}{\PYGZdl{}v\PYGZus{}l\PYGZdl{} / \PYGZdl{}v\PYGZus{}l\PYGZca{}1\PYGZdl{}}\PYG{l+s+s1}{\PYGZsq{}}\PYG{p}{,}\PYG{n}{fontsize}\PYG{o}{=}\PYG{l+m+mi}{18}\PYG{p}{)}\PYG{p}{;}
\PYG{n}{axes}\PYG{o}{.}\PYG{n}{legend}\PYG{p}{(}\PYG{p}{[}\PYG{l+s+s1}{\PYGZsq{}}\PYG{l+s+s1}{Type I}\PYG{l+s+s1}{\PYGZsq{}}\PYG{p}{,}\PYG{l+s+s1}{\PYGZsq{}}\PYG{l+s+s1}{Type II}\PYG{l+s+s1}{\PYGZsq{}}\PYG{p}{,}\PYG{l+s+s1}{\PYGZsq{}}\PYG{l+s+s1}{Type III}\PYG{l+s+s1}{\PYGZsq{}}\PYG{p}{,}\PYG{l+s+s1}{\PYGZsq{}}\PYG{l+s+s1}{Type IV}\PYG{l+s+s1}{\PYGZsq{}}\PYG{p}{,}\PYG{l+s+s1}{\PYGZsq{}}\PYG{l+s+s1}{Type V}\PYG{l+s+s1}{\PYGZsq{}}\PYG{p}{]}\PYG{p}{,} \PYG{n}{fontsize}\PYG{o}{=}\PYG{l+m+mi}{18}\PYG{p}{)}\PYG{p}{;}
\end{sphinxVerbatim}

\end{sphinxuseclass}\end{sphinxVerbatimInput}
\begin{sphinxVerbatimOutput}

\begin{sphinxuseclass}{cell_output}
\noindent\sphinxincludegraphics{{38bfe769f2bdc4d818a61b01faaff21516e66ee829a4eac67a5086569b5fb4ce}.png}

\end{sphinxuseclass}\end{sphinxVerbatimOutput}

\end{sphinxuseclass}
\sphinxstepscope


\section{7. Equilibria in Ion Exchange Systems}
\label{\detokenize{Ion_Exchange:equilibria-in-ion-exchange-systems}}\label{\detokenize{Ion_Exchange::doc}}
\sphinxAtStartPar
In ion exchange processes adsorption takes place through by exchanging the counterions bound to charged ligands immobilised on the stationary phase.
In the following examples the ligands are anionic and monovalent and are indicated with \(R^-\).


\subsection{7.1 Monovalent, symmetric cation exchange.}
\label{\detokenize{Ion_Exchange:monovalent-symmetric-cation-exchange}}
\sphinxAtStartPar
Let us consider an ion exchange system in which the cations that are selectively displaced from the liquid phase are indicated as \(A^+\), while the ligand counterions are indicated with \(B^+\).
The ion exchange reaction considered can be written as:
\begin{equation*}
\begin{split}
A^+ + R^-B^+ + X^- \leftrightarrows B^+ + R^-A^+ + X^- \\
\end{split}
\end{equation*}
\sphinxAtStartPar
where \(X^-\) are spectator anions.

\sphinxAtStartPar
The isotherm in the case of an ion exchange process can be derived starting from the equilibrium constant for the ion exchange reaction.
If we indicate with \(C_A\), \(C_B\) the concentration of cations A and B in the fluid phase and with \(n_A\), \(n_B\) their respective concentration in the adsorbed phase.

\sphinxAtStartPar
The equilibrtium constant is written as:
\begin{equation}\label{equation:Ion_Exchange:IEXeq1}
\begin{split}
K_{AB}=\frac{n_AC_B}{n_BC_A}
\end{split}
\end{equation}
\sphinxAtStartPar
The concentration of surface ligands is related to the concentration in the adsorbed phase of cations \(A^+\) and \(B^+\) by the following expression:
\begin{equation}\label{equation:Ion_Exchange:IEXeq2}
\begin{split}
n_R=n_A+n_B
\end{split}
\end{equation}
\sphinxAtStartPar
The total cationic concentration in the fluid phase \(C_T\) is analogously defined as:
\begin{equation}\label{equation:Ion_Exchange:IEXeq3}
\begin{split}
C_T=C_A+C_B
\end{split}
\end{equation}
\sphinxAtStartPar
hence the equilibrium constant becomes:
\begin{equation}\label{equation:Ion_Exchange:IEXeq4}
\begin{split}
K_{AB}=\frac{n_A(C_T-C_A)}{(n_R-n_A)C_A}
\end{split}
\end{equation}
\sphinxAtStartPar
This expression can be manipulated to provide an explicit relation between the concetration in the adsorbed phase \(n_A\), and the concentration in the fluid phase \(C_A\), i.e. an isotherm for the ion exchange process:
\begin{equation}\label{equation:Ion_Exchange:IEXeq5}
\begin{split}
n_A=\frac{n_RK_{AB}C_A}{C_T+C_A(K_{AB}-1)}
\end{split}
\end{equation}
\sphinxAtStartPar
This expression can be further simplified by introducing the molar fractions in the adsorbed and fluid phases indicated with \(y_A\) and \(x_A\), respectively:
\begin{equation}\label{equation:Ion_Exchange:IEXeq6}
\begin{split}
y_A=\frac{K_{AB}x_A}{1+x_A(K_{AB}-1)}
\end{split}
\end{equation}
\begin{sphinxuseclass}{cell}\begin{sphinxVerbatimInput}

\begin{sphinxuseclass}{cell_input}
\begin{sphinxVerbatim}[commandchars=\\\{\}]
\PYG{k+kn}{import}\PYG{+w}{ }\PYG{n+nn}{numpy}\PYG{+w}{ }\PYG{k}{as}\PYG{+w}{ }\PYG{n+nn}{np}
\PYG{k+kn}{import}\PYG{+w}{ }\PYG{n+nn}{matplotlib}\PYG{n+nn}{.}\PYG{n+nn}{pyplot}\PYG{+w}{ }\PYG{k}{as}\PYG{+w}{ }\PYG{n+nn}{plt} 
\PYG{k+kn}{from}\PYG{+w}{ }\PYG{n+nn}{matplotlib}\PYG{n+nn}{.}\PYG{n+nn}{pyplot}\PYG{+w}{ }\PYG{k+kn}{import} \PYG{n}{cm}

\PYG{c+c1}{\PYGZsh{}Plotting}
\PYG{n}{figure}\PYG{o}{=}\PYG{n}{plt}\PYG{o}{.}\PYG{n}{figure}\PYG{p}{(}\PYG{p}{)}
\PYG{n}{axes} \PYG{o}{=} \PYG{n}{figure}\PYG{o}{.}\PYG{n}{add\PYGZus{}axes}\PYG{p}{(}\PYG{p}{[}\PYG{l+m+mf}{0.1}\PYG{p}{,}\PYG{l+m+mf}{0.1}\PYG{p}{,}\PYG{l+m+mf}{1.2}\PYG{p}{,}\PYG{l+m+mf}{1.2}\PYG{p}{]}\PYG{p}{)}
\PYG{n}{plt}\PYG{o}{.}\PYG{n}{xticks}\PYG{p}{(}\PYG{n}{fontsize}\PYG{o}{=}\PYG{l+m+mi}{14}\PYG{p}{)}
\PYG{n}{plt}\PYG{o}{.}\PYG{n}{yticks}\PYG{p}{(}\PYG{n}{fontsize}\PYG{o}{=}\PYG{l+m+mi}{14}\PYG{p}{)}


\PYG{c+c1}{\PYGZsh{} molar fraction in the feed}
\PYG{n}{z} \PYG{o}{=} \PYG{l+m+mf}{0.3}
\PYG{n}{N} \PYG{o}{=} \PYG{l+m+mi}{500} \PYG{c+c1}{\PYGZsh{}number of points}
\PYG{n}{x} \PYG{o}{=} \PYG{n}{np}\PYG{o}{.}\PYG{n}{linspace}\PYG{p}{(}\PYG{l+m+mi}{0}\PYG{p}{,} \PYG{l+m+mi}{1}\PYG{p}{,} \PYG{n}{N}\PYG{p}{)}
\PYG{n}{kAB} \PYG{o}{=} \PYG{n}{np}\PYG{o}{.}\PYG{n}{array}\PYG{p}{(}\PYG{p}{[}\PYG{l+m+mf}{0.1}\PYG{p}{,} \PYG{l+m+mi}{1}\PYG{p}{,} \PYG{l+m+mi}{2}\PYG{p}{,} \PYG{l+m+mi}{5}\PYG{p}{,} \PYG{l+m+mi}{10}\PYG{p}{,} \PYG{l+m+mi}{50}\PYG{p}{,} \PYG{l+m+mi}{100}\PYG{p}{,} \PYG{l+m+mi}{1000}\PYG{p}{]}\PYG{p}{)}\PYG{p}{;}

\PYG{n}{color}\PYG{o}{=}\PYG{n+nb}{iter}\PYG{p}{(}\PYG{n}{cm}\PYG{o}{.}\PYG{n}{gist\PYGZus{}heat}\PYG{p}{(}\PYG{n}{np}\PYG{o}{.}\PYG{n}{linspace}\PYG{p}{(}\PYG{l+m+mi}{0}\PYG{p}{,}\PYG{l+m+mi}{1}\PYG{p}{,}\PYG{n}{np}\PYG{o}{.}\PYG{n}{size}\PYG{p}{(}\PYG{n}{kAB}\PYG{p}{)}\PYG{p}{)}\PYG{p}{)}\PYG{p}{)}

\PYG{k}{for} \PYG{n}{i} \PYG{o+ow}{in} \PYG{n+nb}{range}\PYG{p}{(}\PYG{l+m+mi}{0}\PYG{p}{,}\PYG{n}{np}\PYG{o}{.}\PYG{n}{size}\PYG{p}{(}\PYG{n}{kAB}\PYG{p}{)}\PYG{p}{)}\PYG{p}{:}     
    \PYG{n}{c}\PYG{o}{=}\PYG{n+nb}{next}\PYG{p}{(}\PYG{n}{color}\PYG{p}{)}
    \PYG{c+c1}{\PYGZsh{}Langmuir isotherm}
    \PYG{n}{y} \PYG{o}{=} \PYG{n}{kAB}\PYG{p}{[}\PYG{n}{i}\PYG{p}{]} \PYG{o}{*} \PYG{n}{x}  \PYG{o}{/} \PYG{p}{(}\PYG{l+m+mi}{1} \PYG{o}{+} \PYG{n}{x}\PYG{o}{*}\PYG{p}{(}\PYG{n}{kAB}\PYG{p}{[}\PYG{n}{i}\PYG{p}{]}\PYG{o}{\PYGZhy{}}\PYG{l+m+mi}{1}\PYG{p}{)}\PYG{p}{)} 
    \PYG{n}{axes}\PYG{o}{.}\PYG{n}{plot}\PYG{p}{(}\PYG{n}{x}\PYG{p}{,}\PYG{n}{y}\PYG{p}{,} \PYG{n}{marker}\PYG{o}{=}\PYG{l+s+s1}{\PYGZsq{}}\PYG{l+s+s1}{ }\PYG{l+s+s1}{\PYGZsq{}} \PYG{p}{,} \PYG{n}{c}\PYG{o}{=}\PYG{n}{c}\PYG{p}{)}

    
\PYG{n}{plt}\PYG{o}{.}\PYG{n}{title}\PYG{p}{(}\PYG{l+s+s1}{\PYGZsq{}}\PYG{l+s+s1}{Ion Exchange}\PYG{l+s+s1}{\PYGZsq{}}\PYG{p}{,} \PYG{n}{fontsize}\PYG{o}{=}\PYG{l+m+mi}{18}\PYG{p}{)}\PYG{p}{;}
\PYG{n}{axes}\PYG{o}{.}\PYG{n}{set\PYGZus{}xlabel}\PYG{p}{(}\PYG{l+s+s1}{\PYGZsq{}}\PYG{l+s+s1}{\PYGZdl{}x\PYGZus{}A\PYGZdl{}}\PYG{l+s+s1}{\PYGZsq{}}\PYG{p}{,} \PYG{n}{fontsize}\PYG{o}{=}\PYG{l+m+mi}{18}\PYG{p}{)}\PYG{p}{;}
\PYG{n}{axes}\PYG{o}{.}\PYG{n}{set\PYGZus{}ylabel}\PYG{p}{(}\PYG{l+s+s1}{\PYGZsq{}}\PYG{l+s+s1}{\PYGZdl{}y\PYGZus{}A\PYGZdl{}}\PYG{l+s+s1}{\PYGZsq{}}\PYG{p}{,}\PYG{n}{fontsize}\PYG{o}{=}\PYG{l+m+mi}{18}\PYG{p}{)}\PYG{p}{;}
\end{sphinxVerbatim}

\end{sphinxuseclass}\end{sphinxVerbatimInput}
\begin{sphinxVerbatimOutput}

\begin{sphinxuseclass}{cell_output}
\noindent\sphinxincludegraphics{{2d8727eefcd18047dc4bf1e6c10905a239ac940e0fd35abfd7aa913929b8c301}.png}

\end{sphinxuseclass}\end{sphinxVerbatimOutput}

\end{sphinxuseclass}
\sphinxAtStartPar
The concentration of surface groups in the statrionary phase \(n_R\) also plays a role in determining the amound adsorbed. In the following we report the concentration in the adsorbed phase, in equilibrium with an arbitrary molar fraction \(x_A'\) in the fluid phase as a function of the affinity towards the stationary phase (captured by the equilibrium constant \(K_{AB}\)) and of the density of surface charged ligands \(n_R\).

\begin{sphinxuseclass}{cell}\begin{sphinxVerbatimInput}

\begin{sphinxuseclass}{cell_input}
\begin{sphinxVerbatim}[commandchars=\\\{\}]
\PYG{k+kn}{import}\PYG{+w}{ }\PYG{n+nn}{numpy}\PYG{+w}{ }\PYG{k}{as}\PYG{+w}{ }\PYG{n+nn}{np}
\PYG{k+kn}{import}\PYG{+w}{ }\PYG{n+nn}{pandas}\PYG{+w}{ }\PYG{k}{as}\PYG{+w}{ }\PYG{n+nn}{pd}
\PYG{k+kn}{from}\PYG{+w}{ }\PYG{n+nn}{matplotlib}\PYG{+w}{ }\PYG{k+kn}{import} \PYG{n}{cm}
\PYG{k+kn}{import}\PYG{+w}{ }\PYG{n+nn}{matplotlib}\PYG{n+nn}{.}\PYG{n+nn}{pyplot}\PYG{+w}{ }\PYG{k}{as}\PYG{+w}{ }\PYG{n+nn}{plt} 
\PYG{k+kn}{from}\PYG{+w}{ }\PYG{n+nn}{mpl\PYGZus{}toolkits}\PYG{n+nn}{.}\PYG{n+nn}{mplot3d}\PYG{+w}{ }\PYG{k+kn}{import} \PYG{n}{Axes3D} 

\PYG{c+c1}{\PYGZsh{} Parameters: }
\PYG{c+c1}{\PYGZsh{} molar fraction in the feed}
\PYG{n}{N} \PYG{o}{=} \PYG{l+m+mi}{100} \PYG{c+c1}{\PYGZsh{}number of points}
\PYG{n}{kAB} \PYG{o}{=} \PYG{n}{np}\PYG{o}{.}\PYG{n}{logspace}\PYG{p}{(}\PYG{o}{\PYGZhy{}}\PYG{l+m+mi}{3}\PYG{p}{,} \PYG{l+m+mf}{1.5}\PYG{p}{,} \PYG{n}{N}\PYG{p}{)}
\PYG{n}{nR} \PYG{o}{=} \PYG{n}{np}\PYG{o}{.}\PYG{n}{logspace}\PYG{p}{(}\PYG{l+m+mi}{1}\PYG{p}{,} \PYG{l+m+mi}{3}\PYG{p}{,} \PYG{n}{N}\PYG{p}{)}
\PYG{n}{x} \PYG{o}{=} \PYG{l+m+mf}{0.1} \PYG{p}{;} \PYG{c+c1}{\PYGZsh{} arbitrary value of C\PYGZsq{}}
\PYG{n}{kAB}\PYG{p}{,} \PYG{n}{nR} \PYG{o}{=} \PYG{n}{np}\PYG{o}{.}\PYG{n}{meshgrid}\PYG{p}{(}\PYG{n}{kAB}\PYG{p}{,} \PYG{n}{nR}\PYG{p}{)} 

\PYG{n}{nA} \PYG{o}{=} \PYG{n}{nR} \PYG{o}{*} \PYG{n}{kAB} \PYG{o}{*} \PYG{n}{x}  \PYG{o}{/} \PYG{p}{(}\PYG{l+m+mi}{1} \PYG{o}{+} \PYG{n}{x}\PYG{o}{*}\PYG{p}{(}\PYG{n}{kAB}\PYG{o}{\PYGZhy{}}\PYG{l+m+mi}{1}\PYG{p}{)}\PYG{p}{)} 


\PYG{c+c1}{\PYGZsh{}Plotting}
\PYG{n}{figure}\PYG{o}{=}\PYG{n}{plt}\PYG{o}{.}\PYG{n}{figure}\PYG{p}{(}\PYG{n}{figsize}\PYG{o}{=}\PYG{p}{(}\PYG{l+m+mi}{10}\PYG{p}{,} \PYG{l+m+mi}{8}\PYG{p}{)}\PYG{p}{)}
\PYG{n}{axes} \PYG{o}{=} \PYG{n}{figure}\PYG{o}{.}\PYG{n}{add\PYGZus{}subplot}\PYG{p}{(}\PYG{n}{projection}\PYG{o}{=}\PYG{l+s+s1}{\PYGZsq{}}\PYG{l+s+s1}{3d}\PYG{l+s+s1}{\PYGZsq{}}\PYG{p}{)}

\PYG{n}{plt}\PYG{o}{.}\PYG{n}{xticks}\PYG{p}{(}\PYG{n}{fontsize}\PYG{o}{=}\PYG{l+m+mi}{14}\PYG{p}{)}
\PYG{n}{plt}\PYG{o}{.}\PYG{n}{yticks}\PYG{p}{(}\PYG{n}{fontsize}\PYG{o}{=}\PYG{l+m+mi}{14}\PYG{p}{)}


\PYG{n}{surf}\PYG{o}{=}\PYG{n}{axes}\PYG{o}{.}\PYG{n}{plot\PYGZus{}surface}\PYG{p}{(}\PYG{n}{nR}\PYG{p}{,}\PYG{n}{kAB}\PYG{p}{,}\PYG{n}{nA}\PYG{p}{,}\PYG{n}{cmap}\PYG{o}{=}\PYG{n}{cm}\PYG{o}{.}\PYG{n}{coolwarm}\PYG{p}{,}
                       \PYG{n}{linewidth}\PYG{o}{=}\PYG{l+m+mi}{0}\PYG{p}{,} \PYG{n}{antialiased}\PYG{o}{=}\PYG{k+kc}{False}\PYG{p}{,}\PYG{n}{alpha}\PYG{o}{=}\PYG{l+m+mf}{0.2}\PYG{p}{)}

\PYG{n}{contour}\PYG{o}{=}\PYG{n}{axes}\PYG{o}{.}\PYG{n}{contour}\PYG{p}{(}\PYG{n}{nR}\PYG{p}{,}\PYG{n}{kAB}\PYG{p}{,}\PYG{n}{nA}\PYG{p}{,}\PYG{n}{np}\PYG{o}{.}\PYG{n}{linspace}\PYG{p}{(}\PYG{l+m+mi}{100}\PYG{p}{,} \PYG{l+m+mi}{800}\PYG{p}{,} \PYG{l+m+mi}{10}\PYG{p}{)}\PYG{p}{,}\PYG{n}{cmap}\PYG{o}{=}\PYG{n}{cm}\PYG{o}{.}\PYG{n}{coolwarm}\PYG{p}{)}

\PYG{n}{plt}\PYG{o}{.}\PYG{n}{title}\PYG{p}{(}\PYG{l+s+s1}{\PYGZsq{}}\PYG{l+s+s1}{Adsorbed concentration vs \PYGZdl{}K\PYGZus{}}\PYG{l+s+si}{\PYGZob{}AB\PYGZcb{}}\PYG{l+s+s1}{\PYGZdl{} / \PYGZdl{}nR\PYGZdl{}}\PYG{l+s+s1}{\PYGZsq{}}\PYG{p}{,} \PYG{n}{fontsize}\PYG{o}{=}\PYG{l+m+mi}{18}\PYG{p}{)}\PYG{p}{;}
\PYG{n}{axes}\PYG{o}{.}\PYG{n}{set\PYGZus{}xlabel}\PYG{p}{(}\PYG{l+s+s1}{\PYGZsq{}}\PYG{l+s+s1}{\PYGZdl{}nR\PYGZdl{}}\PYG{l+s+s1}{\PYGZsq{}}\PYG{p}{,} \PYG{n}{fontsize}\PYG{o}{=}\PYG{l+m+mi}{14}\PYG{p}{)}\PYG{p}{;}
\PYG{n}{axes}\PYG{o}{.}\PYG{n}{set\PYGZus{}ylabel}\PYG{p}{(}\PYG{l+s+s1}{\PYGZsq{}}\PYG{l+s+s1}{\PYGZdl{}k\PYGZus{}}\PYG{l+s+si}{\PYGZob{}AB\PYGZcb{}}\PYG{l+s+s1}{\PYGZdl{}}\PYG{l+s+s1}{\PYGZsq{}}\PYG{p}{,}\PYG{n}{fontsize}\PYG{o}{=}\PYG{l+m+mi}{14}\PYG{p}{)}\PYG{p}{;}
\PYG{n}{axes}\PYG{o}{.}\PYG{n}{set\PYGZus{}zlabel}\PYG{p}{(}\PYG{l+s+s1}{\PYGZsq{}}\PYG{l+s+s1}{\PYGZdl{}nA\PYGZdl{}}\PYG{l+s+s1}{\PYGZsq{}}\PYG{p}{,}\PYG{n}{fontsize}\PYG{o}{=}\PYG{l+m+mi}{14}\PYG{p}{)}\PYG{p}{;}
\PYG{n}{figure}\PYG{o}{.}\PYG{n}{colorbar}\PYG{p}{(}\PYG{n}{surf}\PYG{p}{,} \PYG{n}{shrink}\PYG{o}{=}\PYG{l+m+mf}{0.5}\PYG{p}{,} \PYG{n}{aspect}\PYG{o}{=}\PYG{l+m+mi}{5}\PYG{p}{)}\PYG{p}{;}
\end{sphinxVerbatim}

\end{sphinxuseclass}\end{sphinxVerbatimInput}
\begin{sphinxVerbatimOutput}

\begin{sphinxuseclass}{cell_output}
\noindent\sphinxincludegraphics{{290882bd437184a921882aaae5f0516735c05b46d81d3729b429c6c645426325}.png}

\end{sphinxuseclass}\end{sphinxVerbatimOutput}

\end{sphinxuseclass}

\subsection{7.2 Example problem}
\label{\detokenize{Ion_Exchange:example-problem}}
\sphinxAtStartPar
An aqueous solution containing 10\% by mass of NaCl and an unknown concentration of KCl exchanges cations with a stationary phase until reaching equilibrium. When equilibrated, the resin has 10.5\% of sites occupied by sodium ions. The ion exchange equilibrium constant, KNa+/K+ , is equal to 10.

\sphinxAtStartPar
a)	Write the ion exchange reaction describing the symmetric monovalent exchange taking place in the ion\sphinxhyphen{}exchange column. 			\\
b)	Compute the molar fraction of HNO3 in solution at equilibrium.	\\
c)	If a solution with the same composition, but containing KNO3 instead of NaNO3
(KK+/H+ =9.2) was equilibrated with the same resin, how do you expect the fraction of occupied resin sites to change?

\sphinxstepscope


\section{8. Batch Slurry Adsorption}
\label{\detokenize{Slurry_Adsorption:batch-slurry-adsorption}}\label{\detokenize{Slurry_Adsorption::doc}}
\sphinxAtStartPar
In this notes we will use an example problem to introduce material balances in slurry adsorption processes. We will consider a Batch, process configuration and then its extension to a continuous process configuration.


\subsection{8.1 Problem Statement}
\label{\detokenize{Slurry_Adsorption:problem-statement}}
\sphinxAtStartPar
A water purification process is carried out via batch adsorption. The initial concentration of the pollutant is 0.1 \([mol l^{-1}]\), and it needs to be lowered by at least an order of magnitude. The process is carried out in batches characterised by a water volume \(V_{\ell}\) of 2 \(m^3\), in which \(100\,kg\) of adsorbent are dispersed to form a slurry. The adsorbent is characterised spherical particles with an average radius of 1 \(mm\) and density of \(700 kg\,m^{-3}\). The pollutant adsorption on the surface of the stationary phase particles is well described by a linear isotherm \(q=Hc^*\) with \(H=1\times{10^{-1}} [dm]\).  The liquid\sphinxhyphen{}side mass transfer coefficient is \(k_{\ell}=1\times{10^{-8}}\,[m h^{-1}]\).
\begin{itemize}
\item {} 
\sphinxAtStartPar
Is it possible to carry out the requested purification with the operation parameters described in the problem statment?

\item {} 
\sphinxAtStartPar
If this is the case how much time does it take to lower the pollutant concetration to the target specification?

\end{itemize}


\subsection{8.2 Material Balance in a Batch Slurry adsorption process}
\label{\detokenize{Slurry_Adsorption:material-balance-in-a-batch-slurry-adsorption-process}}
\sphinxAtStartPar
Let us begin by writing the material balance for a batch membrane separator.

\sphinxAtStartPar
Since we have a batch system, no inlet or outlet streams are present and we can write:
\begin{equation}\label{equation:Slurry_Adsorption:slurryeq1}
\begin{split}
V_\ell{c_0}=V_\ell{c}+S_pq
\end{split}
\end{equation}
\sphinxAtStartPar
Where \(V_{\ell}\) is the volume of the liquid phase (it can be considered constant), \(c_0\) is the initial concentration of pollutant in solution, \(c\) is the instantaneous concentration of pollutant in solution,  \(q\) the instantaneous surface concentration of pollutant adsorbed on the stationary phase and \(S_p\) the total surface exposed by the stationary phase particles.

\sphinxAtStartPar
Mass transfer in these cases is dominated by bulk transport on the liquid side, thus the flux \(J\) of the pollutant leaving the fluid phase to be incorporated in the stationary phase can be written as:
\begin{equation}\label{equation:Slurry_Adsorption:slurryeq2}
\begin{split}
J=-k_\ell(c-c^*)
\end{split}
\end{equation}
\sphinxAtStartPar
where \(k_\ell\) is the mass transfer coefficient on the liquid side, and \(c^*\) is the equilibrium concentration, reached at the solid/liquid interface. \(J\) has the units of number of moles per unit time per unit surface.

\sphinxAtStartPar
The differential material balance on the solute concentration can thus be written as:
\begin{equation}\label{equation:Slurry_Adsorption:slurryeq3}
\begin{split}
\frac{dc}{dt}=-k_\ell\frac{S_p}{V_\ell}(c-c^*)=-k_\ell{a}(c-c^*)
\end{split}
\end{equation}
\sphinxAtStartPar
Where \(a\) represents the particles surface area per unit of fluid volume.

\sphinxAtStartPar
The equilibrium of the adsorption process is well captured by a linear isotherm:
\begin{equation}\label{equation:Slurry_Adsorption:slurryeq4}
\begin{split}
c^*=\frac{q}{H}
\label{eq:equilibrium}
\end{split}
\end{equation}
\sphinxAtStartPar
Using Eq. \eqref{equation:Slurry_Adsorption:slurryeq1} and \eqref{equation:Slurry_Adsorption:slurryeq4}, Eq. \eqref{equation:Slurry_Adsorption:slurryeq2} can be rewritten as:
\begin{equation*}
\begin{split}
\frac{dc}{dt}=-k_\ell{a}\left(c-\frac{q}{H}\right) = -k_\ell{a}\left(c-\frac{(c_0-c)V_\ell}{HS_p}\right) = 
\end{split}
\end{equation*}\begin{equation}\label{equation:Slurry_Adsorption:slurryeq5}
\begin{split}
=-k_\ell{a}c\left(1+\frac{V_\ell}{HS_p}\right)+k_\ell{a}c_0\frac{V_\ell}{HS_p}
\end{split}
\end{equation}
\sphinxAtStartPar
Defining: \(\beta=\left(1+\frac{V_\ell}{HS_p}\right)\), and \(\alpha=\frac{V_\ell}{HS_p}\), one gets:
\begin{equation}\label{equation:Slurry_Adsorption:slurryeq6}
\begin{split}
\frac{dc}{dt}+k_\ell{a}\beta\,c=k_\ell{a}c_0\alpha
\end{split}
\end{equation}
\sphinxAtStartPar
which is a first order, non homogeneous ODE in the form \(y^\prime+C_1y=C_2\) with constant \(C_1\) and \(C_2\) coefficients.

\sphinxAtStartPar
In order to solve our problem Eq. \eqref{equation:Slurry_Adsorption:slurryeq6} should either numerically or analytically.
In this case the analytical solution can be obtained from the general solution for ODEs in the form:
\begin{equation}\label{equation:Slurry_Adsorption:slurryeq7}
\begin{split}
y(t)=e^{-A(t)}\left[k_1+\int{be^{A(t)}dt}\right]
\end{split}
\end{equation}
\sphinxAtStartPar
where \(A(t)=\int{C_1dt}\)

\sphinxAtStartPar
In our case:
\begin{itemize}
\item {} 
\sphinxAtStartPar
\(C_1=k_\ell{a\beta}\), thus \(A(t)=k_\ell{a\beta}\,t\)

\item {} 
\sphinxAtStartPar
\(C_2=k_\ell{a\alpha}c_0\),

\end{itemize}

\sphinxAtStartPar
thus:
\begin{equation}\label{equation:Slurry_Adsorption:slurryeq8}
\begin{split}
\int{be^{A(t)}dt}=\int{k_\ell{a\alpha}e^{k_\ell{a\beta}\,t}dt}=\frac{k_\ell{a\alpha}c_0}{k_\ell{a\beta}}e^{k_\ell{a\beta}\,t}=\frac{\alpha{c_0}}{\beta}e^{k_\ell{a\beta}\,t}
\end{split}
\end{equation}
\sphinxAtStartPar
putting everything together:
\begin{equation}\label{equation:Slurry_Adsorption:slurryeq9}
\begin{split}
c(t)=k_1e^{-k_\ell{a\beta}\,t}+\frac{\alpha{c_0}}{\beta}
\end{split}
\end{equation}
\sphinxAtStartPar
The constant \(C_1\) can be computed from the initial conditions: at \(t=0\) \(c(t)=c_0\), therefore:
\begin{equation}\label{equation:Slurry_Adsorption:slurryeq10}
\begin{split}
k_1=c_0\left(1-\frac{\alpha}{\beta}\right)=\frac{c_0}{\beta}
\end{split}
\end{equation}
\sphinxAtStartPar
The final solution for the time\sphinxhyphen{}dependent concentration profile \(c(t)\) is:
\begin{equation}\label{equation:Slurry_Adsorption:slurryeq11}
\begin{split}
c(t)=\frac{c_0}{\beta}\left(e^{-k_\ell{a\beta}\,t}+\alpha\right)
\end{split}
\end{equation}
\sphinxAtStartPar
The performance of a batch slurry process in removing the concentration in the fluid phase is limited, i.e. the concentration has an asymptotic behavior in time \(\lim_{t\to\infty} c(t) = \frac{\alpha{c_0}}{\beta}\).

\begin{sphinxuseclass}{cell}\begin{sphinxVerbatimInput}

\begin{sphinxuseclass}{cell_input}
\begin{sphinxVerbatim}[commandchars=\\\{\}]
\PYG{k+kn}{import}\PYG{+w}{ }\PYG{n+nn}{numpy}\PYG{+w}{ }\PYG{k}{as}\PYG{+w}{ }\PYG{n+nn}{np}
\PYG{k+kn}{import}\PYG{+w}{ }\PYG{n+nn}{matplotlib}\PYG{n+nn}{.}\PYG{n+nn}{pyplot}\PYG{+w}{ }\PYG{k}{as}\PYG{+w}{ }\PYG{n+nn}{plt} 
\PYG{k+kn}{from}\PYG{+w}{ }\PYG{n+nn}{matplotlib}\PYG{n+nn}{.}\PYG{n+nn}{pyplot}\PYG{+w}{ }\PYG{k+kn}{import} \PYG{n}{cm}

\PYG{c+c1}{\PYGZsh{}Plotting}
\PYG{n}{figure}\PYG{o}{=}\PYG{n}{plt}\PYG{o}{.}\PYG{n}{figure}\PYG{p}{(}\PYG{p}{)}
\PYG{n}{axes} \PYG{o}{=} \PYG{n}{figure}\PYG{o}{.}\PYG{n}{add\PYGZus{}axes}\PYG{p}{(}\PYG{p}{[}\PYG{l+m+mf}{0.1}\PYG{p}{,}\PYG{l+m+mf}{0.1}\PYG{p}{,}\PYG{l+m+mf}{1.2}\PYG{p}{,}\PYG{l+m+mf}{1.2}\PYG{p}{]}\PYG{p}{)}
\PYG{n}{plt}\PYG{o}{.}\PYG{n}{xticks}\PYG{p}{(}\PYG{n}{fontsize}\PYG{o}{=}\PYG{l+m+mi}{14}\PYG{p}{)}
\PYG{n}{plt}\PYG{o}{.}\PYG{n}{yticks}\PYG{p}{(}\PYG{n}{fontsize}\PYG{o}{=}\PYG{l+m+mi}{14}\PYG{p}{)}
\PYG{n}{N} \PYG{o}{=} \PYG{l+m+mi}{500} \PYG{c+c1}{\PYGZsh{}number of points}
\PYG{n}{time} \PYG{o}{=} \PYG{n}{np}\PYG{o}{.}\PYG{n}{linspace}\PYG{p}{(}\PYG{l+m+mi}{0}\PYG{p}{,}\PYG{l+m+mi}{15}\PYG{p}{,} \PYG{n}{N}\PYG{p}{)}

\PYG{c+c1}{\PYGZsh{}particles mass}
\PYG{n}{pm}\PYG{o}{=}\PYG{l+m+mi}{100}\PYG{p}{;} \PYG{c+c1}{\PYGZsh{}kg}
\PYG{c+c1}{\PYGZsh{}particles volume}
\PYG{n}{pV}\PYG{o}{=}\PYG{l+m+mi}{100}\PYG{o}{/}\PYG{l+m+mi}{700}\PYG{p}{;} \PYG{c+c1}{\PYGZsh{}m\PYGZca{}3}

\PYG{c+c1}{\PYGZsh{}single particle volume}
\PYG{n}{radius}\PYG{o}{=}\PYG{l+m+mf}{1E\PYGZhy{}3}\PYG{p}{;}
\PYG{n}{uV}\PYG{o}{=}\PYG{p}{(}\PYG{l+m+mi}{4}\PYG{o}{/}\PYG{l+m+mi}{3}\PYG{p}{)}\PYG{o}{*}\PYG{n}{np}\PYG{o}{.}\PYG{n}{pi}\PYG{o}{*}\PYG{n}{np}\PYG{o}{.}\PYG{n}{power}\PYG{p}{(}\PYG{n}{radius}\PYG{p}{,}\PYG{l+m+mi}{3}\PYG{p}{)}\PYG{p}{;}
\PYG{n}{uS}\PYG{o}{=}\PYG{l+m+mi}{4}\PYG{o}{*}\PYG{n}{np}\PYG{o}{.}\PYG{n}{pi}\PYG{o}{*}\PYG{o}{*}\PYG{n}{np}\PYG{o}{.}\PYG{n}{power}\PYG{p}{(}\PYG{n}{radius}\PYG{p}{,}\PYG{l+m+mi}{2}\PYG{p}{)}\PYG{p}{;}

\PYG{c+c1}{\PYGZsh{}number\PYGZus{}of particles}
\PYG{n}{NP}\PYG{o}{=}\PYG{n}{pV}\PYG{o}{/}\PYG{n}{uV}\PYG{p}{;}
\PYG{n}{Sp}\PYG{o}{=}\PYG{n}{NP}\PYG{o}{*}\PYG{n}{uS}\PYG{p}{;} \PYG{c+c1}{\PYGZsh{}particles surface}

\PYG{c+c1}{\PYGZsh{}Batch volume}
\PYG{n}{Vl}\PYG{o}{=}\PYG{l+m+mi}{2}\PYG{p}{;} \PYG{c+c1}{\PYGZsh{}m\PYGZca{}3}
\PYG{c+c1}{\PYGZsh{}Mass transfer}
\PYG{n}{kl}\PYG{o}{=}\PYG{l+m+mf}{1E\PYGZhy{}8}\PYG{p}{;} \PYG{c+c1}{\PYGZsh{}m/h}
\PYG{c+c1}{\PYGZsh{}Equilibrium}
\PYG{n}{H}\PYG{o}{=}\PYG{l+m+mf}{1E\PYGZhy{}1}\PYG{p}{;} \PYG{c+c1}{\PYGZsh{}}

\PYG{c+c1}{\PYGZsh{}Analytical solution 1}
\PYG{n}{C0}\PYG{o}{=}\PYG{l+m+mf}{0.1}\PYG{p}{;} \PYG{c+c1}{\PYGZsh{}mol/l}
\PYG{n}{alpha}\PYG{o}{=}\PYG{n}{Vl}\PYG{o}{/}\PYG{n}{H}\PYG{o}{/}\PYG{n}{Sp}\PYG{p}{;}
\PYG{n}{beta}\PYG{o}{=}\PYG{l+m+mi}{1}\PYG{o}{+}\PYG{n}{alpha}\PYG{p}{;}
\PYG{n}{a}\PYG{o}{=}\PYG{n}{Sp}\PYG{o}{/}\PYG{n}{Vl}\PYG{p}{;}
\PYG{n}{C\PYGZus{}inf}\PYG{o}{=} \PYG{n}{np}\PYG{o}{.}\PYG{n}{ones}\PYG{p}{(}\PYG{n}{np}\PYG{o}{.}\PYG{n}{size}\PYG{p}{(}\PYG{n}{time}\PYG{p}{)}\PYG{p}{)} \PYG{o}{*} \PYG{n}{C0}\PYG{o}{*}\PYG{n}{alpha}\PYG{o}{/}\PYG{n}{beta}
\PYG{n}{C} \PYG{o}{=} \PYG{n}{C0}\PYG{o}{/}\PYG{n}{beta} \PYG{o}{*} \PYG{p}{(}\PYG{n}{np}\PYG{o}{.}\PYG{n}{exp}\PYG{p}{(}\PYG{o}{\PYGZhy{}}\PYG{n}{kl}\PYG{o}{*}\PYG{n}{a}\PYG{o}{*}\PYG{n}{beta}\PYG{o}{*}\PYG{n}{time}\PYG{p}{)}\PYG{o}{+}\PYG{n}{alpha}\PYG{p}{)} 

\PYG{n}{color}\PYG{o}{=}\PYG{n+nb}{iter}\PYG{p}{(}\PYG{n}{cm}\PYG{o}{.}\PYG{n}{coolwarm}\PYG{p}{(}\PYG{n}{np}\PYG{o}{.}\PYG{n}{linspace}\PYG{p}{(}\PYG{l+m+mi}{0}\PYG{p}{,}\PYG{l+m+mi}{1}\PYG{p}{,}\PYG{l+m+mi}{2}\PYG{p}{)}\PYG{p}{)}\PYG{p}{)}
\PYG{n}{c}\PYG{o}{=}\PYG{n+nb}{next}\PYG{p}{(}\PYG{n}{color}\PYG{p}{)}
\PYG{n}{axes}\PYG{o}{.}\PYG{n}{plot}\PYG{p}{(}\PYG{n}{time}\PYG{p}{,}\PYG{n}{C}\PYG{p}{,} \PYG{n}{marker}\PYG{o}{=}\PYG{l+s+s1}{\PYGZsq{}}\PYG{l+s+s1}{ }\PYG{l+s+s1}{\PYGZsq{}}\PYG{p}{,}\PYG{n}{c}\PYG{o}{=}\PYG{n}{c}\PYG{p}{)}
\PYG{n}{c}\PYG{o}{=}\PYG{n+nb}{next}\PYG{p}{(}\PYG{n}{color}\PYG{p}{)}
\PYG{n}{axes}\PYG{o}{.}\PYG{n}{plot}\PYG{p}{(}\PYG{n}{time}\PYG{p}{,}\PYG{n}{C\PYGZus{}inf}\PYG{p}{,}\PYG{n}{marker}\PYG{o}{=}\PYG{l+s+s1}{\PYGZsq{}}\PYG{l+s+s1}{ }\PYG{l+s+s1}{\PYGZsq{}}\PYG{p}{,}\PYG{n}{linestyle}\PYG{o}{=}\PYG{l+s+s1}{\PYGZsq{}}\PYG{l+s+s1}{dotted}\PYG{l+s+s1}{\PYGZsq{}}\PYG{p}{,}\PYG{n}{c}\PYG{o}{=}\PYG{n}{c}\PYG{p}{)}
    
\PYG{n}{plt}\PYG{o}{.}\PYG{n}{title}\PYG{p}{(}\PYG{l+s+s1}{\PYGZsq{}}\PYG{l+s+s1}{\PYGZsq{}}\PYG{p}{,} \PYG{n}{fontsize}\PYG{o}{=}\PYG{l+m+mi}{18}\PYG{p}{)}\PYG{p}{;}
\PYG{n}{axes}\PYG{o}{.}\PYG{n}{set\PYGZus{}xlabel}\PYG{p}{(}\PYG{l+s+s1}{\PYGZsq{}}\PYG{l+s+s1}{\PYGZdl{}time}\PYG{l+s+s1}{\PYGZbs{}}\PYG{l+s+s1}{,}\PYG{l+s+s1}{\PYGZbs{}}\PYG{l+s+s1}{,[hours]\PYGZdl{}}\PYG{l+s+s1}{\PYGZsq{}}\PYG{p}{,} \PYG{n}{fontsize}\PYG{o}{=}\PYG{l+m+mi}{18}\PYG{p}{)}\PYG{p}{;}
\PYG{n}{axes}\PYG{o}{.}\PYG{n}{set\PYGZus{}ylabel}\PYG{p}{(}\PYG{l+s+s1}{\PYGZsq{}}\PYG{l+s+s1}{\PYGZdl{}C}\PYG{l+s+s1}{\PYGZbs{}}\PYG{l+s+s1}{,}\PYG{l+s+s1}{\PYGZbs{}}\PYG{l+s+s1}{,[mol / l]\PYGZdl{}}\PYG{l+s+s1}{\PYGZsq{}}\PYG{p}{,}\PYG{n}{fontsize}\PYG{o}{=}\PYG{l+m+mi}{18}\PYG{p}{)}\PYG{p}{;}
\PYG{n}{axes}\PYG{o}{.}\PYG{n}{legend}\PYG{p}{(}\PYG{p}{[}\PYG{l+s+s1}{\PYGZsq{}}\PYG{l+s+s1}{\PYGZdl{}C(t)\PYGZdl{}}\PYG{l+s+s1}{\PYGZsq{}}\PYG{p}{,}\PYG{l+s+s1}{\PYGZsq{}}\PYG{l+s+s1}{\PYGZdl{}C\PYGZus{}}\PYG{l+s+si}{\PYGZob{}inf\PYGZcb{}}\PYG{l+s+s1}{\PYGZdl{}}\PYG{l+s+s1}{\PYGZsq{}}\PYG{p}{]}\PYG{p}{,} \PYG{n}{fontsize}\PYG{o}{=}\PYG{l+m+mi}{18}\PYG{p}{)}\PYG{p}{;}
\end{sphinxVerbatim}

\end{sphinxuseclass}\end{sphinxVerbatimInput}
\begin{sphinxVerbatimOutput}

\begin{sphinxuseclass}{cell_output}
\noindent\sphinxincludegraphics{{a926497cf9cb7223e455165204a6cfe82e8dd1380d4d18ca33009f90ac4ef85c}.png}

\end{sphinxuseclass}\end{sphinxVerbatimOutput}

\end{sphinxuseclass}

\subsection{8.3 Continuous Slurry Adsorption}
\label{\detokenize{Slurry_Adsorption:continuous-slurry-adsorption}}
\sphinxAtStartPar
The slurry adsorption process can be carried out in a continuous mode. In this mode both the adsorbent and the fluid phase are continuously fed into and removed from a stirred tank.

\sphinxAtStartPar
In a continuous process it is convenient to define \(a\) as:
\begin{equation}\label{equation:Slurry_Adsorption:slurryeq13}
\begin{split}
a={S}\sigma/Q
\end{split}
\end{equation}
\sphinxAtStartPar
where \(Q\) is the volumetric flowrate of the liquid phase, ans \(S\) is the mass flowrate of the solid adsorpbent phase, and \(\sigma\) is the surface area per unit mass of the adsorbent.

\sphinxAtStartPar
In this context the  differential material balance on the solute concentration used to solve the unsteady state batch process should be rewritten as:
\begin{equation}\label{equation:Slurry_Adsorption:slurryeq12}
\begin{split}
\frac{dc}{dt}=-k_\ell{a}(c-c^*) \rightarrow \frac{{c_{in}-c_{out}}}{\tau}=k_\ell{a}(c_{out}-c^*)
\end{split}
\end{equation}
\sphinxAtStartPar
where \(\tau\) is the residence time, \(c_{out}\) the concentration in the liquid phase in the unit, \(c_{in}\) the concentration in the feed, and \(a\) the amount of surface of the stationary phase particles per unit volume of the liquid phase.

\sphinxAtStartPar
This expression should solved with the adsorption isotherm \(q=kc^*\), where \(q\) is the amount adsorbed per kg of stationary phase, \(c\) is the molar concentration in solution and \(k\) is the partition coefficient typically expressed in {[}l kg\(^{-1}\){]}, and with the global material balance at steady state, which reads:
\begin{equation}\label{equation:Slurry_Adsorption:slurryeq14}
\begin{split}
c_{in}Q=qS+c_{out}Q
\end{split}
\end{equation}

\subsection{Contributions}
\label{\detokenize{Slurry_Adsorption:contributions}}
\sphinxAtStartPar
Romain Dujardin, 2 May 2022\\
Salar Rahbari Namin, 8 Feb 2021\\
Nikita Gusev, 9 Feb 2021

\sphinxstepscope


\section{9. Liquid/Solid Equilibrium in multicomponent systems: defining solubility}
\label{\detokenize{Solubility:liquid-solid-equilibrium-in-multicomponent-systems-defining-solubility}}\label{\detokenize{Solubility::doc}}
\sphinxAtStartPar
Let us begin by writing the solid\sphinxhyphen{}liquid equilibrium conditions at temperature \(T\) and pressure \(P\), between a pure crystalline phase and a multicomponent liquid phase characterised by composition \(\vec{x}\).
Equilibrium is stated by the equality of chemical potentials:
\begin{equation*}
\begin{split}
\mu_i^s(T,P)=\mu_i^\ell(T,P,\vec{x})
\end{split}
\end{equation*}
\sphinxAtStartPar
which is equivalent to the equality of the fugacity in the two phases:
\begin{equation*}
\begin{split}
f_i^s(T,P)=\hat{f}_i^\ell(T,P,\vec{x})
\end{split}
\end{equation*}
\sphinxAtStartPar
where \(f_i^s(T,P)\) is the fugacity of a pure crystalline solid at \(T\) and \(P\), while \(\hat{f}_i^\ell(T,P,\vec{x})\), is the fugacity of a \textbackslash{}emph\{real\} mixture at the same temperature and pressure, at composition \(\vec{x}\).


\subsection{Fugacity of a real solution}
\label{\detokenize{Solubility:fugacity-of-a-real-solution}}
\sphinxAtStartPar
Focusing on the right hand of Eq.\textbackslash{}ref\{eq:fugacity\}, the fugacity in the \textbackslash{}textbf\{real solution\} phase can be rewritten as a function of the fugacity of the \textbackslash{}textbf\{pure liquid phase\} as follows:
\begin{equation*}
\begin{split}
\hat{f}_i^\ell(T,P,\vec{x})=f_i^\ell(T,P)x_i\gamma_i(T,P,\vec{x})
\end{split}
\end{equation*}
\sphinxAtStartPar
where \(x_i\) is the molar fraction that realizes the equality between fugacities, i.e. the solubility of component \(i\), \(f_i^\ell(T,P)x_i\) is the fugacity of an \textbackslash{}emph\{ideal\} solution at molar fraction \(x_i\), \(\gamma_i(T,P,\vec{x})\) is the activity coefficient, capturing the non\sphinxhyphen{}idealities of the liquid phase in equilibrium with the solid. By definition of an ideal solution \(f_i^\ell(T,P)\) is the fugacity of the \textbackslash{}textbf\{pure solute in its liquid phase\} at temperature T and pressure P.


\subsection{Fugacity of a pure crystalline solid}
\label{\detokenize{Solubility:fugacity-of-a-pure-crystalline-solid}}
\sphinxAtStartPar
Let us now focus on the left\sphinxhyphen{}hand side of Eq.\textbackslash{}ref\{eq:fugacity\}. Here, we are only concerned with the \textbackslash{}textbf\{physical properties of the pure solute\}.
The fugacity of the pure solid phase can also be written as a function of the fugacity of the pure liquid by introducing the definition of fugacity and integrating between the liquid and solid states \textbackslash{}emph\{\(\ell\)\}\(\rightarrow\)\textbackslash{}emph\{s\}:
\begin{equation*}
\begin{split}
\int^{T,P,s}_{T,P,\ell}d\mu_i=\int^{T,P,s}_{T,P,\ell}RTd\ln{f_i}
\end{split}
\end{equation*}
\sphinxAtStartPar
which, integrated, leads to:
\begin{equation*}
\begin{split}
\Delta{\mu}_{i,\ell\rightarrow{s}}=RTln\left(\frac{f_i^s(T,P)}{f_i^\ell(T,P)}\right)
\end{split}
\end{equation*}
\sphinxAtStartPar
and rearranged gives:
\begin{equation*}
\begin{split}
f_i^s(T,P)=f_i^\ell(T,P)e^{\frac{\Delta{\mu}_{i,\ell\rightarrow{s}}}{RT}}
\end{split}
\end{equation*}
\sphinxAtStartPar
where notably \(\Delta{\mu}_{i,\ell\rightarrow{s}}\), is the change in chemical potential associated with the liquid to solid transition of the \textbackslash{}emph\{pure solid\}.
We note that the change in chemical potential for a pure substance corresponds to the the partial molar change in free energy \(\Delta{g_i(T,P)}\):
\begin{equation*}
\begin{split}
\Delta\mu_{i,\ell\rightarrow{s}}=\Delta{g}_i^s(T,P)=\Delta{h}_i(T,P)-T\Delta{s}_i(T,P)
\end{split}
\end{equation*}
\sphinxAtStartPar
Now we express the molar enthalpy change \textbackslash{}textbf\{for the pure solute\} by considering a thermodynamic cycle, that goes through the melting temperature \(T_f\) at pressure P. This enables the computation of \(\Delta{h}_i(T,P)\) and \(\Delta{s}_i(T,P)\) via an arbitrary but convenient path, taking advantage of the fact that all terms in the free energy expression are \textbackslash{}emph\{state functions\}.


\subsection{Enthalpy of phase transition at T and P for pure solute}
\label{\detokenize{Solubility:enthalpy-of-phase-transition-at-t-and-p-for-pure-solute}}
\sphinxAtStartPar
The enthalpy change associated with the solid\sphinxhyphen{}to\sphinxhyphen{}liquid transition becomes:
\begin{equation*}
\begin{split}
\Delta{h}_i(T,P)=-\Delta{h_{fus}\left(T_f\right)}+\int_T^{T_f}Cp^{\ell}dT-\int_T^{T_f}Cp^{s}dT
\end{split}
\end{equation*}
\sphinxAtStartPar
where:
\begin{itemize}
\item {} 
\sphinxAtStartPar
\(\int_T^{T_f}Cp^{\ell}dT\) is the enthalpy change associated with raising the temperature of the \sphinxstyleemphasis{pure solute liquid} from \(T\) to \(T_m\)

\item {} 
\sphinxAtStartPar
\(-\int_T^{T_f}Cp^{s}dT\) is the enthalpy change associated with decreasing the temperature of the \sphinxstyleemphasis{pure solid} from \(T_m\) to \(T\)

\item {} 
\sphinxAtStartPar
\({-\Delta{h_{fus}\left(T_f\right)}}\) is the change in enthalpy associated with the liquid to solid transition of the \sphinxstyleemphasis{pure solute} at the melting temperature

\end{itemize}

\sphinxAtStartPar
For heat capacities independent of T, the enthalpic contribution reduces to:
\begin{equation*}
\begin{split}
\Delta{h}_i(T,P)\simeq-\Delta{h_{fus}\left(T_f\right)}+\left(Cp^{\ell}-Cp^{s}\right)\left(T_f-T\right)
\end{split}
\end{equation*}

\subsection{Entropy of phase transition at T and P for pure solute}
\label{\detokenize{Solubility:entropy-of-phase-transition-at-t-and-p-for-pure-solute}}
\sphinxAtStartPar
Analogously to the enthalpy term, we can write the molar change in entropy via a the same thermodynamic cycle:
\begin{equation*}
\begin{split}
\Delta{s}_i(T,P)=-\frac{\Delta{h_{fus}\left(T_f\right)}}{T_f}+\int_T^{T_f}\frac{Cp^{\ell}}{T}dT-\int_T^{T_f}\frac{Cp^{s}}{T}dT 
\end{split}
\end{equation*}
\sphinxAtStartPar
where:
\begin{itemize}
\item {} 
\sphinxAtStartPar
\(+\int_T^{T_f}\frac{Cp^{\ell}}{T}dT\) is the entropy change associated with raising the temperature of the \sphinxstyleemphasis{pure solute liquid} from \(T\) to \(T_m\)

\item {} 
\sphinxAtStartPar
\(-\int_T^{T_f}\frac{Cp^{s}}{T}dT \) is the enthalpy change associated with decreasing the temperature of the \sphinxstyleemphasis{pure solid} from \(T_m\) to \(T\)

\item {} 
\sphinxAtStartPar
\(-\frac{\Delta{h_{fus}\left(T_f\right)}}{T_f}\) is the change in entropy associated with the liquid to the solid transition of the \sphinxstyleemphasis{pure solute} at the melting temperature

\end{itemize}

\sphinxAtStartPar
Again, for heat capacities independent of T, the entropy change reduces to:
\begin{equation*}
\begin{split}
\Delta{s}_i(T,P) \simeq -\frac{\Delta{h_{fus}}}{T_f}+\left(Cp^{\ell}-Cp^{s}\right)\ln{\left(\frac{T_f}{T}\right)}
\end{split}
\end{equation*}

\subsection{Free energy of phase transition at T and P for pure solute}
\label{\detokenize{Solubility:free-energy-of-phase-transition-at-t-and-p-for-pure-solute}}
\sphinxAtStartPar
Defining \(\Delta{Cp}_{\ell\rightarrow{s}}=\left(Cp^{\ell}-Cp^{s}\right)\) we can now recover the expression for \(\Delta{g}_i^s(T,P)\), where the only assumption is that \(Cp^{\ell}\) and \(Cp^{s}\) are independent from the temperature T.
\begin{equation*}
\begin{split}
\Delta{g}_i^s(T,P)=-\Delta{h_{fus}\left(1-\frac{T}{T_f}\right)}+\Delta{Cp}_{\ell\rightarrow{s}}\left(T_f-T-T\ln{\left(\frac{T_f}{T}\right)}\right)
\end{split}
\end{equation*}

\subsection{Solubility}
\label{\detokenize{Solubility:solubility}}
\sphinxAtStartPar
Inserting now this expression in Eq. \textbackslash{}ref\{eq:fugacitysolid\} yields:
\begin{equation*}
\begin{split}
f_i^s(T,P)=f_i^\ell(T,P)\exp\left[\frac{{-\Delta{h_{fus}\left(1-\frac{T}{T_f}\right)}+\Delta{Cp}_{\ell\rightarrow{s}}\left(T_f-T-T\ln{\left(\frac{T_f}{T}\right)}\right)}}{RT}\right]
\end{split}
\end{equation*}
\sphinxAtStartPar
which then can be inserted into the equality of fugacities defining equilibrium Eq.\textbackslash{}ref\{eq:fugacity\} giving:
\begin{equation*}
\begin{split}
f_i^\ell(T,P)x_i\gamma_i(T,P,\vec{x})=f_i^\ell(T,P)\exp\left[\frac{{-\Delta{h_{fus}\left(1-\frac{T}{T_f}\right)}+\Delta{Cp}_{\ell\rightarrow{s}}\left(T_f-T-T\ln{\left(\frac{T_f}{T}\right)}\right)}}{RT}\right]
\end{split}
\end{equation*}
\sphinxAtStartPar
which can be rearranged to yield an expression of the solubility \(x_i\) where the only approximation introduced is in considering the heat capacities of both the crystal and the liquid constant in T.
\begin{equation*}
\begin{split}
x_i=\frac{1}{\gamma_i(T,P,\vec{x})}\exp\left[\frac{{-\Delta{h_{fus}\left(1-\frac{T}{T_f}\right)}+\Delta{Cp}_{\ell\rightarrow{s}}\left(T_f-T-T\ln{\left(\frac{T_f}{T}\right)}\right)}}{RT}\right]
\end{split}
\end{equation*}
\sphinxAtStartPar
In this expression, the fugacity of the pure liquid solute, \(f_i^\ell(T,P)\) here cancels out, and the term \(-\Delta{h_{fus}}\) is not approximating a different quantity (like the enthalpy of dissolution). It is actually introduced rigorously by the convenient choice of an appropriate thermodynamic cycle for the calculation of \(\Delta\mu_{i,\ell\rightarrow{s}}\).


\subsection{The Schroeder\sphinxhyphen{}Van Laar Equation}
\label{\detokenize{Solubility:the-schroeder-van-laar-equation}}
\sphinxAtStartPar
The Schroeder\sphinxhyphen{}Van Laar equation for the solubility is derived from \textbackslash{}ref\{eq:solubility\} by assuming that \(\Delta{Cp}_{\ell\rightarrow{s}}\) is negligible.

\sphinxAtStartPar
This simplification yields:
\begin{equation*}
\begin{split}
\Delta{g}_i^s(T,P)\simeq-\Delta{h_{fus}\left(1-\frac{T}{T_f}\right)}
\end{split}
\end{equation*}
\sphinxAtStartPar
Thus leading to the Schroeder\sphinxhyphen{}Van Laar solubility equation for \(x_i\):
\begin{equation*}
\begin{split}
x_i=\frac{1}{\gamma_i(T,P,\vec{x})}\exp\left[{-\frac{\Delta{h_{fus}}}{R}\left(\frac{1}{T}-\frac{1}{T_f}\right)}\right]
\end{split}
\end{equation*}
\sphinxstepscope


\section{10. MSMPR Crystallizer: population balance}
\label{\detokenize{MSMPR_crystallizer:msmpr-crystallizer-population-balance}}\label{\detokenize{MSMPR_crystallizer::doc}}
\sphinxAtStartPar
Let us consider an Mixed Suspension Mixed Product removal crystallizer operating at steady state at constant temperature.
In these notes we will derive the analystical expression of the steady state particle size distribution by introducing a monodimensional population balance, and we will discuss its coupling with the global material balance.


\subsection{10.0 Material Balance}
\label{\detokenize{MSMPR_crystallizer:material-balance}}
\sphinxAtStartPar
Let us begin by considering the global material balance for an MSMPR crystallizer. The operation of an ideal MSMPR crystallizer is characterised by the following set of hypotheses:
\begin{itemize}
\item {} 
\sphinxAtStartPar
The crystalliser operates continuously

\item {} 
\sphinxAtStartPar
Operation is at steady state

\item {} 
\sphinxAtStartPar
The crystalliser is perfectly mixed (both with respect to the liquid and the solid phases)

\item {} 
\sphinxAtStartPar
The outlet has a composition identical in all respects to the crystallizer content

\item {} 
\sphinxAtStartPar
There are no crystals in the feed

\item {} 
\sphinxAtStartPar
The growth rate G=dL/dt is independent from the crystal size (L)

\item {} 
\sphinxAtStartPar
All crystals have the same shape

\item {} 
\sphinxAtStartPar
The system is not subject to agglomeration or breakage

\end{itemize}

\sphinxAtStartPar
At steady state the mass balance for the solute precipitating in an MSMPR crystallizer can be written as:
\begin{equation}\label{equation:MSMPR_crystallizer:MSMPReq0}
\begin{split}
Q_{IN}C_{IN}=Q_{OUT}C_{OUT}+Q_{OUT}M_T
\end{split}
\end{equation}
\sphinxAtStartPar
where \(Q_{IN}\) is the volumetric inlet flowrate of the liquid phase, \(Q_{OUT}\) is the volumetric outlet flowrate of the liquid phase (mother liquor), \(C_{IN}\) is the solute density in the feed, \(C_{OUT}\) is the solute density int he outlet stream, and \(M_T\) is the mass of crystals precipitated per unit volume of liquid phase.
Typically \(Q_{IN}\simeq{Q_{OUT}}=Q\) and thus \(M_T=C_{IN}-C_{OUT}\).

\sphinxAtStartPar
It should be noted that \(M_T\) is determined by the process of precipitation producing a certain particles population and can be computed from the kinetics of crystal nucleation and growth, based on the population balance discussed in the neLt section.


\subsection{10.1 Number and number population density}
\label{\detokenize{MSMPR_crystallizer:number-and-number-population-density}}
\sphinxAtStartPar
Let us begin by defining \(N(L)\) as the \sphinxstylestrong{number} of  crystals of characteristic length \(L\) per unit volume.
The \sphinxstylestrong{number population density} \(n\), or in short number density, is defined as:
\begin{equation}\label{equation:MSMPR_crystallizer:MSMPReq1}
\begin{split}
\lim_{\Delta{L}\to{0}}\frac{\Delta{N(L)}}{\Delta{L}}=\frac{d{N(L)}}{d{L}}=n(L)
\end{split}
\end{equation}
\sphinxAtStartPar
and represents the number of crystals of length \(L\) \sphinxstylestrong{per unit length}, per unit volume.
The number of crystals per unit volume with dimension comprised between \(L_1\) and \(L_2\) can thus be computed as a definite integral of the population density:
\begin{equation}\label{equation:MSMPR_crystallizer:MSMPReq2}
\begin{split}
N_{1,2}=\int_{L_1}^{L_2}n(L)dL
\end{split}
\end{equation}
\sphinxAtStartPar
Similarly the total number of crystals per unit volume is obtained as:
\begin{equation}\label{equation:MSMPR_crystallizer:MSMPReq3}
\begin{split}
N_{T}=\int_{0}^{\infty}n(L)dL
\end{split}
\end{equation}

\subsection{10.2 Population Balance in an MSMPR crystallizer}
\label{\detokenize{MSMPR_crystallizer:population-balance-in-an-msmpr-crystallizer}}
\sphinxAtStartPar
Let us write a steady state population balance for particles with size comprised between \(L_1\) and \(L_2\).
\begin{equation}\label{equation:MSMPR_crystallizer:MSMPReq4}
\begin{split}
{\Delta{N}}V=n_1G_1V\Delta{t}-n_2G_2V\Delta{t}-Q\overline{n}\Delta{L}\Delta{t}
\end{split}
\end{equation}
\sphinxAtStartPar
where \(\Delta{N}V\) is the change in the number of particles within the length interval \(\Delta{L}=L_2-L_1\) per unit volume, accumulated in time \(\Delta{t}\).
The term \(n_1G_1V\Delta{t}\) captures the number of crystals that outgrow length \(L_1\) in time \(\Delta{t}\), where \(G_1\) is the growth rate of particles of length \(L_1\), \(V\) is the total volume.
Similarly, the term \(n_2G_2V\Delta{t}\) takes into account particles leaving the length  interval \(\Delta{L}\), outgrowing length \(L_2\).
Finally the term \(Q\overline{n}\Delta{L}\Delta{t}\) accounts for the removal of particles with size comprised in  \(\Delta{L}\) from the crystalliser due to the magma flowrate with volumetric flowrate \(Q\), where \(\overline{n}\) represents the average population density within the length interval \(\Delta{L}\).

\sphinxAtStartPar
Eq. \eqref{equation:MSMPR_crystallizer:MSMPReq4} can be rearranged to obtain on the left side a discrete accumulation term, i.e. the number of particles with size compriused in \(\delta{L}\) accumulated in time \(\Delta{t}\):
\begin{equation}\label{equation:MSMPR_crystallizer:MSMPReq5}
\begin{split}
\frac{\Delta{N}}{\Delta{t}}=n_1G_1-n_2G_2-\frac{Q}{V}\overline{n}\Delta{L}
\end{split}
\end{equation}
\sphinxAtStartPar
At steady state there is no accumulation and Eq. \eqref{equation:MSMPR_crystallizer:MSMPReq5} can be rewritten as:
\begin{equation}\label{equation:MSMPR_crystallizer:MSMPReq6}
\begin{split}
\frac{n_2G_2-n_1G_1}{\Delta{L}}=-\frac{Q}{V}\overline{n}
\end{split}
\end{equation}
\sphinxAtStartPar
Noting that the residence time \(\tau=V/Q\), and taking the limit for \(\Delta{L}\rightarrow{0}\) we get:
\begin{equation}\label{equation:MSMPR_crystallizer:MSMPReq7}
\begin{split}
\frac{dn}{dL}=\frac{n}{G\tau}
\end{split}
\end{equation}
\sphinxAtStartPar
Which can be solved analytically, yielding the steady state number population density n(L):
\begin{equation}\label{equation:MSMPR_crystallizer:MSMPReq8}
\begin{split}
n(L)=n_0\eLp\left(-\frac{L}{G\tau}\right)
\end{split}
\end{equation}
\sphinxAtStartPar
where \(n_0\) is an integration constant, determined from the boundary conditions, that has the physical meaning of population density of nuclei in the crystallizer at steady state.

\sphinxAtStartPar
In this formulation of the population balance we consider particles to be born as nuclei at length \(L=0\).
The nucleation rate \(J\), i.e. the number of nuclei generated per unit time per unit volume can be written as:
\begin{equation}\label{equation:MSMPR_crystallizer:MSMPReq9}
\begin{split}
J={\frac{dN}{dt}}\vert_{L=0}
\end{split}
\end{equation}
\sphinxAtStartPar
by applying the chain rule we can rewrite the nucleation rate as:
\begin{equation}\label{equation:MSMPR_crystallizer:MSMPReq10}
\begin{split}
J=\frac{dN}{dL}\vert_{L=0}\frac{dL}{dt}\vert_{L=0}=n_0G
\end{split}
\end{equation}
\sphinxAtStartPar
where, by definition \(\frac{dN}{dL}\vert_{L=0}=n_0\), and \(\frac{dL}{dt}=G\). We shall note that here we are considering the growth rate to be independent from the crystal size.

\sphinxAtStartPar
This eLpression allows us to determine the integration constant \(n_0\) based on characteristic rates of crystal nucleation and growth, yielding:
\begin{equation}\label{equation:MSMPR_crystallizer:MSMPReq11}
\begin{split}
n_0=\frac{J}{G}
\end{split}
\end{equation}

\subsection{10.3 Particle size distribution and process parameters: \protect\(\tau\protect\), \protect\(G\protect\), \protect\(J\protect\)}
\label{\detokenize{MSMPR_crystallizer:particle-size-distribution-and-process-parameters-tau-g-j}}
\begin{sphinxuseclass}{cell}\begin{sphinxVerbatimInput}

\begin{sphinxuseclass}{cell_input}
\begin{sphinxVerbatim}[commandchars=\\\{\}]
\PYG{k+kn}{import}\PYG{+w}{ }\PYG{n+nn}{numpy}\PYG{+w}{ }\PYG{k}{as}\PYG{+w}{ }\PYG{n+nn}{np}
\PYG{k+kn}{import}\PYG{+w}{ }\PYG{n+nn}{matplotlib}\PYG{n+nn}{.}\PYG{n+nn}{pyplot}\PYG{+w}{ }\PYG{k}{as}\PYG{+w}{ }\PYG{n+nn}{plt} 
\PYG{k+kn}{from}\PYG{+w}{ }\PYG{n+nn}{matplotlib}\PYG{n+nn}{.}\PYG{n+nn}{pyplot}\PYG{+w}{ }\PYG{k+kn}{import} \PYG{n}{cm}

\PYG{c+c1}{\PYGZsh{}Plotting}
\PYG{n}{figure}\PYG{o}{=}\PYG{n}{plt}\PYG{o}{.}\PYG{n}{figure}\PYG{p}{(}\PYG{p}{)}
\PYG{n}{axes} \PYG{o}{=} \PYG{n}{figure}\PYG{o}{.}\PYG{n}{add\PYGZus{}axes}\PYG{p}{(}\PYG{p}{[}\PYG{l+m+mf}{0.1}\PYG{p}{,}\PYG{l+m+mf}{0.1}\PYG{p}{,}\PYG{l+m+mf}{1.2}\PYG{p}{,}\PYG{l+m+mf}{1.2}\PYG{p}{]}\PYG{p}{)}
\PYG{n}{plt}\PYG{o}{.}\PYG{n}{xticks}\PYG{p}{(}\PYG{n}{fontsize}\PYG{o}{=}\PYG{l+m+mi}{14}\PYG{p}{)}
\PYG{n}{plt}\PYG{o}{.}\PYG{n}{yticks}\PYG{p}{(}\PYG{n}{fontsize}\PYG{o}{=}\PYG{l+m+mi}{14}\PYG{p}{)}

\PYG{c+c1}{\PYGZsh{} Data}
\PYG{n}{G}\PYG{o}{=}\PYG{l+m+mi}{1}
\PYG{n}{J}\PYG{o}{=}\PYG{l+m+mi}{1}
\PYG{n}{tau}\PYG{o}{=}\PYG{n}{np}\PYG{o}{.}\PYG{n}{array}\PYG{p}{(}\PYG{p}{[}\PYG{l+m+mf}{0.1}\PYG{p}{,} \PYG{l+m+mi}{1}\PYG{p}{,} \PYG{l+m+mi}{2}\PYG{p}{,} \PYG{l+m+mi}{5}\PYG{p}{,} \PYG{l+m+mi}{10}\PYG{p}{,} \PYG{l+m+mi}{50}\PYG{p}{,} \PYG{l+m+mi}{100}\PYG{p}{,} \PYG{l+m+mi}{500}\PYG{p}{]}\PYG{p}{)}\PYG{p}{;}
\PYG{n}{N} \PYG{o}{=} \PYG{l+m+mi}{500} \PYG{c+c1}{\PYGZsh{}number of points}
\PYG{n}{L} \PYG{o}{=} \PYG{n}{np}\PYG{o}{.}\PYG{n}{linspace}\PYG{p}{(}\PYG{l+m+mi}{0}\PYG{p}{,}\PYG{l+m+mi}{100}\PYG{p}{,} \PYG{n}{N}\PYG{p}{)}

\PYG{n}{color}\PYG{o}{=}\PYG{n+nb}{iter}\PYG{p}{(}\PYG{n}{cm}\PYG{o}{.}\PYG{n}{coolwarm}\PYG{p}{(}\PYG{n}{np}\PYG{o}{.}\PYG{n}{linspace}\PYG{p}{(}\PYG{l+m+mi}{0}\PYG{p}{,}\PYG{l+m+mi}{1}\PYG{p}{,}\PYG{n}{np}\PYG{o}{.}\PYG{n}{size}\PYG{p}{(}\PYG{n}{tau}\PYG{p}{)}\PYG{p}{)}\PYG{p}{)}\PYG{p}{)}

\PYG{k}{for} \PYG{n}{i} \PYG{o+ow}{in} \PYG{n+nb}{range}\PYG{p}{(}\PYG{l+m+mi}{0}\PYG{p}{,}\PYG{n}{np}\PYG{o}{.}\PYG{n}{size}\PYG{p}{(}\PYG{n}{tau}\PYG{p}{)}\PYG{p}{)}\PYG{p}{:}     
    \PYG{n}{c}\PYG{o}{=}\PYG{n+nb}{next}\PYG{p}{(}\PYG{n}{color}\PYG{p}{)}
    \PYG{n}{n}\PYG{o}{=} \PYG{n}{J} \PYG{o}{/} \PYG{n}{G} \PYG{o}{*} \PYG{n}{np}\PYG{o}{.}\PYG{n}{exp}\PYG{p}{(}\PYG{o}{\PYGZhy{}}\PYG{n}{L}\PYG{o}{/}\PYG{n}{G}\PYG{o}{/}\PYG{n}{tau}\PYG{p}{[}\PYG{n}{i}\PYG{p}{]}\PYG{p}{)}
    \PYG{n}{axes}\PYG{o}{.}\PYG{n}{plot}\PYG{p}{(}\PYG{n}{L}\PYG{p}{,}\PYG{n}{n}\PYG{p}{,} \PYG{n}{marker}\PYG{o}{=}\PYG{l+s+s1}{\PYGZsq{}}\PYG{l+s+s1}{ }\PYG{l+s+s1}{\PYGZsq{}} \PYG{p}{,} \PYG{n}{c}\PYG{o}{=}\PYG{n}{c}\PYG{p}{)}    
    
\PYG{n}{plt}\PYG{o}{.}\PYG{n}{title}\PYG{p}{(}\PYG{l+s+s1}{\PYGZsq{}}\PYG{l+s+s1}{n(L) vs. residence time}\PYG{l+s+s1}{\PYGZsq{}}\PYG{p}{,} \PYG{n}{fontsize}\PYG{o}{=}\PYG{l+m+mi}{18}\PYG{p}{)}\PYG{p}{;}
\PYG{n}{axes}\PYG{o}{.}\PYG{n}{set\PYGZus{}xlabel}\PYG{p}{(}\PYG{l+s+s1}{\PYGZsq{}}\PYG{l+s+s1}{\PYGZdl{}L\PYGZdl{}}\PYG{l+s+s1}{\PYGZsq{}}\PYG{p}{,} \PYG{n}{fontsize}\PYG{o}{=}\PYG{l+m+mi}{18}\PYG{p}{)}\PYG{p}{;}
\PYG{n}{axes}\PYG{o}{.}\PYG{n}{set\PYGZus{}ylabel}\PYG{p}{(}\PYG{l+s+s1}{\PYGZsq{}}\PYG{l+s+s1}{\PYGZdl{}n\PYGZdl{}}\PYG{l+s+s1}{\PYGZsq{}}\PYG{p}{,}\PYG{n}{fontsize}\PYG{o}{=}\PYG{l+m+mi}{18}\PYG{p}{)}\PYG{p}{;}

\PYG{c+c1}{\PYGZsh{}\PYGZsh{}\PYGZsh{}\PYGZsh{}\PYGZsh{}\PYGZsh{}\PYGZsh{}\PYGZsh{}\PYGZsh{}\PYGZsh{}\PYGZsh{}\PYGZsh{}\PYGZsh{}\PYGZsh{}\PYGZsh{}\PYGZsh{}\PYGZsh{}\PYGZsh{}\PYGZsh{}\PYGZsh{}\PYGZsh{}\PYGZsh{}\PYGZsh{}\PYGZsh{}\PYGZsh{}\PYGZsh{}\PYGZsh{}\PYGZsh{}\PYGZsh{}\PYGZsh{}\PYGZsh{}\PYGZsh{}\PYGZsh{}\PYGZsh{}\PYGZsh{}\PYGZsh{}\PYGZsh{}\PYGZsh{}\PYGZsh{}\PYGZsh{}\PYGZsh{}\PYGZsh{}\PYGZsh{}\PYGZsh{}\PYGZsh{}\PYGZsh{}\PYGZsh{}\PYGZsh{}\PYGZsh{}\PYGZsh{}\PYGZsh{}\PYGZsh{}\PYGZsh{}\PYGZsh{}\PYGZsh{}\PYGZsh{}\PYGZsh{}}
\PYG{c+c1}{\PYGZsh{} Plotting}
\PYG{n}{figure}\PYG{o}{=}\PYG{n}{plt}\PYG{o}{.}\PYG{n}{figure}\PYG{p}{(}\PYG{p}{)}
\PYG{n}{axes} \PYG{o}{=} \PYG{n}{figure}\PYG{o}{.}\PYG{n}{add\PYGZus{}axes}\PYG{p}{(}\PYG{p}{[}\PYG{l+m+mf}{0.1}\PYG{p}{,}\PYG{l+m+mf}{0.1}\PYG{p}{,}\PYG{l+m+mf}{1.2}\PYG{p}{,}\PYG{l+m+mf}{1.2}\PYG{p}{]}\PYG{p}{)}
\PYG{n}{plt}\PYG{o}{.}\PYG{n}{xticks}\PYG{p}{(}\PYG{n}{fontsize}\PYG{o}{=}\PYG{l+m+mi}{14}\PYG{p}{)}
\PYG{n}{plt}\PYG{o}{.}\PYG{n}{yticks}\PYG{p}{(}\PYG{n}{fontsize}\PYG{o}{=}\PYG{l+m+mi}{14}\PYG{p}{)}

\PYG{c+c1}{\PYGZsh{} Data}
\PYG{n}{G}\PYG{o}{=}\PYG{n}{np}\PYG{o}{.}\PYG{n}{array}\PYG{p}{(}\PYG{p}{[}\PYG{l+m+mi}{1}\PYG{p}{,} \PYG{l+m+mi}{2}\PYG{p}{,} \PYG{l+m+mi}{5}\PYG{p}{,} \PYG{l+m+mi}{10}\PYG{p}{,} \PYG{l+m+mi}{50}\PYG{p}{,} \PYG{l+m+mi}{100}\PYG{p}{,} \PYG{l+m+mi}{500}\PYG{p}{]}\PYG{p}{)}\PYG{p}{;}
\PYG{n}{J}\PYG{o}{=}\PYG{l+m+mi}{1}\PYG{p}{;}
\PYG{n}{tau}\PYG{o}{=}\PYG{l+m+mi}{5}\PYG{p}{;}
\PYG{n}{N} \PYG{o}{=} \PYG{l+m+mi}{500} \PYG{c+c1}{\PYGZsh{}number of points}
\PYG{n}{L} \PYG{o}{=} \PYG{n}{np}\PYG{o}{.}\PYG{n}{linspace}\PYG{p}{(}\PYG{l+m+mi}{0}\PYG{p}{,}\PYG{l+m+mi}{50}\PYG{p}{,} \PYG{n}{N}\PYG{p}{)}\PYG{p}{;}

\PYG{n}{color}\PYG{o}{=}\PYG{n+nb}{iter}\PYG{p}{(}\PYG{n}{cm}\PYG{o}{.}\PYG{n}{coolwarm}\PYG{p}{(}\PYG{n}{np}\PYG{o}{.}\PYG{n}{linspace}\PYG{p}{(}\PYG{l+m+mi}{0}\PYG{p}{,}\PYG{l+m+mi}{1}\PYG{p}{,}\PYG{n}{np}\PYG{o}{.}\PYG{n}{size}\PYG{p}{(}\PYG{n}{G}\PYG{p}{)}\PYG{p}{)}\PYG{p}{)}\PYG{p}{)}\PYG{p}{;}

\PYG{k}{for} \PYG{n}{i} \PYG{o+ow}{in} \PYG{n+nb}{range}\PYG{p}{(}\PYG{l+m+mi}{0}\PYG{p}{,}\PYG{n}{np}\PYG{o}{.}\PYG{n}{size}\PYG{p}{(}\PYG{n}{G}\PYG{p}{)}\PYG{p}{)}\PYG{p}{:}     
    \PYG{n}{c}\PYG{o}{=}\PYG{n+nb}{next}\PYG{p}{(}\PYG{n}{color}\PYG{p}{)}
    \PYG{n}{n}\PYG{o}{=} \PYG{n}{J} \PYG{o}{/} \PYG{n}{G}\PYG{p}{[}\PYG{n}{i}\PYG{p}{]} \PYG{o}{*} \PYG{n}{np}\PYG{o}{.}\PYG{n}{exp}\PYG{p}{(}\PYG{o}{\PYGZhy{}}\PYG{n}{L}\PYG{o}{/}\PYG{n}{G}\PYG{p}{[}\PYG{n}{i}\PYG{p}{]}\PYG{o}{/}\PYG{n}{tau}\PYG{p}{)}
    \PYG{n}{axes}\PYG{o}{.}\PYG{n}{plot}\PYG{p}{(}\PYG{n}{L}\PYG{p}{,}\PYG{n}{n}\PYG{p}{,} \PYG{n}{marker}\PYG{o}{=}\PYG{l+s+s1}{\PYGZsq{}}\PYG{l+s+s1}{ }\PYG{l+s+s1}{\PYGZsq{}} \PYG{p}{,} \PYG{n}{c}\PYG{o}{=}\PYG{n}{c}\PYG{p}{)}    
    
\PYG{n}{plt}\PYG{o}{.}\PYG{n}{title}\PYG{p}{(}\PYG{l+s+s1}{\PYGZsq{}}\PYG{l+s+s1}{n(L) vs. growth rate}\PYG{l+s+s1}{\PYGZsq{}}\PYG{p}{,} \PYG{n}{fontsize}\PYG{o}{=}\PYG{l+m+mi}{18}\PYG{p}{)}\PYG{p}{;}
\PYG{n}{axes}\PYG{o}{.}\PYG{n}{set\PYGZus{}xlabel}\PYG{p}{(}\PYG{l+s+s1}{\PYGZsq{}}\PYG{l+s+s1}{\PYGZdl{}L\PYGZdl{}}\PYG{l+s+s1}{\PYGZsq{}}\PYG{p}{,} \PYG{n}{fontsize}\PYG{o}{=}\PYG{l+m+mi}{18}\PYG{p}{)}\PYG{p}{;}
\PYG{n}{axes}\PYG{o}{.}\PYG{n}{set\PYGZus{}ylabel}\PYG{p}{(}\PYG{l+s+s1}{\PYGZsq{}}\PYG{l+s+s1}{\PYGZdl{}n\PYGZdl{}}\PYG{l+s+s1}{\PYGZsq{}}\PYG{p}{,}\PYG{n}{fontsize}\PYG{o}{=}\PYG{l+m+mi}{18}\PYG{p}{)}\PYG{p}{;}


\PYG{c+c1}{\PYGZsh{}\PYGZsh{}\PYGZsh{}\PYGZsh{}\PYGZsh{}\PYGZsh{}\PYGZsh{}\PYGZsh{}\PYGZsh{}\PYGZsh{}\PYGZsh{}\PYGZsh{}\PYGZsh{}\PYGZsh{}\PYGZsh{}\PYGZsh{}\PYGZsh{}\PYGZsh{}\PYGZsh{}\PYGZsh{}\PYGZsh{}\PYGZsh{}\PYGZsh{}\PYGZsh{}\PYGZsh{}\PYGZsh{}\PYGZsh{}\PYGZsh{}\PYGZsh{}\PYGZsh{}\PYGZsh{}\PYGZsh{}\PYGZsh{}\PYGZsh{}\PYGZsh{}\PYGZsh{}\PYGZsh{}\PYGZsh{}\PYGZsh{}\PYGZsh{}\PYGZsh{}\PYGZsh{}\PYGZsh{}\PYGZsh{}\PYGZsh{}\PYGZsh{}\PYGZsh{}\PYGZsh{}\PYGZsh{}\PYGZsh{}\PYGZsh{}\PYGZsh{}\PYGZsh{}\PYGZsh{}\PYGZsh{}\PYGZsh{}\PYGZsh{}}
\PYG{c+c1}{\PYGZsh{} Plotting}
\PYG{n}{figure}\PYG{o}{=}\PYG{n}{plt}\PYG{o}{.}\PYG{n}{figure}\PYG{p}{(}\PYG{p}{)}
\PYG{n}{axes} \PYG{o}{=} \PYG{n}{figure}\PYG{o}{.}\PYG{n}{add\PYGZus{}axes}\PYG{p}{(}\PYG{p}{[}\PYG{l+m+mf}{0.1}\PYG{p}{,}\PYG{l+m+mf}{0.1}\PYG{p}{,}\PYG{l+m+mf}{1.2}\PYG{p}{,}\PYG{l+m+mf}{1.2}\PYG{p}{]}\PYG{p}{)}
\PYG{n}{plt}\PYG{o}{.}\PYG{n}{Lticks}\PYG{p}{(}\PYG{n}{fontsize}\PYG{o}{=}\PYG{l+m+mi}{14}\PYG{p}{)}
\PYG{n}{plt}\PYG{o}{.}\PYG{n}{yticks}\PYG{p}{(}\PYG{n}{fontsize}\PYG{o}{=}\PYG{l+m+mi}{14}\PYG{p}{)}

\PYG{c+c1}{\PYGZsh{} Data}
\PYG{n}{G}\PYG{o}{=}\PYG{l+m+mi}{1}\PYG{p}{;}
\PYG{n}{J}\PYG{o}{=}\PYG{n}{np}\PYG{o}{.}\PYG{n}{array}\PYG{p}{(}\PYG{p}{[}\PYG{l+m+mi}{1}\PYG{p}{,} \PYG{l+m+mi}{2}\PYG{p}{,} \PYG{l+m+mi}{5}\PYG{p}{,} \PYG{l+m+mi}{10}\PYG{p}{,} \PYG{l+m+mi}{50}\PYG{p}{,} \PYG{l+m+mi}{100}\PYG{p}{,} \PYG{l+m+mi}{500}\PYG{p}{]}\PYG{p}{)}\PYG{p}{;}
\PYG{n}{tau}\PYG{o}{=}\PYG{l+m+mi}{5}\PYG{p}{;}
\PYG{n}{N} \PYG{o}{=} \PYG{l+m+mi}{500} \PYG{c+c1}{\PYGZsh{}number of points}
\PYG{n}{L} \PYG{o}{=} \PYG{n}{np}\PYG{o}{.}\PYG{n}{linspace}\PYG{p}{(}\PYG{l+m+mi}{0}\PYG{p}{,}\PYG{l+m+mi}{50}\PYG{p}{,} \PYG{n}{N}\PYG{p}{)}\PYG{p}{;}

\PYG{n}{color}\PYG{o}{=}\PYG{n+nb}{iter}\PYG{p}{(}\PYG{n}{cm}\PYG{o}{.}\PYG{n}{coolwarm}\PYG{p}{(}\PYG{n}{np}\PYG{o}{.}\PYG{n}{linspace}\PYG{p}{(}\PYG{l+m+mi}{0}\PYG{p}{,}\PYG{l+m+mi}{1}\PYG{p}{,}\PYG{n}{np}\PYG{o}{.}\PYG{n}{size}\PYG{p}{(}\PYG{n}{J}\PYG{p}{)}\PYG{p}{)}\PYG{p}{)}\PYG{p}{)}\PYG{p}{;}

\PYG{k}{for} \PYG{n}{i} \PYG{o+ow}{in} \PYG{n+nb}{range}\PYG{p}{(}\PYG{l+m+mi}{0}\PYG{p}{,}\PYG{n}{np}\PYG{o}{.}\PYG{n}{size}\PYG{p}{(}\PYG{n}{J}\PYG{p}{)}\PYG{p}{)}\PYG{p}{:}     
    \PYG{n}{c}\PYG{o}{=}\PYG{n+nb}{next}\PYG{p}{(}\PYG{n}{color}\PYG{p}{)}
    \PYG{n}{n}\PYG{o}{=} \PYG{n}{J}\PYG{p}{[}\PYG{n}{i}\PYG{p}{]} \PYG{o}{/} \PYG{n}{G} \PYG{o}{*} \PYG{n}{np}\PYG{o}{.}\PYG{n}{exp}\PYG{p}{(}\PYG{o}{\PYGZhy{}}\PYG{n}{L}\PYG{o}{/}\PYG{n}{G}\PYG{o}{/}\PYG{n}{tau}\PYG{p}{)}
    \PYG{n}{axes}\PYG{o}{.}\PYG{n}{plot}\PYG{p}{(}\PYG{n}{L}\PYG{p}{,}\PYG{n}{n}\PYG{p}{,} \PYG{n}{marker}\PYG{o}{=}\PYG{l+s+s1}{\PYGZsq{}}\PYG{l+s+s1}{ }\PYG{l+s+s1}{\PYGZsq{}} \PYG{p}{,} \PYG{n}{c}\PYG{o}{=}\PYG{n}{c}\PYG{p}{)}    
    
\PYG{n}{plt}\PYG{o}{.}\PYG{n}{title}\PYG{p}{(}\PYG{l+s+s1}{\PYGZsq{}}\PYG{l+s+s1}{n(L) vs. nucleation rate}\PYG{l+s+s1}{\PYGZsq{}}\PYG{p}{,} \PYG{n}{fontsize}\PYG{o}{=}\PYG{l+m+mi}{18}\PYG{p}{)}\PYG{p}{;}
\PYG{n}{axes}\PYG{o}{.}\PYG{n}{set\PYGZus{}xlabel}\PYG{p}{(}\PYG{l+s+s1}{\PYGZsq{}}\PYG{l+s+s1}{\PYGZdl{}L\PYGZdl{}}\PYG{l+s+s1}{\PYGZsq{}}\PYG{p}{,} \PYG{n}{fontsize}\PYG{o}{=}\PYG{l+m+mi}{18}\PYG{p}{)}\PYG{p}{;}
\PYG{n}{axes}\PYG{o}{.}\PYG{n}{set\PYGZus{}ylabel}\PYG{p}{(}\PYG{l+s+s1}{\PYGZsq{}}\PYG{l+s+s1}{\PYGZdl{}n\PYGZdl{}}\PYG{l+s+s1}{\PYGZsq{}}\PYG{p}{,}\PYG{n}{fontsize}\PYG{o}{=}\PYG{l+m+mi}{18}\PYG{p}{)}\PYG{p}{;}
\end{sphinxVerbatim}

\end{sphinxuseclass}\end{sphinxVerbatimInput}
\begin{sphinxVerbatimOutput}

\begin{sphinxuseclass}{cell_output}
\begin{sphinxVerbatim}[commandchars=\\\{\}]
\PYG{g+gt}{\PYGZhy{}\PYGZhy{}\PYGZhy{}\PYGZhy{}\PYGZhy{}\PYGZhy{}\PYGZhy{}\PYGZhy{}\PYGZhy{}\PYGZhy{}\PYGZhy{}\PYGZhy{}\PYGZhy{}\PYGZhy{}\PYGZhy{}\PYGZhy{}\PYGZhy{}\PYGZhy{}\PYGZhy{}\PYGZhy{}\PYGZhy{}\PYGZhy{}\PYGZhy{}\PYGZhy{}\PYGZhy{}\PYGZhy{}\PYGZhy{}\PYGZhy{}\PYGZhy{}\PYGZhy{}\PYGZhy{}\PYGZhy{}\PYGZhy{}\PYGZhy{}\PYGZhy{}\PYGZhy{}\PYGZhy{}\PYGZhy{}\PYGZhy{}\PYGZhy{}\PYGZhy{}\PYGZhy{}\PYGZhy{}\PYGZhy{}\PYGZhy{}\PYGZhy{}\PYGZhy{}\PYGZhy{}\PYGZhy{}\PYGZhy{}\PYGZhy{}\PYGZhy{}\PYGZhy{}\PYGZhy{}\PYGZhy{}\PYGZhy{}\PYGZhy{}\PYGZhy{}\PYGZhy{}\PYGZhy{}\PYGZhy{}\PYGZhy{}\PYGZhy{}\PYGZhy{}\PYGZhy{}\PYGZhy{}\PYGZhy{}\PYGZhy{}\PYGZhy{}\PYGZhy{}\PYGZhy{}\PYGZhy{}\PYGZhy{}\PYGZhy{}\PYGZhy{}}
\PYG{n+ne}{ModuleNotFoundError}\PYG{g+gWhitespace}{                       }Traceback (most recent call last)
\PYG{n}{Cell} \PYG{n}{In}\PYG{p}{[}\PYG{l+m+mi}{1}\PYG{p}{]}\PYG{p}{,} \PYG{n}{line} \PYG{l+m+mi}{1}
\PYG{n+ne}{\PYGZhy{}\PYGZhy{}\PYGZhy{}\PYGZhy{}\PYGZgt{} }\PYG{l+m+mi}{1} \PYG{k+kn}{import}\PYG{+w}{ }\PYG{n+nn}{numpy}\PYG{+w}{ }\PYG{k}{as}\PYG{+w}{ }\PYG{n+nn}{np}
\PYG{g+gWhitespace}{      }\PYG{l+m+mi}{2} \PYG{k+kn}{import}\PYG{+w}{ }\PYG{n+nn}{matplotlib}\PYG{n+nn}{.}\PYG{n+nn}{pyplot}\PYG{+w}{ }\PYG{k}{as}\PYG{+w}{ }\PYG{n+nn}{plt} 
\PYG{g+gWhitespace}{      }\PYG{l+m+mi}{3} \PYG{k+kn}{from}\PYG{+w}{ }\PYG{n+nn}{matplotlib}\PYG{n+nn}{.}\PYG{n+nn}{pyplot}\PYG{+w}{ }\PYG{k+kn}{import} \PYG{n}{cm}

\PYG{n+ne}{ModuleNotFoundError}: No module named \PYGZsq{}numpy\PYGZsq{}
\end{sphinxVerbatim}

\end{sphinxuseclass}\end{sphinxVerbatimOutput}

\end{sphinxuseclass}

\subsection{10.4 Moments of the PSD in a MSMPR crystallizer}
\label{\detokenize{MSMPR_crystallizer:moments-of-the-psd-in-a-msmpr-crystallizer}}
\sphinxAtStartPar
The number population density \(n(L)\) can be used to derive other properties of the population of particles through the calculation of the moments of the distribution.
The \(i^{th}\) order moment of the number population density is defined as:
\begin{equation}\label{equation:MSMPR_crystallizer:MSMPReq12}
\begin{split}
m_i=\int_0^{\infty}L^indL
\end{split}
\end{equation}
\sphinxAtStartPar
Based on this definition we note that the moments from order 0 to 3 have a well defined physical meaning:
\begin{itemize}
\item {} 
\sphinxAtStartPar
\(m_0=\int_0^{\infty}L^0ndL=N_L\) represents the number of particles per unit volume of liquid phase  with size smaller or equal than size L. For \(L\rightarrow{\infty}\) \(N_L\) tends to the total number of particles per unit volume of liquid phase in the crystallizer:

\end{itemize}
\begin{equation}\label{equation:MSMPR_crystallizer:MSMPReq13}
\begin{split}
N_T=n_0G\tau
\end{split}
\end{equation}\begin{itemize}
\item {} 
\sphinxAtStartPar
\(m_1=\int_0^{\infty}L^1ndL=\mathcal{L}\) represents the total length of particles per unit volume of liquid phase with size smaller or equal than size \(L\). For \(L\rightarrow{\infty}\), \(\mathcal{L}\) tends to the total length of particles per unit volume of liquid phase in the crystallizer:

\end{itemize}
\begin{equation}\label{equation:MSMPR_crystallizer:MSMPReq14}
\begin{split}
\mathcal{L}_T=n_0(G\tau)^2
\end{split}
\end{equation}\begin{itemize}
\item {} 
\sphinxAtStartPar
\(m_2=\int_0^{\infty}L^2ndL\propto{A}\) is proportional to the total area of particles per unit volume of liquid phase with size smaller or equal than size \(L\). For \(L\rightarrow{\infty}\), \(A\) tends to the total area of the particles per unit volume of liquid phase contained in the crystallizer:

\end{itemize}
\begin{equation}\label{equation:MSMPR_crystallizer:MSMPReq15}
\begin{split}
A_T=2{\beta}n_0(G\tau)^3
\end{split}
\end{equation}
\sphinxAtStartPar
where \(\beta\) is the surface shape factor of the particles.
\begin{itemize}
\item {} 
\sphinxAtStartPar
\(m_3=\int_0^{\infty}L^3ndL\propto{V}\) is proportional to the total volume of particles with size smaller or equal than size \(L\). For \(L\rightarrow{\infty}\), \(V\) tends to the total volume of the particles per unit volume contained in the crystallizer:

\end{itemize}
\begin{equation}\label{equation:MSMPR_crystallizer:MSMPReq16}
\begin{split}
V_T=6{\alpha}n_0(G\tau)^4
\end{split}
\end{equation}
\sphinxAtStartPar
where \(\alpha\) is the volumetric shape factor of the particles.

\sphinxAtStartPar
It should be noted that the total volume of the particles per unit volume of liquid phase allows to compute the total mass of particles per unit volume of liquid phase in the system:
\begin{equation}\label{equation:MSMPR_crystallizer:MSMPReq17}
\begin{split}
M_T=\rho_c{V_T}=6\rho_c{\alpha}n_0(G\tau)^4
\end{split}
\end{equation}
\sphinxAtStartPar
where \(\rho_c\) is the density of the crystalline phase.

\sphinxAtStartPar
This expression for \(M_T\) allows to couple the population balance and the gobal material balance in an MSMPR crystallizer since \(M_T=C_{in}-C_{out}\), where \(C_{in}\) and \(C_{out}\) are the inlet and outlet concentrations of solute in the liquid phase.


\subsection{10.5 Adimensional Mass distribution and dominant crystal size}
\label{\detokenize{MSMPR_crystallizer:adimensional-mass-distribution-and-dominant-crystal-size}}
\sphinxAtStartPar
In order to obtain a general relationship between the crystallization process and an average indicator of the crystalline particles size obtained in an MSMPR crystallizer it is convenient to write the cumulative mass distribution in a dimensionless form. To this aim we introduce the dimensionless length \(x=L/G\tau\), where \(G\) is the growth rate and \(\tau\) is the residence time.

\sphinxAtStartPar
The cumulative mass distribution can thus be written as:
\begin{equation}\label{equation:MSMPR_crystallizer:MSMPReq18}
\begin{split}
M(x)=\frac{M}{M_T}=\frac{\int_0^{x}x^3\exp(-x)dx}{\int_0^{\infty}x^3\exp(-x)dx}=1-\left(1+x+\frac{1}{2}x^2+\frac{1}{6}x^3\right)\exp(-x)
\end{split}
\end{equation}
\sphinxAtStartPar
The corresponding mass population density is computed differentiating \(M(x)\)with respect to \(x\):
\begin{equation}\label{equation:MSMPR_crystallizer:MSMPReq19}
\begin{split}
m(x)=\frac{dM(x)}{dx}=\frac{1}{6}x^3\exp(-x)
\end{split}
\end{equation}
\sphinxAtStartPar
The mass\sphinxhyphen{}based population density m(x) has a maximum for \(x=3\), hence the most probable \sphinxhyphen{} often referred to as the modal or dominant particle size can be computed as:
\begin{equation}\label{equation:MSMPR_crystallizer:MSMPReq20}
\begin{split}
L_D=3G\tau
\end{split}
\end{equation}
\begin{sphinxuseclass}{cell}\begin{sphinxVerbatimInput}

\begin{sphinxuseclass}{cell_input}
\begin{sphinxVerbatim}[commandchars=\\\{\}]
\PYG{k+kn}{import}\PYG{+w}{ }\PYG{n+nn}{numpy}\PYG{+w}{ }\PYG{k}{as}\PYG{+w}{ }\PYG{n+nn}{np}
\PYG{k+kn}{import}\PYG{+w}{ }\PYG{n+nn}{matplotlib}\PYG{n+nn}{.}\PYG{n+nn}{pyplot}\PYG{+w}{ }\PYG{k}{as}\PYG{+w}{ }\PYG{n+nn}{plt} 
\PYG{k+kn}{from}\PYG{+w}{ }\PYG{n+nn}{matplotlib}\PYG{n+nn}{.}\PYG{n+nn}{pyplot}\PYG{+w}{ }\PYG{k+kn}{import} \PYG{n}{cm}

\PYG{c+c1}{\PYGZsh{}Plotting}
\PYG{n}{figure}\PYG{o}{=}\PYG{n}{plt}\PYG{o}{.}\PYG{n}{figure}\PYG{p}{(}\PYG{p}{)}
\PYG{n}{axes} \PYG{o}{=} \PYG{n}{figure}\PYG{o}{.}\PYG{n}{add\PYGZus{}axes}\PYG{p}{(}\PYG{p}{[}\PYG{l+m+mf}{0.1}\PYG{p}{,}\PYG{l+m+mf}{0.1}\PYG{p}{,}\PYG{l+m+mf}{1.2}\PYG{p}{,}\PYG{l+m+mf}{1.2}\PYG{p}{]}\PYG{p}{)}
\PYG{n}{plt}\PYG{o}{.}\PYG{n}{xticks}\PYG{p}{(}\PYG{n}{fontsize}\PYG{o}{=}\PYG{l+m+mi}{14}\PYG{p}{)}
\PYG{n}{plt}\PYG{o}{.}\PYG{n}{yticks}\PYG{p}{(}\PYG{n}{fontsize}\PYG{o}{=}\PYG{l+m+mi}{14}\PYG{p}{)}

\PYG{c+c1}{\PYGZsh{} Data}
\PYG{n}{x} \PYG{o}{=} \PYG{n}{np}\PYG{o}{.}\PYG{n}{linspace}\PYG{p}{(}\PYG{l+m+mi}{0}\PYG{p}{,}\PYG{l+m+mi}{20}\PYG{p}{,} \PYG{n}{N}\PYG{p}{)}

\PYG{n}{color}\PYG{o}{=}\PYG{n+nb}{iter}\PYG{p}{(}\PYG{n}{cm}\PYG{o}{.}\PYG{n}{coolwarm}\PYG{p}{(}\PYG{n}{np}\PYG{o}{.}\PYG{n}{linspace}\PYG{p}{(}\PYG{l+m+mi}{0}\PYG{p}{,}\PYG{l+m+mi}{1}\PYG{p}{,}\PYG{l+m+mi}{2}\PYG{p}{)}\PYG{p}{)}\PYG{p}{)}
  
\PYG{n}{c}\PYG{o}{=}\PYG{n+nb}{next}\PYG{p}{(}\PYG{n}{color}\PYG{p}{)}
\PYG{n}{M}\PYG{o}{=}\PYG{l+m+mi}{1}\PYG{o}{\PYGZhy{}}\PYG{p}{(}\PYG{l+m+mi}{1}\PYG{o}{+}\PYG{n}{x}\PYG{o}{+}\PYG{l+m+mf}{0.5}\PYG{o}{*}\PYG{n}{np}\PYG{o}{.}\PYG{n}{power}\PYG{p}{(}\PYG{n}{x}\PYG{p}{,}\PYG{l+m+mi}{2}\PYG{p}{)}\PYG{o}{+}\PYG{l+m+mi}{1}\PYG{o}{/}\PYG{l+m+mi}{6}\PYG{o}{*}\PYG{n}{np}\PYG{o}{.}\PYG{n}{power}\PYG{p}{(}\PYG{n}{x}\PYG{p}{,}\PYG{l+m+mi}{3}\PYG{p}{)}\PYG{p}{)}\PYG{o}{*}\PYG{n}{np}\PYG{o}{.}\PYG{n}{exp}\PYG{p}{(}\PYG{o}{\PYGZhy{}}\PYG{n}{x}\PYG{p}{)}
\PYG{n}{axes}\PYG{o}{.}\PYG{n}{plot}\PYG{p}{(}\PYG{n}{x}\PYG{p}{,}\PYG{n}{M}\PYG{p}{,} \PYG{n}{marker}\PYG{o}{=}\PYG{l+s+s1}{\PYGZsq{}}\PYG{l+s+s1}{ }\PYG{l+s+s1}{\PYGZsq{}} \PYG{p}{,} \PYG{n}{c}\PYG{o}{=}\PYG{n}{c}\PYG{p}{)} 
\PYG{n}{c}\PYG{o}{=}\PYG{n+nb}{next}\PYG{p}{(}\PYG{n}{color}\PYG{p}{)}
\PYG{n}{m}\PYG{o}{=}\PYG{l+m+mi}{1}\PYG{o}{/}\PYG{l+m+mi}{6}\PYG{o}{*}\PYG{n}{np}\PYG{o}{.}\PYG{n}{power}\PYG{p}{(}\PYG{n}{x}\PYG{p}{,}\PYG{l+m+mi}{3}\PYG{p}{)}\PYG{o}{*}\PYG{n}{np}\PYG{o}{.}\PYG{n}{exp}\PYG{p}{(}\PYG{o}{\PYGZhy{}}\PYG{n}{x}\PYG{p}{)}
\PYG{n}{axes}\PYG{o}{.}\PYG{n}{plot}\PYG{p}{(}\PYG{n}{x}\PYG{p}{,}\PYG{n}{m}\PYG{p}{,} \PYG{n}{marker}\PYG{o}{=}\PYG{l+s+s1}{\PYGZsq{}}\PYG{l+s+s1}{ }\PYG{l+s+s1}{\PYGZsq{}} \PYG{p}{,} \PYG{n}{c}\PYG{o}{=}\PYG{n}{c}\PYG{p}{)}    
\PYG{n}{axes}\PYG{o}{.}\PYG{n}{plot}\PYG{p}{(}\PYG{n}{np}\PYG{o}{.}\PYG{n}{array}\PYG{p}{(}\PYG{p}{[}\PYG{l+m+mi}{3}\PYG{p}{,}\PYG{l+m+mi}{3}\PYG{p}{]}\PYG{p}{)}\PYG{p}{,}\PYG{n}{np}\PYG{o}{.}\PYG{n}{array}\PYG{p}{(}\PYG{p}{[}\PYG{l+m+mi}{0}\PYG{p}{,} \PYG{l+m+mi}{1}\PYG{o}{/}\PYG{l+m+mi}{6}\PYG{o}{*}\PYG{n}{np}\PYG{o}{.}\PYG{n}{power}\PYG{p}{(}\PYG{l+m+mi}{3}\PYG{p}{,}\PYG{l+m+mi}{3}\PYG{p}{)}\PYG{o}{*}\PYG{n}{np}\PYG{o}{.}\PYG{n}{exp}\PYG{p}{(}\PYG{o}{\PYGZhy{}}\PYG{l+m+mi}{3}\PYG{p}{)}\PYG{p}{]}\PYG{p}{)}\PYG{p}{,}\PYG{l+s+s1}{\PYGZsq{}}\PYG{l+s+s1}{\PYGZhy{}\PYGZhy{}}\PYG{l+s+s1}{\PYGZsq{}}\PYG{p}{,}\PYG{p}{)}

\PYG{n}{axes}\PYG{o}{.}\PYG{n}{set\PYGZus{}xlabel}\PYG{p}{(}\PYG{l+s+s1}{\PYGZsq{}}\PYG{l+s+s1}{\PYGZdl{}x\PYGZdl{}}\PYG{l+s+s1}{\PYGZsq{}}\PYG{p}{,} \PYG{n}{fontsize}\PYG{o}{=}\PYG{l+m+mi}{18}\PYG{p}{)}\PYG{p}{;}
\PYG{n}{axes}\PYG{o}{.}\PYG{n}{set\PYGZus{}ylabel}\PYG{p}{(}\PYG{l+s+s1}{\PYGZsq{}}\PYG{l+s+s1}{\PYGZdl{}M(x)/m(x)\PYGZdl{}}\PYG{l+s+s1}{\PYGZsq{}}\PYG{p}{,}\PYG{n}{fontsize}\PYG{o}{=}\PYG{l+m+mi}{18}\PYG{p}{)}\PYG{p}{;}
\PYG{n}{axes}\PYG{o}{.}\PYG{n}{legend}\PYG{p}{(}\PYG{p}{[}\PYG{l+s+s1}{\PYGZsq{}}\PYG{l+s+s1}{M(x)}\PYG{l+s+s1}{\PYGZsq{}}\PYG{p}{,}\PYG{l+s+s1}{\PYGZsq{}}\PYG{l+s+s1}{m(x)}\PYG{l+s+s1}{\PYGZsq{}}\PYG{p}{]}\PYG{p}{,} \PYG{n}{fontsize}\PYG{o}{=}\PYG{l+m+mi}{18}\PYG{p}{)}\PYG{p}{;}
\end{sphinxVerbatim}

\end{sphinxuseclass}\end{sphinxVerbatimInput}
\begin{sphinxVerbatimOutput}

\begin{sphinxuseclass}{cell_output}
\noindent\sphinxincludegraphics{{c8a0374e9fc935dc9057894c4ae75a3b0f244034899660da8f410e6f51506708}.png}

\end{sphinxuseclass}\end{sphinxVerbatimOutput}

\end{sphinxuseclass}

\subsection{Contributions}
\label{\detokenize{MSMPR_crystallizer:contributions}}
\sphinxAtStartPar
Guilherme Pizarro Werner, 23 February 2021

\sphinxAtStartPar
Ahmad Eblagh, 13 March 2021

\sphinxstepscope


\section{Week 2 \sphinxhyphen{} Exercise}
\label{\detokenize{Exercise_W2_2025:week-2-exercise}}\label{\detokenize{Exercise_W2_2025::doc}}

\subsection{Problem Statement}
\label{\detokenize{Exercise_W2_2025:problem-statement}}
\sphinxAtStartPar
A perfectly mixed gas permeation module is used to separate carbon dioxide from nitrogen using a poly
(2,6\sphinxhyphen{}dimethylphenylene oxide) membrane	(\(P_{CO_2}=0.034\times10^{-13}\) \(P_{N_2}=0.089\times10^{-14} [cm^3 (STP) / cm s Pa]\)). The process carried out at a temperature of 25\(^oC\).
The feed flowrate is 20.0 mol\% carbon dioxide.
The module has 15.0 \(m^2\) of membrane. The module is operated with a retentate pressure of 5.5 atm
and a permeate pressure of 1.01 atm.

\sphinxAtStartPar
The permeate is enriched to 38.0 mol\% in carbon dioxide. The membrane thickness is \(\ell = 2.0 \times 10^{-4}\) cm.
\begin{itemize}
\item {} 
\sphinxAtStartPar
Draw a scheme of the unit clearly labelling all the streams and relevant variables.

\item {} 
\sphinxAtStartPar
Completely characterise the unit, by computing feed, permeate and retentate flow rates, cut, and retentate composition.

\end{itemize}


\subsection{Solution trace}
\label{\detokenize{Exercise_W2_2025:solution-trace}}
\begin{figure}[htbp]
\centering
\capstart

\noindent\sphinxincludegraphics[height=300\sphinxpxdimen]{{scheme}.png}
\caption{Schematics of a membrane process with all variables defined.}\label{\detokenize{Exercise_W2_2025:scheme}}\end{figure}



\sphinxAtStartPar
\sphinxstylestrong{Unknowns}:
\(x\),\(F_{p}\),\(F_{in}\), \(F_{r}\), \(\theta\)

\sphinxAtStartPar
We can start computing the ideal separation factor and the composition in the retentate:
\begin{equation*}
\begin{split}
\alpha=P_{CO_2}/P_{N_2}
\end{split}
\end{equation*}
\sphinxAtStartPar
and then using the rate transfer equation to computing \(x\):
\begin{equation*}
\begin{split}
x=\frac{y\left[1+\frac{p_p}{p_r}(1-y)\left(\alpha_{AB}-1\right)\right]}{\alpha_{AB}-\left(\alpha_{AB}-1\right)y}
\end{split}
\end{equation*}
\sphinxAtStartPar
Then we can compute the cut:
\begin{equation*}
\begin{split}
\theta=\frac{z-x}{y-x}
\end{split}
\end{equation*}
\sphinxAtStartPar
and the permeate flowrate from the flux equation:
\begin{equation*}
\begin{split}
Fp=\frac{A\,P_{CO_2}\rho_{STP}}{l y}\left(xP_r-yP_p\right)
\end{split}
\end{equation*}
\sphinxAtStartPar
where \(\rho_{STP}\) is the molar density of an ideal gas at standard T and P.

\sphinxAtStartPar
Finally, the retentate flowrate can be computed as:
\begin{equation*}
\begin{split}
F_r=F_{in}-F_p
\end{split}
\end{equation*}
\begin{sphinxuseclass}{cell}
\begin{sphinxuseclass}{tag_hide_input}\begin{sphinxVerbatimInput}

\begin{sphinxuseclass}{cell_input}
\begin{sphinxVerbatim}[commandchars=\\\{\}]
\PYG{c+c1}{\PYGZsh{}\PYGZsh{} Solution trace}
\PYG{c+c1}{\PYGZsh{} Variables}
\PYG{n}{A}\PYG{o}{=}\PYG{l+m+mi}{15} \PYG{o}{*} \PYG{l+m+mf}{1E4} \PYG{c+c1}{\PYGZsh{} [m\PYGZca{}2 * cm\PYGZca{}2 / m\PYGZca{}2]}
\PYG{n}{P\PYGZus{}r}\PYG{o}{=}\PYG{l+m+mf}{5.5} \PYG{o}{*} \PYG{l+m+mi}{101325} \PYG{c+c1}{\PYGZsh{} [atm] * [Pa / atm]}
\PYG{n}{P\PYGZus{}p}\PYG{o}{=}\PYG{l+m+mf}{1.01} \PYG{o}{*} \PYG{l+m+mi}{101325} \PYG{c+c1}{\PYGZsh{} [atm] * [Pa / atm] }
\PYG{n}{l}\PYG{o}{=}\PYG{l+m+mf}{2.0}\PYG{o}{*}\PYG{l+m+mf}{1E\PYGZhy{}4} \PYG{c+c1}{\PYGZsh{} cm}
\PYG{n}{rho\PYGZus{}STP}\PYG{o}{=}\PYG{l+m+mi}{1}\PYG{o}{/}\PYG{l+m+mf}{22.4E3} \PYG{c+c1}{\PYGZsh{} / [mol/cm3(stp)]}
\PYG{n}{PCO2}\PYG{o}{=}\PYG{l+m+mf}{0.034E\PYGZhy{}13} \PYG{o}{*} \PYG{l+m+mi}{3600} \PYG{c+c1}{\PYGZsh{} [cm3(STP) * cm/cm2/s/Pa] * [s/h]}
\PYG{n}{PN2}\PYG{o}{=}\PYG{l+m+mf}{0.089E\PYGZhy{}14} \PYG{o}{*} \PYG{l+m+mi}{3600}\PYG{c+c1}{\PYGZsh{}  [cm3(STP) * cm/cm2/s/Pa] * [s/h]}
\PYG{c+c1}{\PYGZsh{} Permeate Composition}
\PYG{n}{y}\PYG{o}{=}\PYG{l+m+mf}{0.38}
\PYG{c+c1}{\PYGZsh{} Feed Composition}
\PYG{n}{z} \PYG{o}{=} \PYG{l+m+mf}{0.20}

\PYG{c+c1}{\PYGZsh{} 1. compute the retentate composition: }
\PYG{n}{alpha}\PYG{o}{=}\PYG{n}{PCO2}\PYG{o}{/}\PYG{n}{PN2}
\PYG{n}{x}\PYG{o}{=} \PYG{n}{y} \PYG{o}{*} \PYG{p}{(}\PYG{l+m+mi}{1} \PYG{o}{+} \PYG{n}{P\PYGZus{}p}\PYG{o}{/}\PYG{n}{P\PYGZus{}r} \PYG{o}{*} \PYG{p}{(}\PYG{l+m+mi}{1}\PYG{o}{\PYGZhy{}}\PYG{n}{y}\PYG{p}{)} \PYG{o}{*} \PYG{p}{(}\PYG{n}{alpha} \PYG{o}{\PYGZhy{}} \PYG{l+m+mi}{1}\PYG{p}{)}\PYG{p}{)} \PYG{o}{/} \PYG{p}{(}\PYG{n}{alpha} \PYG{o}{\PYGZhy{}} \PYG{p}{(}\PYG{n}{alpha} \PYG{o}{\PYGZhy{}} \PYG{l+m+mi}{1}\PYG{p}{)} \PYG{o}{*} \PYG{n}{y}\PYG{p}{)}

\PYG{c+c1}{\PYGZsh{} 2. compute the cut: }
\PYG{n}{theta}\PYG{o}{=}\PYG{p}{(}\PYG{n}{z}\PYG{o}{\PYGZhy{}}\PYG{n}{x}\PYG{p}{)}\PYG{o}{/}\PYG{p}{(}\PYG{n}{y}\PYG{o}{\PYGZhy{}}\PYG{n}{x}\PYG{p}{)}

\PYG{c+c1}{\PYGZsh{} 3. compute the permeate flowrate:}
\PYG{n}{Fp}\PYG{o}{=}\PYG{n}{A}\PYG{o}{*}\PYG{n}{PCO2}\PYG{o}{*}\PYG{n}{rho\PYGZus{}STP} \PYG{o}{/} \PYG{n}{y} \PYG{o}{/} \PYG{n}{l} \PYG{o}{*} \PYG{p}{(}\PYG{n}{x}\PYG{o}{*}\PYG{n}{P\PYGZus{}r}\PYG{o}{\PYGZhy{}}\PYG{n}{y}\PYG{o}{*}\PYG{n}{P\PYGZus{}p}\PYG{p}{)}

\PYG{c+c1}{\PYGZsh{} 4. compute the feed flowrate: }
\PYG{n}{Fin}\PYG{o}{=}\PYG{n}{Fp}\PYG{o}{/}\PYG{n}{theta}

\PYG{c+c1}{\PYGZsh{} 5 compute the retentate flowrate: }
\PYG{n}{Fr}\PYG{o}{=}\PYG{n}{Fin}\PYG{o}{\PYGZhy{}}\PYG{n}{Fp}

\PYG{n+nb}{print}\PYG{p}{(}\PYG{l+s+s2}{\PYGZdq{}}\PYG{l+s+se}{\PYGZbs{}n}\PYG{l+s+s2}{Retentate composition: }\PYG{l+s+s2}{\PYGZdq{}}\PYG{p}{,} \PYG{l+s+sa}{f}\PYG{l+s+s2}{\PYGZdq{}}\PYG{l+s+si}{\PYGZob{}}\PYG{n}{x}\PYG{l+s+si}{:}\PYG{l+s+s2}{.4}\PYG{l+s+si}{\PYGZcb{}}\PYG{l+s+s2}{\PYGZdq{}}\PYG{p}{,} \PYG{l+s+s2}{\PYGZdq{}}\PYG{l+s+s2}{ [\PYGZhy{}]}\PYG{l+s+s2}{\PYGZdq{}}\PYG{p}{)}
\PYG{n+nb}{print}\PYG{p}{(}\PYG{l+s+s2}{\PYGZdq{}}\PYG{l+s+s2}{Cut: }\PYG{l+s+s2}{\PYGZdq{}}\PYG{p}{,} \PYG{l+s+sa}{f}\PYG{l+s+s2}{\PYGZdq{}}\PYG{l+s+si}{\PYGZob{}}\PYG{n}{theta}\PYG{l+s+si}{:}\PYG{l+s+s2}{.4}\PYG{l+s+si}{\PYGZcb{}}\PYG{l+s+s2}{\PYGZdq{}}\PYG{p}{,} \PYG{l+s+s2}{\PYGZdq{}}\PYG{l+s+s2}{[\PYGZhy{}]}\PYG{l+s+s2}{\PYGZdq{}}\PYG{p}{)}
\PYG{n+nb}{print}\PYG{p}{(}\PYG{l+s+s2}{\PYGZdq{}}\PYG{l+s+s2}{Permeate flowrate: }\PYG{l+s+s2}{\PYGZdq{}}\PYG{p}{,} \PYG{l+s+sa}{f}\PYG{l+s+s2}{\PYGZdq{}}\PYG{l+s+si}{\PYGZob{}}\PYG{n}{Fp}\PYG{l+s+si}{:}\PYG{l+s+s2}{.4}\PYG{l+s+si}{\PYGZcb{}}\PYG{l+s+s2}{\PYGZdq{}}\PYG{p}{,} \PYG{l+s+s2}{\PYGZdq{}}\PYG{l+s+s2}{ [mol/h]}\PYG{l+s+s2}{\PYGZdq{}}\PYG{p}{)}
\PYG{n+nb}{print}\PYG{p}{(}\PYG{l+s+s2}{\PYGZdq{}}\PYG{l+s+s2}{Feed flowrate: }\PYG{l+s+s2}{\PYGZdq{}}\PYG{p}{,} \PYG{l+s+sa}{f}\PYG{l+s+s2}{\PYGZdq{}}\PYG{l+s+si}{\PYGZob{}}\PYG{n}{Fin}\PYG{l+s+si}{:}\PYG{l+s+s2}{.4}\PYG{l+s+si}{\PYGZcb{}}\PYG{l+s+s2}{\PYGZdq{}}\PYG{p}{,}  \PYG{l+s+s2}{\PYGZdq{}}\PYG{l+s+s2}{ [mol/h]}\PYG{l+s+s2}{\PYGZdq{}}\PYG{p}{)}
\PYG{n+nb}{print}\PYG{p}{(}\PYG{l+s+s2}{\PYGZdq{}}\PYG{l+s+s2}{Retentate flowrate: }\PYG{l+s+s2}{\PYGZdq{}}\PYG{p}{,} \PYG{l+s+sa}{f}\PYG{l+s+s2}{\PYGZdq{}}\PYG{l+s+si}{\PYGZob{}}\PYG{n}{Fr}\PYG{l+s+si}{:}\PYG{l+s+s2}{.4}\PYG{l+s+si}{\PYGZcb{}}\PYG{l+s+s2}{\PYGZdq{}}\PYG{p}{,} \PYG{l+s+s2}{\PYGZdq{}}\PYG{l+s+s2}{ [mol/h]}\PYG{l+s+se}{\PYGZbs{}n}\PYG{l+s+s2}{\PYGZdq{}}\PYG{p}{)}
\end{sphinxVerbatim}

\end{sphinxuseclass}\end{sphinxVerbatimInput}
\begin{sphinxVerbatimOutput}

\begin{sphinxuseclass}{cell_output}
\begin{sphinxVerbatim}[commandchars=\\\{\}]
Retentate composition:  0.1826  [\PYGZhy{}]
Cut:  0.08792 [\PYGZhy{}]
Permeate flowrate:  0.06784  [mol/h]
Feed flowrate:  0.7715  [mol/h]
Retentate flowrate:  0.7037  [mol/h]
\end{sphinxVerbatim}

\end{sphinxuseclass}\end{sphinxVerbatimOutput}

\end{sphinxuseclass}
\end{sphinxuseclass}
\sphinxAtStartPar
\sphinxstylestrong{The brute force way}

\begin{sphinxuseclass}{cell}
\begin{sphinxuseclass}{tag_hide_input}\begin{sphinxVerbatimInput}

\begin{sphinxuseclass}{cell_input}
\begin{sphinxVerbatim}[commandchars=\\\{\}]
\PYG{c+c1}{\PYGZsh{}\PYGZsh{} Solution trace}
\PYG{k+kn}{import}\PYG{+w}{ }\PYG{n+nn}{scipy}\PYG{+w}{ }\PYG{k}{as}\PYG{+w}{ }\PYG{n+nn}{sp}
\PYG{c+c1}{\PYGZsh{}\PYGZsh{}from sp.optimize import fsolve}

\PYG{c+c1}{\PYGZsh{} Variables}
\PYG{n}{A}\PYG{o}{=}\PYG{l+m+mi}{15} \PYG{o}{*} \PYG{l+m+mf}{1E4} \PYG{c+c1}{\PYGZsh{} [m\PYGZca{}2 * cm\PYGZca{}2 / m\PYGZca{}2]}
\PYG{n}{P\PYGZus{}r}\PYG{o}{=}\PYG{l+m+mf}{5.5} \PYG{o}{*} \PYG{l+m+mi}{101325} \PYG{c+c1}{\PYGZsh{} [atm] * [Pa / atm]}
\PYG{n}{P\PYGZus{}p}\PYG{o}{=}\PYG{l+m+mf}{1.01} \PYG{o}{*} \PYG{l+m+mi}{101325} \PYG{c+c1}{\PYGZsh{} [atm] * [Pa / atm] }
\PYG{n}{l}\PYG{o}{=}\PYG{l+m+mf}{2.0}\PYG{o}{*}\PYG{l+m+mf}{1E\PYGZhy{}4} \PYG{c+c1}{\PYGZsh{} cm}
\PYG{n}{rho\PYGZus{}STP}\PYG{o}{=}\PYG{l+m+mi}{1}\PYG{o}{/}\PYG{l+m+mf}{22.4E3} \PYG{c+c1}{\PYGZsh{} / [mol/cm3(stp)]}
\PYG{n}{PCO2}\PYG{o}{=}\PYG{l+m+mf}{0.034E\PYGZhy{}13} \PYG{o}{*} \PYG{l+m+mi}{3600} \PYG{c+c1}{\PYGZsh{} [cm3(STP) * cm/cm2/s/Pa] * [s/h]}
\PYG{n}{PN2}\PYG{o}{=}\PYG{l+m+mf}{0.089E\PYGZhy{}14} \PYG{o}{*} \PYG{l+m+mi}{3600}\PYG{c+c1}{\PYGZsh{}  [cm3(STP) * cm/cm2/s/Pa] * [s/h]}
\PYG{c+c1}{\PYGZsh{} Permeate Composition}
\PYG{n}{y}\PYG{o}{=}\PYG{l+m+mf}{0.38}
\PYG{c+c1}{\PYGZsh{} Feed Composition}
\PYG{n}{z} \PYG{o}{=} \PYG{l+m+mf}{0.20}

\PYG{n}{alpha}\PYG{o}{=}\PYG{n}{PCO2}\PYG{o}{/}\PYG{n}{PN2}

\PYG{k}{def}\PYG{+w}{ }\PYG{n+nf}{equations}\PYG{p}{(}\PYG{n+nb}{vars}\PYG{p}{)}\PYG{p}{:}
    \PYG{n}{x}\PYG{p}{,} \PYG{n}{Fin}\PYG{p}{,} \PYG{n}{Fp}\PYG{p}{,} \PYG{n}{Fr}\PYG{p}{,} \PYG{n}{theta} \PYG{o}{=} \PYG{n+nb}{vars}
    \PYG{n}{eq1} \PYG{o}{=} \PYG{p}{(}\PYG{p}{(}\PYG{n}{theta} \PYG{o}{\PYGZhy{}} \PYG{l+m+mi}{1}\PYG{p}{)} \PYG{o}{/} \PYG{n}{theta}\PYG{p}{)} \PYG{o}{*} \PYG{n}{x} \PYG{o}{+} \PYG{p}{(}\PYG{n}{z} \PYG{o}{/} \PYG{n}{theta}\PYG{p}{)} \PYG{o}{\PYGZhy{}} \PYG{n}{y}
    \PYG{n}{eq2} \PYG{o}{=} \PYG{n}{x} \PYG{o}{\PYGZhy{}} \PYG{p}{(}\PYG{n}{y} \PYG{o}{*} \PYG{p}{(}\PYG{l+m+mi}{1} \PYG{o}{+} \PYG{n}{P\PYGZus{}p}\PYG{o}{/}\PYG{n}{P\PYGZus{}r} \PYG{o}{*} \PYG{p}{(}\PYG{l+m+mi}{1}\PYG{o}{\PYGZhy{}}\PYG{n}{y}\PYG{p}{)} \PYG{o}{*} \PYG{p}{(}\PYG{n}{alpha} \PYG{o}{\PYGZhy{}} \PYG{l+m+mi}{1}\PYG{p}{)}\PYG{p}{)} \PYG{o}{/} \PYG{p}{(}\PYG{n}{alpha} \PYG{o}{\PYGZhy{}} \PYG{p}{(}\PYG{n}{alpha} \PYG{o}{\PYGZhy{}} \PYG{l+m+mi}{1}\PYG{p}{)} \PYG{o}{*} \PYG{n}{y}\PYG{p}{)}\PYG{p}{)}
    \PYG{n}{eq3} \PYG{o}{=} \PYG{n}{Fin} \PYG{o}{\PYGZhy{}} \PYG{n}{Fr} \PYG{o}{\PYGZhy{}} \PYG{n}{Fp}
    \PYG{n}{eq4} \PYG{o}{=} \PYG{n}{Fp}\PYG{o}{\PYGZhy{}}\PYG{n}{A}\PYG{o}{*}\PYG{n}{PCO2}\PYG{o}{*}\PYG{n}{rho\PYGZus{}STP} \PYG{o}{/} \PYG{n}{y} \PYG{o}{/} \PYG{n}{l} \PYG{o}{*} \PYG{p}{(}\PYG{n}{x}\PYG{o}{*}\PYG{n}{P\PYGZus{}r}\PYG{o}{\PYGZhy{}}\PYG{n}{y}\PYG{o}{*}\PYG{n}{P\PYGZus{}p}\PYG{p}{)}
    \PYG{n}{eq5} \PYG{o}{=} \PYG{n}{Fin} \PYG{o}{\PYGZhy{}} \PYG{p}{(}\PYG{n}{Fp}\PYG{o}{/}\PYG{n}{theta}\PYG{p}{)}
    \PYG{k}{return} \PYG{p}{[}\PYG{n}{eq1}\PYG{p}{,} \PYG{n}{eq2}\PYG{p}{,} \PYG{n}{eq3}\PYG{p}{,} \PYG{n}{eq4}\PYG{p}{,} \PYG{n}{eq5}\PYG{p}{]}

\PYG{n}{x}\PYG{p}{,} \PYG{n}{Fin}\PYG{p}{,} \PYG{n}{Fp}\PYG{p}{,} \PYG{n}{Fr}\PYG{p}{,} \PYG{n}{theta} \PYG{o}{=}  \PYG{n}{sp}\PYG{o}{.}\PYG{n}{optimize}\PYG{o}{.}\PYG{n}{fsolve}\PYG{p}{(}\PYG{n}{equations}\PYG{p}{,} \PYG{p}{(}\PYG{l+m+mf}{0.1}\PYG{p}{,} \PYG{l+m+mf}{0.1}\PYG{p}{,} \PYG{l+m+mf}{0.1}\PYG{p}{,} \PYG{l+m+mf}{0.1}\PYG{p}{,} \PYG{l+m+mf}{0.1}\PYG{p}{)}\PYG{p}{)}

\PYG{n+nb}{print}\PYG{p}{(}\PYG{l+s+s2}{\PYGZdq{}}\PYG{l+s+se}{\PYGZbs{}n}\PYG{l+s+s2}{Retentate composition: }\PYG{l+s+s2}{\PYGZdq{}}\PYG{p}{,} \PYG{l+s+sa}{f}\PYG{l+s+s2}{\PYGZdq{}}\PYG{l+s+si}{\PYGZob{}}\PYG{n}{x}\PYG{l+s+si}{:}\PYG{l+s+s2}{.4}\PYG{l+s+si}{\PYGZcb{}}\PYG{l+s+s2}{\PYGZdq{}}\PYG{p}{,} \PYG{l+s+s2}{\PYGZdq{}}\PYG{l+s+s2}{ [\PYGZhy{}]}\PYG{l+s+s2}{\PYGZdq{}}\PYG{p}{)}
\PYG{n+nb}{print}\PYG{p}{(}\PYG{l+s+s2}{\PYGZdq{}}\PYG{l+s+s2}{Cut: }\PYG{l+s+s2}{\PYGZdq{}}\PYG{p}{,} \PYG{l+s+sa}{f}\PYG{l+s+s2}{\PYGZdq{}}\PYG{l+s+si}{\PYGZob{}}\PYG{n}{theta}\PYG{l+s+si}{:}\PYG{l+s+s2}{.4}\PYG{l+s+si}{\PYGZcb{}}\PYG{l+s+s2}{\PYGZdq{}}\PYG{p}{,} \PYG{l+s+s2}{\PYGZdq{}}\PYG{l+s+s2}{[\PYGZhy{}]}\PYG{l+s+s2}{\PYGZdq{}}\PYG{p}{)}
\PYG{n+nb}{print}\PYG{p}{(}\PYG{l+s+s2}{\PYGZdq{}}\PYG{l+s+s2}{Permeate flowrate: }\PYG{l+s+s2}{\PYGZdq{}}\PYG{p}{,} \PYG{l+s+sa}{f}\PYG{l+s+s2}{\PYGZdq{}}\PYG{l+s+si}{\PYGZob{}}\PYG{n}{Fp}\PYG{l+s+si}{:}\PYG{l+s+s2}{.4}\PYG{l+s+si}{\PYGZcb{}}\PYG{l+s+s2}{\PYGZdq{}}\PYG{p}{,} \PYG{l+s+s2}{\PYGZdq{}}\PYG{l+s+s2}{ [mol/h]}\PYG{l+s+s2}{\PYGZdq{}}\PYG{p}{)}
\PYG{n+nb}{print}\PYG{p}{(}\PYG{l+s+s2}{\PYGZdq{}}\PYG{l+s+s2}{Feed flowrate: }\PYG{l+s+s2}{\PYGZdq{}}\PYG{p}{,} \PYG{l+s+sa}{f}\PYG{l+s+s2}{\PYGZdq{}}\PYG{l+s+si}{\PYGZob{}}\PYG{n}{Fin}\PYG{l+s+si}{:}\PYG{l+s+s2}{.4}\PYG{l+s+si}{\PYGZcb{}}\PYG{l+s+s2}{\PYGZdq{}}\PYG{p}{,}  \PYG{l+s+s2}{\PYGZdq{}}\PYG{l+s+s2}{ [mol/h]}\PYG{l+s+s2}{\PYGZdq{}}\PYG{p}{)}
\PYG{n+nb}{print}\PYG{p}{(}\PYG{l+s+s2}{\PYGZdq{}}\PYG{l+s+s2}{Retentate flowrate: }\PYG{l+s+s2}{\PYGZdq{}}\PYG{p}{,} \PYG{l+s+sa}{f}\PYG{l+s+s2}{\PYGZdq{}}\PYG{l+s+si}{\PYGZob{}}\PYG{n}{Fr}\PYG{l+s+si}{:}\PYG{l+s+s2}{.4}\PYG{l+s+si}{\PYGZcb{}}\PYG{l+s+s2}{\PYGZdq{}}\PYG{p}{,} \PYG{l+s+s2}{\PYGZdq{}}\PYG{l+s+s2}{ [mol/h]}\PYG{l+s+se}{\PYGZbs{}n}\PYG{l+s+s2}{\PYGZdq{}}\PYG{p}{)}
\end{sphinxVerbatim}

\end{sphinxuseclass}\end{sphinxVerbatimInput}
\begin{sphinxVerbatimOutput}

\begin{sphinxuseclass}{cell_output}
\begin{sphinxVerbatim}[commandchars=\\\{\}]
Retentate composition:  0.1826  [\PYGZhy{}]
Cut:  0.08792 [\PYGZhy{}]
Permeate flowrate:  0.06784  [mol/h]
Feed flowrate:  0.7715  [mol/h]
Retentate flowrate:  0.7037  [mol/h]
\end{sphinxVerbatim}

\end{sphinxuseclass}\end{sphinxVerbatimOutput}

\end{sphinxuseclass}
\end{sphinxuseclass}

\subsection{Contributions}
\label{\detokenize{Exercise_W2_2025:contributions}}\begin{itemize}
\item {} 
\sphinxAtStartPar
Divyansh Vashishtha, 19 Feb 2024

\end{itemize}

\sphinxstepscope


\section{Week 3 \sphinxhyphen{} Exercise}
\label{\detokenize{Exercise_W3_2021:week-3-exercise}}\label{\detokenize{Exercise_W3_2021::doc}}

\subsection{Problem Statement}
\label{\detokenize{Exercise_W3_2021:problem-statement}}
\sphinxAtStartPar
Compute the total membrane area necessary to produce 10 l/h purified water (0.03 wt\% NaCl) through a single stage Reverse Osmosis separation of seawater (water 4 wt\% NaCl).
The feed pressure is 100 bar and the composition in the retentate stream is 7 wt\%.\\
The permeate stream is at ambient pressure (1 bar), all streams are at 25\(^o\) C, and the osmotic pressure in bar is captured by the semiempirical expression:
\begin{equation*}
\begin{split}
\pi=\frac{1.12}{14.5}T\sum_{i=1}^Nc_i
\end{split}
\end{equation*}
\sphinxAtStartPar
Where T is the temperature in K,  \(N\) is the total number of ionic species in solution, and ci is the concentration expressed in \([mol/l]\).

\sphinxAtStartPar
The asymmetric membrane used for the separation has a permeance of \(10^{-5}\,[m h^{-1} bar^{-1}]\).

\sphinxAtStartPar
Once you have solved the problem consider the following questions:
\begin{itemize}
\item {} 
\sphinxAtStartPar
Can you evaluate the salt rejection in this process, and salt passage?

\item {} 
\sphinxAtStartPar
How would you modify the operation of the system to increase the productivity of the process?

\end{itemize}


\subsection{Solution trace}
\label{\detokenize{Exercise_W3_2021:solution-trace}}
\sphinxAtStartPar
\sphinxstylestrong{Data:}

\sphinxAtStartPar
NaCl weigth percentage:
\(w_{NaCl,P}=0.03\)
\(w_{NaCl,F}=4\)
\(w_{NaCl,R}=7\)

\sphinxAtStartPar
\sphinxstylestrong{Water density:}
\(\rho_{H_2O}\)=1\(\times\)10\(^3\) {[}g/l{]}

\sphinxAtStartPar
\sphinxstylestrong{Working expression:}
\begin{equation*}
\begin{split}
c_{NaCl,i}=\frac{w_{NaCl,i}}{100}\times\frac{\rho_{Solution,i}}{MM_{NaCl}}
\end{split}
\end{equation*}
\sphinxAtStartPar
Where: \(i\) indicates the stream we are referring to, \(\rho_{Solution,i}\)  is the mass density of the NaCl solution in each stream, and \(MM_{NaCl}\) is the molar mass of NaCl. 
The mass density of the solution \(\rho_{Solution}\) can be computed known the NaCl weight percentage and the density of  water as:
\begin{equation*}
\begin{split}
\rho_{solution}=\frac{\rho_{H_2O}}{1-\frac{w_{NaCl,i}}{100}}
\end{split}
\end{equation*}
\sphinxAtStartPar
\sphinxstylestrong{Computing the Osmotic pressure difference across the membrane}
The expression provided in the exercise assignment allows to compute the osmotic pressure associated with an assigned concentration of NaCl in a solution of assigned NaCl concentration.
\sphinxstyleemphasis{Note that the concentration appearing in such expression is the total ionic concentration, which can be computed from the NaCl concentration knowing the unit formula of NaCl.}
\begin{equation*}
\begin{split}
c_{TOT,i}=c_{Na^+,i}+c_{Cl^-,i}=2c_{NaCl,i}
\end{split}
\end{equation*}
\sphinxAtStartPar
Note that in this case it is typically reasonable to assume the concentration on the feed/retentate side of the membrane to be the average between concentrations in the feed and retentate.
\begin{equation*}
\begin{split}
c_{TOT,F/R}=2 \left({\frac{c_{NaCl,F}+c_{NaCl,R}}{2}}\right)
\end{split}
\end{equation*}
\sphinxAtStartPar
Alternatively one can consider the feed/retentate side to be perfectly mixed and its composition constant and equal to the retentate one.

\sphinxAtStartPar
The difference in the osmotic pressure can thus be computed as:
\begin{equation*}
\begin{split}
\Delta{\pi}=\frac{1.12}{14.5}T\left(c_{TOT,F/R}-c_{TOT,P}\right)
\end{split}
\end{equation*}
\sphinxAtStartPar
\sphinxstylestrong{Computing the membrane area}
The flow rate through the membrane is computed by the expression:
\begin{equation*}
\begin{split}
F_{p,wat}=J\times{A}=A\,P_{wat}\,\left(\Delta{P}-\Delta{\pi}\right)
\end{split}
\end{equation*}
\sphinxAtStartPar
hence:
\begin{equation*}
\begin{split}
A=\frac{F_{p,wat}}{P_{wat}\,\left(\Delta{P}-\Delta{\pi}\right)}
\end{split}
\end{equation*}
\sphinxAtStartPar
\sphinxstylestrong{Salt rejection}
Given the concentrations computed earlier salt rejection can be computed from its definition:
\begin{equation*}
\begin{split}
R_s=1-\frac{c_{NaCl,P}}{c_{NaCl,R}}
\end{split}
\end{equation*}
\sphinxAtStartPar
Salt passage is given by:
\begin{equation*}
\begin{split}
P_s=\frac{c_{NaCl,P}}{c_{NaCl,R}}
\end{split}
\end{equation*}
\begin{sphinxuseclass}{cell}\begin{sphinxVerbatimInput}

\begin{sphinxuseclass}{cell_input}
\begin{sphinxVerbatim}[commandchars=\\\{\}]
\PYG{k+kn}{import}\PYG{+w}{ }\PYG{n+nn}{numpy}\PYG{+w}{ }\PYG{k}{as}\PYG{+w}{ }\PYG{n+nn}{np}
\PYG{k+kn}{import}\PYG{+w}{ }\PYG{n+nn}{matplotlib}\PYG{n+nn}{.}\PYG{n+nn}{pyplot}\PYG{+w}{ }\PYG{k}{as}\PYG{+w}{ }\PYG{n+nn}{plt} 
\PYG{k+kn}{from}\PYG{+w}{ }\PYG{n+nn}{matplotlib}\PYG{n+nn}{.}\PYG{n+nn}{pyplot}\PYG{+w}{ }\PYG{k+kn}{import} \PYG{n}{cm}

\PYG{c+c1}{\PYGZsh{}}
\PYG{n}{PM}\PYG{o}{=}\PYG{l+m+mf}{1E\PYGZhy{}5}\PYG{p}{;} \PYG{c+c1}{\PYGZsh{}[m/h/bar]}
\PYG{n}{MMNaCl}\PYG{o}{=}\PYG{l+m+mf}{58.5}\PYG{p}{;} \PYG{c+c1}{\PYGZsh{}[g/mol]}

\PYG{c+c1}{\PYGZsh{}Feed Retentate Permeate}
\PYG{n}{wtp}\PYG{o}{=}\PYG{n}{np}\PYG{o}{.}\PYG{n}{array}\PYG{p}{(}\PYG{p}{[}\PYG{l+m+mi}{4}\PYG{p}{,} \PYG{l+m+mi}{7}\PYG{p}{,} \PYG{l+m+mf}{0.03}\PYG{p}{]}\PYG{p}{)}\PYG{o}{/}\PYG{l+m+mi}{100}\PYG{p}{;} \PYG{c+c1}{\PYGZsh{}mass fractions}
\PYG{n}{water\PYGZus{}density}\PYG{o}{=}\PYG{l+m+mf}{1E3}\PYG{p}{;} \PYG{c+c1}{\PYGZsh{}[g/l]}

\PYG{n}{c}\PYG{o}{=}\PYG{n}{wtp}\PYG{o}{*}\PYG{n}{water\PYGZus{}density}\PYG{o}{/}\PYG{p}{(}\PYG{n}{MMNaCl}\PYG{o}{*}\PYG{p}{(}\PYG{l+m+mi}{1}\PYG{o}{\PYGZhy{}}\PYG{n}{wtp}\PYG{p}{)}\PYG{p}{)}\PYG{p}{;}  \PYG{c+c1}{\PYGZsh{}[mol/l]}

\PYG{n}{C}\PYG{o}{=}\PYG{l+m+mf}{1.12}\PYG{o}{/}\PYG{l+m+mf}{14.5}\PYG{p}{;}
\PYG{n}{T}\PYG{o}{=}\PYG{l+m+mi}{298}\PYG{p}{;}
\PYG{n}{pi\PYGZus{}P}\PYG{o}{=}\PYG{l+m+mi}{2}\PYG{o}{*}\PYG{n}{C}\PYG{o}{*}\PYG{n}{T}\PYG{o}{*}\PYG{n}{c}\PYG{p}{[}\PYG{l+m+mi}{2}\PYG{p}{]}\PYG{p}{;}
\PYG{n}{pi\PYGZus{}R}\PYG{o}{=}\PYG{l+m+mi}{2}\PYG{o}{*}\PYG{n}{C}\PYG{o}{*}\PYG{n}{T}\PYG{o}{*}\PYG{n}{np}\PYG{o}{.}\PYG{n}{mean}\PYG{p}{(}\PYG{n}{c}\PYG{p}{[}\PYG{l+m+mi}{0}\PYG{p}{:}\PYG{l+m+mi}{2}\PYG{p}{]}\PYG{p}{)}\PYG{p}{;}
\PYG{n}{Dpi}\PYG{o}{=}\PYG{n}{pi\PYGZus{}R}\PYG{o}{\PYGZhy{}}\PYG{n}{pi\PYGZus{}P}\PYG{p}{;} \PYG{c+c1}{\PYGZsh{} [bar]}

\PYG{n}{DP}\PYG{o}{=}\PYG{l+m+mi}{100}\PYG{o}{\PYGZhy{}}\PYG{l+m+mi}{1}\PYG{p}{;} \PYG{c+c1}{\PYGZsh{} [bar]}
\PYG{n}{Fpwat}\PYG{o}{=}\PYG{l+m+mf}{10E\PYGZhy{}3}\PYG{p}{;} \PYG{c+c1}{\PYGZsh{} [m3/h]}

\PYG{n}{A}\PYG{o}{=}\PYG{n}{Fpwat}\PYG{o}{/}\PYG{n}{PM}\PYG{o}{/}\PYG{p}{(}\PYG{n}{DP}\PYG{o}{\PYGZhy{}}\PYG{n}{Dpi}\PYG{p}{)}

\PYG{n}{SP}\PYG{o}{=}\PYG{n}{c}\PYG{p}{[}\PYG{l+m+mi}{2}\PYG{p}{]}\PYG{o}{/}\PYG{n}{c}\PYG{p}{[}\PYG{l+m+mi}{1}\PYG{p}{]}
\PYG{n}{SR}\PYG{o}{=}\PYG{p}{(}\PYG{l+m+mi}{1}\PYG{o}{\PYGZhy{}}\PYG{n}{SP}\PYG{p}{)}

\PYG{n+nb}{print}\PYG{p}{(}\PYG{l+s+s2}{\PYGZdq{}}\PYG{l+s+se}{\PYGZbs{}n}\PYG{l+s+s2}{Permeate concentration: }\PYG{l+s+s2}{\PYGZdq{}}\PYG{p}{,} \PYG{l+s+sa}{f}\PYG{l+s+s2}{\PYGZdq{}}\PYG{l+s+si}{\PYGZob{}}\PYG{n}{c}\PYG{p}{[}\PYG{l+m+mi}{2}\PYG{p}{]}\PYG{l+s+si}{:}\PYG{l+s+s2}{.4}\PYG{l+s+si}{\PYGZcb{}}\PYG{l+s+s2}{\PYGZdq{}}\PYG{p}{,} \PYG{l+s+s2}{\PYGZdq{}}\PYG{l+s+s2}{ [mol/l]}\PYG{l+s+s2}{\PYGZdq{}}\PYG{p}{)}
\PYG{n+nb}{print}\PYG{p}{(}\PYG{l+s+s2}{\PYGZdq{}}\PYG{l+s+s2}{Feed concentration: }\PYG{l+s+s2}{\PYGZdq{}}\PYG{p}{,} \PYG{l+s+sa}{f}\PYG{l+s+s2}{\PYGZdq{}}\PYG{l+s+si}{\PYGZob{}}\PYG{n}{c}\PYG{p}{[}\PYG{l+m+mi}{0}\PYG{p}{]}\PYG{l+s+si}{:}\PYG{l+s+s2}{.4}\PYG{l+s+si}{\PYGZcb{}}\PYG{l+s+s2}{\PYGZdq{}}\PYG{p}{,} \PYG{l+s+s2}{\PYGZdq{}}\PYG{l+s+s2}{ [mol/l]}\PYG{l+s+s2}{\PYGZdq{}}\PYG{p}{)}
\PYG{n+nb}{print}\PYG{p}{(}\PYG{l+s+s2}{\PYGZdq{}}\PYG{l+s+s2}{Retentate concentration: }\PYG{l+s+s2}{\PYGZdq{}}\PYG{p}{,} \PYG{l+s+sa}{f}\PYG{l+s+s2}{\PYGZdq{}}\PYG{l+s+si}{\PYGZob{}}\PYG{n}{c}\PYG{p}{[}\PYG{l+m+mi}{1}\PYG{p}{]}\PYG{l+s+si}{:}\PYG{l+s+s2}{.4}\PYG{l+s+si}{\PYGZcb{}}\PYG{l+s+s2}{\PYGZdq{}}\PYG{p}{,} \PYG{l+s+s2}{\PYGZdq{}}\PYG{l+s+s2}{ [mol/l]}\PYG{l+s+s2}{\PYGZdq{}}\PYG{p}{)}

\PYG{n+nb}{print}\PYG{p}{(}\PYG{l+s+s2}{\PYGZdq{}}\PYG{l+s+se}{\PYGZbs{}n}\PYG{l+s+s2}{Osmotic pressure: }\PYG{l+s+s2}{\PYGZdq{}}\PYG{p}{,} \PYG{l+s+sa}{f}\PYG{l+s+s2}{\PYGZdq{}}\PYG{l+s+si}{\PYGZob{}}\PYG{n}{Dpi}\PYG{l+s+si}{:}\PYG{l+s+s2}{.4}\PYG{l+s+si}{\PYGZcb{}}\PYG{l+s+s2}{\PYGZdq{}}\PYG{p}{,} \PYG{l+s+s2}{\PYGZdq{}}\PYG{l+s+s2}{ [bar]}\PYG{l+s+s2}{\PYGZdq{}}\PYG{p}{)}

\PYG{n+nb}{print}\PYG{p}{(}\PYG{l+s+s2}{\PYGZdq{}}\PYG{l+s+se}{\PYGZbs{}n}\PYG{l+s+s2}{Membrane area: }\PYG{l+s+s2}{\PYGZdq{}}\PYG{p}{,} \PYG{l+s+sa}{f}\PYG{l+s+s2}{\PYGZdq{}}\PYG{l+s+si}{\PYGZob{}}\PYG{n}{A}\PYG{l+s+si}{:}\PYG{l+s+s2}{.4}\PYG{l+s+si}{\PYGZcb{}}\PYG{l+s+s2}{\PYGZdq{}}\PYG{p}{,} \PYG{l+s+s2}{\PYGZdq{}}\PYG{l+s+s2}{ [m\PYGZca{}2]}\PYG{l+s+s2}{\PYGZdq{}}\PYG{p}{)}
\PYG{n+nb}{print}\PYG{p}{(}\PYG{l+s+s2}{\PYGZdq{}}\PYG{l+s+s2}{Salt rejection: }\PYG{l+s+s2}{\PYGZdq{}}\PYG{p}{,} \PYG{l+s+sa}{f}\PYG{l+s+s2}{\PYGZdq{}}\PYG{l+s+si}{\PYGZob{}}\PYG{n}{SR}\PYG{l+s+si}{:}\PYG{l+s+s2}{.4}\PYG{l+s+si}{\PYGZcb{}}\PYG{l+s+s2}{\PYGZdq{}}\PYG{p}{,} \PYG{l+s+s2}{\PYGZdq{}}\PYG{l+s+s2}{[\PYGZhy{}]}\PYG{l+s+s2}{\PYGZdq{}}\PYG{p}{)}
\PYG{n+nb}{print}\PYG{p}{(}\PYG{l+s+s2}{\PYGZdq{}}\PYG{l+s+s2}{Salt passage: }\PYG{l+s+s2}{\PYGZdq{}}\PYG{p}{,} \PYG{l+s+sa}{f}\PYG{l+s+s2}{\PYGZdq{}}\PYG{l+s+si}{\PYGZob{}}\PYG{n}{SP}\PYG{l+s+si}{:}\PYG{l+s+s2}{.2}\PYG{l+s+si}{\PYGZcb{}}\PYG{l+s+s2}{\PYGZdq{}}\PYG{p}{,} \PYG{l+s+s2}{\PYGZdq{}}\PYG{l+s+s2}{[\PYGZhy{}]}\PYG{l+s+s2}{\PYGZdq{}}\PYG{p}{)}
\end{sphinxVerbatim}

\end{sphinxuseclass}\end{sphinxVerbatimInput}
\begin{sphinxVerbatimOutput}

\begin{sphinxuseclass}{cell_output}
\begin{sphinxVerbatim}[commandchars=\\\{\}]
Permeate concentration:  0.00513  [mol/l]
Feed concentration:  0.7123  [mol/l]
Retentate concentration:  1.287  [mol/l]

Osmotic pressure:  45.77  [bar]

Membrane area:  18.79  [m\PYGZca{}2]
Salt rejection:  0.996 [\PYGZhy{}]
Salt passage:  0.004 [\PYGZhy{}]
\end{sphinxVerbatim}

\end{sphinxuseclass}\end{sphinxVerbatimOutput}

\end{sphinxuseclass}

\subsection{Contributions}
\label{\detokenize{Exercise_W3_2021:contributions}}
\sphinxAtStartPar
Debugged by Nikita Gusev, 29 Jan 2021

\sphinxstepscope


\section{Week 4 \sphinxhyphen{} Exercise}
\label{\detokenize{Exercise_W4_2024_updated:week-4-exercise}}\label{\detokenize{Exercise_W4_2024_updated::doc}}

\subsection{Problem Statement}
\label{\detokenize{Exercise_W4_2024_updated:problem-statement}}
\sphinxAtStartPar
The downstream processing of an aqueous stream containing therapeutic proteins produced in a bioreactor is carried out through an ultrafiltration membrane process.
The volumetric flow to be processed is 200 \(\mathrm{m^3\,day^{-1}}\), with a suspended protein concentration of 0.5 \(\mathrm{kg\,m^{-3}}\).
In order to be used in subsequent formulation stages the protein stream needs to be concentrated to 20 \(\mathrm{kg\,m^{-3}}\).

\sphinxAtStartPar
Tubular membrane modules with an exchange surface of 5 \(\mathrm{m^2}\) are used to carry out the ultrafiltration process.

\sphinxAtStartPar
Due to fouling, the water flow through these membranes decreases when protein concentration (C) increases following the empirical law:
\begin{equation*}
\begin{split}
J(C)=3.8\times10^{-2}\ln\left(\frac{145}{C}\right) \mathrm{[m\,h^{-1}]}
\end{split}
\end{equation*}

\subsection{Requests}
\label{\detokenize{Exercise_W4_2024_updated:requests}}\begin{itemize}
\item {} 
\sphinxAtStartPar
Based on steady state calculations compute the number of membrane modules necessary for the operation of this process in a single\sphinxhyphen{}stage feed\sphinxhyphen{}and\sphinxhyphen{}bleed configuration.

\item {} 
\sphinxAtStartPar
Consider now a cascade of two stages. Compute the optimal number of membrane modules and their arrangement in stages in order to achieve the specifics required. Compute all the concentrations and volumetric flows in your process configuration of choice.  Define clearly your chosen criterion for optimality.

\end{itemize}


\subsection{Solution:}
\label{\detokenize{Exercise_W4_2024_updated:solution}}
\sphinxAtStartPar
\sphinxstyleemphasis{{[}contributed by Max Shone, February 2024{]}}


\subsection{Request 1}
\label{\detokenize{Exercise_W4_2024_updated:request-1}}
\sphinxAtStartPar
To compute the mumber of membrane modules necessary for the operation of this process in a single\sphinxhyphen{}stage feed\sphinxhyphen{}and\sphinxhyphen{}bleed configuration, a mass balance and volume balance can be set up:


\subsubsection{Mass Balance:}
\label{\detokenize{Exercise_W4_2024_updated:mass-balance}}\begin{equation*}
\begin{split}
Q_{in}\times C_{in}=Q_{out}\times C_{out}
\end{split}
\end{equation*}

\subsubsection{Volume Balance:}
\label{\detokenize{Exercise_W4_2024_updated:volume-balance}}\begin{equation*}
\begin{split}
Q_{in}=Q_{out}+J(C)A
\end{split}
\end{equation*}
\sphinxAtStartPar
To solve this problem, the following steps are taken:
\begin{itemize}
\item {} 
\sphinxAtStartPar
Solve mass balance for Qin.

\item {} 
\sphinxAtStartPar
Solve empirical formula for flux through membrane.

\item {} 
\sphinxAtStartPar
Rearrange volume balance for area and solve, rounding the final answer up.

\end{itemize}

\begin{sphinxuseclass}{cell}\begin{sphinxVerbatimInput}

\begin{sphinxuseclass}{cell_input}
\begin{sphinxVerbatim}[commandchars=\\\{\}]
\PYG{k+kn}{import}\PYG{+w}{ }\PYG{n+nn}{numpy}\PYG{+w}{ }\PYG{k}{as}\PYG{+w}{ }\PYG{n+nn}{np}
\PYG{n}{Qin} \PYG{o}{=} \PYG{l+m+mi}{200}\PYG{o}{/}\PYG{l+m+mi}{24} \PYG{c+c1}{\PYGZsh{}m\PYGZca{}3/hr}
\PYG{n}{Cin} \PYG{o}{=} \PYG{l+m+mf}{0.5} \PYG{c+c1}{\PYGZsh{}kg/m\PYGZca{}3}
\PYG{n}{Cout} \PYG{o}{=} \PYG{l+m+mi}{20} \PYG{c+c1}{\PYGZsh{}kg/m\PYGZca{}3}

\PYG{c+c1}{\PYGZsh{}equations}

\PYG{c+c1}{\PYGZsh{}mass balance}
\PYG{n}{Qout} \PYG{o}{=} \PYG{n}{Qin}\PYG{o}{*}\PYG{n}{Cin}\PYG{o}{/}\PYG{n}{Cout}

\PYG{c+c1}{\PYGZsh{}Finding flux going through membrane}
\PYG{n}{Jc} \PYG{o}{=} \PYG{p}{(}\PYG{l+m+mf}{3.8}\PYG{o}{*}\PYG{l+m+mi}{10}\PYG{o}{*}\PYG{o}{*}\PYG{p}{(}\PYG{o}{\PYGZhy{}}\PYG{l+m+mi}{2}\PYG{p}{)}\PYG{o}{*}\PYG{n}{np}\PYG{o}{.}\PYG{n}{log}\PYG{p}{(}\PYG{l+m+mi}{145}\PYG{o}{/}\PYG{n}{Cout}\PYG{p}{)}\PYG{p}{)}

\PYG{c+c1}{\PYGZsh{}Volume balance rearranged for A}
\PYG{n}{A} \PYG{o}{=} \PYG{p}{(}\PYG{n}{Qin}\PYG{o}{\PYGZhy{}}\PYG{n}{Qout}\PYG{p}{)}\PYG{o}{/}\PYG{n}{Jc}
\PYG{n}{n}\PYG{o}{=}\PYG{n}{np}\PYG{o}{.}\PYG{n}{ceil}\PYG{p}{(}\PYG{n}{A}\PYG{o}{/}\PYG{l+m+mi}{5}\PYG{p}{)}
\PYG{n+nb}{print}\PYG{p}{(}\PYG{l+s+s2}{\PYGZdq{}}\PYG{l+s+se}{\PYGZbs{}n}\PYG{l+s+s2}{Flow rate out: }\PYG{l+s+s2}{\PYGZdq{}}\PYG{p}{,} \PYG{l+s+sa}{f}\PYG{l+s+s2}{\PYGZdq{}}\PYG{l+s+si}{\PYGZob{}}\PYG{n}{Qout}\PYG{l+s+si}{:}\PYG{l+s+s2}{.4}\PYG{l+s+si}{\PYGZcb{}}\PYG{l+s+s2}{\PYGZdq{}}\PYG{p}{,} \PYG{l+s+s2}{\PYGZdq{}}\PYG{l+s+s2}{ [m\PYGZca{}3/h]}\PYG{l+s+s2}{\PYGZdq{}}\PYG{p}{)}
\PYG{n+nb}{print}\PYG{p}{(}\PYG{l+s+s2}{\PYGZdq{}}\PYG{l+s+se}{\PYGZbs{}n}\PYG{l+s+s2}{Flux: }\PYG{l+s+s2}{\PYGZdq{}}\PYG{p}{,} \PYG{l+s+sa}{f}\PYG{l+s+s2}{\PYGZdq{}}\PYG{l+s+si}{\PYGZob{}}\PYG{n}{Jc}\PYG{l+s+si}{:}\PYG{l+s+s2}{.4}\PYG{l+s+si}{\PYGZcb{}}\PYG{l+s+s2}{\PYGZdq{}}\PYG{p}{,} \PYG{l+s+s2}{\PYGZdq{}}\PYG{l+s+s2}{ [m/h]}\PYG{l+s+s2}{\PYGZdq{}}\PYG{p}{)}
\PYG{n+nb}{print}\PYG{p}{(}\PYG{l+s+s2}{\PYGZdq{}}\PYG{l+s+se}{\PYGZbs{}n}\PYG{l+s+s2}{Membrane Area: }\PYG{l+s+s2}{\PYGZdq{}}\PYG{p}{,} \PYG{l+s+sa}{f}\PYG{l+s+s2}{\PYGZdq{}}\PYG{l+s+si}{\PYGZob{}}\PYG{n}{A}\PYG{l+s+si}{:}\PYG{l+s+s2}{.4}\PYG{l+s+si}{\PYGZcb{}}\PYG{l+s+s2}{\PYGZdq{}}\PYG{p}{,} \PYG{l+s+s2}{\PYGZdq{}}\PYG{l+s+s2}{ [m\PYGZca{}2]}\PYG{l+s+s2}{\PYGZdq{}}\PYG{p}{)}
\PYG{n+nb}{print}\PYG{p}{(}\PYG{l+s+s2}{\PYGZdq{}}\PYG{l+s+se}{\PYGZbs{}n}\PYG{l+s+s2}{Number of stages: }\PYG{l+s+s2}{\PYGZdq{}}\PYG{p}{,} \PYG{l+s+sa}{f}\PYG{l+s+s2}{\PYGZdq{}}\PYG{l+s+si}{\PYGZob{}}\PYG{n}{n}\PYG{l+s+si}{\PYGZcb{}}\PYG{l+s+s2}{\PYGZdq{}}\PYG{p}{,} \PYG{l+s+s2}{\PYGZdq{}}\PYG{l+s+s2}{ [\PYGZhy{}]}\PYG{l+s+s2}{\PYGZdq{}}\PYG{p}{)}
\end{sphinxVerbatim}

\end{sphinxuseclass}\end{sphinxVerbatimInput}
\begin{sphinxVerbatimOutput}

\begin{sphinxuseclass}{cell_output}
\begin{sphinxVerbatim}[commandchars=\\\{\}]
Flow rate out:  0.2083  [m\PYGZca{}3/h]

Flux:  0.07528  [m/h]

Membrane Area:  107.9  [m\PYGZca{}2]

Number of stages:  22.0  [\PYGZhy{}]
\end{sphinxVerbatim}

\end{sphinxuseclass}\end{sphinxVerbatimOutput}

\end{sphinxuseclass}

\subsection{Request 2}
\label{\detokenize{Exercise_W4_2024_updated:request-2}}
\sphinxAtStartPar
To find the optimal membrane configuration for a cascade of two stages, we would like to find the least number of total membranes that will result in an outlet concentration of 20 \(\mathrm{kg\,m^{-3}}\).

\sphinxAtStartPar
This can be done through completing a mass and volume balance over the first membrane, assuming the outlet concentration (C1) value, followed by completing the same calculations over the second membrane, with an outlet concentration of 20 \(\mathrm{kg\,m^{-3}}\). The following graph shows all possible values of C1 between 0.5 and 20 \(\mathrm{kg\,m^{-3}}\).

\begin{sphinxuseclass}{cell}\begin{sphinxVerbatimInput}

\begin{sphinxuseclass}{cell_input}
\begin{sphinxVerbatim}[commandchars=\\\{\}]
\PYG{k+kn}{import}\PYG{+w}{ }\PYG{n+nn}{numpy}\PYG{+w}{ }\PYG{k}{as}\PYG{+w}{ }\PYG{n+nn}{np}
\PYG{k+kn}{import}\PYG{+w}{ }\PYG{n+nn}{matplotlib}\PYG{n+nn}{.}\PYG{n+nn}{pyplot}\PYG{+w}{ }\PYG{k}{as}\PYG{+w}{ }\PYG{n+nn}{plt}

\PYG{c+c1}{\PYGZsh{}Question 2 Open exercsie}
\PYG{c+c1}{\PYGZsh{}Initialising variables}
\PYG{n}{Q0} \PYG{o}{=} \PYG{l+m+mi}{200}\PYG{o}{/}\PYG{l+m+mi}{24}\PYG{p}{;} \PYG{c+c1}{\PYGZsh{}[m\PYGZca{}3/day] * [day/hour]}
\PYG{n}{C0} \PYG{o}{=} \PYG{l+m+mf}{0.5}\PYG{p}{;} \PYG{c+c1}{\PYGZsh{}kg/m\PYGZca{}3}
\PYG{n}{C1} \PYG{o}{=} \PYG{n}{np}\PYG{o}{.}\PYG{n}{linspace}\PYG{p}{(}\PYG{l+m+mf}{0.5}\PYG{p}{,} \PYG{l+m+mi}{20}\PYG{p}{,} \PYG{n}{num}\PYG{o}{=}\PYG{l+m+mi}{400}\PYG{p}{)} \PYG{c+c1}{\PYGZsh{}kg/m\PYGZca{}3}
\PYG{n}{C2} \PYG{o}{=} \PYG{l+m+mi}{20}\PYG{p}{;} \PYG{c+c1}{\PYGZsh{}kg/m\PYGZca{}3}
\PYG{n}{JC} \PYG{o}{=} \PYG{l+m+mf}{3.8}\PYG{o}{*}\PYG{l+m+mi}{10}\PYG{o}{*}\PYG{o}{*}\PYG{p}{(}\PYG{o}{\PYGZhy{}}\PYG{l+m+mi}{2}\PYG{p}{)}\PYG{o}{*}\PYG{n}{np}\PYG{o}{.}\PYG{n}{log}\PYG{p}{(}\PYG{l+m+mf}{145.}\PYG{o}{/}\PYG{n}{C1}\PYG{p}{)}

\PYG{c+c1}{\PYGZsh{}Membrane Surface Area}
\PYG{n}{MSA} \PYG{o}{=} \PYG{l+m+mi}{5} 

\PYG{c+c1}{\PYGZsh{}Calculations over first membrane}
\PYG{n}{Q1} \PYG{o}{=} \PYG{n}{Q0}\PYG{o}{*}\PYG{n}{C0}\PYG{o}{/}\PYG{n}{C1}
\PYG{n}{A1} \PYG{o}{=} \PYG{p}{(}\PYG{n}{Q0}\PYG{o}{\PYGZhy{}}\PYG{n}{Q1}\PYG{p}{)}\PYG{o}{/}\PYG{n}{JC}
\PYG{n}{n1} \PYG{o}{=} \PYG{n}{np}\PYG{o}{.}\PYG{n}{ceil}\PYG{p}{(}\PYG{n}{A1}\PYG{o}{/}\PYG{n}{MSA}\PYG{p}{)}

\PYG{c+c1}{\PYGZsh{}Calculations over second membrane}
\PYG{n}{Q2} \PYG{o}{=} \PYG{n}{Q1}\PYG{o}{*}\PYG{n}{C1}\PYG{o}{/}\PYG{n}{C2}
\PYG{n}{A2} \PYG{o}{=} \PYG{p}{(}\PYG{n}{Q1}\PYG{o}{\PYGZhy{}}\PYG{n}{Q2}\PYG{p}{)}\PYG{o}{/}\PYG{n}{JC}\PYG{p}{[}\PYG{o}{\PYGZhy{}}\PYG{l+m+mi}{1}\PYG{p}{]}
\PYG{n}{n2} \PYG{o}{=} \PYG{n}{np}\PYG{o}{.}\PYG{n}{ceil}\PYG{p}{(}\PYG{n}{A2}\PYG{o}{/}\PYG{n}{MSA}\PYG{p}{)}

\PYG{c+c1}{\PYGZsh{}Total membranes}
\PYG{n}{n} \PYG{o}{=} \PYG{n}{n1}\PYG{o}{+}\PYG{n}{n2}

\PYG{c+c1}{\PYGZsh{} Plotting graphs}
\PYG{n}{plt}\PYG{o}{.}\PYG{n}{plot}\PYG{p}{(}\PYG{n}{C1}\PYG{p}{,} \PYG{n}{n}\PYG{p}{,} \PYG{n}{label}\PYG{o}{=}\PYG{l+s+s2}{\PYGZdq{}}\PYG{l+s+s2}{Total membranes}\PYG{l+s+s2}{\PYGZdq{}}\PYG{p}{)}
\PYG{n}{plt}\PYG{o}{.}\PYG{n}{plot}\PYG{p}{(}\PYG{n}{C1}\PYG{p}{,} \PYG{n}{n1}\PYG{p}{,} \PYG{n}{label}\PYG{o}{=}\PYG{l+s+s2}{\PYGZdq{}}\PYG{l+s+s2}{Membrane 1}\PYG{l+s+s2}{\PYGZdq{}}\PYG{p}{)}
\PYG{n}{plt}\PYG{o}{.}\PYG{n}{plot}\PYG{p}{(}\PYG{n}{C1}\PYG{p}{,} \PYG{n}{n2}\PYG{p}{,} \PYG{n}{label}\PYG{o}{=}\PYG{l+s+s2}{\PYGZdq{}}\PYG{l+s+s2}{Membrane 2}\PYG{l+s+s2}{\PYGZdq{}}\PYG{p}{)}

\PYG{c+c1}{\PYGZsh{} Adding legend, labels, and title}
\PYG{n}{plt}\PYG{o}{.}\PYG{n}{legend}\PYG{p}{(}\PYG{p}{)}
\PYG{n}{plt}\PYG{o}{.}\PYG{n}{xlabel}\PYG{p}{(}\PYG{l+s+s2}{\PYGZdq{}}\PYG{l+s+s2}{C1}\PYG{l+s+s2}{\PYGZdq{}}\PYG{p}{)}
\PYG{n}{plt}\PYG{o}{.}\PYG{n}{ylabel}\PYG{p}{(}\PYG{l+s+s2}{\PYGZdq{}}\PYG{l+s+s2}{Number of Membranes}\PYG{l+s+s2}{\PYGZdq{}}\PYG{p}{)}
\PYG{n}{plt}\PYG{o}{.}\PYG{n}{title}\PYG{p}{(}\PYG{l+s+s2}{\PYGZdq{}}\PYG{l+s+s2}{Number of Membranes vs C1}\PYG{l+s+s2}{\PYGZdq{}}\PYG{p}{)}


\PYG{c+c1}{\PYGZsh{} Show plot}
\PYG{n}{plt}\PYG{o}{.}\PYG{n}{show}\PYG{p}{(}\PYG{p}{)}

\PYG{c+c1}{\PYGZsh{}Finding index of minimum number of membranes}
\PYG{n}{minN\PYGZus{}index} \PYG{o}{=} \PYG{n}{np}\PYG{o}{.}\PYG{n}{argmin}\PYG{p}{(}\PYG{n}{n}\PYG{p}{)}

\PYG{c+c1}{\PYGZsh{}Finding values at minimum number of membranes}
\PYG{n}{minN} \PYG{o}{=} \PYG{n}{n}\PYG{p}{[}\PYG{n}{minN\PYGZus{}index}\PYG{p}{]}
\PYG{n}{minN1} \PYG{o}{=} \PYG{n}{n1}\PYG{p}{[}\PYG{n}{minN\PYGZus{}index}\PYG{p}{]}
\PYG{n}{minN2} \PYG{o}{=} \PYG{n}{n2}\PYG{p}{[}\PYG{n}{minN\PYGZus{}index}\PYG{p}{]}
\PYG{n}{minC1} \PYG{o}{=} \PYG{n}{C1}\PYG{p}{[}\PYG{n}{minN\PYGZus{}index}\PYG{p}{]}

\PYG{n+nb}{print}\PYG{p}{(}\PYG{l+s+s2}{\PYGZdq{}}\PYG{l+s+se}{\PYGZbs{}n}\PYG{l+s+s2}{Stage 1 outlet concentration: }\PYG{l+s+s2}{\PYGZdq{}}\PYG{p}{,} \PYG{l+s+sa}{f}\PYG{l+s+s2}{\PYGZdq{}}\PYG{l+s+si}{\PYGZob{}}\PYG{n}{minC1}\PYG{l+s+si}{:}\PYG{l+s+s2}{.4}\PYG{l+s+si}{\PYGZcb{}}\PYG{l+s+s2}{\PYGZdq{}}\PYG{p}{,} \PYG{l+s+s2}{\PYGZdq{}}\PYG{l+s+s2}{ [kg/m\PYGZca{}3]}\PYG{l+s+s2}{\PYGZdq{}}\PYG{p}{)}
\PYG{n+nb}{print}\PYG{p}{(}\PYG{l+s+s2}{\PYGZdq{}}\PYG{l+s+se}{\PYGZbs{}n}\PYG{l+s+s2}{Minimum number of membranes for stage 1: }\PYG{l+s+s2}{\PYGZdq{}}\PYG{p}{,} \PYG{l+s+sa}{f}\PYG{l+s+s2}{\PYGZdq{}}\PYG{l+s+si}{\PYGZob{}}\PYG{n}{minN1}\PYG{l+s+si}{\PYGZcb{}}\PYG{l+s+s2}{\PYGZdq{}}\PYG{p}{,} \PYG{l+s+s2}{\PYGZdq{}}\PYG{l+s+s2}{ [\PYGZhy{}]}\PYG{l+s+s2}{\PYGZdq{}}\PYG{p}{)}
\PYG{n+nb}{print}\PYG{p}{(}\PYG{l+s+s2}{\PYGZdq{}}\PYG{l+s+se}{\PYGZbs{}n}\PYG{l+s+s2}{Minimum number of membranes for stage 2: }\PYG{l+s+s2}{\PYGZdq{}}\PYG{p}{,} \PYG{l+s+sa}{f}\PYG{l+s+s2}{\PYGZdq{}}\PYG{l+s+si}{\PYGZob{}}\PYG{n}{minN2}\PYG{l+s+si}{\PYGZcb{}}\PYG{l+s+s2}{\PYGZdq{}}\PYG{p}{,} \PYG{l+s+s2}{\PYGZdq{}}\PYG{l+s+s2}{ [\PYGZhy{}]}\PYG{l+s+s2}{\PYGZdq{}}\PYG{p}{)}
\PYG{n+nb}{print}\PYG{p}{(}\PYG{l+s+s2}{\PYGZdq{}}\PYG{l+s+se}{\PYGZbs{}n}\PYG{l+s+s2}{Minimum number of total membranes: }\PYG{l+s+s2}{\PYGZdq{}}\PYG{p}{,} \PYG{l+s+sa}{f}\PYG{l+s+s2}{\PYGZdq{}}\PYG{l+s+si}{\PYGZob{}}\PYG{n}{minN}\PYG{l+s+si}{\PYGZcb{}}\PYG{l+s+s2}{\PYGZdq{}}\PYG{p}{,} \PYG{l+s+s2}{\PYGZdq{}}\PYG{l+s+s2}{ [\PYGZhy{}]}\PYG{l+s+s2}{\PYGZdq{}}\PYG{p}{)}
\end{sphinxVerbatim}

\end{sphinxuseclass}\end{sphinxVerbatimInput}
\begin{sphinxVerbatimOutput}

\begin{sphinxuseclass}{cell_output}
\noindent\sphinxincludegraphics{{2422669e8ec62b36a5e7ba0267c0434d7fe2a3521d173581fd75a92e9fa3f449}.png}

\begin{sphinxVerbatim}[commandchars=\\\{\}]
Stage 1 outlet concentration:  2.015  [kg/m\PYGZca{}3]

Minimum number of membranes for stage 1:  8.0  [\PYGZhy{}]

Minimum number of membranes for stage 2:  5.0  [\PYGZhy{}]

Minimum number of total membranes:  13.0  [\PYGZhy{}]
\end{sphinxVerbatim}

\end{sphinxuseclass}\end{sphinxVerbatimOutput}

\end{sphinxuseclass}
\sphinxAtStartPar
This can also be solved with GAMS:

\begin{sphinxVerbatim}[commandchars=\\\{\}]
\PYG{c+c1}{\PYGZsh{}Note: This code cannot be run in Jupyter notebook, copy paste into GAMS.}

\PYG{n}{Parameters}
\PYG{n}{Q0}\PYG{p}{,}\PYG{n}{C0}\PYG{p}{,}\PYG{n}{C2}\PYG{p}{;}
\PYG{n}{Q0}\PYG{o}{=}\PYG{l+m+mi}{200}\PYG{o}{/}\PYG{l+m+mi}{24}\PYG{p}{;}
\PYG{n}{C0}\PYG{o}{=}\PYG{l+m+mf}{0.5}\PYG{p}{;}
\PYG{n}{C2}\PYG{o}{=}\PYG{l+m+mi}{20}\PYG{p}{;}

\PYG{n}{Variables}
\PYG{n}{n}\PYG{p}{,} \PYG{n}{Q1}\PYG{p}{,} \PYG{n}{C1}\PYG{p}{,} \PYG{n}{Q2}\PYG{p}{,} \PYG{n}{A1}\PYG{p}{,} \PYG{n}{A2}
\PYG{p}{;}

\PYG{n}{Equations}
\PYG{n}{mb1}\PYG{p}{,} \PYG{n}{mb2}\PYG{p}{,} \PYG{n}{area}\PYG{p}{,} \PYG{n}{vb1}\PYG{p}{,} \PYG{n}{vb2}
\PYG{p}{;}

\PYG{n}{mb1}\PYG{o}{.}\PYG{o}{.} \PYG{n}{Q0}\PYG{o}{=}\PYG{n}{e}\PYG{o}{=}\PYG{p}{(}\PYG{l+m+mf}{3.8}\PYG{o}{*}\PYG{l+m+mi}{10}\PYG{o}{*}\PYG{o}{*}\PYG{p}{(}\PYG{o}{\PYGZhy{}}\PYG{l+m+mi}{2}\PYG{p}{)}\PYG{p}{)}\PYG{o}{*}\PYG{n}{log}\PYG{p}{(}\PYG{l+m+mi}{145}\PYG{o}{/}\PYG{n}{C1}\PYG{p}{)}\PYG{o}{*}\PYG{n}{A1}\PYG{o}{+}\PYG{n}{Q1}\PYG{p}{;}
\PYG{n}{vb1}\PYG{o}{.}\PYG{o}{.} \PYG{n}{Q1}\PYG{o}{=}\PYG{n}{e}\PYG{o}{=}\PYG{n}{Q0}\PYG{o}{*}\PYG{n}{C0}\PYG{o}{/}\PYG{n}{C1}\PYG{p}{;}
\PYG{n}{mb2}\PYG{o}{.}\PYG{o}{.} \PYG{n}{Q1}\PYG{o}{=}\PYG{n}{e}\PYG{o}{=}\PYG{p}{(}\PYG{l+m+mf}{3.8}\PYG{o}{*}\PYG{l+m+mi}{10}\PYG{o}{*}\PYG{o}{*}\PYG{p}{(}\PYG{o}{\PYGZhy{}}\PYG{l+m+mi}{2}\PYG{p}{)}\PYG{p}{)}\PYG{o}{*}\PYG{n}{log}\PYG{p}{(}\PYG{l+m+mi}{145}\PYG{o}{/}\PYG{l+m+mi}{20}\PYG{p}{)}\PYG{o}{*}\PYG{n}{A2}\PYG{o}{+}\PYG{n}{Q2}\PYG{p}{;}
\PYG{n}{vb2}\PYG{o}{.}\PYG{o}{.} \PYG{n}{Q2}\PYG{o}{=}\PYG{n}{e}\PYG{o}{=}\PYG{n}{Q1}\PYG{o}{*}\PYG{n}{C1}\PYG{o}{/}\PYG{n}{C2}\PYG{p}{;}
\PYG{n}{area}\PYG{o}{.}\PYG{o}{.} \PYG{n}{n}\PYG{o}{=}\PYG{n}{e}\PYG{o}{=}\PYG{n}{A1}\PYG{o}{/}\PYG{l+m+mi}{5}\PYG{o}{+}\PYG{n}{A2}\PYG{o}{/}\PYG{l+m+mi}{5}\PYG{p}{;}

\PYG{n}{Model} \PYG{n}{process} \PYG{o}{/}\PYG{n+nb}{all}\PYG{o}{/}\PYG{p}{;}
\PYG{n}{c1}\PYG{o}{.}\PYG{n}{l}\PYG{o}{=}\PYG{l+m+mi}{5}\PYG{p}{;}
\PYG{n}{Solve} \PYG{n}{process} \PYG{n}{using} \PYG{n}{NLP} \PYG{n}{minimising} \PYG{n}{n}\PYG{p}{;}
\end{sphinxVerbatim}


\subsection{Refinement of the GAMS Solution:}
\label{\detokenize{Exercise_W4_2024_updated:refinement-of-the-gams-solution}}
\sphinxAtStartPar
\sphinxstyleemphasis{{[}Contributed by Divyansh Vashishtha, March 2024{]}}

\sphinxAtStartPar
In this version only an integer number of membrane modules can be adopted in every stage. In this version, only an integer number of membrane modules can be adopted in every stage.

\begin{sphinxVerbatim}[commandchars=\\\{\}]
\PYG{n}{Parameters}
\PYG{n}{Q0} \PYG{n}{intial} \PYG{n}{volumetric} \PYG{n}{flowrate} \PYG{n}{coming} \PYG{n}{into} \PYG{n}{stage} \PYG{l+m+mi}{1} \PYG{p}{[}\PYG{n}{m}\PYG{o}{\PYGZca{}}\PYG{l+m+mi}{3} \PYG{n}{h}\PYG{o}{\PYGZca{}}\PYG{o}{\PYGZhy{}}\PYG{l+m+mi}{1}\PYG{p}{]}\PYG{p}{,}
\PYG{n}{C0} \PYG{n}{initial} \PYG{n}{protein} \PYG{n}{concentration} \PYG{n}{coming} \PYG{n}{into} \PYG{n}{stage} \PYG{l+m+mi}{1}  \PYG{p}{[}\PYG{n}{kg} \PYG{n}{m}\PYG{o}{\PYGZca{}}\PYG{o}{\PYGZhy{}}\PYG{l+m+mi}{3} \PYG{p}{]}\PYG{p}{,}\PYG{n}{C2} \PYG{n}{final} \PYG{n}{protein} \PYG{n}{concentration} \PYG{n}{leaving} \PYG{n}{stage} \PYG{l+m+mi}{2} \PYG{p}{[}\PYG{n}{kg} \PYG{n}{m}\PYG{o}{\PYGZca{}}\PYG{o}{\PYGZhy{}}\PYG{l+m+mi}{3}\PYG{p}{]}\PYG{p}{;}
\PYG{o}{*}\PYG{n}{Initialisation} \PYG{n}{of} \PYG{n}{Parameters} \PYG{k+kn}{from}\PYG{+w}{ }\PYG{n+nn}{the} \PYG{n}{Process} \PYG{n}{Description}\PYG{p}{:} \PYG{n}{Q0}\PYG{p}{,} \PYG{n}{C0}\PYG{p}{,} \PYG{n}{C2}\PYG{o}{.}
\PYG{n}{Q0}\PYG{o}{=}\PYG{l+m+mi}{200}\PYG{o}{/}\PYG{l+m+mi}{24}\PYG{p}{;}
\PYG{n}{C0}\PYG{o}{=}\PYG{l+m+mf}{0.5}\PYG{p}{;}
\PYG{n}{C2}\PYG{o}{=}\PYG{l+m+mi}{20}\PYG{p}{;}

\PYG{n}{Variables}
\PYG{n}{n} \PYG{n}{total} \PYG{n}{number} \PYG{n}{of} \PYG{n}{membrane} \PYG{n}{modules} \PYG{n}{across} \PYG{n}{the} \PYG{n}{two} \PYG{n}{stages} \PYG{p}{[}\PYG{n}{decimal}\PYG{p}{]}\PYG{p}{,}
\PYG{n}{Q1} \PYG{n}{volumetric} \PYG{n}{flowrate} \PYG{n}{leaving} \PYG{n}{the} \PYG{n}{first} \PYG{n}{stage} \PYG{p}{[}\PYG{n}{m}\PYG{o}{\PYGZca{}}\PYG{l+m+mi}{3} \PYG{n}{h}\PYG{o}{\PYGZca{}}\PYG{o}{\PYGZhy{}}\PYG{l+m+mi}{1}\PYG{p}{]}\PYG{p}{,}
\PYG{n}{C1} \PYG{n}{concentration} \PYG{n}{of} \PYG{n}{protein} \PYG{n}{leaving} \PYG{n}{the} \PYG{n}{first} \PYG{n}{stage} \PYG{p}{[}\PYG{n}{kg} \PYG{n}{m}\PYG{o}{\PYGZca{}}\PYG{o}{\PYGZhy{}}\PYG{l+m+mi}{3}\PYG{p}{]}\PYG{p}{,}
\PYG{n}{Q2} \PYG{n}{volumetric} \PYG{n}{flowrate} \PYG{n}{leaving} \PYG{n}{the} \PYG{n}{second} \PYG{n}{stage} \PYG{p}{[}\PYG{n}{m}\PYG{o}{\PYGZca{}}\PYG{l+m+mi}{3} \PYG{n}{h}\PYG{o}{\PYGZca{}}\PYG{o}{\PYGZhy{}}\PYG{l+m+mi}{1}\PYG{p}{]}\PYG{p}{,}
\PYG{n}{A1} \PYG{n}{Total} \PYG{n}{area} \PYG{n}{of} \PYG{n}{membranes} \PYG{n}{required} \PYG{o+ow}{in} \PYG{n}{stage} \PYG{l+m+mi}{1} \PYG{p}{[}\PYG{n}{m}\PYG{o}{\PYGZca{}}\PYG{l+m+mi}{2}\PYG{p}{]}\PYG{p}{,}
\PYG{n}{A2} \PYG{n}{Total} \PYG{n}{area} \PYG{n}{of} \PYG{n}{membranes} \PYG{n}{required} \PYG{o+ow}{in} \PYG{n}{stage} \PYG{l+m+mi}{2} \PYG{p}{[}\PYG{n}{m}\PYG{o}{\PYGZca{}}\PYG{l+m+mi}{2}\PYG{p}{]}\PYG{p}{,}
\PYG{n}{n1} \PYG{n}{number} \PYG{n}{of} \PYG{n}{membrane} \PYG{n}{modules} \PYG{n}{required} \PYG{k}{for} \PYG{n}{stage} \PYG{l+m+mi}{1} \PYG{p}{[}\PYG{n}{decimal}\PYG{p}{]}\PYG{p}{,}
\PYG{n}{n1\PYGZus{}rounded} \PYG{n}{minimum} \PYG{n}{number} \PYG{n}{of} \PYG{n}{membrane} \PYG{n}{modules} \PYG{n}{required} \PYG{k}{for} \PYG{n}{stage} \PYG{l+m+mi}{1} \PYG{p}{[}\PYG{n}{integer}\PYG{p}{]}\PYG{p}{,}
\PYG{n}{n2} \PYG{n}{number} \PYG{n}{of} \PYG{n}{membrane} \PYG{n}{modules} \PYG{n}{required} \PYG{k}{for} \PYG{n}{stage} \PYG{l+m+mi}{2} \PYG{p}{[}\PYG{n}{decimal}\PYG{p}{]}\PYG{p}{,}
\PYG{n}{n2\PYGZus{}rounded} \PYG{n}{minimum} \PYG{n}{number} \PYG{n}{of} \PYG{n}{membrane} \PYG{n}{modules} \PYG{n}{required} \PYG{k}{for} \PYG{n}{stage} \PYG{l+m+mi}{2} \PYG{p}{[}\PYG{n}{integer}\PYG{p}{]}\PYG{p}{,}
\PYG{n}{nx} \PYG{n}{rounded} \PYG{n}{up} \PYG{n}{value} \PYG{n}{of} \PYG{n}{the} \PYG{n}{total} \PYG{n}{membranes} \PYG{n}{required} \PYG{p}{[}\PYG{n}{integer}\PYG{p}{]}\PYG{o}{.}
\PYG{p}{;}

\PYG{n}{Equations}
\PYG{n}{mb1} \PYG{n}{equation} \PYG{k}{for} \PYG{n}{the} \PYG{n}{mass} \PYG{n}{balane} \PYG{n}{of} \PYG{n}{the} \PYG{n}{solute} \PYG{n}{over} \PYG{n}{stage} \PYG{l+m+mi}{1}\PYG{p}{,}
\PYG{n}{mb2} \PYG{n}{equation} \PYG{k}{for} \PYG{n}{the} \PYG{n}{mass} \PYG{n}{balance} \PYG{n}{of} \PYG{n}{the} \PYG{n}{solute} \PYG{n}{over} \PYG{n}{stage} \PYG{l+m+mi}{2}\PYG{p}{,}
\PYG{n}{area1} \PYG{n}{equation} \PYG{k}{for} \PYG{n}{the} \PYG{n}{area} \PYG{n}{of} \PYG{n}{membrane} \PYG{n}{required} \PYG{n}{over} \PYG{n}{the} \PYG{n}{first} \PYG{n}{stage}\PYG{p}{,}
\PYG{n}{vb1} \PYG{n}{equation} \PYG{k}{for} \PYG{n}{the} \PYG{n}{volume} \PYG{n}{balance} \PYG{n}{over} \PYG{n}{stage} \PYG{l+m+mi}{1}\PYG{p}{,}
\PYG{n}{vb2} \PYG{n}{equation} \PYG{k}{for} \PYG{n}{the} \PYG{n}{volume} \PYG{n}{balance} \PYG{n}{over} \PYG{n}{stage} \PYG{l+m+mi}{2}\PYG{p}{,}
\PYG{n}{area2} \PYG{n}{equation} \PYG{k}{for} \PYG{n}{the} \PYG{n}{area} \PYG{n}{of} \PYG{n}{membrane} \PYG{n}{required} \PYG{n}{over} \PYG{n}{the} \PYG{n}{second} \PYG{n}{stage}\PYG{p}{,}
\PYG{n}{roundn1} \PYG{n}{equation} \PYG{n}{that} \PYG{n}{rounds} \PYG{n}{up} \PYG{n}{the} \PYG{n}{decimal} \PYG{n}{value} \PYG{n}{of} \PYG{n}{the} \PYG{n}{number} \PYG{n}{of} \PYG{n}{membrane} \PYG{n}{modules} \PYG{n}{obtained} \PYG{o+ow}{in} \PYG{n}{the} \PYG{n}{first} \PYG{n}{stage} \PYG{p}{,}
\PYG{n}{roundn2} \PYG{n}{equation} \PYG{n}{that} \PYG{n}{rounds} \PYG{n}{up} \PYG{n}{the} \PYG{n}{decimal} \PYG{n}{value} \PYG{n}{of} \PYG{n}{the} \PYG{n}{number} \PYG{n}{of} \PYG{n}{membrane} \PYG{n}{modules} \PYG{n}{obtained} \PYG{o+ow}{in} \PYG{n}{the} \PYG{n}{second} \PYG{n}{stage}\PYG{p}{,}
\PYG{n}{total\PYGZus{}n} \PYG{n}{equation} \PYG{n}{that} \PYG{n}{determines} \PYG{n}{the} \PYG{n}{total} \PYG{n}{number} \PYG{n}{of} \PYG{n}{membrane} \PYG{n}{modules} \PYG{n}{obtained} \PYG{n}{over} \PYG{n}{the} \PYG{n}{two} \PYG{n}{stages} \PYG{p}{[}\PYG{n}{returns} \PYG{n}{a} \PYG{n}{decimal}\PYG{p}{]}\PYG{p}{,}
\PYG{n}{roundN} \PYG{n}{equation} \PYG{n}{that} \PYG{n}{displays} \PYG{n}{the} \PYG{n}{rounded} \PYG{n}{number} \PYG{n}{of} \PYG{n}{total} \PYG{n}{membrane} \PYG{n}{modules} \PYG{n}{over} \PYG{n}{the} \PYG{n}{two} \PYG{n}{stages}\PYG{p}{;}

\PYG{n}{vb1}\PYG{o}{.}\PYG{o}{.} \PYG{n}{Q0}\PYG{o}{=}\PYG{n}{e}\PYG{o}{=}\PYG{p}{(}\PYG{l+m+mf}{3.8}\PYG{o}{*}\PYG{l+m+mi}{10}\PYG{o}{*}\PYG{o}{*}\PYG{p}{(}\PYG{o}{\PYGZhy{}}\PYG{l+m+mi}{2}\PYG{p}{)}\PYG{p}{)}\PYG{o}{*}\PYG{n}{log}\PYG{p}{(}\PYG{l+m+mi}{145}\PYG{o}{/}\PYG{n}{C1}\PYG{p}{)}\PYG{o}{*}\PYG{n}{A1}\PYG{o}{+}\PYG{n}{Q1}\PYG{p}{;}
\PYG{n}{mb1}\PYG{o}{.}\PYG{o}{.} \PYG{n}{Q1}\PYG{o}{=}\PYG{n}{e}\PYG{o}{=}\PYG{n}{Q0}\PYG{o}{*}\PYG{n}{C0}\PYG{o}{/}\PYG{n}{C1}\PYG{p}{;}
\PYG{n}{vb2}\PYG{o}{.}\PYG{o}{.} \PYG{n}{Q1}\PYG{o}{=}\PYG{n}{e}\PYG{o}{=}\PYG{p}{(}\PYG{l+m+mf}{3.8}\PYG{o}{*}\PYG{l+m+mi}{10}\PYG{o}{*}\PYG{o}{*}\PYG{p}{(}\PYG{o}{\PYGZhy{}}\PYG{l+m+mi}{2}\PYG{p}{)}\PYG{p}{)}\PYG{o}{*}\PYG{n}{log}\PYG{p}{(}\PYG{l+m+mi}{145}\PYG{o}{/}\PYG{n}{C2}\PYG{p}{)}\PYG{o}{*}\PYG{n}{A2}\PYG{o}{+}\PYG{n}{Q2}\PYG{p}{;}
\PYG{n}{mb2}\PYG{o}{.}\PYG{o}{.} \PYG{n}{Q2}\PYG{o}{=}\PYG{n}{e}\PYG{o}{=}\PYG{n}{Q1}\PYG{o}{*}\PYG{n}{C1}\PYG{o}{/}\PYG{n}{C2}\PYG{p}{;}
\PYG{n}{area1}\PYG{o}{.}\PYG{o}{.} \PYG{n}{n1} \PYG{o}{=}\PYG{n}{e}\PYG{o}{=} \PYG{n}{A1}\PYG{o}{/}\PYG{l+m+mi}{5}\PYG{p}{;}
\PYG{n}{area2}\PYG{o}{.}\PYG{o}{.} \PYG{n}{n2} \PYG{o}{=}\PYG{n}{e}\PYG{o}{=} \PYG{n}{A2}\PYG{o}{/}\PYG{l+m+mi}{5}\PYG{p}{;}
\PYG{n}{roundn1}\PYG{o}{.}\PYG{o}{.} \PYG{n}{n1\PYGZus{}rounded} \PYG{o}{=}\PYG{n}{e}\PYG{o}{=} \PYG{n}{ceil}\PYG{p}{(}\PYG{n}{n1}\PYG{p}{)}\PYG{p}{;}
\PYG{n}{roundn2}\PYG{o}{.}\PYG{o}{.} \PYG{n}{n2\PYGZus{}rounded} \PYG{o}{=}\PYG{n}{e}\PYG{o}{=} \PYG{n}{ceil}\PYG{p}{(}\PYG{n}{n2}\PYG{p}{)}\PYG{p}{;}
\PYG{n}{total\PYGZus{}n}\PYG{o}{.}\PYG{o}{.} \PYG{n}{n} \PYG{o}{=}\PYG{n}{e}\PYG{o}{=} \PYG{n}{n1} \PYG{o}{+} \PYG{n}{n2}\PYG{p}{;}
\PYG{n}{roundN}\PYG{o}{.}\PYG{o}{.} \PYG{n}{nx} \PYG{o}{=}\PYG{n}{e}\PYG{o}{=} \PYG{n}{ceil}\PYG{p}{(}\PYG{n}{n}\PYG{p}{)}\PYG{p}{;}

\PYG{n}{Model} \PYG{n}{process} \PYG{o}{/}\PYG{n+nb}{all}\PYG{o}{/}\PYG{p}{;}
\PYG{n}{c1}\PYG{o}{.}\PYG{n}{l}\PYG{o}{=}\PYG{l+m+mi}{5}\PYG{p}{;}
\PYG{n}{Solve} \PYG{n}{process} \PYG{n}{using} \PYG{n}{DNLP} \PYG{n}{minimising} \PYG{n}{n}\PYG{p}{;}
\PYG{n}{display} \PYG{n}{n}\PYG{o}{.}\PYG{n}{l}\PYG{p}{,} \PYG{n}{n1\PYGZus{}rounded}\PYG{o}{.}\PYG{n}{l}\PYG{p}{,} \PYG{n}{n2\PYGZus{}rounded}\PYG{o}{.}\PYG{n}{l}\PYG{p}{,} \PYG{n}{nx}\PYG{o}{.}\PYG{n}{l}\PYG{p}{;}
\end{sphinxVerbatim}

\sphinxstepscope


\section{Week 5 \sphinxhyphen{} Exercise}
\label{\detokenize{Exercise_W5_2021:week-5-exercise}}\label{\detokenize{Exercise_W5_2021::doc}}

\subsection{Problem Statement}
\label{\detokenize{Exercise_W5_2021:problem-statement}}
\sphinxAtStartPar
A protein mixture has been separated through an ion\sphinxhyphen{}exchange chromatography process. The residence time of the column in non\sphinxhyphen{}adsorbing conditons is \(t_0=1\) min.
The chromatogram resulting from the separation is reported in the following:

\begin{figure}[htbp]
\centering
\capstart

\noindent\sphinxincludegraphics[height=300\sphinxpxdimen]{{chromatogram}.png}
\caption{Experimental chromatogram.}\label{\detokenize{Exercise_W5_2021:scheme}}\end{figure}


\subsubsection{Determine:}
\label{\detokenize{Exercise_W5_2021:determine}}\begin{itemize}
\item {} 
\sphinxAtStartPar
How many components in the mixture can you identify from the chromatogram?

\item {} 
\sphinxAtStartPar
What is the selectivity of this column for each couple of components?

\item {} 
\sphinxAtStartPar
Assuming a linear adsorption isotherm holds in all cases, what conclusions can you draw about adsorption thermodynamics of each component in the column?

\item {} 
\sphinxAtStartPar
What is the number of theoretical plates associated to the elution of each component you have identified? If the coulmn is 10 cm long, what is the height equivalent of a theoretical plate for each component of the mixture?

\item {} 
\sphinxAtStartPar
Compute the resolution relative to each couple of components in the chromatogram. Does it depend only on the adsorption thermodynamics?

\item {} 
\sphinxAtStartPar
Comment on the separation performances of the column. Does it allow for a sharp resolution of all components?

\end{itemize}


\subsubsection{Solution}
\label{\detokenize{Exercise_W5_2021:solution}}\begin{enumerate}
\sphinxsetlistlabels{\arabic}{enumi}{enumii}{}{.}%
\item {} 
\sphinxAtStartPar
The chromatogram clearly indicates three components, let us indicate them with A, B, and C.
The retention times for these components are:

\end{enumerate}
\begin{itemize}
\item {} 
\sphinxAtStartPar
\(t_{r,A}=7\)

\item {} 
\sphinxAtStartPar
\(t_{r,B}=10\)

\item {} 
\sphinxAtStartPar
\(t_{r,C}=15\)

\end{itemize}
\begin{enumerate}
\sphinxsetlistlabels{\arabic}{enumi}{enumii}{}{.}%
\setcounter{enumi}{1}
\item {} 
\sphinxAtStartPar
The selectivity can be computed for every pair of components as:

\end{enumerate}
\begin{equation*}
\begin{split}
S_{i,j}=\frac{t_{r,i}-t_0}{t_{r,j}-t_0}
\end{split}
\end{equation*}
\begin{sphinxuseclass}{cell}\begin{sphinxVerbatimInput}

\begin{sphinxuseclass}{cell_input}
\begin{sphinxVerbatim}[commandchars=\\\{\}]
\PYG{k+kn}{import}\PYG{+w}{ }\PYG{n+nn}{numpy}\PYG{+w}{ }\PYG{k}{as}\PYG{+w}{ }\PYG{n+nn}{np}

\PYG{n}{tr}\PYG{o}{=}\PYG{n}{np}\PYG{o}{.}\PYG{n}{array}\PYG{p}{(}\PYG{p}{[}\PYG{l+m+mi}{7}\PYG{p}{,}\PYG{l+m+mi}{10}\PYG{p}{,}\PYG{l+m+mi}{15}\PYG{p}{]}\PYG{p}{)}\PYG{p}{;} 
\PYG{n}{t0}\PYG{o}{=}\PYG{l+m+mi}{1}
\PYG{n}{SEL}\PYG{o}{=}\PYG{n}{np}\PYG{o}{.}\PYG{n}{zeros}\PYG{p}{(}\PYG{p}{(}\PYG{n}{np}\PYG{o}{.}\PYG{n}{size}\PYG{p}{(}\PYG{n}{tr}\PYG{p}{)}\PYG{p}{,}\PYG{n}{np}\PYG{o}{.}\PYG{n}{size}\PYG{p}{(}\PYG{n}{tr}\PYG{p}{)}\PYG{p}{)}\PYG{p}{)}\PYG{p}{;}

\PYG{k}{for} \PYG{n}{i} \PYG{o+ow}{in} \PYG{n+nb}{range}\PYG{p}{(}\PYG{l+m+mi}{0}\PYG{p}{,}\PYG{n}{np}\PYG{o}{.}\PYG{n}{size}\PYG{p}{(}\PYG{n}{tr}\PYG{p}{)}\PYG{p}{)}\PYG{p}{:}
    \PYG{k}{for} \PYG{n}{j} \PYG{o+ow}{in} \PYG{n+nb}{range}\PYG{p}{(}\PYG{l+m+mi}{0}\PYG{p}{,}\PYG{n}{np}\PYG{o}{.}\PYG{n}{size}\PYG{p}{(}\PYG{n}{tr}\PYG{p}{)}\PYG{p}{)}\PYG{p}{:}
        \PYG{n}{SEL}\PYG{p}{[}\PYG{n}{i}\PYG{p}{,}\PYG{n}{j}\PYG{p}{]}\PYG{o}{=}\PYG{p}{(}\PYG{n}{tr}\PYG{p}{[}\PYG{n}{i}\PYG{p}{]}\PYG{o}{\PYGZhy{}}\PYG{n}{t0}\PYG{p}{)}\PYG{o}{/}\PYG{p}{(}\PYG{n}{tr}\PYG{p}{[}\PYG{n}{j}\PYG{p}{]}\PYG{o}{\PYGZhy{}}\PYG{n}{t0}\PYG{p}{)}

\PYG{n+nb}{print}\PYG{p}{(}\PYG{l+s+s1}{\PYGZsq{}}\PYG{l+s+s1}{Selectivity i,j}\PYG{l+s+se}{\PYGZbs{}n}\PYG{l+s+s1}{\PYGZsq{}}\PYG{p}{,}\PYG{n}{SEL}\PYG{p}{)}
\end{sphinxVerbatim}

\end{sphinxuseclass}\end{sphinxVerbatimInput}
\begin{sphinxVerbatimOutput}

\begin{sphinxuseclass}{cell_output}
\begin{sphinxVerbatim}[commandchars=\\\{\}]
Selectivity i,j
 [[1.         0.66666667 0.42857143]
 [1.5        1.         0.64285714]
 [2.33333333 1.55555556 1.        ]]
\end{sphinxVerbatim}

\end{sphinxuseclass}\end{sphinxVerbatimOutput}

\end{sphinxuseclass}\begin{enumerate}
\sphinxsetlistlabels{\arabic}{enumi}{enumii}{}{.}%
\setcounter{enumi}{2}
\item {} 
\sphinxAtStartPar
Since \(S_{i,j}=H_i/H_j\) we can determine the relative affinity of components A, B and C for the startionary phase. Component C shows the highest affinity for the stationary phase, followed by B and then C.

\end{enumerate}

\sphinxAtStartPar
4/5. The number of theoretical plates associated to the elution of components A, B, and C can be obtained applying the relationship \(N=16*(tr/tw)^2\). The resolution for every couple of components is computed as: \(R_{i,j}=\frac{t_{r,i}-t_{r,j}}{\frac{1}{2}(t_{w,i}+t_{w,j})}\)

\begin{sphinxuseclass}{cell}\begin{sphinxVerbatimInput}

\begin{sphinxuseclass}{cell_input}
\begin{sphinxVerbatim}[commandchars=\\\{\}]
\PYG{k+kn}{import}\PYG{+w}{ }\PYG{n+nn}{numpy}\PYG{+w}{ }\PYG{k}{as}\PYG{+w}{ }\PYG{n+nn}{np}

\PYG{n}{tr}\PYG{o}{=}\PYG{n}{np}\PYG{o}{.}\PYG{n}{array}\PYG{p}{(}\PYG{p}{[}\PYG{l+m+mi}{7}\PYG{p}{,}\PYG{l+m+mi}{10}\PYG{p}{,}\PYG{l+m+mi}{15}\PYG{p}{]}\PYG{p}{)}\PYG{p}{;} 
\PYG{n}{tw}\PYG{o}{=}\PYG{n}{np}\PYG{o}{.}\PYG{n}{array}\PYG{p}{(}\PYG{p}{[}\PYG{l+m+mi}{1}\PYG{p}{,}\PYG{l+m+mi}{2}\PYG{p}{,}\PYG{l+m+mf}{1.5}\PYG{p}{]}\PYG{p}{)}\PYG{p}{;} 
\PYG{n}{L}\PYG{o}{=}\PYG{l+m+mi}{10} \PYG{c+c1}{\PYGZsh{}cm}

\PYG{n}{N}\PYG{o}{=}\PYG{p}{(}\PYG{l+m+mi}{16}\PYG{o}{*}\PYG{p}{(}\PYG{n}{tr}\PYG{o}{/}\PYG{n}{tw}\PYG{p}{)}\PYG{o}{*}\PYG{o}{*}\PYG{l+m+mi}{2}\PYG{p}{)}

\PYG{n}{H}\PYG{o}{=}\PYG{n}{L}\PYG{o}{/}\PYG{n}{N}

\PYG{n}{R}\PYG{o}{=}\PYG{n}{np}\PYG{o}{.}\PYG{n}{zeros}\PYG{p}{(}\PYG{p}{(}\PYG{n}{np}\PYG{o}{.}\PYG{n}{size}\PYG{p}{(}\PYG{n}{tr}\PYG{p}{)}\PYG{p}{,}\PYG{n}{np}\PYG{o}{.}\PYG{n}{size}\PYG{p}{(}\PYG{n}{tr}\PYG{p}{)}\PYG{p}{)}\PYG{p}{)}\PYG{p}{;}

\PYG{k}{for} \PYG{n}{i} \PYG{o+ow}{in} \PYG{n+nb}{range}\PYG{p}{(}\PYG{l+m+mi}{0}\PYG{p}{,}\PYG{n}{np}\PYG{o}{.}\PYG{n}{size}\PYG{p}{(}\PYG{n}{tr}\PYG{p}{)}\PYG{p}{)}\PYG{p}{:}
    \PYG{k}{for} \PYG{n}{j} \PYG{o+ow}{in} \PYG{n+nb}{range}\PYG{p}{(}\PYG{l+m+mi}{0}\PYG{p}{,}\PYG{n}{np}\PYG{o}{.}\PYG{n}{size}\PYG{p}{(}\PYG{n}{tr}\PYG{p}{)}\PYG{p}{)}\PYG{p}{:}
        \PYG{n}{R}\PYG{p}{[}\PYG{n}{i}\PYG{p}{,}\PYG{n}{j}\PYG{p}{]}\PYG{o}{=}\PYG{n+nb}{abs}\PYG{p}{(}\PYG{p}{(}\PYG{n}{tr}\PYG{p}{[}\PYG{n}{i}\PYG{p}{]}\PYG{o}{\PYGZhy{}}\PYG{n}{tr}\PYG{p}{[}\PYG{n}{j}\PYG{p}{]}\PYG{p}{)}\PYG{o}{/}\PYG{p}{(}\PYG{l+m+mi}{2}\PYG{o}{*}\PYG{p}{(}\PYG{n}{tw}\PYG{p}{[}\PYG{n}{i}\PYG{p}{]}\PYG{o}{+}\PYG{n}{tw}\PYG{p}{[}\PYG{n}{j}\PYG{p}{]}\PYG{p}{)}\PYG{p}{)}\PYG{p}{)}


\PYG{n+nb}{print}\PYG{p}{(}\PYG{l+s+s1}{\PYGZsq{}}\PYG{l+s+se}{\PYGZbs{}n}\PYG{l+s+s1}{Height of a theoretical plates }\PYG{l+s+s1}{\PYGZsq{}}\PYG{p}{,}\PYG{n}{H}\PYG{p}{,}\PYG{l+s+s1}{\PYGZsq{}}\PYG{l+s+s1}{ [\PYGZhy{}]}\PYG{l+s+s1}{\PYGZsq{}}\PYG{p}{)}
\PYG{n+nb}{print}\PYG{p}{(}\PYG{l+s+s1}{\PYGZsq{}}\PYG{l+s+s1}{Number of theoretical plates }\PYG{l+s+s1}{\PYGZsq{}}\PYG{p}{,}\PYG{n}{H}\PYG{p}{,}\PYG{l+s+s1}{\PYGZsq{}}\PYG{l+s+s1}{ cm}\PYG{l+s+s1}{\PYGZsq{}}\PYG{p}{)}
\PYG{n+nb}{print}\PYG{p}{(}\PYG{l+s+s1}{\PYGZsq{}}\PYG{l+s+s1}{Resolution i,j}\PYG{l+s+se}{\PYGZbs{}n}\PYG{l+s+s1}{\PYGZsq{}}\PYG{p}{,}\PYG{n}{R}\PYG{p}{)}
\end{sphinxVerbatim}

\end{sphinxuseclass}\end{sphinxVerbatimInput}
\begin{sphinxVerbatimOutput}

\begin{sphinxuseclass}{cell_output}
\begin{sphinxVerbatim}[commandchars=\\\{\}]
Height of a theoretical plates  [0.0127551 0.025     0.00625  ]  [\PYGZhy{}]
Number of theoretical plates  [0.0127551 0.025     0.00625  ]  cm
Resolution i,j
 [[0.         0.5        1.6       ]
 [0.5        0.         0.71428571]
 [1.6        0.71428571 0.        ]]
\end{sphinxVerbatim}

\end{sphinxuseclass}\end{sphinxVerbatimOutput}

\end{sphinxuseclass}
\sphinxstepscope


\section{Week 6 \sphinxhyphen{} Exercise}
\label{\detokenize{Exercise_W6_2021:week-6-exercise}}\label{\detokenize{Exercise_W6_2021::doc}}

\subsection{Problem Statement}
\label{\detokenize{Exercise_W6_2021:problem-statement}}
\sphinxAtStartPar
Continuous slurry adsorption is carried out to reduce the pollutant content in an aqueous stream from an inlet concentration of 1x10\(^{-2}\) mol l\(^{-1}\) . The flowrate of the aqueous stream is 50 m\(^3\) h\(^{-1}\) .
The adsorbent flowrate is 750 kg h\(^{-1}\). The residence time is 3 h.

\sphinxAtStartPar
Adsorption is controlled by external mass transfer with a coefficient k\(_L\)= 0.5 m h\(^{-1}\).
The adsorbent particles have a surface area of 10 m\(^2\) kg\(^{-1}\) and a density of 1850 kg m\(^{-3}\).
Consider a linear isotherm of the form \(q=kc\), where \(q\) {[}mol kg\(^{-1}\){]} is the amount of phenol adsorbed per kg of adsorbent, \(c\) is the molar concentration of pollutant in solution and \(k=1\times{10}^4\) {[}l kg\(^{-1}\){]} is the partition coefficient.

\sphinxAtStartPar
Determine the steady state concentration of pollutant in the aqueous stream leaving the unit.

\sphinxstepscope


\section{Week 8 \sphinxhyphen{} Exercise}
\label{\detokenize{Exercise_W8_2021:week-8-exercise}}\label{\detokenize{Exercise_W8_2021::doc}}

\subsection{Problem Statement 1}
\label{\detokenize{Exercise_W8_2021:problem-statement-1}}\begin{itemize}
\item {} 
\sphinxAtStartPar
In a system with two monotropic polymorphs, is it true that the polymorph with the lowest solubility is thermodynamically more stable?

\item {} 
\sphinxAtStartPar
Can you demonstrate it using the definition of solubility?

\end{itemize}


\subsection{Solution Trace}
\label{\detokenize{Exercise_W8_2021:solution-trace}}
\sphinxAtStartPar
\sphinxstyleemphasis{{[}contributed by Nikita Gusev, 9/March/2021{]}}
\begin{itemize}
\item {} 
\sphinxAtStartPar
Crystal polymorphs are defined as substances that are chemically identical but exist in more than one crystal form. A polymorphic transition is a reversible transition of a solid crystalline phase at a certain temperature and pressure (the inversion point) to another phase of the same chemical composition with a different crystal structure.
For a monotropic system, plots of the solubility of the various polymorphs against temperature do not cross, meaning that different polymorphs always have different solubilities, and different stability.
Solution of a crystal is favored if the energy of dissolution is greater than the lattice energy, i.e., by dissolving, a compound forms a more thermodynamically stable state. In that sense, if the solubility of a polymorph is greater, means that this polymorph is less stable in its
crystalline state, and vice versa, the polymorph with the lowest solubility is thermodynamically more stable in its crystalline state.

\item {} 
\sphinxAtStartPar
The chemical potential of a solid crystalline phase is given by:

\end{itemize}
\begin{equation*}
\begin{split}
\mu_{i}^s\simeq\mu_0+RT\ln{x_i^*}
\end{split}
\end{equation*}
\sphinxAtStartPar
where \(x_i^*\) is the solubility.

\sphinxAtStartPar
From the equation above it can be seen than the lower is the solubility, the lower is the chemical potential and we know that the phase with the lowest chemical potential at any temperature is the most stable phase at that temperature. This again shows that a less soluble polymorph is more stable.


\subsection{Problem Statement 2}
\label{\detokenize{Exercise_W8_2021:problem-statement-2}}
\sphinxAtStartPar
An aqueous feed of 10000 kg h\(^{-1}\), saturated with \(BaCl_2\) at 100\(^o\)C, enters a crystallizer that can be simulated with the MSMPR model. The slurry leaves the crystallizer at 20\(^o\)C with crystals of the dihydrate. The crystallizer has a volume of 2.0 \(m^3\). From laboratory experiments, the crystal growth rate is essentially constant at \(4.0\times{10^{-7}}\) m/s. The solubility of \(BaCl_2\) is 58.3 g/100 g water at 100\(^o\) C and 35.7 g/100 g water at 20\(^o\) C. The molar mass of \(BaCl_2\) is 208.27 g/mol, the the dihydrate, \(BaCl_22H_2O\) has a molar mass of 244.31 g/mol. The density of the crystals is 3097 \(kg/m^3\), the density of the solution in the crystallizer is 1279 \(kg/m^3\)
\begin{itemize}
\item {} 
\sphinxAtStartPar
Determine the mass flow rate of crystals in the product stream.

\item {} 
\sphinxAtStartPar
Compute the predominant crystal size \(L_D\) in mm.

\end{itemize}


\subsection{Solution Trace}
\label{\detokenize{Exercise_W8_2021:id1}}
\sphinxAtStartPar
\sphinxstyleemphasis{{[}Nikita Gusev and Guilherme Pizarro Werner, 9/March/2021{]}}

\sphinxAtStartPar
A total material balance around the crystallizer (assuming steady state):
\begin{equation*}
\begin{split}
F=V+L+S
\end{split}
\end{equation*}
\sphinxAtStartPar
where:
\begin{itemize}
\item {} 
\sphinxAtStartPar
\(F\) = total feed rate (kg/h)

\item {} 
\sphinxAtStartPar
\(V\) = evaporation rate

\item {} 
\sphinxAtStartPar
\(L\) = Liquid phase flow rate kg/h of liquid

\item {} 
\sphinxAtStartPar
\(S\) = Crystals flow rate kg/h of crystals

\end{itemize}

\sphinxAtStartPar
Assuming evaporation is negligible and rearranging for \(L\):
\begin{equation*}
\begin{split}L=F-S\end{split}
\end{equation*}
\sphinxAtStartPar
Material balance for \(BaCl_2\):
\begin{equation*}
\begin{split}
x_F{F} = x_LL + x_SS
\end{split}
\end{equation*}\begin{itemize}
\item {} 
\sphinxAtStartPar
where \(x_F\), \(x_L\), \(x_S\) are the mass fractions of \(BaCl_2\) in feed, liquid and crystals, respectively.

\end{itemize}

\sphinxAtStartPar
F is known, and mass fractions can be found from the data, so we have two equations and two unknowns (L and S), and so there is 0 degrees of freedom and the problem can be solved.

\sphinxAtStartPar
Solving simultaneously using the expressions above:
\begin{equation*}
\begin{split}
x_FF= x_L(F - S) + x_SS
\end{split}
\end{equation*}
\sphinxAtStartPar
Rearranging:
\begin{equation*}
\begin{split}
S = F \frac{x_F-x_L}{x_S-x_L}
\end{split}
\end{equation*}
\sphinxAtStartPar
Now finding the mass fractions:
\begin{equation*}
\begin{split}
x_{F,L}=\frac{w_{F,L}}{w_{F,L} + 100}
\end{split}
\end{equation*}
\sphinxAtStartPar
where \(w_{F,L}\) is the solubility mass of solute per 100 g of water.
\begin{equation*}
\begin{split}
x_S=\frac{MM_{BaCl_2}}{MM_{BaCl_22H_2O}}
\end{split}
\end{equation*}
\begin{sphinxuseclass}{cell}\begin{sphinxVerbatimInput}

\begin{sphinxuseclass}{cell_input}
\begin{sphinxVerbatim}[commandchars=\\\{\}]
\PYG{c+c1}{\PYGZsh{}\PYGZsh{} Solution trace}
\PYG{c+c1}{\PYGZsh{} Variables}
\PYG{n}{wF}\PYG{o}{=}\PYG{l+m+mf}{58.3} \PYG{c+c1}{\PYGZsh{}g/100 g water }
\PYG{n}{wL}\PYG{o}{=}\PYG{l+m+mf}{35.7} \PYG{c+c1}{\PYGZsh{}g/100 g water }
\PYG{n}{MMBaCl2}\PYG{o}{=} \PYG{l+m+mf}{208.27} \PYG{c+c1}{\PYGZsh{}g/mol}
\PYG{n}{MMBaCl22H2O}\PYG{o}{=}\PYG{l+m+mf}{244.31} \PYG{c+c1}{\PYGZsh{}g/mol}
\PYG{n}{F}\PYG{o}{=}\PYG{l+m+mf}{1E4} \PYG{c+c1}{\PYGZsh{}kg/h}


\PYG{n}{xF}\PYG{o}{=}\PYG{n}{wF}\PYG{o}{/}\PYG{p}{(}\PYG{n}{wF}\PYG{o}{+}\PYG{l+m+mi}{100}\PYG{p}{)}\PYG{p}{;}
\PYG{n}{xL}\PYG{o}{=}\PYG{n}{wL}\PYG{o}{/}\PYG{p}{(}\PYG{n}{wL}\PYG{o}{+}\PYG{l+m+mi}{100}\PYG{p}{)}\PYG{p}{;}
\PYG{n}{xS}\PYG{o}{=}\PYG{n}{MMBaCl2}\PYG{o}{/}\PYG{n}{MMBaCl22H2O}\PYG{p}{;}

\PYG{n}{S}\PYG{o}{=}\PYG{n}{F}\PYG{o}{*}\PYG{p}{(}\PYG{n}{xF}\PYG{o}{\PYGZhy{}}\PYG{n}{xL}\PYG{p}{)}\PYG{o}{/}\PYG{p}{(}\PYG{n}{xS}\PYG{o}{\PYGZhy{}}\PYG{n}{xL}\PYG{p}{)}

\PYG{n}{L}\PYG{o}{=}\PYG{n}{F}\PYG{o}{\PYGZhy{}}\PYG{n}{S}

\PYG{n+nb}{print}\PYG{p}{(}\PYG{l+s+s2}{\PYGZdq{}}\PYG{l+s+se}{\PYGZbs{}n}\PYG{l+s+s2}{Feed mass fraction }\PYG{l+s+s2}{\PYGZdq{}}\PYG{p}{,} \PYG{l+s+sa}{f}\PYG{l+s+s2}{\PYGZdq{}}\PYG{l+s+si}{\PYGZob{}}\PYG{n}{xF}\PYG{l+s+si}{:}\PYG{l+s+s2}{.4}\PYG{l+s+si}{\PYGZcb{}}\PYG{l+s+s2}{\PYGZdq{}}\PYG{p}{,} \PYG{l+s+s2}{\PYGZdq{}}\PYG{l+s+s2}{ [\PYGZhy{}]}\PYG{l+s+s2}{\PYGZdq{}}\PYG{p}{)}
\PYG{n+nb}{print}\PYG{p}{(}\PYG{l+s+s2}{\PYGZdq{}}\PYG{l+s+s2}{Liquid phase mass fraction }\PYG{l+s+s2}{\PYGZdq{}}\PYG{p}{,} \PYG{l+s+sa}{f}\PYG{l+s+s2}{\PYGZdq{}}\PYG{l+s+si}{\PYGZob{}}\PYG{n}{xL}\PYG{l+s+si}{:}\PYG{l+s+s2}{.4}\PYG{l+s+si}{\PYGZcb{}}\PYG{l+s+s2}{\PYGZdq{}}\PYG{p}{,} \PYG{l+s+s2}{\PYGZdq{}}\PYG{l+s+s2}{[\PYGZhy{}]}\PYG{l+s+s2}{\PYGZdq{}}\PYG{p}{)}
\PYG{n+nb}{print}\PYG{p}{(}\PYG{l+s+s2}{\PYGZdq{}}\PYG{l+s+s2}{Crystal mass fraction }\PYG{l+s+s2}{\PYGZdq{}}\PYG{p}{,} \PYG{l+s+sa}{f}\PYG{l+s+s2}{\PYGZdq{}}\PYG{l+s+si}{\PYGZob{}}\PYG{n}{xS}\PYG{l+s+si}{:}\PYG{l+s+s2}{.4}\PYG{l+s+si}{\PYGZcb{}}\PYG{l+s+s2}{\PYGZdq{}}\PYG{p}{,} \PYG{l+s+s2}{\PYGZdq{}}\PYG{l+s+s2}{[\PYGZhy{}]}\PYG{l+s+s2}{\PYGZdq{}}\PYG{p}{)}
\PYG{n+nb}{print}\PYG{p}{(}\PYG{l+s+s2}{\PYGZdq{}}\PYG{l+s+s2}{Crystals flowrate: }\PYG{l+s+s2}{\PYGZdq{}}\PYG{p}{,} \PYG{l+s+sa}{f}\PYG{l+s+s2}{\PYGZdq{}}\PYG{l+s+si}{\PYGZob{}}\PYG{n}{S}\PYG{l+s+si}{:}\PYG{l+s+s2}{.4}\PYG{l+s+si}{\PYGZcb{}}\PYG{l+s+s2}{\PYGZdq{}}\PYG{p}{,} \PYG{l+s+s2}{\PYGZdq{}}\PYG{l+s+s2}{ [kg/h]}\PYG{l+s+s2}{\PYGZdq{}}\PYG{p}{)}
\PYG{n+nb}{print}\PYG{p}{(}\PYG{l+s+s2}{\PYGZdq{}}\PYG{l+s+s2}{Liquid phase flowrate: }\PYG{l+s+s2}{\PYGZdq{}}\PYG{p}{,} \PYG{l+s+sa}{f}\PYG{l+s+s2}{\PYGZdq{}}\PYG{l+s+si}{\PYGZob{}}\PYG{n}{L}\PYG{l+s+si}{:}\PYG{l+s+s2}{.4}\PYG{l+s+si}{\PYGZcb{}}\PYG{l+s+s2}{\PYGZdq{}}\PYG{p}{,}  \PYG{l+s+s2}{\PYGZdq{}}\PYG{l+s+s2}{ [kg/h]}\PYG{l+s+s2}{\PYGZdq{}}\PYG{p}{)}
\end{sphinxVerbatim}

\end{sphinxuseclass}\end{sphinxVerbatimInput}
\begin{sphinxVerbatimOutput}

\begin{sphinxuseclass}{cell_output}
\begin{sphinxVerbatim}[commandchars=\\\{\}]
Feed mass fraction  0.3683  [\PYGZhy{}]
Liquid phase mass fraction  0.2631 [\PYGZhy{}]
Crystal mass fraction  0.8525 [\PYGZhy{}]
Crystals flowrate:  1.785e+03  [kg/h]
Liquid phase flowrate:  8.215e+03  [kg/h]
\end{sphinxVerbatim}

\end{sphinxuseclass}\end{sphinxVerbatimOutput}

\end{sphinxuseclass}
\sphinxAtStartPar
The modal or dominant particle size can be computed as:
\begin{equation*}
\begin{split}
L=3G\tau
\end{split}
\end{equation*}
\sphinxAtStartPar
where:
\begin{itemize}
\item {} 
\sphinxAtStartPar
\(\tau\) = average residence time

\item {} 
\sphinxAtStartPar
\(G\) = growth rate

\end{itemize}

\sphinxAtStartPar
We need to find the average residence time:
\begin{equation*}
\begin{split}
\tau = \frac{V}{Q} 
\end{split}
\end{equation*}
\sphinxAtStartPar
where:
\begin{itemize}
\item {} 
\sphinxAtStartPar
\(V\)  = volume of the crystallizer

\item {} 
\sphinxAtStartPar
\(Q\)  = outlet volumetric flowrate

\end{itemize}

\sphinxAtStartPar
But first we need to calculate the outlet volumetric flowrate:
\begin{equation*}
\begin{split}
Q=\frac{S}{\rho_S}+\frac{L}{\rho_L}
\end{split}
\end{equation*}
\begin{sphinxuseclass}{cell}\begin{sphinxVerbatimInput}

\begin{sphinxuseclass}{cell_input}
\begin{sphinxVerbatim}[commandchars=\\\{\}]
\PYG{c+c1}{\PYGZsh{}\PYGZsh{} Solution trace}
\PYG{c+c1}{\PYGZsh{} Variables}
\PYG{n}{rho\PYGZus{}C}\PYG{o}{=}\PYG{l+m+mi}{3097} 
\PYG{n}{rho\PYGZus{}L}\PYG{o}{=}\PYG{l+m+mi}{1290}
\PYG{n}{V}\PYG{o}{=}\PYG{l+m+mi}{2} \PYG{c+c1}{\PYGZsh{}m3}
\PYG{n}{G}\PYG{o}{=}\PYG{l+m+mf}{4E\PYGZhy{}7} \PYG{c+c1}{\PYGZsh{}m/s}

\PYG{n}{Q}\PYG{o}{=}\PYG{n}{L}\PYG{o}{/}\PYG{n}{rho\PYGZus{}L}\PYG{o}{+}\PYG{n}{S}\PYG{o}{/}\PYG{n}{rho\PYGZus{}C} \PYG{c+c1}{\PYGZsh{}m3/h}

\PYG{n}{tau}\PYG{o}{=}\PYG{n}{V}\PYG{o}{/}\PYG{n}{Q}\PYG{o}{*}\PYG{l+m+mi}{3600} \PYG{c+c1}{\PYGZsh{}s}

\PYG{n}{LD}\PYG{o}{=}\PYG{l+m+mi}{3}\PYG{o}{*}\PYG{n}{G}\PYG{o}{*}\PYG{n}{tau}\PYG{o}{*}\PYG{l+m+mf}{1E3} \PYG{c+c1}{\PYGZsh{}mm}

\PYG{n+nb}{print}\PYG{p}{(}\PYG{l+s+s2}{\PYGZdq{}}\PYG{l+s+se}{\PYGZbs{}n}\PYG{l+s+s2}{The volumetric flowrate }\PYG{l+s+s2}{\PYGZdq{}}\PYG{p}{,} \PYG{l+s+sa}{f}\PYG{l+s+s2}{\PYGZdq{}}\PYG{l+s+si}{\PYGZob{}}\PYG{n}{Q}\PYG{l+s+si}{:}\PYG{l+s+s2}{.4}\PYG{l+s+si}{\PYGZcb{}}\PYG{l+s+s2}{\PYGZdq{}}\PYG{p}{,} \PYG{l+s+s2}{\PYGZdq{}}\PYG{l+s+s2}{[m3/h]}\PYG{l+s+s2}{\PYGZdq{}}\PYG{p}{)}
\PYG{n+nb}{print}\PYG{p}{(}\PYG{l+s+s2}{\PYGZdq{}}\PYG{l+s+s2}{The residence time }\PYG{l+s+s2}{\PYGZdq{}}\PYG{p}{,} \PYG{l+s+sa}{f}\PYG{l+s+s2}{\PYGZdq{}}\PYG{l+s+si}{\PYGZob{}}\PYG{n}{tau}\PYG{l+s+si}{:}\PYG{l+s+s2}{.4}\PYG{l+s+si}{\PYGZcb{}}\PYG{l+s+s2}{\PYGZdq{}}\PYG{p}{,} \PYG{l+s+s2}{\PYGZdq{}}\PYG{l+s+s2}{[s]}\PYG{l+s+s2}{\PYGZdq{}}\PYG{p}{)}
\PYG{n+nb}{print}\PYG{p}{(}\PYG{l+s+s2}{\PYGZdq{}}\PYG{l+s+s2}{The predominant crystal size is }\PYG{l+s+s2}{\PYGZdq{}}\PYG{p}{,} \PYG{l+s+sa}{f}\PYG{l+s+s2}{\PYGZdq{}}\PYG{l+s+si}{\PYGZob{}}\PYG{n}{LD}\PYG{l+s+si}{:}\PYG{l+s+s2}{.4}\PYG{l+s+si}{\PYGZcb{}}\PYG{l+s+s2}{\PYGZdq{}}\PYG{p}{,} \PYG{l+s+s2}{\PYGZdq{}}\PYG{l+s+s2}{ [mm]}\PYG{l+s+s2}{\PYGZdq{}}\PYG{p}{)}
\end{sphinxVerbatim}

\end{sphinxuseclass}\end{sphinxVerbatimInput}
\begin{sphinxVerbatimOutput}

\begin{sphinxuseclass}{cell_output}
\begin{sphinxVerbatim}[commandchars=\\\{\}]
The volumetric flowrate  6.945 [m3/h]
The residence time  1.037e+03 [s]
The predominant crystal size is  1.244  [mm]
\end{sphinxVerbatim}

\end{sphinxuseclass}\end{sphinxVerbatimOutput}

\end{sphinxuseclass}
\sphinxstepscope


\section{Week 9 \sphinxhyphen{} Exercise}
\label{\detokenize{Exercise_W9_2021:week-9-exercise}}\label{\detokenize{Exercise_W9_2021::doc}}

\subsection{Problem Statement}
\label{\detokenize{Exercise_W9_2021:problem-statement}}
\sphinxAtStartPar
The cumulative particle size distribution \(N(L)\) has been measured for two different MSMPR crystallization processes (A and B) and is reported in Fig. 1. Both processes are producing the same compact crystals, under different residence times and temperatures. The growth rate is known to be independent from both temperature and composition (\(G=const\)). The nucleation rate \(J\) instead depends exponentially on the inverse of the temperature through an Arrhenius\sphinxhyphen{}like expression (\(J=A\exp{-B/T}\),  where \(A\) and \(B\) are constants).




\subsubsection{Tasks:}
\label{\detokenize{Exercise_W9_2021:tasks}}\begin{itemize}
\item {} 
\sphinxAtStartPar
Compute and plot the number density, length, area, and normalised mass distribution for these two processes. Download the dataset: \sphinxhref{https://raw.githubusercontent.com/msalvalaglio/PSSP\_notebook/master/EX\_W9\_data.txt}{plain text}, \sphinxhref{https://raw.githubusercontent.com/msalvalaglio/PSSP\_notebook/master/EX\_W9\_data.xlsx}{excel xlsx format}.

\item {} 
\sphinxAtStartPar
Compute the dominant length \(L_D\) obtained from both processes.

\item {} 
\sphinxAtStartPar
Can you say anything about the residence time in the two crystallisers?

\item {} 
\sphinxAtStartPar
Which process is operating at the highest temperature?

\end{itemize}


\subsubsection{Solution}
\label{\detokenize{Exercise_W9_2021:solution}}\begin{itemize}
\item {} 
\sphinxAtStartPar
\sphinxhref{https://raw.githubusercontent.com/msalvalaglio/PSSP\_notebook/master/PSSP\_Solution\_Proposal\_PizarroWerner.pdf}{Contributed by Guilherme Pizarro Werner}

\item {} 
\sphinxAtStartPar
\sphinxhref{https://raw.githubusercontent.com/msalvalaglio/PSSP\_notebook/master/Week\_9\_Exercise\_Gusev.pdf}{Contributed by Nikita Gusev}

\end{itemize}

\sphinxstepscope


\section{Week 10 \sphinxhyphen{} Exercise}
\label{\detokenize{Exercise_W10_2021:week-10-exercise}}\label{\detokenize{Exercise_W10_2021::doc}}

\subsection{Problem Statement}
\label{\detokenize{Exercise_W10_2021:problem-statement}}
\sphinxAtStartPar
Consider a drying process evolving through two stages of \sphinxstyleemphasis{i} constant drying rate, and \sphinxstyleemphasis{ii} falling rate (first order kinetics).

\sphinxAtStartPar
During stage \sphinxstyleemphasis{i} the drying rate \sphinxstyleemphasis{per unit area} is constant, and indicated with \(R_c=m\times{f_c}=const\).

\sphinxAtStartPar
Where:
\begin{itemize}
\item {} 
\sphinxAtStartPar
\(m\) is a drying rate constant

\item {} 
\sphinxAtStartPar
\(f_c=w_c-w_e\)

\item {} 
\sphinxAtStartPar
\(w_c\) is the critical moisture content

\item {} 
\sphinxAtStartPar
\(w_e\) is the equilibrium moisture content

\item {} 
\sphinxAtStartPar
\(w_0\) is the initial moisture content

\end{itemize}

\sphinxAtStartPar
During stage \sphinxstyleemphasis{ii} the rate of drying is proportional to the free moisture content (\(f=w-w_e\)). The falling rate regime stops at the rquilibrium moisture content \(w_e\).


\subsection{Tasks}
\label{\detokenize{Exercise_W10_2021:tasks}}\begin{enumerate}
\sphinxsetlistlabels{\arabic}{enumi}{enumii}{}{.}%
\item {} 
\sphinxAtStartPar
Derive an expression for the time necessary to reduce the moisture content from \(w_0\) to \(w_c\).

\item {} 
\sphinxAtStartPar
Derive an expression for the time necessary to reduce the moisture content from \(w_c\) to \(w_e\) in the (linear) falling regime.

\item {} 
\sphinxAtStartPar
Derive an expression for the time necessary to reduce the moisture content from \(w_0\) to \(w_e<w<w_c\).

\item {} 
\sphinxAtStartPar
If a system, characterised by critical moisture content of 18\% and an equilibrium moisture content of 2\% takes 7300 seconds to reach a final moisture content of 5\% from an initial moisture content of 20\%, how long will it take, in the same conditions, to dry a system from 30\% to 3\%.

\item {} 
\sphinxAtStartPar
For the same system produce a plot of the drying time as a function of the kinetic constant m.

\end{enumerate}


\subsection{Solution}
\label{\detokenize{Exercise_W10_2021:solution}}

\subsubsection{Task 1.}
\label{\detokenize{Exercise_W10_2021:task-1}}
\sphinxAtStartPar
In order to derive an expression for the time necessary to obtain a reduction in the moisture content from an initial value \(w_0\) to the critical moisture content \(w_c\) we shall start by realising that for \(w\geq{w_c}\) the rate of drying (per unit area) \(R=const=R_c=m(w_c-w_e)\).

\sphinxAtStartPar
Hence we can write the rate of reduction in moisture content \(\frac{dw}{dt}\) as:
\begin{equation*}
\begin{split}
\frac{dw}{dt}=-R_cA
\end{split}
\end{equation*}
\sphinxAtStartPar
which can be solved by separating the variables
\begin{equation*}
\begin{split}
-\frac{dw}{R_cA}=dt
\end{split}
\end{equation*}
\sphinxAtStartPar
and integrating:
\begin{equation*}
\begin{split}
-\frac{1}{AR_c}\int_{w_0}^{w_c}dw=\int_{0}^{t_c}dt
\end{split}
\end{equation*}
\sphinxAtStartPar
leads to:
\begin{equation*}
\begin{split}
t_c=\frac{w_0-w_c}{AR_c}=\frac{w_0-w_c}{mA(w_c-w_e)}
\end{split}
\end{equation*}

\subsubsection{Task 2.}
\label{\detokenize{Exercise_W10_2021:task-2}}
\sphinxAtStartPar
In order to derive an expression for the time necessary to obtain a reduction in the moisture content from the critical value  \(w_c\) to the equilibrium moisture content let us begin first by deriving an expression for a generic \(w_e\geq{w}<w_c\).
In this case we know that the rate is linearly dependent on \(w\) accordig to the expression given in the text \(R=m(w-w_e)\).
The rate of reduction in moisture content \(\frac{dw}{dt}\) now becomes:
\begin{equation*}
\begin{split}
\frac{dw}{dt}=-mA(w-w_e)
\end{split}
\end{equation*}
\sphinxAtStartPar
which can be solved again by separating the variables
\begin{equation*}
\begin{split}
-\frac{1}{mA}\frac{dw}{w-w_e}=dt
\end{split}
\end{equation*}
\sphinxAtStartPar
and integrating:
\begin{equation*}
\begin{split}
-\frac{1}{mA}\int_{w_c}^{w^\prime}\frac{dw}{w-w_e}=\int_0^{t_w}dt
\end{split}
\end{equation*}
\sphinxAtStartPar
leads to:
\begin{equation*}
\begin{split}
t_w=\frac{1}{mA}\ln\left(\frac{w_c-w_e}{w-w_e}\right)
\end{split}
\end{equation*}
\sphinxAtStartPar
we note that since \({w-w_e}\) is at the denominator,
\begin{equation*}
\begin{split}
\lim_{w\rightarrow{w_e}}\frac{1}{mA}\ln\left(\frac{w_c-w_e}{w-w_e}\right)=+\infty
\end{split}
\end{equation*}
\sphinxAtStartPar
so for the moisture content approaching iuts equilibrium value the time necessary diverges.


\subsubsection{Task 3.}
\label{\detokenize{Exercise_W10_2021:task-3}}
\sphinxAtStartPar
The total time necessary to dry from an initial moiusture content \(w_0>w_c\) to a final moisture content \(w<w_c\) is equal to:
\begin{equation*}
\begin{split}
t_{tot}=\frac{1}{mA}\left(\frac{w_0-w_c}{w_c-w_e}+\ln\frac{w_c-w_e}{w-w_e}\right)
\end{split}
\end{equation*}

\subsubsection{Task 4.}
\label{\detokenize{Exercise_W10_2021:task-4}}
\sphinxAtStartPar
Using the expression derived earlier one can compute first
\begin{equation*}
\begin{split}
mA=\frac{1}{t_{tot,1}}\left(\frac{w_{0,1}-w_c}{w_c-w_e}+\ln\frac{w_c-w_e}{w_1-w_e}\right)
\end{split}
\end{equation*}
\sphinxAtStartPar
and then once \(mA\) is known,
\begin{equation*}
\begin{split}
t_{tot,2}=\frac{1}{mA}\left(\frac{w_{0,2}-w_c}{w_c-w_e}+\ln\frac{w_c-w_e}{w_2-w_e}\right)
\end{split}
\end{equation*}
\begin{sphinxuseclass}{cell}\begin{sphinxVerbatimInput}

\begin{sphinxuseclass}{cell_input}
\begin{sphinxVerbatim}[commandchars=\\\{\}]
\PYG{k+kn}{import}\PYG{+w}{ }\PYG{n+nn}{numpy}\PYG{+w}{ }\PYG{k}{as}\PYG{+w}{ }\PYG{n+nn}{np}

\PYG{n}{w\PYGZus{}e}\PYG{o}{=}\PYG{l+m+mf}{0.02}
\PYG{n}{w\PYGZus{}c}\PYG{o}{=}\PYG{l+m+mf}{0.18}

\PYG{n}{w0\PYGZus{}1}\PYG{o}{=}\PYG{l+m+mf}{0.2}
\PYG{n}{w\PYGZus{}1}\PYG{o}{=}\PYG{l+m+mf}{0.05}
\PYG{n}{t\PYGZus{}1}\PYG{o}{=}\PYG{l+m+mi}{7300}


\PYG{n}{mA}\PYG{o}{=}\PYG{p}{(}\PYG{l+m+mi}{1}\PYG{o}{/}\PYG{n}{t\PYGZus{}1}\PYG{p}{)}\PYG{o}{*}\PYG{p}{(}\PYG{p}{(}\PYG{n}{w0\PYGZus{}1}\PYG{o}{\PYGZhy{}}\PYG{n}{w\PYGZus{}c}\PYG{p}{)}\PYG{o}{/}\PYG{p}{(}\PYG{n}{w\PYGZus{}c}\PYG{o}{\PYGZhy{}}\PYG{n}{w\PYGZus{}e}\PYG{p}{)}\PYG{o}{+}\PYG{n}{np}\PYG{o}{.}\PYG{n}{log}\PYG{p}{(}\PYG{p}{(}\PYG{n}{w\PYGZus{}c}\PYG{o}{\PYGZhy{}}\PYG{n}{w\PYGZus{}e}\PYG{p}{)}\PYG{o}{/}\PYG{p}{(}\PYG{n}{w\PYGZus{}1}\PYG{o}{\PYGZhy{}}\PYG{n}{w\PYGZus{}e}\PYG{p}{)}\PYG{p}{)}\PYG{p}{)}
            
\PYG{n}{w0\PYGZus{}2} \PYG{o}{=} \PYG{l+m+mf}{0.3}
\PYG{n}{w\PYGZus{}2} \PYG{o}{=} \PYG{l+m+mf}{0.03}
            
\PYG{n}{t\PYGZus{}2}\PYG{o}{=}\PYG{p}{(}\PYG{l+m+mi}{1}\PYG{o}{/}\PYG{n}{mA}\PYG{p}{)}\PYG{o}{*}\PYG{p}{(}\PYG{p}{(}\PYG{n}{w0\PYGZus{}2}\PYG{o}{\PYGZhy{}}\PYG{n}{w\PYGZus{}c}\PYG{p}{)}\PYG{o}{/}\PYG{p}{(}\PYG{n}{w\PYGZus{}c}\PYG{o}{\PYGZhy{}}\PYG{n}{w\PYGZus{}e}\PYG{p}{)}\PYG{o}{+}\PYG{n}{np}\PYG{o}{.}\PYG{n}{log}\PYG{p}{(}\PYG{p}{(}\PYG{n}{w\PYGZus{}c}\PYG{o}{\PYGZhy{}}\PYG{n}{w\PYGZus{}e}\PYG{p}{)}\PYG{o}{/}\PYG{p}{(}\PYG{n}{w\PYGZus{}2}\PYG{o}{\PYGZhy{}}\PYG{n}{w\PYGZus{}e}\PYG{p}{)}\PYG{p}{)}\PYG{p}{)}
            
\PYG{n+nb}{print}\PYG{p}{(}\PYG{l+s+s2}{\PYGZdq{}}\PYG{l+s+se}{\PYGZbs{}n}\PYG{l+s+s2}{The drying time is}\PYG{l+s+s2}{\PYGZdq{}}\PYG{p}{,} \PYG{l+s+sa}{f}\PYG{l+s+s2}{\PYGZdq{}}\PYG{l+s+si}{\PYGZob{}}\PYG{n}{t\PYGZus{}2}\PYG{l+s+si}{:}\PYG{l+s+s2}{.4}\PYG{l+s+si}{\PYGZcb{}}\PYG{l+s+s2}{\PYGZdq{}}\PYG{p}{,} \PYG{l+s+s2}{\PYGZdq{}}\PYG{l+s+s2}{[s]}\PYG{l+s+s2}{\PYGZdq{}}\PYG{p}{)} 
\end{sphinxVerbatim}

\end{sphinxuseclass}\end{sphinxVerbatimInput}
\begin{sphinxVerbatimOutput}

\begin{sphinxuseclass}{cell_output}
\begin{sphinxVerbatim}[commandchars=\\\{\}]
The drying time is 1.429e+04 [s]
\end{sphinxVerbatim}

\end{sphinxuseclass}\end{sphinxVerbatimOutput}

\end{sphinxuseclass}

\subsubsection{Task 5.}
\label{\detokenize{Exercise_W10_2021:task-5}}
\sphinxAtStartPar
One cannot plot the dependence of the time on m alone, but the term mA can be used instead.

\begin{sphinxuseclass}{cell}\begin{sphinxVerbatimInput}

\begin{sphinxuseclass}{cell_input}
\begin{sphinxVerbatim}[commandchars=\\\{\}]
\PYG{k+kn}{import}\PYG{+w}{ }\PYG{n+nn}{numpy}\PYG{+w}{ }\PYG{k}{as}\PYG{+w}{ }\PYG{n+nn}{np}
\PYG{k+kn}{import}\PYG{+w}{ }\PYG{n+nn}{matplotlib}\PYG{n+nn}{.}\PYG{n+nn}{pyplot}\PYG{+w}{ }\PYG{k}{as}\PYG{+w}{ }\PYG{n+nn}{plt} 

\PYG{c+c1}{\PYGZsh{} Data}
\PYG{n}{N}\PYG{o}{=}\PYG{l+m+mi}{1000}
\PYG{n}{mA\PYGZus{}v} \PYG{o}{=} \PYG{n}{np}\PYG{o}{.}\PYG{n}{linspace}\PYG{p}{(}\PYG{n}{mA}\PYG{o}{/}\PYG{l+m+mi}{10}\PYG{p}{,}\PYG{l+m+mi}{10}\PYG{o}{*}\PYG{n}{mA}\PYG{p}{,} \PYG{n}{N}\PYG{p}{)}
\PYG{n}{t\PYGZus{}v}\PYG{o}{=}\PYG{p}{(}\PYG{l+m+mi}{1}\PYG{o}{/}\PYG{n}{mA\PYGZus{}v}\PYG{p}{)}\PYG{o}{*}\PYG{p}{(}\PYG{p}{(}\PYG{n}{w0\PYGZus{}2}\PYG{o}{\PYGZhy{}}\PYG{n}{w\PYGZus{}c}\PYG{p}{)}\PYG{o}{/}\PYG{p}{(}\PYG{n}{w\PYGZus{}c}\PYG{o}{\PYGZhy{}}\PYG{n}{w\PYGZus{}e}\PYG{p}{)}\PYG{o}{+}\PYG{n}{np}\PYG{o}{.}\PYG{n}{log}\PYG{p}{(}\PYG{p}{(}\PYG{n}{w\PYGZus{}c}\PYG{o}{\PYGZhy{}}\PYG{n}{w\PYGZus{}e}\PYG{p}{)}\PYG{o}{/}\PYG{p}{(}\PYG{n}{w\PYGZus{}2}\PYG{o}{\PYGZhy{}}\PYG{n}{w\PYGZus{}e}\PYG{p}{)}\PYG{p}{)}\PYG{p}{)}

\PYG{c+c1}{\PYGZsh{}Plotting}
\PYG{n}{figure}\PYG{o}{=}\PYG{n}{plt}\PYG{o}{.}\PYG{n}{figure}\PYG{p}{(}\PYG{p}{)}
\PYG{n}{axes} \PYG{o}{=} \PYG{n}{figure}\PYG{o}{.}\PYG{n}{add\PYGZus{}axes}\PYG{p}{(}\PYG{p}{[}\PYG{l+m+mf}{0.1}\PYG{p}{,}\PYG{l+m+mf}{0.1}\PYG{p}{,}\PYG{l+m+mi}{1}\PYG{p}{,}\PYG{l+m+mi}{1}\PYG{p}{]}\PYG{p}{)}
\PYG{n}{plt}\PYG{o}{.}\PYG{n}{xticks}\PYG{p}{(}\PYG{n}{fontsize}\PYG{o}{=}\PYG{l+m+mi}{14}\PYG{p}{)}
\PYG{n}{plt}\PYG{o}{.}\PYG{n}{yticks}\PYG{p}{(}\PYG{n}{fontsize}\PYG{o}{=}\PYG{l+m+mi}{14}\PYG{p}{)}
\PYG{n}{axes}\PYG{o}{.}\PYG{n}{plot}\PYG{p}{(}\PYG{n}{mA\PYGZus{}v}\PYG{p}{,}\PYG{n}{t\PYGZus{}v}\PYG{p}{,} \PYG{n}{marker}\PYG{o}{=}\PYG{l+s+s1}{\PYGZsq{}}\PYG{l+s+s1}{ }\PYG{l+s+s1}{\PYGZsq{}}\PYG{p}{)} 
\PYG{n}{axes}\PYG{o}{.}\PYG{n}{plot}\PYG{p}{(}\PYG{n}{mA}\PYG{p}{,}\PYG{n}{t\PYGZus{}2}\PYG{p}{,}\PYG{n}{marker}\PYG{o}{=}\PYG{l+s+s1}{\PYGZsq{}}\PYG{l+s+s1}{o}\PYG{l+s+s1}{\PYGZsq{}}\PYG{p}{,}\PYG{n}{markersize}\PYG{o}{=}\PYG{l+m+mi}{10}\PYG{p}{)}
\PYG{n}{axes}\PYG{o}{.}\PYG{n}{set\PYGZus{}xlabel}\PYG{p}{(}\PYG{l+s+s1}{\PYGZsq{}}\PYG{l+s+s1}{\PYGZdl{}mA\PYGZdl{} [\PYGZdl{}s\PYGZca{}}\PYG{l+s+s1}{\PYGZob{}}\PYG{l+s+s1}{\PYGZhy{}1\PYGZcb{}\PYGZdl{}]}\PYG{l+s+s1}{\PYGZsq{}}\PYG{p}{,} \PYG{n}{fontsize}\PYG{o}{=}\PYG{l+m+mi}{18}\PYG{p}{)}\PYG{p}{;}
\PYG{n}{axes}\PYG{o}{.}\PYG{n}{set\PYGZus{}ylabel}\PYG{p}{(}\PYG{l+s+s1}{\PYGZsq{}}\PYG{l+s+s1}{\PYGZdl{}t\PYGZdl{} [s]}\PYG{l+s+s1}{\PYGZsq{}}\PYG{p}{,}\PYG{n}{fontsize}\PYG{o}{=}\PYG{l+m+mi}{18}\PYG{p}{)}\PYG{p}{;}
\end{sphinxVerbatim}

\end{sphinxuseclass}\end{sphinxVerbatimInput}
\begin{sphinxVerbatimOutput}

\begin{sphinxuseclass}{cell_output}
\noindent\sphinxincludegraphics{{ff76b6e4ce0f1a546e41636192571507bb14a4aea2d3ec4ed649b92020511752}.png}

\end{sphinxuseclass}\end{sphinxVerbatimOutput}

\end{sphinxuseclass}

\subsection{Additional info}
\label{\detokenize{Exercise_W10_2021:additional-info}}\begin{enumerate}
\sphinxsetlistlabels{\arabic}{enumi}{enumii}{}{.}%
\item {} 
\sphinxAtStartPar
Asymptotic behaviour of the time for the moisture content reaching the equilibrium moisture content:

\end{enumerate}

\begin{sphinxuseclass}{cell}\begin{sphinxVerbatimInput}

\begin{sphinxuseclass}{cell_input}
\begin{sphinxVerbatim}[commandchars=\\\{\}]
\PYG{k+kn}{import}\PYG{+w}{ }\PYG{n+nn}{numpy}\PYG{+w}{ }\PYG{k}{as}\PYG{+w}{ }\PYG{n+nn}{np}
\PYG{k+kn}{import}\PYG{+w}{ }\PYG{n+nn}{matplotlib}\PYG{n+nn}{.}\PYG{n+nn}{pyplot}\PYG{+w}{ }\PYG{k}{as}\PYG{+w}{ }\PYG{n+nn}{plt} 

\PYG{c+c1}{\PYGZsh{} Data}
\PYG{n}{N}\PYG{o}{=}\PYG{l+m+mi}{1000}
\PYG{n}{ww} \PYG{o}{=} \PYG{n}{np}\PYG{o}{.}\PYG{n}{linspace}\PYG{p}{(}\PYG{n}{w\PYGZus{}e}\PYG{o}{*}\PYG{l+m+mf}{1.01}\PYG{p}{,}\PYG{n}{w\PYGZus{}c}\PYG{o}{*}\PYG{l+m+mf}{0.99}\PYG{p}{,} \PYG{n}{N}\PYG{p}{)}
\PYG{n}{t\PYGZus{}div}\PYG{o}{=}\PYG{p}{(}\PYG{l+m+mi}{1}\PYG{o}{/}\PYG{n}{mA}\PYG{p}{)}\PYG{o}{*}\PYG{p}{(}\PYG{p}{(}\PYG{n}{w0\PYGZus{}2}\PYG{o}{\PYGZhy{}}\PYG{n}{w\PYGZus{}c}\PYG{p}{)}\PYG{o}{/}\PYG{p}{(}\PYG{n}{w\PYGZus{}c}\PYG{o}{\PYGZhy{}}\PYG{n}{w\PYGZus{}e}\PYG{p}{)}\PYG{o}{+}\PYG{n}{np}\PYG{o}{.}\PYG{n}{log}\PYG{p}{(}\PYG{p}{(}\PYG{n}{w\PYGZus{}c}\PYG{o}{\PYGZhy{}}\PYG{n}{w\PYGZus{}e}\PYG{p}{)}\PYG{o}{/}\PYG{p}{(}\PYG{n}{ww}\PYG{o}{\PYGZhy{}}\PYG{n}{w\PYGZus{}e}\PYG{p}{)}\PYG{p}{)}\PYG{p}{)}


\PYG{c+c1}{\PYGZsh{}Plotting}
\PYG{n}{figure}\PYG{o}{=}\PYG{n}{plt}\PYG{o}{.}\PYG{n}{figure}\PYG{p}{(}\PYG{p}{)}
\PYG{n}{axes} \PYG{o}{=} \PYG{n}{figure}\PYG{o}{.}\PYG{n}{add\PYGZus{}axes}\PYG{p}{(}\PYG{p}{[}\PYG{l+m+mf}{0.1}\PYG{p}{,}\PYG{l+m+mf}{0.1}\PYG{p}{,}\PYG{l+m+mf}{1.2}\PYG{p}{,}\PYG{l+m+mf}{1.2}\PYG{p}{]}\PYG{p}{)}
\PYG{n}{plt}\PYG{o}{.}\PYG{n}{xticks}\PYG{p}{(}\PYG{n}{fontsize}\PYG{o}{=}\PYG{l+m+mi}{14}\PYG{p}{)}
\PYG{n}{plt}\PYG{o}{.}\PYG{n}{yticks}\PYG{p}{(}\PYG{n}{fontsize}\PYG{o}{=}\PYG{l+m+mi}{14}\PYG{p}{)}
\PYG{n}{axes}\PYG{o}{.}\PYG{n}{plot}\PYG{p}{(}\PYG{n}{ww}\PYG{p}{,}\PYG{n}{t\PYGZus{}div}\PYG{p}{,} \PYG{n}{marker}\PYG{o}{=}\PYG{l+s+s1}{\PYGZsq{}}\PYG{l+s+s1}{ }\PYG{l+s+s1}{\PYGZsq{}}\PYG{p}{)}
\PYG{n}{axes}\PYG{o}{.}\PYG{n}{plot}\PYG{p}{(}\PYG{n}{np}\PYG{o}{.}\PYG{n}{array}\PYG{p}{(}\PYG{p}{[}\PYG{n}{w\PYGZus{}e}\PYG{p}{,} \PYG{n}{w\PYGZus{}e}\PYG{p}{]}\PYG{p}{)}\PYG{p}{,}\PYG{n}{np}\PYG{o}{.}\PYG{n}{array}\PYG{p}{(}\PYG{p}{[}\PYG{l+m+mi}{0}\PYG{p}{,} \PYG{l+m+mi}{44000}\PYG{p}{]}\PYG{p}{)}\PYG{p}{)}
\PYG{n}{axes}\PYG{o}{.}\PYG{n}{set\PYGZus{}xlabel}\PYG{p}{(}\PYG{l+s+s1}{\PYGZsq{}}\PYG{l+s+s1}{\PYGZdl{}w\PYGZdl{} [\PYGZhy{}]}\PYG{l+s+s1}{\PYGZsq{}}\PYG{p}{,} \PYG{n}{fontsize}\PYG{o}{=}\PYG{l+m+mi}{18}\PYG{p}{)}\PYG{p}{;}
\PYG{n}{axes}\PYG{o}{.}\PYG{n}{set\PYGZus{}ylabel}\PYG{p}{(}\PYG{l+s+s1}{\PYGZsq{}}\PYG{l+s+s1}{\PYGZdl{}t\PYGZdl{} [s]}\PYG{l+s+s1}{\PYGZsq{}}\PYG{p}{,}\PYG{n}{fontsize}\PYG{o}{=}\PYG{l+m+mi}{18}\PYG{p}{)}\PYG{p}{;}
\end{sphinxVerbatim}

\end{sphinxuseclass}\end{sphinxVerbatimInput}
\begin{sphinxVerbatimOutput}

\begin{sphinxuseclass}{cell_output}
\noindent\sphinxincludegraphics{{5e643180601f325462e9e0168b8df5da546b363632f3c86dd560a204247a9d9e}.png}

\end{sphinxuseclass}\end{sphinxVerbatimOutput}

\end{sphinxuseclass}

\subsection{Contributions}
\label{\detokenize{Exercise_W10_2021:contributions}}
\sphinxAtStartPar
Nikita Gusev, 29 March 2021

\sphinxstepscope


\section{Week 11 \sphinxhyphen{} Exercise}
\label{\detokenize{Exercise_W11_2021:week-11-exercise}}\label{\detokenize{Exercise_W11_2021::doc}}

\subsection{Problem Statement}
\label{\detokenize{Exercise_W11_2021:problem-statement}}
\sphinxAtStartPar
Water (density 1000 kg m\(^{-3}\), viscosity 1 mN s m\(^{-2}\)) is used to fluidise a bed of packed spheres characterised by a diameter of 0.5 mm and a density of 1500 kg m\(^{-3}\).


\subsubsection{Tasks}
\label{\detokenize{Exercise_W11_2021:tasks}}\begin{itemize}
\item {} 
\sphinxAtStartPar
Demonstrate that, in laminar regime, the minimum fluidization velocity is always smaller than the terminal falling velocity.

\item {} 
\sphinxAtStartPar
Compute the minimum fluidisation velocity and the terminal falling velocity of the bed particles.

\item {} 
\sphinxAtStartPar
Sketch a plot of the \(\Delta{P}\) as a function of the fluid velocity, clearly indicating the minimum fluidisation velocity.

\item {} 
\sphinxAtStartPar
How would your result change if particles were twice as large?

\item {} 
\sphinxAtStartPar
What does the terminal falling velocity represent in the context of fluidisation?

\end{itemize}


\subsubsection{Solution}
\label{\detokenize{Exercise_W11_2021:solution}}

\paragraph{Task 1}
\label{\detokenize{Exercise_W11_2021:task-1}}
\sphinxAtStartPar
In the laminar regime, for spherical particles of diameter \(d\) the minimum fluidization velocity is given by:
\begin{equation*}
\begin{split}
v_{mf}=\frac{1}{180}\frac{e^3}{(1-e)}\frac{d^2(\rho_s-\rho_f)g}{\mu}
\end{split}
\end{equation*}
\sphinxAtStartPar
where, \(e\) is the void fraction, \(d\) the diameter of the particles, \(\rho_s\) is the density of the particles,  \(\rho_f\) is the density of the fluidising fluid, \(g\) is the gravity acceleration constant, and \(\mu\) the viscosity of the fluid.

\sphinxAtStartPar
The terminal falling velocity, conversely is determined by the steady state force balance on a single particle:
\begin{equation*}
\begin{split}
v_t^2=Cd^{-1}\frac{(\rho_s-\rho_f)}{\rho_f}\frac{4d}{3}g
\end{split}
\end{equation*}
\sphinxAtStartPar
where \(Cd\) is the drag coefficient.

\sphinxAtStartPar
If we introduce the definition of the drag coefficient in laminar regime given by the Stokes relation:
\begin{equation*}
\begin{split}
Cd=\frac{24}{Re_p}=\frac{24\mu}{\rho_fdv}
\end{split}
\end{equation*}
\sphinxAtStartPar
where \(Re_p\) is the Reynolds number and \(v\) the velocity of the fluid, we get:
\begin{equation*}
\begin{split}
v_t=\frac{\rho_fd}{24\mu}\frac{(\rho_s-\rho_f)}{\rho_f}\frac{4d}{3}g=\frac{1}{18}\frac{d^2(\rho_s-\rho_f)g}{\mu}
\end{split}
\end{equation*}
\sphinxAtStartPar
Now, if we take the ratio between \(v_{mf}\) and \(v_t\) we get:
\begin{equation*}
\begin{split}
\frac{v_{mf}}{v_t}=\frac{\frac{1}{180}\frac{e^3}{(1-e)}\frac{d^2(\rho_s-\rho_f)g}{\mu}}{\frac{1}{18}\frac{d^2(\rho_s-\rho_f)g}{\mu}}=\frac{1}{10}\frac{e^3}{(1-e)}
\end{split}
\end{equation*}
\sphinxAtStartPar
Since the void fraction \(0<e<1\) by definition, with a lower theoretical limit at \(\approx{0.36}\) (corresponding to a close\sphinxhyphen{}packing of equal spheres),  the term \(\frac{e^3}{(1-e)}\) is of order \(e^2\) and smaller than 1.

\sphinxAtStartPar
Hence
\begin{equation*}
\begin{split}
\frac{v_{mf}}{v_t}<1
\end{split}
\end{equation*}
\sphinxAtStartPar
and the minimum fluidization velocity is always smaller than the terminal Falling velocity, irrespective of \(d\), \(\rho_p\), \(\rho_f\) or \(\mu\).


\paragraph{Task 2 / 4}
\label{\detokenize{Exercise_W11_2021:task-2-4}}
\sphinxAtStartPar
In order to solve task 2 we need to estimate the void fraction. A typical approach whne the actual void fraction is not provided is to compute it as the ratio between the volume of a sphere of diameter \(d\) and a cube with the side of lenghth \(d\).
\begin{equation*}
\begin{split}
e=1-\frac{\pi}{6}
\end{split}
\end{equation*}
\begin{sphinxuseclass}{cell}\begin{sphinxVerbatimInput}

\begin{sphinxuseclass}{cell_input}
\begin{sphinxVerbatim}[commandchars=\\\{\}]
\PYG{k+kn}{import}\PYG{+w}{ }\PYG{n+nn}{numpy}\PYG{+w}{ }\PYG{k}{as}\PYG{+w}{ }\PYG{n+nn}{np}
\PYG{k+kn}{import}\PYG{+w}{ }\PYG{n+nn}{matplotlib}\PYG{n+nn}{.}\PYG{n+nn}{pyplot}\PYG{+w}{ }\PYG{k}{as}\PYG{+w}{ }\PYG{n+nn}{plt} 
\PYG{k+kn}{from}\PYG{+w}{ }\PYG{n+nn}{matplotlib}\PYG{n+nn}{.}\PYG{n+nn}{pyplot}\PYG{+w}{ }\PYG{k+kn}{import} \PYG{n}{cm}

\PYG{n}{rho\PYGZus{}f}\PYG{o}{=}\PYG{l+m+mi}{1000} \PYG{c+c1}{\PYGZsh{}}
\PYG{n}{rho\PYGZus{}s}\PYG{o}{=}\PYG{l+m+mi}{1500} \PYG{c+c1}{\PYGZsh{} kg m\PYGZdl{}\PYGZca{}\PYGZob{}\PYGZhy{}3\PYGZcb{}}

\PYG{n}{mu} \PYG{o}{=} \PYG{l+m+mf}{1E\PYGZhy{}3} \PYG{c+c1}{\PYGZsh{} N s m\PYGZdl{}\PYGZca{}\PYGZob{}\PYGZhy{}2\PYGZcb{}\PYGZdl{}}
\PYG{n}{d} \PYG{o}{=} \PYG{l+m+mf}{0.5} \PYG{o}{*} \PYG{l+m+mf}{1E\PYGZhy{}3} \PYG{c+c1}{\PYGZsh{}m}
\PYG{n}{g}\PYG{o}{=}\PYG{l+m+mf}{9.81} \PYG{c+c1}{\PYGZsh{} m /s\PYGZca{}2}
\PYG{n}{e}\PYG{o}{=}\PYG{l+m+mi}{1}\PYG{o}{\PYGZhy{}}\PYG{n}{np}\PYG{o}{.}\PYG{n}{pi}\PYG{o}{/}\PYG{l+m+mi}{6}
\PYG{n}{common}\PYG{o}{=}\PYG{n}{d}\PYG{o}{*}\PYG{o}{*}\PYG{l+m+mi}{2}\PYG{o}{*}\PYG{p}{(}\PYG{n}{rho\PYGZus{}s}\PYG{o}{\PYGZhy{}}\PYG{n}{rho\PYGZus{}f}\PYG{p}{)}\PYG{o}{*}\PYG{n}{g}\PYG{o}{/}\PYG{n}{mu} \PYG{c+c1}{\PYGZsh{}m/s}
\PYG{n}{v\PYGZus{}mf}\PYG{o}{=}\PYG{l+m+mi}{1}\PYG{o}{/}\PYG{l+m+mi}{180}\PYG{o}{*}\PYG{p}{(}\PYG{n}{e}\PYG{o}{*}\PYG{o}{*}\PYG{l+m+mi}{3}\PYG{o}{/}\PYG{p}{(}\PYG{l+m+mi}{1}\PYG{o}{\PYGZhy{}}\PYG{n}{e}\PYG{p}{)}\PYG{p}{)}\PYG{o}{*}\PYG{n}{common}
\PYG{n}{v\PYGZus{}t}\PYG{o}{=}\PYG{l+m+mi}{1}\PYG{o}{/}\PYG{l+m+mi}{18}\PYG{o}{*}\PYG{n}{common}

\PYG{n+nb}{print}\PYG{p}{(}\PYG{l+s+s2}{\PYGZdq{}}\PYG{l+s+se}{\PYGZbs{}n}\PYG{l+s+s2}{Diameter: }\PYG{l+s+s2}{\PYGZdq{}}\PYG{p}{,} \PYG{l+s+sa}{f}\PYG{l+s+s2}{\PYGZdq{}}\PYG{l+s+si}{\PYGZob{}}\PYG{n}{d}\PYG{l+s+si}{:}\PYG{l+s+s2}{.4}\PYG{l+s+si}{\PYGZcb{}}\PYG{l+s+s2}{\PYGZdq{}}\PYG{p}{,} \PYG{l+s+s2}{\PYGZdq{}}\PYG{l+s+s2}{ [mm]}\PYG{l+s+s2}{\PYGZdq{}}\PYG{p}{)}
\PYG{n+nb}{print}\PYG{p}{(}\PYG{l+s+s2}{\PYGZdq{}}\PYG{l+s+se}{\PYGZbs{}n}\PYG{l+s+s2}{Terminal Falling velocity: }\PYG{l+s+s2}{\PYGZdq{}}\PYG{p}{,} \PYG{l+s+sa}{f}\PYG{l+s+s2}{\PYGZdq{}}\PYG{l+s+si}{\PYGZob{}}\PYG{n}{v\PYGZus{}t}\PYG{l+s+si}{:}\PYG{l+s+s2}{.4}\PYG{l+s+si}{\PYGZcb{}}\PYG{l+s+s2}{\PYGZdq{}}\PYG{p}{,} \PYG{l+s+s2}{\PYGZdq{}}\PYG{l+s+s2}{ [m/s]}\PYG{l+s+s2}{\PYGZdq{}}\PYG{p}{)}
\PYG{n+nb}{print}\PYG{p}{(}\PYG{l+s+s2}{\PYGZdq{}}\PYG{l+s+s2}{Minimum Fluidisation velocity: }\PYG{l+s+s2}{\PYGZdq{}}\PYG{p}{,} \PYG{l+s+sa}{f}\PYG{l+s+s2}{\PYGZdq{}}\PYG{l+s+si}{\PYGZob{}}\PYG{n}{v\PYGZus{}mf}\PYG{l+s+si}{:}\PYG{l+s+s2}{.4}\PYG{l+s+si}{\PYGZcb{}}\PYG{l+s+s2}{\PYGZdq{}}\PYG{p}{,} \PYG{l+s+s2}{\PYGZdq{}}\PYG{l+s+s2}{ [m/s]}\PYG{l+s+s2}{\PYGZdq{}}\PYG{p}{)}

\PYG{n}{common}\PYG{o}{=}\PYG{p}{(}\PYG{l+m+mi}{2}\PYG{o}{*}\PYG{n}{d}\PYG{p}{)}\PYG{o}{*}\PYG{o}{*}\PYG{l+m+mi}{2}\PYG{o}{*}\PYG{p}{(}\PYG{n}{rho\PYGZus{}s}\PYG{o}{\PYGZhy{}}\PYG{n}{rho\PYGZus{}f}\PYG{p}{)}\PYG{o}{*}\PYG{n}{g}\PYG{o}{/}\PYG{n}{mu} \PYG{c+c1}{\PYGZsh{}m/s}
\PYG{n}{v\PYGZus{}mf}\PYG{o}{=}\PYG{l+m+mi}{1}\PYG{o}{/}\PYG{l+m+mi}{180}\PYG{o}{*}\PYG{p}{(}\PYG{n}{e}\PYG{o}{*}\PYG{o}{*}\PYG{l+m+mi}{3}\PYG{o}{/}\PYG{p}{(}\PYG{l+m+mi}{1}\PYG{o}{\PYGZhy{}}\PYG{n}{e}\PYG{p}{)}\PYG{p}{)}\PYG{o}{*}\PYG{n}{common}
\PYG{n}{v\PYGZus{}t}\PYG{o}{=}\PYG{l+m+mi}{1}\PYG{o}{/}\PYG{l+m+mi}{18}\PYG{o}{*}\PYG{n}{common}

\PYG{n+nb}{print}\PYG{p}{(}\PYG{l+s+s2}{\PYGZdq{}}\PYG{l+s+se}{\PYGZbs{}n}\PYG{l+s+s2}{Diameter: }\PYG{l+s+s2}{\PYGZdq{}}\PYG{p}{,} \PYG{l+s+sa}{f}\PYG{l+s+s2}{\PYGZdq{}}\PYG{l+s+si}{\PYGZob{}}\PYG{l+m+mi}{2}\PYG{o}{*}\PYG{n}{d}\PYG{l+s+si}{:}\PYG{l+s+s2}{.4}\PYG{l+s+si}{\PYGZcb{}}\PYG{l+s+s2}{\PYGZdq{}}\PYG{p}{,} \PYG{l+s+s2}{\PYGZdq{}}\PYG{l+s+s2}{ [mm]}\PYG{l+s+s2}{\PYGZdq{}}\PYG{p}{)}
\PYG{n+nb}{print}\PYG{p}{(}\PYG{l+s+s2}{\PYGZdq{}}\PYG{l+s+se}{\PYGZbs{}n}\PYG{l+s+s2}{Terminal Falling velocity: }\PYG{l+s+s2}{\PYGZdq{}}\PYG{p}{,} \PYG{l+s+sa}{f}\PYG{l+s+s2}{\PYGZdq{}}\PYG{l+s+si}{\PYGZob{}}\PYG{n}{v\PYGZus{}t}\PYG{l+s+si}{:}\PYG{l+s+s2}{.4}\PYG{l+s+si}{\PYGZcb{}}\PYG{l+s+s2}{\PYGZdq{}}\PYG{p}{,} \PYG{l+s+s2}{\PYGZdq{}}\PYG{l+s+s2}{ [m/s]}\PYG{l+s+s2}{\PYGZdq{}}\PYG{p}{)}
\PYG{n+nb}{print}\PYG{p}{(}\PYG{l+s+s2}{\PYGZdq{}}\PYG{l+s+s2}{Minimum Fluidisation velocity: }\PYG{l+s+s2}{\PYGZdq{}}\PYG{p}{,} \PYG{l+s+sa}{f}\PYG{l+s+s2}{\PYGZdq{}}\PYG{l+s+si}{\PYGZob{}}\PYG{n}{v\PYGZus{}mf}\PYG{l+s+si}{:}\PYG{l+s+s2}{.4}\PYG{l+s+si}{\PYGZcb{}}\PYG{l+s+s2}{\PYGZdq{}}\PYG{p}{,} \PYG{l+s+s2}{\PYGZdq{}}\PYG{l+s+s2}{ [m/s]}\PYG{l+s+se}{\PYGZbs{}n}\PYG{l+s+s2}{\PYGZdq{}}\PYG{p}{)}
\end{sphinxVerbatim}

\end{sphinxuseclass}\end{sphinxVerbatimInput}
\begin{sphinxVerbatimOutput}

\begin{sphinxuseclass}{cell_output}
\begin{sphinxVerbatim}[commandchars=\\\{\}]
Diameter:  0.0005  [mm]

Terminal Falling velocity:  0.06813  [m/s]
Minimum Fluidisation velocity:  0.001407  [m/s]

Diameter:  0.001  [mm]

Terminal Falling velocity:  0.2725  [m/s]
Minimum Fluidisation velocity:  0.005627  [m/s]
\end{sphinxVerbatim}

\end{sphinxuseclass}\end{sphinxVerbatimOutput}

\end{sphinxuseclass}

\paragraph{Task 3 / 4}
\label{\detokenize{Exercise_W11_2021:task-3-4}}
\sphinxAtStartPar
Without specifying the total lenght of the packed bed it is only possible to plot the pressure gradient \(\Delta{P}/l\).

\begin{sphinxuseclass}{cell}\begin{sphinxVerbatimInput}

\begin{sphinxuseclass}{cell_input}
\begin{sphinxVerbatim}[commandchars=\\\{\}]
\PYG{k+kn}{import}\PYG{+w}{ }\PYG{n+nn}{numpy}\PYG{+w}{ }\PYG{k}{as}\PYG{+w}{ }\PYG{n+nn}{np}
\PYG{k+kn}{import}\PYG{+w}{ }\PYG{n+nn}{matplotlib}\PYG{n+nn}{.}\PYG{n+nn}{pyplot}\PYG{+w}{ }\PYG{k}{as}\PYG{+w}{ }\PYG{n+nn}{plt} 
\PYG{k+kn}{from}\PYG{+w}{ }\PYG{n+nn}{matplotlib}\PYG{n+nn}{.}\PYG{n+nn}{pyplot}\PYG{+w}{ }\PYG{k+kn}{import} \PYG{n}{cm}

\PYG{c+c1}{\PYGZsh{} Data}
\PYG{n}{N}\PYG{o}{=}\PYG{l+m+mi}{1000}

\PYG{n}{v}\PYG{o}{=}\PYG{n}{np}\PYG{o}{.}\PYG{n}{linspace}\PYG{p}{(}\PYG{l+m+mi}{0}\PYG{p}{,}\PYG{n}{v\PYGZus{}mf}\PYG{p}{,} \PYG{n}{N}\PYG{p}{)}
\PYG{n}{color}\PYG{o}{=}\PYG{n+nb}{iter}\PYG{p}{(}\PYG{n}{cm}\PYG{o}{.}\PYG{n}{gist\PYGZus{}heat}\PYG{p}{(}\PYG{n}{np}\PYG{o}{.}\PYG{n}{linspace}\PYG{p}{(}\PYG{l+m+mi}{0}\PYG{p}{,}\PYG{l+m+mi}{1}\PYG{p}{,}\PYG{l+m+mi}{3}\PYG{p}{)}\PYG{p}{)}\PYG{p}{)}


\PYG{c+c1}{\PYGZsh{}Plotting}
\PYG{n}{figure}\PYG{o}{=}\PYG{n}{plt}\PYG{o}{.}\PYG{n}{figure}\PYG{p}{(}\PYG{p}{)}
\PYG{n}{axes} \PYG{o}{=} \PYG{n}{figure}\PYG{o}{.}\PYG{n}{add\PYGZus{}axes}\PYG{p}{(}\PYG{p}{[}\PYG{l+m+mf}{0.1}\PYG{p}{,}\PYG{l+m+mf}{0.1}\PYG{p}{,}\PYG{l+m+mf}{1.2}\PYG{p}{,}\PYG{l+m+mf}{1.2}\PYG{p}{]}\PYG{p}{)}
\PYG{n}{plt}\PYG{o}{.}\PYG{n}{xticks}\PYG{p}{(}\PYG{n}{fontsize}\PYG{o}{=}\PYG{l+m+mi}{14}\PYG{p}{)}
\PYG{n}{plt}\PYG{o}{.}\PYG{n}{yticks}\PYG{p}{(}\PYG{n}{fontsize}\PYG{o}{=}\PYG{l+m+mi}{14}\PYG{p}{)}
\PYG{n}{mDeltaP\PYGZus{}l\PYGZus{}mf}\PYG{o}{=}\PYG{n}{v\PYGZus{}mf}\PYG{o}{*}\PYG{l+m+mi}{180}\PYG{o}{*}\PYG{p}{(}\PYG{l+m+mi}{1}\PYG{o}{\PYGZhy{}}\PYG{n}{e}\PYG{p}{)}\PYG{o}{*}\PYG{o}{*}\PYG{l+m+mi}{2}\PYG{o}{/}\PYG{n}{e}\PYG{o}{*}\PYG{o}{*}\PYG{l+m+mi}{3}\PYG{o}{*}\PYG{n}{mu}\PYG{o}{/}\PYG{n}{d}\PYG{o}{*}\PYG{o}{*}\PYG{l+m+mi}{2} \PYG{c+c1}{\PYGZsh{} N/m}
\PYG{n}{mDeltaP\PYGZus{}l}\PYG{o}{=}\PYG{n}{v}\PYG{o}{*}\PYG{l+m+mi}{180}\PYG{o}{*}\PYG{p}{(}\PYG{l+m+mi}{1}\PYG{o}{\PYGZhy{}}\PYG{n}{e}\PYG{p}{)}\PYG{o}{*}\PYG{o}{*}\PYG{l+m+mi}{2}\PYG{o}{/}\PYG{n}{e}\PYG{o}{*}\PYG{o}{*}\PYG{l+m+mi}{3}\PYG{o}{*}\PYG{n}{mu}\PYG{o}{/}\PYG{n}{d}\PYG{o}{*}\PYG{o}{*}\PYG{l+m+mi}{2} \PYG{c+c1}{\PYGZsh{} N/m}
\PYG{n}{c}\PYG{o}{=}\PYG{n+nb}{next}\PYG{p}{(}\PYG{n}{color}\PYG{p}{)}
\PYG{n}{axes}\PYG{o}{.}\PYG{n}{plot}\PYG{p}{(}\PYG{n}{v}\PYG{p}{,}\PYG{n}{mDeltaP\PYGZus{}l}\PYG{p}{,} \PYG{n}{marker}\PYG{o}{=}\PYG{l+s+s1}{\PYGZsq{}}\PYG{l+s+s1}{ }\PYG{l+s+s1}{\PYGZsq{}}\PYG{p}{,}\PYG{n}{c}\PYG{o}{=}\PYG{n}{c}\PYG{p}{)} 
\PYG{n}{axes}\PYG{o}{.}\PYG{n}{plot}\PYG{p}{(}\PYG{n}{v\PYGZus{}mf}\PYG{p}{,}\PYG{n}{mDeltaP\PYGZus{}l\PYGZus{}mf}\PYG{p}{,}\PYG{n}{marker}\PYG{o}{=}\PYG{l+s+s1}{\PYGZsq{}}\PYG{l+s+s1}{o}\PYG{l+s+s1}{\PYGZsq{}}\PYG{p}{,}\PYG{n}{markersize}\PYG{o}{=}\PYG{l+m+mi}{10}\PYG{p}{,}\PYG{n}{c}\PYG{o}{=}\PYG{n}{c}\PYG{p}{)}


\PYG{c+c1}{\PYGZsh{} Double particle diameter}
\PYG{n}{c}\PYG{o}{=}\PYG{n+nb}{next}\PYG{p}{(}\PYG{n}{color}\PYG{p}{)}
\PYG{n}{mDeltaP\PYGZus{}l\PYGZus{}mf}\PYG{o}{=}\PYG{n}{v\PYGZus{}mf}\PYG{o}{*}\PYG{l+m+mi}{180}\PYG{o}{*}\PYG{p}{(}\PYG{l+m+mi}{1}\PYG{o}{\PYGZhy{}}\PYG{n}{e}\PYG{p}{)}\PYG{o}{*}\PYG{o}{*}\PYG{l+m+mi}{2}\PYG{o}{/}\PYG{n}{e}\PYG{o}{*}\PYG{o}{*}\PYG{l+m+mi}{3}\PYG{o}{*}\PYG{n}{mu}\PYG{o}{/}\PYG{p}{(}\PYG{l+m+mi}{2}\PYG{o}{*}\PYG{n}{d}\PYG{p}{)}\PYG{o}{*}\PYG{o}{*}\PYG{l+m+mi}{2} \PYG{c+c1}{\PYGZsh{} N/m}
\PYG{n}{mDeltaP\PYGZus{}l}\PYG{o}{=}\PYG{n}{v}\PYG{o}{*}\PYG{l+m+mi}{180}\PYG{o}{*}\PYG{p}{(}\PYG{l+m+mi}{1}\PYG{o}{\PYGZhy{}}\PYG{n}{e}\PYG{p}{)}\PYG{o}{*}\PYG{o}{*}\PYG{l+m+mi}{2}\PYG{o}{/}\PYG{n}{e}\PYG{o}{*}\PYG{o}{*}\PYG{l+m+mi}{3}\PYG{o}{*}\PYG{n}{mu}\PYG{o}{/}\PYG{p}{(}\PYG{l+m+mi}{2}\PYG{o}{*}\PYG{n}{d}\PYG{p}{)}\PYG{o}{*}\PYG{o}{*}\PYG{l+m+mi}{2} \PYG{c+c1}{\PYGZsh{} N/m}
\PYG{n}{axes}\PYG{o}{.}\PYG{n}{plot}\PYG{p}{(}\PYG{n}{v}\PYG{p}{,}\PYG{n}{mDeltaP\PYGZus{}l}\PYG{p}{,} \PYG{n}{marker}\PYG{o}{=}\PYG{l+s+s1}{\PYGZsq{}}\PYG{l+s+s1}{ }\PYG{l+s+s1}{\PYGZsq{}}\PYG{p}{,}\PYG{n}{c}\PYG{o}{=}\PYG{n}{c}\PYG{p}{)} 
\PYG{n}{axes}\PYG{o}{.}\PYG{n}{plot}\PYG{p}{(}\PYG{n}{v\PYGZus{}mf}\PYG{p}{,}\PYG{n}{mDeltaP\PYGZus{}l\PYGZus{}mf}\PYG{p}{,}\PYG{n}{marker}\PYG{o}{=}\PYG{l+s+s1}{\PYGZsq{}}\PYG{l+s+s1}{o}\PYG{l+s+s1}{\PYGZsq{}}\PYG{p}{,}\PYG{n}{markersize}\PYG{o}{=}\PYG{l+m+mi}{10}\PYG{p}{,}\PYG{n}{c}\PYG{o}{=}\PYG{n}{c}\PYG{p}{)}
\PYG{n}{axes}\PYG{o}{.}\PYG{n}{set\PYGZus{}xlabel}\PYG{p}{(}\PYG{l+s+s1}{\PYGZsq{}}\PYG{l+s+s1}{\PYGZdl{}v\PYGZdl{} [\PYGZdl{}m s\PYGZca{}}\PYG{l+s+s1}{\PYGZob{}}\PYG{l+s+s1}{\PYGZhy{}1\PYGZcb{}\PYGZdl{}]}\PYG{l+s+s1}{\PYGZsq{}}\PYG{p}{,} \PYG{n}{fontsize}\PYG{o}{=}\PYG{l+m+mi}{18}\PYG{p}{)}\PYG{p}{;}
\PYG{n}{axes}\PYG{o}{.}\PYG{n}{set\PYGZus{}ylabel}\PYG{p}{(}\PYG{l+s+s1}{\PYGZsq{}}\PYG{l+s+s1}{\PYGZdl{}\PYGZhy{}}\PYG{l+s+s1}{\PYGZbs{}}\PYG{l+s+s1}{Delta}\PYG{l+s+si}{\PYGZob{}P\PYGZcb{}}\PYG{l+s+s1}{\PYGZdl{} \PYGZdl{}}\PYG{l+s+s1}{\PYGZob{}}\PYG{l+s+s1}{l\PYGZca{}}\PYG{l+s+s1}{\PYGZob{}}\PYG{l+s+s1}{\PYGZhy{}1\PYGZcb{}\PYGZcb{}\PYGZdl{} [N/m\PYGZdl{}\PYGZca{}3\PYGZdl{}]}\PYG{l+s+s1}{\PYGZsq{}}\PYG{p}{,}\PYG{n}{fontsize}\PYG{o}{=}\PYG{l+m+mi}{18}\PYG{p}{)}\PYG{p}{;}
\end{sphinxVerbatim}

\end{sphinxuseclass}\end{sphinxVerbatimInput}
\begin{sphinxVerbatimOutput}

\begin{sphinxuseclass}{cell_output}
\noindent\sphinxincludegraphics{{318d1eab550381c92380dd1661832518ffa4bf993816f85350fccdbb68ffaec9}.png}

\end{sphinxuseclass}\end{sphinxVerbatimOutput}

\end{sphinxuseclass}
\sphinxAtStartPar
We can treat the bed length as a parameter, to show the behaviour of the \(\Delta{P}\) as a function of velocity:

\begin{sphinxuseclass}{cell}\begin{sphinxVerbatimInput}

\begin{sphinxuseclass}{cell_input}
\begin{sphinxVerbatim}[commandchars=\\\{\}]
\PYG{k+kn}{import}\PYG{+w}{ }\PYG{n+nn}{numpy}\PYG{+w}{ }\PYG{k}{as}\PYG{+w}{ }\PYG{n+nn}{np}
\PYG{k+kn}{import}\PYG{+w}{ }\PYG{n+nn}{matplotlib}\PYG{n+nn}{.}\PYG{n+nn}{pyplot}\PYG{+w}{ }\PYG{k}{as}\PYG{+w}{ }\PYG{n+nn}{plt} 
\PYG{k+kn}{from}\PYG{+w}{ }\PYG{n+nn}{matplotlib}\PYG{n+nn}{.}\PYG{n+nn}{pyplot}\PYG{+w}{ }\PYG{k+kn}{import} \PYG{n}{cm}


\PYG{c+c1}{\PYGZsh{} Data}
\PYG{n}{N}\PYG{o}{=}\PYG{l+m+mi}{1000}
\PYG{c+c1}{\PYGZsh{}Plotting}
\PYG{n}{figure}\PYG{o}{=}\PYG{n}{plt}\PYG{o}{.}\PYG{n}{figure}\PYG{p}{(}\PYG{p}{)}
\PYG{n}{axes} \PYG{o}{=} \PYG{n}{figure}\PYG{o}{.}\PYG{n}{add\PYGZus{}axes}\PYG{p}{(}\PYG{p}{[}\PYG{l+m+mf}{0.1}\PYG{p}{,}\PYG{l+m+mf}{0.1}\PYG{p}{,}\PYG{l+m+mf}{1.2}\PYG{p}{,}\PYG{l+m+mf}{1.2}\PYG{p}{]}\PYG{p}{)}
\PYG{n}{plt}\PYG{o}{.}\PYG{n}{xticks}\PYG{p}{(}\PYG{n}{fontsize}\PYG{o}{=}\PYG{l+m+mi}{14}\PYG{p}{)}
\PYG{n}{plt}\PYG{o}{.}\PYG{n}{yticks}\PYG{p}{(}\PYG{n}{fontsize}\PYG{o}{=}\PYG{l+m+mi}{14}\PYG{p}{)}

\PYG{c+c1}{\PYGZsh{} Parameter: bed length}
\PYG{n}{lengths}\PYG{o}{=}\PYG{n}{np}\PYG{o}{.}\PYG{n}{array}\PYG{p}{(}\PYG{p}{[}\PYG{l+m+mf}{0.5}\PYG{p}{,} \PYG{l+m+mi}{1}\PYG{p}{,} \PYG{l+m+mi}{3}\PYG{p}{,} \PYG{l+m+mi}{5}\PYG{p}{]}\PYG{p}{)}

\PYG{n}{color}\PYG{o}{=}\PYG{n+nb}{iter}\PYG{p}{(}\PYG{n}{cm}\PYG{o}{.}\PYG{n}{gist\PYGZus{}heat}\PYG{p}{(}\PYG{n}{np}\PYG{o}{.}\PYG{n}{linspace}\PYG{p}{(}\PYG{l+m+mi}{0}\PYG{p}{,}\PYG{l+m+mi}{1}\PYG{p}{,}\PYG{n}{np}\PYG{o}{.}\PYG{n}{size}\PYG{p}{(}\PYG{n}{lengths}\PYG{p}{)}\PYG{o}{+}\PYG{l+m+mi}{1}\PYG{p}{)}\PYG{p}{)}\PYG{p}{)}

\PYG{k}{for} \PYG{n}{l} \PYG{o+ow}{in} \PYG{n}{lengths}\PYG{p}{:} 
    \PYG{n}{c}\PYG{o}{=}\PYG{n+nb}{next}\PYG{p}{(}\PYG{n}{color}\PYG{p}{)}
    \PYG{n}{mDeltaP\PYGZus{}mf}\PYG{o}{=}\PYG{n}{mDeltaP\PYGZus{}l\PYGZus{}mf}\PYG{o}{*}\PYG{n}{l} \PYG{c+c1}{\PYGZsh{} N/m}
    \PYG{n}{mDeltaP}\PYG{o}{=}\PYG{n}{mDeltaP\PYGZus{}l}\PYG{o}{*}\PYG{n}{l} \PYG{c+c1}{\PYGZsh{} N/m}
    \PYG{n}{axes}\PYG{o}{.}\PYG{n}{plot}\PYG{p}{(}\PYG{n}{v}\PYG{p}{,}\PYG{n}{mDeltaP}\PYG{p}{,} \PYG{n}{marker}\PYG{o}{=}\PYG{l+s+s1}{\PYGZsq{}}\PYG{l+s+s1}{ }\PYG{l+s+s1}{\PYGZsq{}}\PYG{p}{,}\PYG{n}{c}\PYG{o}{=}\PYG{n}{c}\PYG{p}{)} 
    \PYG{n}{axes}\PYG{o}{.}\PYG{n}{plot}\PYG{p}{(}\PYG{n}{np}\PYG{o}{.}\PYG{n}{array}\PYG{p}{(}\PYG{p}{[}\PYG{n}{v\PYGZus{}mf}\PYG{p}{,} \PYG{l+m+mi}{2}\PYG{o}{*}\PYG{n}{v\PYGZus{}mf}\PYG{p}{]}\PYG{p}{)}\PYG{p}{,}\PYG{n}{mDeltaP\PYGZus{}mf}\PYG{o}{*}\PYG{n}{np}\PYG{o}{.}\PYG{n}{array}\PYG{p}{(}\PYG{p}{[}\PYG{l+m+mi}{1}\PYG{p}{,}\PYG{l+m+mi}{1}\PYG{p}{]}\PYG{p}{)}\PYG{p}{,} \PYG{n}{marker}\PYG{o}{=}\PYG{l+s+s1}{\PYGZsq{}}\PYG{l+s+s1}{ }\PYG{l+s+s1}{\PYGZsq{}}\PYG{p}{,}\PYG{n}{c}\PYG{o}{=}\PYG{n}{c}\PYG{p}{)} 
    \PYG{n}{axes}\PYG{o}{.}\PYG{n}{plot}\PYG{p}{(}\PYG{n}{v\PYGZus{}mf}\PYG{p}{,}\PYG{n}{mDeltaP\PYGZus{}mf}\PYG{p}{,}\PYG{n}{marker}\PYG{o}{=}\PYG{l+s+s1}{\PYGZsq{}}\PYG{l+s+s1}{s}\PYG{l+s+s1}{\PYGZsq{}}\PYG{p}{,}\PYG{n}{markersize}\PYG{o}{=}\PYG{l+m+mi}{10}\PYG{p}{,}\PYG{n}{c}\PYG{o}{=}\PYG{n}{c}\PYG{p}{)}

\PYG{n}{axes}\PYG{o}{.}\PYG{n}{set\PYGZus{}xlabel}\PYG{p}{(}\PYG{l+s+s1}{\PYGZsq{}}\PYG{l+s+s1}{\PYGZdl{}v\PYGZdl{} [\PYGZdl{}m s\PYGZca{}}\PYG{l+s+s1}{\PYGZob{}}\PYG{l+s+s1}{\PYGZhy{}1\PYGZcb{}\PYGZdl{}]}\PYG{l+s+s1}{\PYGZsq{}}\PYG{p}{,} \PYG{n}{fontsize}\PYG{o}{=}\PYG{l+m+mi}{18}\PYG{p}{)}\PYG{p}{;}
\PYG{n}{axes}\PYG{o}{.}\PYG{n}{set\PYGZus{}ylabel}\PYG{p}{(}\PYG{l+s+s1}{\PYGZsq{}}\PYG{l+s+s1}{\PYGZdl{}\PYGZhy{}}\PYG{l+s+s1}{\PYGZbs{}}\PYG{l+s+s1}{Delta}\PYG{l+s+si}{\PYGZob{}P\PYGZcb{}}\PYG{l+s+s1}{\PYGZdl{} [N/m\PYGZdl{}\PYGZca{}2\PYGZdl{}]}\PYG{l+s+s1}{\PYGZsq{}}\PYG{p}{,}\PYG{n}{fontsize}\PYG{o}{=}\PYG{l+m+mi}{18}\PYG{p}{)}\PYG{p}{;}
\end{sphinxVerbatim}

\end{sphinxuseclass}\end{sphinxVerbatimInput}
\begin{sphinxVerbatimOutput}

\begin{sphinxuseclass}{cell_output}
\noindent\sphinxincludegraphics{{6aee70493ba8c03ea8f83bce659a203b9bdf40317b6c677655adb76298e28526}.png}

\end{sphinxuseclass}\end{sphinxVerbatimOutput}

\end{sphinxuseclass}

\paragraph{Task 5}
\label{\detokenize{Exercise_W11_2021:task-5}}
\sphinxAtStartPar
The terminal falling velocity corresponds to the velocity at whixh thew drag froces counterbalance the net weight of the particles, over coming thisb velocity leads to net upward acceleration of the particles and corresponds to the instauration of the transport regime.







\renewcommand{\indexname}{Index}
\printindex
\end{document}